% Generated mechanically from frozen input p04/ja/varieties.tex.
% Do not edit this generated file; edit the bound locale source instead.

% OK, start here.
%

\setcounter{chapter}{32}
\chapter{多様体}



\phantomsection
\label{varieties-section-phantom}


\section{序論}
\label{varieties-section-introduction}

\noindent
本章では、多様体および、より一般に体上のスキームの研究を始める。
基本的な参考文献は \cite{EGA} である。








\section{記号}
\label{varieties-section-notation}

\noindent
本章を通して、文字 $k$ で基礎体を表す。









\section{多様体}
\label{varieties-section-varieties}

\noindent
Stacks プロジェクトでは、多様体の定義として次を用いる。

\begin{definition}
\label{varieties-definition-variety}
$k$ を体とする。{\it 多様体}とは、$X$ が $k$ 上のスキームであって、
$X$ が整であり、かつ構造射
$X \to \Spec(k)$ が分離的かつ有限型であるものをいう。
\end{definition}

\noindent
この定義には次の欠点がある。$k'/k$ を体の拡大とし、$X$ を
$k$ 上の多様体とする。このとき基底変換
$X_{k'} = X \times_{\Spec(k)} \Spec(k')$ は、必ずしも
$k'$ 上の多様体ではない。この現象は、より一般的な形で後続の各節において
詳しく論じる。二つの多様体の積も多様体とは限らない
（これは本質的には同じ現象である）。次に例を挙げる。

\begin{example}
\label{varieties-example-product-not-a-variety}
$k = \mathbf{Q}$ とする。$X = \Spec(\mathbf{Q}(i))$ および
$Y = \Spec(\mathbf{Q}(i))$ とする。このとき積
$X \times_{\Spec(k)} Y$、すなわち多様体 $X$ と $Y$ の積は可約なので、
多様体ではない。（これは
$X$ 二つの不交和と同型である。）
\end{example}

\noindent
しかし基礎体が代数閉ならば、多様体の積は多様体である。これは代数の章の
結果から従うが、そこでははるかに一般的な状況を扱っている。
ここでは簡単な直接証明も与える。

\begin{lemma}
\label{varieties-lemma-product-varieties}
\begin{slogan}
代数閉体上では、多様体の積は多様体である。
\end{slogan}
$k$ を代数閉体とする。$X$, $Y$ を $k$ 上の多様体とする。このとき
$X \times_{\Spec(k)} Y$ は $k$ 上の多様体である。
\end{lemma}

\begin{proof}
射 $X \times_{\Spec(k)} Y \to \Spec(k)$ は、分離的かつ有限型な射
$X \times_{\Spec(k)} Y \to Y \to \Spec(k)$ の合成であるから、
有限型かつ分離的である。Morphisms, Lemmas
\ref{morphisms-lemma-base-change-finite-type} および
\ref{morphisms-lemma-composition-finite-type}、ならびに
Schemes, Lemma \ref{schemes-lemma-separated-permanence} を参照せよ。
証明を終えるには、$X \times_{\Spec(k)} Y$ が整であることを示せば十分である。
$X = \bigcup_{i = 1, \ldots, n} U_i$、
$Y = \bigcup_{j = 1, \ldots, m} V_j$ を有限アフィン開被覆とする。
各 $U_i \times_{\Spec(k)} V_j$ が整であることを示せれば、
Properties, Lemmas \ref{properties-lemma-characterize-reduced}、
\ref{properties-lemma-characterize-irreducible}、および
\ref{properties-lemma-characterize-integral} により証明は完了する。
従ってアフィンの場合に帰着する。

\medskip\noindent
アフィンの場合は、次の代数的主張に移される：$A$, $B$ を整域である有限生成
$k$-代数とする。このとき $A \otimes_k B$ は整域である。矛盾を導くため、
$$
(\sum\nolimits_{i = 1, \ldots, n} a_i \otimes b_i)
(\sum\nolimits_{j = 1, \ldots, m} c_j \otimes d_j) = 0
$$
が $A \otimes_k B$ において成り立ち、両因子が $A \otimes_k B$ で非零であると
仮定する。$b_1, \ldots, b_n$ は $k$-線型独立な $B$ の元であり、
$d_1, \ldots, d_m$ も $k$-線型独立な $B$ の元であると仮定してよい。
もちろん、$a_1$ と $c_1$ が $A$ で非零であると仮定してもよい。
従って $D(a_1c_1) \subset \Spec(A)$ は空でない。Hilbert の零点定理
（Algebra, Theorem \ref{algebra-theorem-nullstellensatz}）により、極大イデアル
$\mathfrak m \subset A$ で $D(a_1c_1)$ に含まれ、かつ
$A/\mathfrak m = k$ となるものを取れる。最後の等式は $k$ が代数閉だからである。
$\overline{a}_i, \overline{c}_j$ で、$a_i, c_j$ の
$A/\mathfrak m = k$ における剰余類を表す。上の等式は
$$
(\sum\nolimits_{i = 1, \ldots, n} \overline{a}_i b_i)
(\sum\nolimits_{j = 1, \ldots, m} \overline{c}_j d_j) = 0
$$
となる。これは $\mathfrak m \in D(a_1c_1)$、
$b_1, \ldots, b_n$ と $d_1, \ldots, d_m$ の線型独立性、および
$B$ が整域であることに矛盾する。
\end{proof}







\section{多様体と有理写像}
\label{varieties-section-varieties-rational-maps}

\noindent
$k$ を体とする。$X$ と $Y$ を $k$ 上の多様体とする。
{\it $X$ から $Y$ への多様体の有理写像}という語句は、Morphisms,
Definition \ref{morphisms-definition-rational-map} で定義された
$\Spec(k)$-有理写像、すなわちスキーム $X$ からスキーム $Y$ への有理写像を
意味するものとする。慣例どおり、
「多様体の有理写像」という語句は、多様体の（共通の）基礎体を指示しない。
ただし一般のスキームについては、有理写像と、与えられた基底上の有理写像とを区別する。

\medskip\noindent
本節の表題は、次の基本定理を指している。

\begin{theorem}
\label{varieties-theorem-varieties-rational-maps}
$k$ を体とする。多様体と支配的有理写像の圏は、有限生成体拡大
$K/k$ の圏と同値である。
\end{theorem}

\begin{proof}
$X$ と $Y$ を多様体とし、その生成点をそれぞれ $x \in X$ と $y \in Y$ とする。
$X$ から $Y$ への支配的有理写像とは、ちょうど $x$ を $y$ に写す有理写像であることを
思い出そう（Morphisms, Definition \ref{morphisms-definition-dominant-rational}
およびその後の議論）。従って、支配的有理写像 $X \supset U \to Y$ が与えられると、
函数体の写像
$$
k(Y) = \kappa(y) = \mathcal{O}_{Y, y}
\longrightarrow
\mathcal{O}_{X, x} = \kappa(x) = k(X)
$$
を得る。逆に、このような $k$-代数写像は、始域と終域が体なので自動的に局所的であり、
Morphisms, Lemma \ref{morphisms-lemma-rational-map-finite-presentation} により
一意な支配的有理写像を定める。このようにして充満忠実な関手を得る。
証明を終えるには、有限生成体拡大 $K/k$ がすべて本質像に属することを示せば十分である。
$K/k$ は有限生成なので、有限型 $k$-代数 $A \subset K$ で、$K$ が $A$ の
分数体となるものが存在する。このとき $X = \Spec(A)$ は、函数体が $K$ である多様体である。
\end{proof}

\noindent
$k$ を体とする。$X$ と $Y$ を $k$ 上の多様体とする。
{\it $X$ と $Y$ は双有理な多様体である}という語句は、Morphisms,
Definition \ref{morphisms-definition-birational} で定義される意味で、$X$ と $Y$ が
$\Spec(k)$-双有理であることを意味するものとする。慣例どおり、「双有理な多様体」
という語句は、多様体の（共通の）基礎体を指示しない。ただし一般の既約スキームについては、
双有理であることと、与えられた基底上で双有理であることとを区別する。

\begin{lemma}
\label{varieties-lemma-birational-varieties}
$X$ と $Y$ を体 $k$ 上の多様体とする。次は同値である。
\begin{enumerate}
\item $X$ と $Y$ は双有理な多様体である。
\item 函数体 $k(X)$ と $k(Y)$ は同型である。
\item $X$ と $Y$ には、多様体として同型な空でない開集合が存在する。
\item 開集合 $U \subset X$ と、多様体の双有理射 $U \to Y$ が存在する。
\end{enumerate}
\end{lemma}

\begin{proof}
これは Morphisms, Lemma
\ref{morphisms-lemma-criterion-birational-finite-presentation} の特別な場合である。
\end{proof}







\section{体の変更と局所環}
\label{varieties-section-local-rings}

\noindent
基礎体の拡大の下で局所環に何が起こるかについて、いくつかの準備的結果を与える。

\begin{lemma}
\label{varieties-lemma-change-fields-flat}
$K/k$ を体の拡大とする。$X$ を $k$ 上のスキームとし、$Y = X_K$ と置く。
$y \in Y$ の像を $x \in X$ とすると、次が成り立つ。
\begin{enumerate}
\item $\mathcal{O}_{X, x} \to \mathcal{O}_{Y, y}$ は忠実平坦な局所環準同型である。
\item $\mathfrak p_0 = \Ker(\kappa(x) \otimes_k K \to \kappa(y))$ と置けば、
$\kappa(y) = \kappa(\mathfrak p_0)$ である。
\item $\mathcal{O}_{Y, y} = (\mathcal{O}_{X, x} \otimes_k K)_\mathfrak p$ である。ここで
$\mathfrak p \subset \mathcal{O}_{X, x} \otimes_k K$ は $\mathfrak p_0$ の逆像である。
\item
$\mathcal{O}_{Y, y}/\mathfrak m_x\mathcal{O}_{Y, y} =
(\kappa(x) \otimes_k K)_{\mathfrak p_0}$ である。
\end{enumerate}
\end{lemma}

\begin{proof}
$X = \Spec(A)$ がアフィンであると仮定してよい。このとき
$Y = \Spec(A \otimes_k K)$ である。$K$ は $k$ 上平坦なので、
$A \to A \otimes_k K$ は平坦である。従って $Y \to X$ は平坦であり、
さらに Algebra, Lemma \ref{algebra-lemma-local-flat-ff} を用いれば第 1 の主張を得る。
第 2 の主張は Schemes, Lemma
\ref{schemes-lemma-points-fibre-product} から従う。さて $y$ は素イデアル
$\mathfrak q \subset A \otimes_k K$ に、$x$ は
$\mathfrak r = A \cap \mathfrak q$ に対応する。このとき $\mathfrak p_0$ は、誘導される写像
$\kappa(\mathfrak r) \otimes_k K \to \kappa(\mathfrak q)$ の核である。
局所環上の写像は
$$
A_\mathfrak r \longrightarrow (A \otimes_k K)_\mathfrak q
$$
である。この写像を
$A_\mathfrak r \otimes_k K = (A \otimes_k K)_{\mathfrak r}$ を経由して分解すると、
$$
A_\mathfrak r \longrightarrow A_\mathfrak r \otimes_k K
\longrightarrow (A \otimes_k K)_\mathfrak q
$$
を得る。このとき第 2 の矢印は、ある素イデアルにおける局所化である。この素イデアルは
$\mathfrak p_0$ の逆像であり（詳細は省略する）、これで (3) が証明された。
(4) を見るには、(3) と、局所化および $- \otimes_k K$ が完全関手であることを用いる。
\end{proof}

\begin{lemma}
\label{varieties-lemma-change-fields-algebraic-dim}
記号は Lemma \ref{varieties-lemma-change-fields-flat} のとおりとする。
$X$ が $k$ 上局所有限型であると仮定する。このとき
$$
\dim(\mathcal{O}_{Y, y}/\mathfrak m_x\mathcal{O}_{Y, y}) =
\text{trdeg}_k(\kappa(x)) - \text{trdeg}_K(\kappa(y)) =
\dim(\mathcal{O}_{Y, y}) - \dim(\mathcal{O}_{X, x})
$$
である。
\end{lemma}

\begin{proof}
これは Algebra, Lemma
\ref{algebra-lemma-inequalities-under-field-extension} の言い換えである。
\end{proof}

\begin{lemma}
\label{varieties-lemma-change-fields-algebraic-unramified}
記号は Lemma \ref{varieties-lemma-change-fields-flat} のとおりとする。
$X$ が $k$ 上局所有限型であり、
$\dim(\mathcal{O}_{X, x}) = \dim(\mathcal{O}_{Y, y})$ であり、かつ
$\kappa(x) \otimes_k K$ が被約であると仮定する
（例えば $\kappa(x)/k$ が分離的であるか、$K/k$ が分離的ならばよい）。
このとき $\mathfrak m_x \mathcal{O}_{Y, y} = \mathfrak m_y$ である。
\end{lemma}

\begin{proof}
（括弧内の主張は Algebra, Lemma
\ref{algebra-lemma-separable-extension-preserves-reducedness} から従う。）
Lemmas \ref{varieties-lemma-change-fields-flat} と
\ref{varieties-lemma-change-fields-algebraic-dim} を合わせると、
$\mathcal{O}_{Y, y}/\mathfrak m_x \mathcal{O}_{Y, y}$ は次元 $0$ で被約である。
従ってこれは体である。
\end{proof}





\section{幾何学的被約スキーム}
\label{varieties-section-geometrically-reduced}

\noindent
$X$ が体上の被約スキームであっても、基礎体を拡大すると $X$ が非被約になることがある。
幾何学的被約スキームでは、このことは起こらない。

\begin{definition}
\label{varieties-definition-geometrically-reduced}
$k$ を体とする。$X$ を $k$ 上のスキームとする。
\begin{enumerate}
\item $x \in X$ を点とする。任意の体拡大を施した後でも、$X$ は
{\it $x$ において幾何学的被約}であるという。すなわち、任意の体拡大 $k'/k$ と、
任意の点 $x' \in X_{k'}$ で $x$ の上にあるものに対し、局所環
$\mathcal{O}_{X_{k'}, x'}$ が被約であることをいう。
\item $X$ は $k$ 上 {\it 幾何学的被約}であるという。すなわち、$X$ が
$X$ のすべての点で幾何学的被約であることをいう。
\end{enumerate}
\end{definition}

\noindent
これは初めには少し不可解に見えるかもしれないが、実際には代数の章で論じた概念と
同じである。次のいくつかの基本結果がその関係を説明する。
Lemma \ref{varieties-lemma-geometrically-reduced-dense-smooth-open} も参照せよ。
この補題から、体上の多様体について「幾何学的被約」であることと
「生成的に滑らか」であることが同値だと従う。

\begin{lemma}
\label{varieties-lemma-geometrically-reduced-at-point}
\begin{slogan}
幾何学的被約性は局所環上で判定できる。
\end{slogan}
$k$ を体とする。$X$ を $k$ 上のスキームとし、$x \in X$ とする。
次は同値である。
\begin{enumerate}
\item $X$ は $x$ において幾何学的被約である。
\item 環 $\mathcal{O}_{X, x}$ は $k$ 上幾何学的被約である（Algebra,
Definition \ref{algebra-definition-geometrically-reduced} を参照）。
\end{enumerate}
\end{lemma}

\begin{proof}
(1) を仮定する。これは特に $\mathcal{O}_{X, x}$ が被約であることを意味する。
$k'/k$ を有限純非分離体拡大とする。環
$\mathcal{O}_{X, x} \otimes_k k'$ を考える。Algebra, Lemma
\ref{algebra-lemma-p-ring-map} により、そのスペクトルは
$\mathcal{O}_{X, x}$ のスペクトルと同じである。従ってこれも局所環である
（Algebra, Lemma \ref{algebra-lemma-characterize-local-ring}）。従って、一意な点
$x' \in X_{k'}$ が $x$ の上にあり、Lemma \ref{varieties-lemma-change-fields-flat} により
$\mathcal{O}_{X_{k'}, x'} \cong \mathcal{O}_{X, x} \otimes_k k'$ である。
仮定により、これは被約環である。従って Algebra, Lemma
\ref{algebra-lemma-geometrically-reduced-finite-purely-inseparable-extension}
により (2) を得る。

\medskip\noindent
(2) を仮定する。$k'/k$ を体拡大とする。
$\Spec(k') \to \Spec(k)$ は全射なので、
$X_{k'} \to X$ も全射である
（Morphisms, Lemma \ref{morphisms-lemma-base-change-surjective}）。
任意の点 $x' \in X_{k'}$ で $x$ の上にあるものを取る。局所環
$\mathcal{O}_{X_{k'}, x'}$ は、Lemma \ref{varieties-lemma-change-fields-flat} により
$\mathcal{O}_{X, x} \otimes_k k'$ の局所化である。
従って仮定により被約であり、(1) が証明された。
\end{proof}

\noindent
標数零では、この概念は興味をもたない。

\begin{lemma}
\label{varieties-lemma-perfect-reduced}
$X$ を完全体 $k$ 上のスキームとする（例えば $k$ の標数が零である場合）。
$x \in X$ とする。$\mathcal{O}_{X, x}$ が被約ならば、$X$ は $x$ において
幾何学的被約である。$X$ が被約ならば、$X$ は $k$ 上幾何学的被約である。
\end{lemma}

\begin{proof}
第 1 の主張は Lemma \ref{varieties-lemma-geometrically-reduced-at-point}、
Algebra, Lemma \ref{algebra-lemma-separable-extension-preserves-reducedness}、
および完全体の定義（Algebra, Definition \ref{algebra-definition-perfect}）から従う。
第 2 の主張は第 1 の主張から従う。
\end{proof}

\begin{lemma}
\label{varieties-lemma-geometrically-reduced}
$k$ を標数 $p > 0$ の体とする。$X$ を $k$ 上のスキームとする。
次は同値である。
\begin{enumerate}
\item $X$ は幾何学的被約である。
\item $X_{k'}$ は任意の体拡大 $k'/k$ に対し被約である。
\item $X_{k'}$ は任意の有限純非分離体拡大 $k'/k$ に対し被約である。
\item $X_{k^{1/p}}$ は被約である。
\item $X_{k^{perf}}$ は被約である。
\item $X_{\bar k}$ は被約である。
\item 任意のアフィン開集合 $U \subset X$ に対し、環 $\mathcal{O}_X(U)$ は
幾何学的被約である（Algebra, Definition
\ref{algebra-definition-geometrically-reduced} を参照）。
\end{enumerate}
\end{lemma}

\begin{proof}
(1) を仮定する。このとき任意の体拡大 $k'/k$ および任意の点
$x' \in X_{k'}$ に対し、$X_{k'}$ の $x'$ における局所環は被約である。
言い換えれば $X_{k'}$ は被約である。従って (2) が成り立つ。

\medskip\noindent
(2) を仮定する。$U \subset X$ をアフィン開集合とする。このとき任意の体拡大
$k'/k$ に対してスキーム $X_{k'}$ は被約であり、従って
$U_{k'} = \Spec(\mathcal{O}(U)\otimes_k k')$ は被約であり、従って
$\mathcal{O}(U)\otimes_k k'$ は被約である（Properties,
Section \ref{properties-section-integral} を参照）。言い換えれば
$\mathcal{O}(U)$ は幾何学的被約なので、(7) が成り立つ。

\medskip\noindent
(7) を仮定する。任意の体拡大 $k'/k$ に対し、基底変換 $X_{k'}$ は、
環 $\mathcal{O}_X(U) \otimes_k k'$（ここで $U$ は $X$ のアフィン開集合）の
スペクトルを貼り合わせることで得られる（Schemes, Section
\ref{schemes-section-fibre-products} を参照）。従って $X_{k'}$ は被約である。
ゆえに (1) が成り立つ。

\medskip\noindent
これで (1)、(2)、(7) が同値であることが示された。これらは (3)、(4)、(5)、(6) とも
同値である。実際、$\mathcal{O}_X(U)$ に、アフィン開集合 $U \subset X$ について
Algebra, Lemma
\ref{algebra-lemma-geometrically-reduced-finite-purely-inseparable-extension}
を適用すればよい。
\end{proof}

\begin{lemma}
\label{varieties-lemma-check-only-finite-inseparable-extensions}
$k$ を標数 $p > 0$ の体とする。$X$ を $k$ 上のスキームとする。
$x \in X$ とする。次は同値である。
\begin{enumerate}
\item $X$ は $x$ において幾何学的被約である。
\item $\mathcal{O}_{X_{k'}, x'}$ は、$k'$ が $k$ の任意の有限純非分離体拡大で、
$x' \in X_{k'}$ が $x$ の上にある一意な点であるとき被約である。
\item $\mathcal{O}_{X_{k^{1/p}}, x'}$ は、$x' \in X_{k^{1/p}}$ が
$x$ の上にある一意な点であるとき被約である。
\item $\mathcal{O}_{X_{k^{perf}}, x'}$ は、$x' \in X_{k^{perf}}$ が
$x$ の上にある一意な点であるとき被約である。
\end{enumerate}
\end{lemma}

\begin{proof}
$k'/k$ が純非分離ならば、$X_{k'} \to X$ は台となる位相空間上の
同相写像を誘導することに注意せよ。Algebra, Lemma
\ref{algebra-lemma-p-ring-map} を参照せよ。従って、主張中で述べた
$x'$（$x$ の上にあるもの）の一意性が従う。さらに、この場合には
$\mathcal{O}_{X_{k'}, x'} = \mathcal{O}_{X, x} \otimes_k k'$ である。
従って補題は、上の Lemma \ref{varieties-lemma-geometrically-reduced-at-point} と Algebra, Lemma
\ref{algebra-lemma-geometrically-reduced-finite-purely-inseparable-extension}
から従う。
\end{proof}

\begin{lemma}
\label{varieties-lemma-geometrically-reduced-upstairs}
$k$ を体とする。$X$ を $k$ 上のスキームとする。$k'/k$ を体拡大とする。
$x \in X$ を点とし、点 $x' \in X_{k'}$ を $x$ の上に取る。
次は同値である。
\begin{enumerate}
\item $X$ は $x$ において幾何学的被約である。
\item $X_{k'}$ は $x'$ において幾何学的被約である。
\end{enumerate}
特に、$X$ が $k$ 上幾何学的被約であるための必要十分条件は、$X_{k'}$ が
$k'$ 上幾何学的被約であることである。
\end{lemma}

\begin{proof}
(1) が (2) を含意することは明らかである。(2) を仮定する。
$k''/k$ を有限純非分離体拡大とし、点 $x'' \in X_{k''}$ を $x$ の上に取る
（実際、これは一意である）。Fields, Lemma
\ref{fields-lemma-common-extension-field} により、$k'''/k$ を $k''/k$ と $k'/k$ の
共通体拡大として選ぶ。点 $x''' \in X_{k'''}$ を $x'$ と $x''$ の両方の上に取る。
局所環の写像
$$
\mathcal{O}_{X_{k''}, x''} \longrightarrow \mathcal{O}_{X_{k'''}, x'''}.
$$
を考える。これは平坦な局所環準同型であり、従って忠実平坦である。
(2) により、右辺の局所環は被約である。従って Algebra, Lemma
\ref{algebra-lemma-descent-reduced} により、
$\mathcal{O}_{X_{k''}, x''}$ は被約であると結論できる。従って Lemma
\ref{varieties-lemma-check-only-finite-inseparable-extensions} により、$X$ は $x$ において
幾何学的被約であると結論できる。
\end{proof}

\begin{lemma}
\label{varieties-lemma-geometrically-reduced-any-base-change}
$k$ を体とする。$X$, $Y$ を $k$ 上のスキームとする。
\begin{enumerate}
\item $X$ が $x$ において幾何学的被約であり、$Y$ が被約ならば、
$X \times_k Y$ は $x$ の上にあるすべての点で被約である。
\item $X$ が $k$ 上幾何学的被約であり、$Y$ が被約ならば、
$X \times_k Y$ は被約である。
\item $X$ と $Y$ が $k$ 上幾何学的被約ならば、$X \times_k Y$ は
幾何学的被約である。
\item $k$ が完全体であり、$X$ と $Y$ が被約ならば、$X \times_k Y$ は
被約である。
\item ここにさらに追加する。
\end{enumerate}
\end{lemma}

\begin{proof}
(1) を証明するには、Lemma \ref{varieties-lemma-geometrically-reduced-at-point} と Algebra, Lemma
\ref{algebra-lemma-geometrically-reduced-any-reduced-base-change} を合わせる。
(2) を証明するには、Lemma \ref{varieties-lemma-geometrically-reduced} と Algebra, Lemma
\ref{algebra-lemma-geometrically-reduced-any-reduced-base-change} を合わせる。
(3) を証明するには、
$(X \times_k Y)_{\overline{k}} =
X_{\overline{k}} \times_{\overline{k}} Y_{\overline{k}}$
に注意し、(2) および Lemma \ref{varieties-lemma-geometrically-reduced} を用いる。
(4) を証明するには、(3) と Lemma \ref{varieties-lemma-perfect-reduced} を合わせる。
\end{proof}

\begin{lemma}
\label{varieties-lemma-generic-points-geometrically-reduced}
$k$ を体とする。$X$ を $k$ 上のスキームとする。
\begin{enumerate}
\item $x' \leadsto x$ が特殊化であり、$X$ が $x$ において幾何学的被約ならば、
$X$ は $x'$ において幾何学的被約である。
\item $x \in X$ が次を満たすとする：(a) $\mathcal{O}_{X, x}$ は被約である。
(b) $x' \leadsto x$ なる各特殊化で、$x'$ が $X$ の既約成分の生成点であるものに対し、
スキーム $X$ は $x'$ において幾何学的被約である。このとき $X$ は $x$ において
幾何学的被約である。
\item $X$ が被約であり、$X$ の既約成分のすべての生成点において幾何学的被約ならば、
$X$ は幾何学的被約である。
\item $X$ が $k$ 上の多様体ならば、$X$ が幾何学的被約であるための必要十分条件は、
その函数体が $k$ の分離拡大であることである。
\end{enumerate}
\end{lemma}

\begin{proof}
(1) は Lemma \ref{varieties-lemma-geometrically-reduced-at-point} と次の事実から従う：
$A$ が幾何学的被約な $k$-代数ならば、$S^{-1}A$ は幾何学的被約な
$k$-代数である。ただし $S$ は $A$ の任意の積閉部分集合である。Algebra, Lemma
\ref{algebra-lemma-geometrically-reduced-permanence} を参照せよ。

\medskip\noindent
$A = \mathcal{O}_{X, x}$ と置く。(2) の仮定 (a) と (b) から、$A$ は被約であり、
$A_{\mathfrak q}$ は $k$ 上、各極小素イデアル $\mathfrak q$（$A$ のもの）に対して
幾何学的被約である。従って $A$ は $k$ 上幾何学的被約である。Algebra, Lemma
\ref{algebra-lemma-generic-points-geometrically-reduced} を参照せよ。
従って Lemma \ref{varieties-lemma-geometrically-reduced-at-point} により、$X$ は $x$ において
幾何学的被約である。

\medskip\noindent
(3) は (2) から直ちに従う。最後に (4) は、Algebra, Lemma
\ref{algebra-lemma-characterize-separable-field-extensions} により (3) から従う。
\end{proof}

\begin{lemma}
\label{varieties-lemma-Noetherian-geometrically-reduced-at-point}
$k$ を体とする。$X$ を $k$ 上のスキームとし、$x \in X$ とする。
$X$ が局所 Noether かつ $x$ において幾何学的被約であると仮定する。このとき、
開近傍 $U \subset X$ で $x$ を含み、$k$ 上幾何学的被約なものが存在する。
\end{lemma}

\begin{proof}
$X$ が局所 Noether かつ $x$ において幾何学的被約であると仮定する。
Properties, Lemma
\ref{properties-lemma-ring-affine-open-injective-local-ring} により、アフィン開近傍
$U \subset X$ で $x$ を含み、$R = \mathcal{O}_X(U) \to \mathcal{O}_{X, x}$ が単射となるものを取れる。
Lemma \ref{varieties-lemma-geometrically-reduced-at-point} により、仮定は
$\mathcal{O}_{X, x}$ が $k$ 上幾何学的被約であることを意味する。
Algebra, Lemma \ref{algebra-lemma-subalgebra-separable} により、これは $R$ が $k$ 上
幾何学的被約であることを含意し、従って $U$ は幾何学的被約である。
\end{proof}

\begin{example}
\label{varieties-example-not-geometrically-reduced}
$k = \mathbf{F}_p(s, t)$、すなわち素体の純超越拡大とする。多様体
$X = \Spec(k[x, y]/(1 + sx^p + ty^p))$ を考える。$k'/k$ を、$s$ と $t$ の双方が
$p$ 乗根を $k'$ にもつような任意の拡大とする。このとき基底変換 $X_{k'}$ は被約でない。
実際、環 $k'[x, y]/(1 + s x^p + ty^p)$ は元
$1 + s^{1/p}x + t^{1/p}y$ を含み、その $p$ 乗は零であるが、この元は零ではない
（イデアル $(1 + sx^p + ty^p)$ は次数 $< p$ の非零元を明らかに含まないからである）。
\end{example}

\begin{lemma}
\label{varieties-lemma-finite-extension-geometrically-reduced}
$k$ を体とする。$X \to \Spec(k)$ は局所有限型であるとする。
$X$ が有限個の既約成分をもつと仮定する。このとき有限純非分離拡大 $k'/k$ で、
$(X_{k'})_{red}$ が $k'$ 上幾何学的被約となるものが存在する。
\end{lemma}

\begin{proof}
この補題を証明するため、$X$ をその被約化 $X_{red}$ で置き換えてよい。
従って $X$ が被約であり、$k$ 上局所有限型であると仮定してよい。
$x_1, \ldots, x_n \in X$ を $X$ の既約成分の生成点とする。
任意の純非分離代数拡大 $k'/k$ に対し、射
$(X_{k'})_{red} \to X$ は同相写像であることに注意せよ。Algebra, Lemma
\ref{algebra-lemma-p-ring-map} を参照せよ。従って点 $x'_1, \ldots, x'_n$ は、
$x_1, \ldots, x_n$ の上にあり、$(X_{k'})_{red}$ の既約成分の生成点である。
$X$ は被約なので、局所環 $K_i = \mathcal{O}_{X, x_i}$ は体である。Algebra, Lemma
\ref{algebra-lemma-minimal-prime-reduced-ring} を参照せよ。
$X$ は $k$ 上局所有限型なので、体拡大 $K_i/k$ は有限生成である。
最後に、局所環 $\mathcal{O}_{(X_{k'})_{red}, x'_i}$ は体
$(K_i \otimes_k k')_{red}$ である。Algebra, Lemma
\ref{algebra-lemma-make-separable} により、有限純非分離拡大 $k'/k$ で、
$(K_i \otimes_k k')_{red}$ が $k'$ の分離体拡大となるものを取れる。
特に、Algebra, Lemma
\ref{algebra-lemma-characterize-separable-field-extensions} により、各
$(K_i \otimes_k k')_{red}$ は $k'$ 上幾何学的被約である。
ここで Lemma \ref{varieties-lemma-generic-points-geometrically-reduced} の (3) から、
$(X_{k'})_{red}$ は幾何学的被約であると従う。
\end{proof}







\section{幾何学的連結スキーム}
\label{varieties-section-geometrically-connected}

\noindent
$X$ が体上の連結スキームであっても、基礎体を拡大すると $X$ が非連結になることがある。
幾何学的連結スキームでは、このことは起こらない。

\begin{definition}
\label{varieties-definition-geometrically-connected}
$X$ を体 $k$ 上のスキームとする。$X$ は $k$ 上 {\it 幾何学的連結}であるというのは、
スキーム $X_{k'}$ が任意の体拡大 $k'$（$k$ の拡大）に対して連結であることをいう。
\end{definition}

\noindent
慣例により、連結な位相空間は空でない。従って、なおさら幾何学的連結スキームは
空でない。次は、幾何学的連結でない多様体の例である。

\begin{example}
\label{varieties-example-not-geometrically-irreducible}
$k = \mathbf{Q}$ とする。スキーム $X = \Spec(\mathbf{Q}(i))$ は
$\Spec(\mathbf{Q})$ 上の多様体である。しかし基底変換 $X_{\mathbf{C}}$ は
$\mathbf{C} \otimes_{\mathbf{Q}} \mathbf{Q}(i) \cong
\mathbf{C} \times \mathbf{C}$ のスペクトルであり、これは
$\Spec(\mathbf{C})$ 二つの不交和である。従って実際、これは幾何学的連結でない
多様体の例である。
\end{example}

\begin{lemma}
\label{varieties-lemma-geometrically-connected-check-after-extension}
$X$ を体 $k$ 上のスキームとする。$k'/k$ を体拡大とする。このとき、
$X$ が $k$ 上幾何学的連結であるための必要十分条件は、$X_{k'}$ が
$k'$ 上幾何学的連結であることである。
\end{lemma}

\begin{proof}
$X$ が $k$ 上幾何学的連結ならば、$X_{k'}$ が $k'$ 上幾何学的連結であることは
明らかである。逆に、任意の体拡大 $k''/k$ に対し、共通体拡大
$k'''/k'$ および $k'''/k''$ が存在する。Fields, Lemma
\ref{fields-lemma-common-extension-field} を参照せよ。射
$X_{k'''} \to X_{k''}$ は全射な体のスペクトル間の射の基底変換として全射なので、
$X_{k'''}$ の連結性から $X_{k''}$ の連結性が従う。従って、$X_{k'}$ が
$k'$ 上幾何学的連結ならば、$X$ は $k$ 上幾何学的連結である。
\end{proof}

\begin{lemma}
\label{varieties-lemma-bijection-connected-components}
$k$ を体とする。$X$, $Y$ を $k$ 上のスキームとする。
$X$ が $k$ 上幾何学的連結であると仮定する。このとき射影
$$
p : X \times_k Y \longrightarrow Y
$$
は連結成分の間の全単射を誘導する。
\end{lemma}

\begin{proof}
$p$ のスキーム論的ファイバーは、幾何学的連結スキーム $X$ の体拡大による
基底変換なので連結である。さらにスキーム論的ファイバーは集合論的ファイバーと
同相である。Schemes, Lemma \ref{schemes-lemma-fibre-topological} を参照せよ。
Morphisms, Lemma \ref{morphisms-lemma-scheme-over-field-universally-open} により、
写像 $p$ は開である。従って Topology, Lemma
\ref{topology-lemma-connected-fibres-connected-components} を適用すれば結論を得る。
\end{proof}

\begin{lemma}
\label{varieties-lemma-affine-geometrically-connected}
$k$ を体とし、$A$ を $k$-代数とする。このとき
$X = \Spec(A)$ が $k$ 上幾何学的連結であるための必要十分条件は、$A$ が
$k$ 上幾何学的連結であることである（Algebra, Definition
\ref{algebra-definition-geometrically-connected} を参照）。
\end{lemma}

\begin{proof}
定義から直ちに従う。
\end{proof}

\begin{lemma}
\label{varieties-lemma-separably-closed-field-connected-components}
$k'/k$ を体拡大とする。$X$ を $k$ 上のスキームとする。
$k$ が分離代数閉であると仮定する。このとき射 $X_{k'} \to X$ は連結成分の
全単射を誘導する。特に、$X$ が $k$ 上幾何学的連結であるための必要十分条件は、
$X$ が連結であることである。
\end{lemma}

\begin{proof}
$k$ は分離代数閉なので、$k'$ は $k$ 上幾何学的連結である。Algebra,
Lemma \ref{algebra-lemma-separably-closed-connected-implies-geometric} を参照せよ。
従って $Z = \Spec(k')$ は、上の Lemma
\ref{varieties-lemma-affine-geometrically-connected} により $k$ 上幾何学的連結である。
$X_{k'} = Z \times_k X$ なので、結果は Lemma
\ref{varieties-lemma-bijection-connected-components} の特別な場合である。
\end{proof}

\begin{lemma}
\label{varieties-lemma-characterize-geometrically-connected}
$k$ を体とする。$X$ を $k$ 上のスキームとする。$\overline{k}$ を $k$ の
分離代数閉包とする。このとき、$X$ が幾何学的連結であるための必要十分条件は、
基底変換 $X_{\overline{k}}$ が連結であることである。
\end{lemma}

\begin{proof}
$X_{\overline{k}}$ が連結であると仮定する。$k'/k$ を体拡大とする。
体拡大 $\overline{k}'/\overline{k}$ で、$k'$ が $\overline{k}'$ に
$k$ の拡大として埋め込まれるものが存在する。Lemma
\ref{varieties-lemma-separably-closed-field-connected-components} により、
$X_{\overline{k}'}$ は連結である。$X_{\overline{k}'} \to X_{k'}$ は全射なので、
$X_{k'}$ は連結であると結論でき、所望の結果を得る。
\end{proof}

\begin{lemma}
\label{varieties-lemma-descend-open}
$k$ を体とする。$X$ を $k$ 上のスキームとする。$A$ を $k$-代数とする。
$V \subset X_A$ を準コンパクト開集合とする。このとき有限生成 $k$-部分代数
$A' \subset A$ と準コンパクト開集合 $V' \subset X_{A'}$ で、
$V = V'_A$ となるものが存在する。
\end{lemma}

\begin{proof}
$X$ がさらに準分離的ならば、これは Limits, Lemma
\ref{limits-lemma-descend-opens} から従うことを注意しておく。
$U_1, \ldots, U_n$ を $X$ の有限個のアフィン開集合で、
$V \subset \bigcup U_{i, A}$ となるものとする。$U_i = \Spec(R_i)$ と書く。
$V$ は準コンパクトなので、
有限個の $f_{ij} \in R_i \otimes_k A$（$j = 1, \ldots, n_i$）で、
$V = \bigcup_i \bigcup_{j = 1, \ldots, n_i} D(f_{ij})$ となるものを取れる。
ここで $D(f_{ij}) \subset U_{i, A}$ は対応する標準開集合である。（
$V \cap U_{i, A}$ が $D(f_{ij})$（$j = 1, \ldots, n_i$）の和であるとは
主張していない。）有限生成 $k$-部分代数 $A' \subset A$ で、$f_{ij}$ がある
$f_{ij}' \in R_i \otimes_k A'$ の像となるものを取れることは明らかである。
$V' = \bigcup D(f_{ij}')$ と置く。これは $X_{A'}$ の準コンパクト開集合である。
$\pi : X_A \to X_{A'}$ を標準射とする。
$\pi(V) \subset V'$ である。実際、$\pi(D(f_{ij})) \subset D(f_{ij}')$ だからである。
$x \in X_A$ が $\pi(x) \in V'$ を満たすならば、
$\pi(x) \in D(f_{ij}')$ となるある $i, j$ があり、
$x \in D(f_{ij})$ であることが分かる。実際、$f_{ij}'$ は $f_{ij}$ に写る。
従って所望のとおり $V = \pi^{-1}(V')$ である。
\end{proof}

\noindent
$k$ を体とする。$\overline{k}/k$ を（無限でもよい）Galois 拡大とする。
例えば $\overline{k}$ は $k$ の分離代数閉包でもよい。
任意の $\sigma \in \text{Gal}(\overline{k}/k)$ に対し、対応する自己同型
$
\Spec(\sigma) :
\Spec(\overline{k})
\longrightarrow
\Spec(\overline{k})
$
を得る。
$\Spec(\sigma) \circ \Spec(\tau) = \Spec(\tau \circ \sigma)$
であることに注意せよ。従って、反対群のスキーム $\Spec(\overline{k})$ への作用
$$
\text{Gal}(\overline{k}/k)^{opp} \times \Spec(\overline{k})
\longrightarrow
\Spec(\overline{k})
$$
を得る。$X$ を $k$ 上のスキームとする。定義により
$X_{\overline{k}} =
\Spec(\overline{k}) \times_{\Spec(k)} X$
なので、上の作用は標準的作用
\begin{equation}
\label{varieties-equation-galois-action-base-change-kbar}
\text{Gal}(\overline{k}/k)^{opp} \times X_{\overline{k}}
\longrightarrow
X_{\overline{k}}.
\end{equation}
を誘導する。

\begin{lemma}
\label{varieties-lemma-Galois-action-quasi-compact-open}
$k$ を体とする。$X$ を $k$ 上のスキームとする。$\overline{k}$ を $k$ の
（無限でもよい）Galois 拡大とする。$V \subset X_{\overline{k}}$ を準コンパクト開集合とする。
このとき、次が成り立つ。
\begin{enumerate}
\item 有限部分拡大 $\overline{k}/k'/k$ と準コンパクト開集合
$V' \subset X_{k'}$ で、$V = (V')_{\overline{k}}$ となるものが存在する。
\item 開部分群 $H \subset \text{Gal}(\overline{k}/k)$ で、
$\sigma(V) = V$ がすべての $\sigma \in H$ に対して成り立つものが存在する。
\end{enumerate}
\end{lemma}

\begin{proof}
Lemma \ref{varieties-lemma-descend-open} により、有限部分拡大
$k'/k \subset \overline{k}$ と、開集合 $V' \subset X_{k'}$ で $V$ に
引き戻されるものが存在する。これで (1) が証明された。
$\text{Gal}(\overline{k}/k')$ は $\text{Gal}(\overline{k}/k)$ の開部分群なので、
(2) も明らかである。
\end{proof}







\begin{lemma}
\label{varieties-lemma-closed-fixed-by-Galois}
$k$ を体とする。$\overline{k}/k$ を（無限でもよい）Galois 拡大とする。
$X$ を $k$ 上のスキームとする。$\overline{T} \subset X_{\overline{k}}$ が
次の性質をもつとする。
\begin{enumerate}
\item $\overline{T}$ は $X_{\overline{k}}$ の閉部分集合である。
\item 任意の $\sigma \in \text{Gal}(\overline{k}/k)$ に対し
$\sigma(\overline{T}) = \overline{T}$ である。
\end{enumerate}
このとき閉部分集合 $T \subset X$ で、その $X_{\overline{k}}$ における逆像が
$\overline{T}$ となるものが存在する。
\end{lemma}

\begin{proof}
この補題は直ちに $X = \Spec(A)$ がアフィンである場合に帰着する。この場合、
$\overline{I} \subset A \otimes_k \overline{k}$ を $\overline{T}$ に対応する根基イデアルとする。
仮定 (2) から、$\sigma(\overline{I}) = \overline{I}$ がすべての
$\sigma \in \text{Gal}(\overline{k}/k)$ に対して成り立つ。$x \in \overline{I}$ を取る。
有限 Galois 拡大 $k'/k$ で $\overline{k}$ に含まれ、
$x \in A \otimes_k k'$ となるものが存在する。$G = \text{Gal}(k'/k)$ と置き、
$$
P(T) = \prod\nolimits_{\sigma \in G} (T - \sigma(x)) \in (A \otimes_k k')[T]
$$
と置く。$P(T)$ がモニックであり、実際に
$(A \otimes_k k')^G[T] = A[T]$ の元であることは明らかである
（Galois 理論の基本事項による）。さらに
$P(T) = T^d + a_1T^{d - 1} + \ldots + a_d$ と書けば、
$a_i \in I := A \cap \overline{I}$ であることが分かる。
$P(x) = 0$ と $a_i \in I$ を合わせると、
$x^d = - a_1 x^{d - 1} - \ldots  - a_d \in I(A \otimes_k \overline{k})$ を得る。
従って $x$ は $I(A \otimes_k \overline{k})$ の根基に含まれる。
ゆえに $\overline{I}$ は $I(A \otimes_k \overline{k})$ の根基であり、
$T = V(I)$ と置けば解を得る。
\end{proof}

\begin{lemma}
\label{varieties-lemma-characterize-geometrically-disconnected}
$k$ を体とする。$X$ を $k$ 上のスキームとする。次は同値である。
\begin{enumerate}
\item $X$ は幾何学的連結である。
\item 任意の有限分離体拡大 $k'/k$ に対し、スキーム $X_{k'}$ は連結である。
\end{enumerate}
\end{lemma}

\begin{proof}
(1) が (2) を含意することは定義から直ちに従う。$X$ が幾何学的連結でないと仮定する。
$k \subset \overline{k}$ を $k$ の分離代数閉包とする。Lemma
\ref{varieties-lemma-characterize-geometrically-connected} により、$X_{\overline{k}}$ は非連結である。
$X_{\overline{k}} = \overline{U} \amalg \overline{V}$ と書く。ここで
$\overline{U}$ と $\overline{V}$ は開、閉、かつ空でない。

\medskip\noindent
$W \subset X$ を任意の準コンパクト開集合とする。このとき
$W_{\overline{k}} \cap \overline{U}$ と $W_{\overline{k}} \cap \overline{V}$ は
$W_{\overline{k}}$ において開かつ閉である。特に
$W_{\overline{k}} \cap \overline{U}$ と $W_{\overline{k}} \cap \overline{V}$ は
準コンパクトであり、Lemma \ref{varieties-lemma-Galois-action-quasi-compact-open} により、
$W_{\overline{k}} \cap \overline{U}$ と $W_{\overline{k}} \cap \overline{V}$ はともに
有限部分拡大上で定義され、$\text{Gal}(\overline{k}/k)$ の開部分群の下で不変である。
以下ではこれを断りなく用いる。

\medskip\noindent
$W_0 \subset X$ を、$W_{0, \overline{k}} \cap \overline{U}$ と
$W_{0, \overline{k}} \cap \overline{V}$ がともに空でない準コンパクト開集合として取る。
有限部分拡大 $\overline{k}/k'/k$ と、開かつ閉な部分集合への分解
$W_{0, k'} = U_0' \amalg V_0'$ で、
$W_{0, \overline{k}} \cap \overline{U} = (U'_0)_{\overline{k}}$ および
$W_{0, \overline{k}} \cap \overline{V} = (V'_0)_{\overline{k}}$ となるものを選ぶ。
$H = \text{Gal}(\overline{k}/k') \subset \text{Gal}(\overline{k}/k)$ と置く。特に
$\sigma(W_{0, \overline{k}} \cap \overline{U}) =
W_{0, \overline{k}} \cap \overline{U}$ であり、$\overline{V}$ についても同様である。

\medskip\noindent
上のように $W_0$, $k'$ を選んだ後、任意の準コンパクト開集合 $W \subset X$ に対して
$$
U_W =
\bigcap\nolimits_{\sigma \in H} \sigma(W_{\overline{k}} \cap \overline{U}),
\quad
V_W =
\bigcup\nolimits_{\sigma \in H} \sigma(W_{\overline{k}} \cap \overline{V}).
$$
と置く。$W_{\overline{k}} \cap \overline{U}$ と
$W_{\overline{k}} \cap \overline{V}$ は $\text{Gal}(\overline{k}/k)$ の開部分群で
固定されるので、上の和と共通部分は有限である。従って $U_W$ と $V_W$ はともに
開かつ閉である。また構成により $W_{\bar k} = U_W \amalg V_W$ である。

\medskip\noindent
$W \subset W' \subset X$ が準コンパクト開集合ならば、
$W_{\overline{k}} \cap U_{W'} = U_W$ および
$W_{\overline{k}} \cap V_{W'} = V_W$ であると主張する。検証は省略する。
従って $U = \bigcup_{W \subset X} U_W$ および
$V = \bigcup_{W \subset X} V_W$ と定めると、
$X_{\overline{k}} = U \amalg V$ は開かつ閉な部分集合の不交和となる。
$V$ は（局所的な）和を取って構成されるので、空でないことは明らかである。
他方、$U$ は構成により $W_0 \cap \overline{U}$ を含むので空でない。
最後に、$U, V \subset X_{\bar k}$ は構成により閉かつ $H$-不変である。
従って Lemma \ref{varieties-lemma-closed-fixed-by-Galois} により、
$U = (U')_{\bar k}$ および $V = (V')_{\bar k}$ が、ある閉部分集合
$U', V' \subset X_{k'}$ に対して成り立つ。明らかに $X_{k'} = U' \amalg V'$ であり、所望のとおり
$X_{k'}$ は非連結である。
\end{proof}

\begin{lemma}
\label{varieties-lemma-tricky}
$k$ を体とする。$\overline{k}/k$ を（無限でもよい）Galois 拡大とする。
$f : T \to X$ を $k$ 上のスキームの射とする。$T_{\overline{k}}$ が連結で、
$X_{\overline{k}}$ が非連結であると仮定する。このとき $X$ は非連結である。
\end{lemma}

\begin{proof}
$X_{\overline{k}} = \overline{U} \amalg \overline{V}$ と書く。ここで
$\overline{U}$ と $\overline{V}$ は開かつ閉である。
$\overline{f} : T_{\overline{k}} \to X_{\overline{k}}$ を $f$ の基底変換とする。
$T_{\overline{k}}$ は連結なので、$T_{\overline{k}}$ は
$\overline{f}^{-1}(\overline{U})$ または $\overline{f}^{-1}(\overline{V})$ のいずれかに含まれる。
$T_{\overline{k}} \subset \overline{f}^{-1}(\overline{U})$ とする。

\medskip\noindent
準コンパクト開集合 $W \subset X$ を固定する。有限 Galois 部分拡大
$\overline{k}/k'/k$ で、$\overline{U} \cap W_{\overline{k}}$ と
$\overline{V} \cap W_{\overline{k}}$ が準コンパクト開集合
$U', V' \subset W_{k'}$ から来るものが存在する。このとき
$W_{k'} = U' \amalg V'$ でもある。次を考える。
$$
U'' = \bigcap\nolimits_{\sigma \in \text{Gal}(k'/k)} \sigma(U'),
\quad
V'' = \bigcup\nolimits_{\sigma \in \text{Gal}(k'/k)} \sigma(V').
$$
これらは Galois 不変で開かつ閉であり、$W_{k'} = U'' \amalg V''$ である。
Lemma \ref{varieties-lemma-closed-fixed-by-Galois} により、開かつ閉な部分集合
$U_W, V_W \subset W$ で、$U'' = (U_W)_{k'}$、$V'' = (V_W)_{k'}$ および
$W = U_W \amalg V_W$ となるものを得る。

\medskip\noindent
$W \subset W' \subset X$ が準コンパクト開集合ならば、
$W \cap U_{W'} = U_W$ および $W \cap V_{W'} = V_W$ であると主張する。
検証は省略する。従って $U = \bigcup_{W \subset X} U_W$ および
$V = \bigcup_{W \subset X} V_W$ と定めると、$X = U \amalg V$ を得る。
$V$ は（局所的な）和を取って構成されるので、空でないことは明らかである。
他方、$U$ は構成により $f(T)$ を含むので空でない。
\end{proof}

\begin{lemma}
\label{varieties-lemma-geometrically-connected-criterion}
\begin{reference}
\cite[IV Corollary 4.5.13.1(i)]{EGA}
\end{reference}
$k$ を体とする。$T \to X$ を $k$ 上のスキームの射とする。
$T$ が幾何学的連結であり、$X$ が連結であると仮定する。このとき
$X$ は幾何学的連結である。
\end{lemma}

\begin{proof}
これは Lemma \ref{varieties-lemma-tricky} の言い換えである。
\end{proof}



\begin{lemma}
\label{varieties-lemma-geometrically-connected-if-connected-and-point}
$k$ を体とする。$X$ を $k$ 上のスキームとする。$X$ は連結であり、
点 $x$ をもち、その点について $k$ が $\kappa(x)$ において代数的に閉であると仮定する。
このとき $X$ は幾何学的連結である。特に、$X$ が $k$-有理点をもち、$X$ が連結ならば、
$X$ は幾何学的連結である。
\end{lemma}

\begin{proof}
$T = \Spec(\kappa(x))$ と置く。$\overline{k}$ を $k$ の分離代数閉包とする。
$\kappa(x)/k$ に関する仮定から $T_{\overline{k}}$ は既約である。Algebra, Lemma
\ref{algebra-lemma-field-extension-geometrically-irreducible} を参照せよ。
従って Lemma \ref{varieties-lemma-geometrically-connected-criterion} により、
$X_{\overline{k}}$ は連結である。Lemma
\ref{varieties-lemma-characterize-geometrically-connected} により、$X$ は幾何学的連結であると結論する。
\end{proof}

\begin{lemma}
\label{varieties-lemma-inverse-image-connected-component}
$K/k$ を体拡大とする。$X$ を $k$ 上のスキームとする。
$T$ を $X$ の任意の連結成分とする。このとき逆像 $T_K \subset X_K$ は
$X_K$ の連結成分の和である。
\end{lemma}

\begin{proof}
これは純粋に位相的な主張である。$p : X_K \to X$ を射影とする。
$T \subset X$ を $X$ の連結成分とする。$t \in T_K = p^{-1}(T)$ とし、
$C \subset X_K$ を $t$ を含む連結成分とする。このとき $p(C)$ は $X$ の連結部分集合で
$T$ と交わるので、$p(C) \subset T$ である。従って $C \subset T_K$ である。
\end{proof}

\noindent
次の補題は、後のより強い Lemma \ref{varieties-lemma-image-connected-component} により置き換えられる。

\begin{lemma}
\label{varieties-lemma-image-connected-component-finite-extension}
$K/k$ を有限体拡大とし、$X$ を $k$ 上のスキームとする。
$p : X_K \to X$ を射影とする。$T$ を $X_K$ の任意の連結成分とすると、像
$p(T)$ は $X$ の連結成分である。
\end{lemma}

\begin{proof}
像 $p(T)$ は、ある $X'$（これは $X$ の連結成分である）に含まれる。$X'$ を任意の仕方で $X$ の
閉部分スキームとみなす。このとき $T$ は $X'_K = p^{-1}(X')$ の連結成分でもあるから、
$X$ が連結であると仮定してよい。射 $p$ は開であり
（Morphisms, Lemma \ref{morphisms-lemma-scheme-over-field-universally-open}）、
閉であり（Morphisms, Lemma \ref{morphisms-lemma-integral-universally-closed}）、
$p$ のファイバーは有限集合である
（Morphisms, Lemma \ref{morphisms-lemma-finite-quasi-finite}）。従って
Topology, Lemma \ref{topology-lemma-finite-fibre-connected-components} を適用して結論を得る。
\end{proof}

\begin{lemma}[Gabber]
\label{varieties-lemma-image-connected-component}
\begin{reference}
Email from Ofer Gabber dated June 4, 2016
\end{reference}
$K/k$ を体拡大とする。$X$ を $k$ 上のスキームとする。
$p : X_K \to X$ を射影とする。$\overline{T} \subset X_K$ を連結成分とする。
このとき $p(\overline{T})$ は $X$ の連結成分である。
\end{lemma}

\begin{proof}
$K/k$ が有限ならば、これは Lemma
\ref{varieties-lemma-image-connected-component-finite-extension} である。一般には証明はより難しい。

\medskip\noindent
$T \subset X$ を、$X$ の連結成分であって $\overline{T}$ の像を含むものとする。
$X$ を $T$ で置き換えてよい（誘導された被約部分スキーム構造を入れる）。
従って $X$ が連結であると仮定してよい。$A = H^0(X, \mathcal{O}_X)$ と置く。
$L \subset A$ を極大弱 \'etale $k$-部分代数とする。More on Algebra, Lemma
\ref{more-algebra-lemma-max-weakly-etale-subalgebra} を参照せよ。
$A$ は非自明な冪等元をもたないので、More on Algebra, Lemma
\ref{more-algebra-lemma-class-weakly-etale-over-field} により、$L$ は体であり、
$k$ の分離代数拡大である。$L$ はまた極大弱 \'etale $L$-部分代数であり、それは $A$ に含まれることを
注意しておく（実際、More on Algebra, Lemma
\ref{more-algebra-lemma-composition-weakly-etale} により、任意の弱 \'etale $L$-代数は
$k$ 上弱 \'etale である）。Schemes, Lemma
\ref{schemes-lemma-morphism-into-affine} により、構造射の分解
$X \to \Spec(L) \to \Spec(k)$ を得る。

\medskip\noindent
$L'/L$ を有限分離拡大とする。Cohomology of Schemes, Lemma
\ref{coherent-lemma-finite-locally-free-base-change-cohomology} により
$$
A \otimes_L L' =
H^0(X \times_{\Spec(L)} \Spec(L'), \mathcal{O}_{X \times_{\Spec(L)} \Spec(L')})
$$
である。極大弱 \'etale $L'$-部分代数を $A \otimes_L L'$ の中で考えると、More on Algebra, Lemma
\ref{more-algebra-lemma-change-fields-max-weakly-etale-subalgebra} により
$L \otimes_L L' = L'$ である。特に $A \otimes_L L'$ は非自明な冪等元をもたない
（そのような冪等元は弱 \'etale 部分代数を生成する）。従って
$X \times_{\Spec(L)} \Spec(L')$ は連結である。Lemma
\ref{varieties-lemma-characterize-geometrically-disconnected} により、$X$ は $L$ 上
幾何学的連結であると結論する。

\medskip\noindent
$\overline{T}$ に誘導される被約スキーム構造を入れ、合成
$$
\overline{T} \xrightarrow{i} X_K = X \times_{\Spec(k)} \Spec(K)
\xrightarrow{\pi}
\Spec(L \otimes_k K)
$$
を考える。像は $\Spec(L \otimes_k K)$ のある連結成分に含まれる。
$K \to L \otimes_k K$ は整なので、$\Spec(L \otimes_k K)$ の連結成分は点であり、
すべての点は閉である。Algebra, Lemma \ref{algebra-lemma-integral-over-field} を参照せよ。
従って商体 $L \otimes_k K \to E$ で、$\overline{T}$ が
$\Spec(E) \subset \Spec(L \otimes_k K)$ に写るものを得る。
従って $i(\overline{T}) \subset \pi^{-1}(\Spec(E))$ である。しかし
$$
\pi^{-1}(\Spec(E)) =
(X \times_{\Spec(k)} \Spec(K)) \times_{\Spec(L \otimes_k K)} \Spec(E) =
X \times_{\Spec(L)} \Spec(E)
$$
であり、$X$ は $L$ 上幾何学的連結なので、これは連結である。
従って集合論的に
$\overline{T} = X \times_{\Spec(L)} \Spec(E)$ という等式を得る。
ゆえに所望のとおり $\overline{T} \to X$ は全射である。
\end{proof}

\noindent
$X$ をスキームとする。$\pi_0(X)$ で $X$ の連結成分の集合を表す。

\begin{lemma}
\label{varieties-lemma-galois-action-connected-components}
$k$ を体とし、その分離代数閉包を $\overline{k}$ とする。$X$ を $k$ 上のスキームとする。
次の作用が存在する。
$$
\text{Gal}(\overline{k}/k)^{opp} \times \pi_0(X_{\overline{k}})
\longrightarrow
\pi_0(X_{\overline{k}})
$$
この作用は次の性質をもつ。
\begin{enumerate}
\item 元 $\overline{T} \in \pi_0(X_{\overline{k}})$ が作用により固定されるための
必要十分条件は、連結成分 $T \subset X$ であって $k$ 上幾何学的連結であり、
$T_{\overline{k}} = \overline{T}$ となるものが存在することである。
\item 体拡大 $k'/k$ とその分離代数閉包 $\overline{k}'$ に対して、図式
$$
\xymatrix{
\text{Gal}(\overline{k}'/k') \times \pi_0(X_{\overline{k}'})
\ar[r] \ar[d] &
\pi_0(X_{\overline{k}'}) \ar[d] \\
\text{Gal}(\overline{k}/k) \times \pi_0(X_{\overline{k}})
\ar[r] &
\pi_0(X_{\overline{k}})
}
$$
は可換である（右の縦矢印は Lemma
\ref{varieties-lemma-separably-closed-field-connected-components} により全単射である）。
\end{enumerate}
\end{lemma}

\begin{proof}
$\text{Gal}(\overline{k}/k)$ の $X_{\overline{k}}$ への作用
(\ref{varieties-equation-galois-action-base-change-kbar}) は、その連結成分への作用を誘導する。
連結成分は常に閉である（Topology, Lemma
\ref{topology-lemma-connected-components}）。従って $\overline{T}$ が (1) のとおりならば、
Lemma \ref{varieties-lemma-closed-fixed-by-Galois} により、閉部分集合 $T \subset X$ で
$\overline{T} = T_{\overline{k}}$ となるものが存在する。$T$ は $k$ 上幾何学的連結である。
Lemma \ref{varieties-lemma-characterize-geometrically-connected} を参照せよ。
$T$ が $X$ の連結成分であることを見るため、$T \subset T'$、
$T \not = T'$ で、$T'$ が $X$ の連結成分であると仮定する。この場合、
$T'_{k'}$ は $\overline{T}$ を真に含み、従って非連結である。Lemma
\ref{varieties-lemma-tricky} により、これは $T'$ が非連結であることを意味する。矛盾である。

\medskip\noindent
(2) における関手性の証明は省略する。
\end{proof}

\begin{lemma}
\label{varieties-lemma-galois-action-connected-components-continuous}
$k$ を体とし、その分離代数閉包を $\overline{k}$ とする。$X$ を $k$ 上のスキームとする。
次を仮定する。
\begin{enumerate}
\item $X$ は準コンパクトである。
\item $X_{\overline{k}}$ の連結成分は開である。
\end{enumerate}
このとき次が成り立つ。
\begin{enumerate}
\item[(a)] $\pi_0(X_{\overline{k}})$ は有限である。
\item[(b)] $\text{Gal}(\overline{k}/k)$ の $\pi_0(X_{\overline{k}})$ への作用は連続である。
\end{enumerate}
さらに、$X$ が $k$ 上有限型ならば仮定 (1)、(2) は満たされる。
\end{lemma}

\begin{proof}
連結成分は開であり $X_{\overline{k}}$ を被覆し（Topology, Lemma
\ref{topology-lemma-connected-components}）、かつ $X_{\overline{k}}$ は準コンパクトなので、
その個数は有限であると結論できる。従って (a) が成り立つ。
Lemma \ref{varieties-lemma-descend-open} により、これらの連結成分はそれぞれ
$\overline{k}/k$ の有限部分拡大上で定義され、(b) を得る。
$X$ が $k$ 上有限型ならば、$X_{\overline{k}}$ は $\overline{k}$ 上有限型である
（Morphisms, Lemma \ref{morphisms-lemma-base-change-finite-type}）。
従って $X_{\overline{k}}$ は Noether スキームである
（Morphisms, Lemma \ref{morphisms-lemma-finite-type-noetherian}）。
ゆえに $X_{\overline{k}}$ は有限個の既約成分をもち
（Properties, Lemma \ref{properties-lemma-Noetherian-irreducible-components}）、
なおさら有限個の連結成分をもつ（従ってそれらは開である）。
\end{proof}







\section{幾何学的既約スキーム}
\label{varieties-section-geometrically-irreducible}

\noindent
体上の既約スキーム $X$ は、基礎体を拡大すると $X$ が可約になることがある。
幾何学的既約スキームではこれは起こらない。

\begin{definition}
\label{varieties-definition-geometrically-irreducible}
$X$ を体 $k$ 上のスキームとする。スキーム $X$ が $k$ 上{\it 幾何学的既約}であるとは、
$X_{k'}$ が既約\footnote{既約空間は空でない。}であることが、任意の体拡大 $k'$ を
$k$ の拡大としてとるたびに成り立つことをいう。
\end{definition}

\begin{lemma}
\label{varieties-lemma-geometrically-irreducible-check-after-extension}
$X$ を体 $k$ 上のスキームとし、$k'/k$ を体拡大とする。
$X$ が $k$ 上幾何学的既約であるための必要十分条件は、$X_{k'}$ が $k'$ 上
幾何学的既約であることである。
\end{lemma}

\begin{proof}
$X$ が $k$ 上幾何学的既約ならば、$X_{k'}$ が $k'$ 上幾何学的既約であることは
明らかである。逆に、任意の体拡大 $k''/k$ に対し、共通の体拡大
$k'''/k'$ および $k'''/k''$ が存在する（Fields, Lemma
\ref{fields-lemma-common-extension-field}）。射 $X_{k'''} \to X_{k''}$ は全射である
（体のスペクトル間の全射の基底変換である）から、$X_{k'''}$ の既約性は
$X_{k''}$ の既約性を導く。従って $X_{k'}$ が $k'$ 上幾何学的既約ならば、
$X$ は $k$ 上幾何学的既約である。
\end{proof}

\begin{lemma}
\label{varieties-lemma-separably-closed-irreducible}
$X$ を分離閉体 $k$ 上のスキームとする。$X$ が既約ならば、$X_K$ は任意の体拡大
$K/k$ に対して既約である。すなわち、$X$ は $k$ 上幾何学的既約である。
\end{lemma}

\begin{proof}
Properties, Lemma \ref{properties-lemma-characterize-irreducible} および
Algebra, Lemma \ref{algebra-lemma-separably-closed-irreducible} を用いればよい。
\end{proof}

\begin{lemma}
\label{varieties-lemma-bijection-irreducible-components}
$k$ を体とし、$X$, $Y$ を $k$ 上のスキームとする。$X$ は $k$ 上幾何学的既約で
あると仮定する。このとき射影
$$
p : X \times_k Y \longrightarrow Y
$$
は既約成分の間の全単射を誘導する。
\end{lemma}

\begin{proof}
まず、$p$ のスキーム論的ファイバーは幾何学的既約スキーム $X$ の体拡大による
基底変換なので既約である。さらに、スキーム論的ファイバーは集合論的ファイバーと
同相である（Schemes, Lemma \ref{schemes-lemma-fibre-topological}）。Morphisms, Lemma
\ref{morphisms-lemma-scheme-over-field-universally-open} により写像 $p$ は開である。
従って Topology, Lemma
\ref{topology-lemma-irreducible-fibres-irreducible-components} を適用すれば結論を得る。
\end{proof}

\begin{lemma}
\label{varieties-lemma-geometrically-irreducible-local}
\begin{slogan}
連結性を除けば、幾何学的既約性は Zariski 局所的である。
\end{slogan}
$k$ を体とし、$X$ を $k$ 上のスキームとする。次は同値である。
\begin{enumerate}
\item $X$ は $k$ 上幾何学的既約である。
\item 任意の空でないアフィン開集合 $U$ に対して、$k$-代数 $\mathcal{O}_X(U)$ は
$k$ 上幾何学的既約である（Algebra, Definition
\ref{algebra-definition-geometrically-irreducible} を参照）。
\item $X$ は既約であり、アフィン開被覆 $X = \bigcup U_i$ が存在して、各 $k$-代数
$\mathcal{O}_X(U_i)$ は幾何学的既約である。
\item 開被覆 $X = \bigcup_{i \in I} X_i$ が存在し、$I \not = \emptyset$ であり、
$X_i$ は各 $i$ に対して幾何学的既約で、$X_i \cap X_j \not = \emptyset$ が
すべての $i, j \in I$ に対して成り立つ。
\end{enumerate}
さらに、$X$ が幾何学的既約ならば、$X$ の空でない任意の開部分スキームも
幾何学的既約である。
\end{lemma}

\begin{proof}
$\Spec(A)$ が $k$ 上の幾何学的既約なアフィンスキームであるための必要十分条件は、
$A$ が $k$ 上幾何学的既約であることであり、これは定義から直ちに従う。
スキームが既約ならば、その空でない任意の開部分スキーム $X$ も既約であり、
空でない二つの開部分集合は空でない交わりをもつことを思い出そう。また、すべての
アフィン開集合が既約ならば、そのスキームは既約である（Properties, Lemma
\ref{properties-lemma-characterize-irreducible}）。従って補足の主張、および含意
(1) $\Rightarrow$ (2)、(2) $\Rightarrow$ (3)、(3) $\Rightarrow$ (4) は明らかである。
(4) が成り立つならば、任意の体拡大 $k'/k$ に対しスキーム $X_{k'}$ は、二つずつ
交わる既約開集合による被覆をもつ。従って $X_{k'}$ は既約である。ゆえに (4) は
(1) を導く。
\end{proof}

\begin{lemma}
\label{varieties-lemma-geometrically-irreducible-function-field}
$X$ を体 $k$ 上の既約スキームとし、$\xi \in X$ をその生成点とする。次は同値である。
\begin{enumerate}
\item $X$ は $k$ 上幾何学的既約である。
\item $\kappa(\xi)$ は $k$ 上幾何学的既約である。
\item $k$ の $\kappa(\xi)$ における分離代数閉包は $k$ に等しい。
\end{enumerate}
\end{lemma}

\begin{proof}
Algebra, Lemma
\ref{algebra-lemma-geometrically-irreducible-separable-elements} により、(2) と (3) は
同値である。

\medskip\noindent
(1) を仮定する。$\mathcal{O}_{X, \xi}$ は $\mathcal{O}_X(U)$ のフィルター余極限であり、
ここで $U$ は $X$ の空でないアフィン開部分スキームを走ることを思い出そう。
Lemma \ref{varieties-lemma-geometrically-irreducible-local} と Algebra, Lemma
\ref{algebra-lemma-subalgebra-geometrically-irreducible} を合わせると、
$\mathcal{O}_{X, \xi}$ は $k$ 上幾何学的既約である。写像
$\mathcal{O}_{X, \xi} \to \kappa(\xi)$ は局所冪零な核をもつ全射であるから
（Algebra, Lemma \ref{algebra-lemma-minimal-prime-reduced-ring}）、
$\kappa(\xi)$ は幾何学的既約である（Algebra, Lemma
\ref{algebra-lemma-p-ring-map}）。

\medskip\noindent
(2) を仮定する。$X$ は被約であるとしてよい。$U \subset X$ を空でないアフィン開集合とする。
このとき $U = \Spec(A)$ であり、$A$ は分数体 $\kappa(\xi)$ をもつ整域である。
従って $A$ は幾何学的既約な $k$-代数の $k$-部分代数である。Algebra, Lemma
\ref{algebra-lemma-subalgebra-geometrically-irreducible} により、$A$ は $k$ 上
幾何学的既約である。Lemma \ref{varieties-lemma-geometrically-irreducible-local} により、
$X$ は $k$ 上幾何学的既約である。
\end{proof}

\begin{lemma}
\label{varieties-lemma-geometrically-irreducible-self-product}
スキーム $X$ を体 $k$ 上のものとする。これが $k$ 上幾何学的既約であるための必要十分条件は、
$X\times_k X$ が既約であることである。
\end{lemma}

\begin{proof}
$X$ が $k$ 上幾何学的既約ならば、Lemma
\ref{varieties-lemma-bijection-irreducible-components} により $X\times_k X$ は既約である。

\medskip\noindent
逆に、$X \times_k X$ が既約であると仮定する。$X$ をその台となる被約スキームで
置き換えても問題は変わらないので、$X$ は被約であるとしてよい。
第一射影 $X \times_k X \to X$ は全射であるから、$X$ も既約であり、従って整である。
元 $\alpha$ が函数体 $k(X)$ に存在し、$k$ 上分離代数的であるが $k$ には属さないと
仮定する。$E = k(\alpha)$ とおく。

\medskip\noindent
このとき、空でないアフィン開部分スキーム $U$ を $X$ の中にとり、
$k \subset E \subset O(U)$ を満たすものが存在し、従って
$E \otimes_k E \subset
O(U)\otimes_k O(U) = O(U\times_k U)$ である。ゆえに支配的射
$U \times_k U \to \Spec(E\otimes_k E)$ を得る。$X\times_k X$ は既約なので、
$U \times_k U$ も既約であり、従って $\Spec(E \otimes_k E)$ は連結である。
$E$ が $k$ より真に大きければ、$E \otimes_k E$ は $E$ を一つの因子として含む
体の積である。これは次元 $(\dim_k(E))^2>\dim_k(E)$ をもつ $k$-ベクトル空間であり、
従って $E\otimes_k E$ は少なくとも二つの体の積である。
これは $\Spec(E \otimes_k E)$ の連結性に矛盾する。ゆえに実際 $E$ は $k$ に等しい。
Lemma \ref{varieties-lemma-geometrically-irreducible-function-field} により、$X$ は $k$ 上
幾何学的既約である。
\end{proof}

\begin{lemma}
\label{varieties-lemma-separably-closed-field-irreducible-components}
$k'/k$ を体拡大とし、$X$ を $k$ 上のスキームとする。$X' = X_{k'}$ とおく。
$k$ は分離代数閉であると仮定する。このとき射 $X' \to X$ は既約成分の全単射を誘導する。
\end{lemma}

\begin{proof}
$k$ は分離代数閉であるから、Algebra, Lemma
\ref{algebra-lemma-separably-closed-irreducible-implies-geometric} により $k'$ は $k$ 上
幾何学的既約である。従って $Z = \Spec(k')$ は $k$ 上幾何学的既約であることが、上の
Lemma \ref{varieties-lemma-geometrically-irreducible-local} から従う。
$X' = Z \times_k X$ なので、これは Lemma
\ref{varieties-lemma-bijection-irreducible-components} の特別な場合である。
\end{proof}

\begin{lemma}
\label{varieties-lemma-characterize-geometrically-irreducible}
\begin{slogan}
幾何学的既約性は、基礎体の分離代数閉包上で判定できる。
\end{slogan}
$k$ を体とし、$X$ を $k$ 上のスキームとする。次は同値である。
\begin{enumerate}
\item $X$ は $k$ 上幾何学的既約である。
\item 任意の有限分離体拡大 $k'/k$ に対してスキーム $X_{k'}$ は既約である。
\item $X_{\overline{k}}$ は既約である。ここで $k \subset \overline{k}$ は $k$ の
分離代数閉包である。
\end{enumerate}
\end{lemma}

\begin{proof}
$X_{\overline{k}}$ は既約、すなわち (3) を仮定する。$k'/k$ を体拡大とする。
体拡大 $\overline{k}'/\overline{k}$ が存在して、$k'$ は $\overline{k}'$ に $k$ の拡大として
埋め込まれる。
Lemma \ref{varieties-lemma-separably-closed-field-irreducible-components} により
$X_{\overline{k}'}$ は既約である。$X_{\overline{k}'} \to X_{k'}$ は全射なので、
$X_{k'}$ は既約であると結論できる。従って (1) が成り立つ。

\medskip\noindent
$k \subset \overline{k}$ を $k$ の分離代数閉包とする。(3) が成り立たない、すなわち
$X_{\overline{k}}$ が可約であると仮定する。ある有限部分拡大について $X_{k'}$ も可約で
あることを示したい。その部分拡大を $\overline{k}/k'/k$ と記す。
アフィン開被覆 $X = \bigcup_{i \in I} U_i$ をとり、各 $U_i$ は空でないとする。
ある $i$ に対してスキーム $U_i$ が可約であるか、またはある組 $i \not = j$ に対して
交わり $U_i \cap U_j$ が空ならば、$X$ は可約であり（Properties, Lemma
\ref{properties-lemma-characterize-irreducible}）、証明は終わる。
特に、$U_{i, \overline{k}} \cap U_{j, \overline{k}}$ はすべての $i, j \in I$ について
空でないと仮定してよく、$U_{i, \overline{k}}$ はある $i$ について可約でなければならない。
Algebra, Lemma
\ref{algebra-lemma-geometrically-irreducible} によれば、これはある有限分離体拡大
$U_{i, k'}$ が可約となる $k'/k$ が存在することを意味する。従って $X_{k'}$ も可約である。
ゆえに (2) は (3) を導く。

\medskip\noindent
含意 (1) $\Rightarrow$ (2) は直ちに従う。これで補題が証明された。
\end{proof}

\begin{lemma}
\label{varieties-lemma-inverse-image-irreducible}
$K/k$ を体拡大とし、$X$ を $k$ 上のスキームとする。$T$ を $X$ の任意の既約成分とすると、
逆像 $T_K \subset X_K$ は $X_K$ の既約成分の合併である。
\end{lemma}

\begin{proof}
$T \subset X$ を $X$ の既約成分とする。射 $T_K \to T$ は平坦なので、一般化は
$T_K \to T$ に沿って持ち上がる。従って、任意の $\xi \in T_K$ で $T_K$ の既約成分の
生成点であるものは、ある生成点 $\eta$ に写り、それは $T$ の生成点である。特殊化 $\xi' \leadsto \xi$ が
$X_K$ にあるならば、$\xi'$ は $\eta$ に写る。実際、$\eta$ に特殊化する点は $X$ に
存在しない。従って $\xi' \in T_K$ であり、$\xi = \xi'$ と結論できる。
言い換えれば、$\xi$ は $X_K$ のある既約成分の生成点である。これは $T_K$ の既約成分が
すべて $X_K$ の既約成分であることを意味する。
\end{proof}

\noindent
スキーム $X$ に対し、$\text{IrredComp}(X)$ で $X$ の既約成分の集合を表す。

\begin{lemma}
\label{varieties-lemma-image-irreducible}
$K/k$ を体拡大とし、$X$ を $k$ 上のスキームとする。任意の既約成分
$\overline{T} \subset X_K$ に対し、$\overline{T}$ の $X$ における像は $X$ の既約成分である。
これにより標準的な全射
$$
\text{IrredComp}(X_K)
\longrightarrow
\text{IrredComp}(X)
$$
が定まる。
\end{lemma}

\begin{proof}
図式
$$
\xymatrix{
X_K \ar[d] & X_{\overline{K}} \ar[d] \ar[l] \\
X & X_{\overline{k}} \ar[l]
}
$$
を考える。ここで $\overline{K}$ は $K$ の分離代数閉包、$\overline{k}$ は $k$ の
分離代数閉包である。Lemma \ref{varieties-lemma-separably-closed-field-irreducible-components} により、
射 $X_{\overline{K}} \to X_{\overline{k}}$ は既約成分の間の全単射を誘導する。従って、
射 $X_{\overline{k}} \to X$ および $X_{\overline{K}} \to X_K$ について補題を示せば十分である。
言い換えれば、$K = \overline{k}$ と仮定してよい。

\medskip\noindent
射 $p : X_{\overline{k}} \to X$ は整、平坦かつ全射である。平坦性により一般化は $p$ に沿って
持ち上がる（Morphisms, Lemma \ref{morphisms-lemma-generalizations-lift-flat}）。従って
$X_{\overline{k}}$ の既約成分の生成点は、$X$ の既約成分の生成点に写る。整性により $p$ は
普遍閉である（Morphisms, Lemma \ref{morphisms-lemma-integral-universally-closed}）。従って、
既約成分の像 $p(\overline{T})$ は、$X$ のある既約成分の生成点を含む閉既約部分集合であり、
ゆえに $p(\overline{T})$ は $X$ の既約成分である。これで第一の主張が証明された。
$T \subset X$ が既約成分ならば、$p^{-1}(T) =T_K$ は空でない既約成分の合併である
（Lemma \ref{varieties-lemma-inverse-image-irreducible}）。第一の部分により、その各成分は必ず $T$ 上へ
写る。従って写像は全射である。
\end{proof}

\begin{lemma}
\label{varieties-lemma-irreducible-dense-rational-points-geometrically-irreducible}
$k$ を体とし、$X$ を $k$ 上のスキームとする。$X$ が既約で、$k$-有理点の稠密な集合を
もつならば、$X$ は幾何学的既約である。
\end{lemma}

\begin{proof}
$k'/k$ を有限体拡大とし、$Z, Z' \subset X_{k'}$ を既約成分とする。Lemma
\ref{varieties-lemma-characterize-geometrically-irreducible} により、$Z = Z'$ を示せば十分である。
Lemma \ref{varieties-lemma-image-irreducible} により $p(Z) = p(Z') = X$ であり、ここで
$p : X_{k'} \to X$ は射影である。$Z \not = Z'$ ならば、$Z \cap Z'$ は $X_{k'}$ の中で
疎であり、従って $p(Z \cap Z')$ は稠密でない（Morphisms, Lemma
\ref{morphisms-lemma-image-nowhere-dense-finite}）。ここでは、$p$ が有限射
$\Spec(k') \to \Spec(k)$ の基底変換として有限射であることも用いている
（Morphisms, Lemma \ref{morphisms-lemma-base-change-finite}）。従って、
$k$-有理点 $x \in X$ で $x \not \in p(Z \cap Z')$ を満たすものを選べる。
$x$ の剰余体は $k$ なので、$p^{-1}(\{x\}) = \{x'\}$ であり、ここで
$x' \in X_{k'}$ は剰余体 $k'$ をもつ点である。$x \in p(Z) = p(Z')$ であるから、
$x' \in Z \cap Z'$ と結論でき、これは求める矛盾である。
\end{proof}

\begin{lemma}
\label{varieties-lemma-galois-action-irreducible-components}
$k$ を体とし、その分離代数閉包を $\overline{k}$ とする。$X$ を $k$ 上のスキームとする。
次の作用が存在する。
$$
\text{Gal}(\overline{k}/k)^{opp} \times \text{IrredComp}(X_{\overline{k}})
\longrightarrow
\text{IrredComp}(X_{\overline{k}})
$$
これは次の性質をもつ。
\begin{enumerate}
\item 元 $\overline{T} \in \text{IrredComp}(X_{\overline{k}})$ が作用で固定されるための
必要十分条件は、既約成分 $T \subset X$ であって $k$ 上幾何学的既約であり、
$T_{\overline{k}} = \overline{T}$ を満たすものが存在することである。
\item 任意の体拡大 $k'/k$ とその分離代数閉包 $\overline{k}'$ に対して、図式
$$
\xymatrix{
\text{Gal}(\overline{k}'/k') \times \text{IrredComp}(X_{\overline{k}'})
\ar[r] \ar[d] &
\text{IrredComp}(X_{\overline{k}'}) \ar[d] \\
\text{Gal}(\overline{k}/k) \times \text{IrredComp}(X_{\overline{k}})
\ar[r] &
\text{IrredComp}(X_{\overline{k}})
}
$$
は可換である（右の垂直矢印は Lemma
\ref{varieties-lemma-separably-closed-field-irreducible-components} により全単射である）。
\end{enumerate}
\end{lemma}

\begin{proof}
$\text{Gal}(\overline{k}/k)$ の $X_{\overline{k}}$ 上の作用
(\ref{varieties-equation-galois-action-base-change-kbar}) は、その既約成分への作用を誘導する。
既約成分は常に閉である（Topology, Lemma
\ref{topology-lemma-connected-components}）。従って $\overline{T}$ が (1) のとおりならば、
Lemma \ref{varieties-lemma-closed-fixed-by-Galois} により、閉部分集合 $T \subset X$ で
$\overline{T} = T_{\overline{k}}$ を満たすものが存在する。$T$ は $k$ 上幾何学的既約であることに注意せよ
（Lemma \ref{varieties-lemma-characterize-geometrically-irreducible}）。$T$ が $X$ の既約成分であることを
見るため、$T \subset T'$、$T \not = T'$ で、$T'$ が $X$ の既約成分であると仮定する。
$\overline{\eta}$ を $\overline{T}$ の生成点とする。これはある生成点 $\eta$ に写り、その点は $T$ の生成点である。
このとき生成点 $\xi \in T'$ は $\eta$ に特殊化する。$X_{\overline{k}} \to X$ は平坦なので、
点 $\overline{\xi} \in X_{\overline{k}}$ が存在し、これは $\xi$ に写って $\overline{\eta}$ に特殊化する。
従って単集合 $\{\overline{\xi}\}$ の閉包は、$X_{\overline{\xi}}$ の既約閉部分集合であり、
$\overline{T}$ を真に含む。これは求める矛盾である。

\medskip\noindent
(2) における関手性の証明は省略する。
\end{proof}

\begin{lemma}
\label{varieties-lemma-orbit-irreducible-components}
$k$ を体とし、その分離代数閉包を $\overline{k}$ とする。$X$ を $k$ 上のスキームとする。
Lemma \ref{varieties-lemma-image-irreducible} の写像
$$
\text{IrredComp}(X_{\overline{k}})
\longrightarrow
\text{IrredComp}(X)
$$
のファイバーは、Lemma \ref{varieties-lemma-galois-action-irreducible-components} の作用による
$\text{Gal}(\overline{k}/k)$ の軌道にちょうど等しい。
\end{lemma}

\begin{proof}
$T \subset X$ を $X$ の既約成分とし、$\eta \in T$ をその生成点とする。Lemmas
\ref{varieties-lemma-inverse-image-irreducible} および \ref{varieties-lemma-image-irreducible} により、
$\overline{T}$ の既約成分の生成点のうち $T$ の中へ写るものは $\eta$ に写る。Algebra, Lemma
\ref{algebra-lemma-Galois-orbit} により、Galois 群は $X_{\overline{k}}$ の点のうち
$\eta$ に写るすべての点に推移的に作用する。従って補題が従う。
\end{proof}

\begin{lemma}
\label{varieties-lemma-galois-action-irreducible-components-locally-finite-type}
$k$ を体とする。$X \to \Spec(k)$ は局所有限型であると仮定する。この場合、
\begin{enumerate}
\item 作用
$$
\text{Gal}(\overline{k}/k)^{opp} \times \text{IrredComp}(X_{\overline{k}})
\longrightarrow
\text{IrredComp}(X_{\overline{k}})
$$
は、$\text{IrredComp}(X_{\overline{k}})$ に離散位相を入れると連続である。
\item $X_{\overline{k}}$ の各既約成分は $k$ の有限拡大上で定義できる。
\item 任意の既約成分 $T \subset X$ に対し、スキーム $T_{\overline{k}}$ は
$X_{\overline{k}}$ の有限個の既約成分の合併であり、それらはすべて同じ
$\text{Gal}(\overline{k}/k)$-軌道に属する。
\end{enumerate}
\end{lemma}

\begin{proof}
$\overline{T}$ を $X_{\overline{k}}$ の既約成分とする。
$U \subset X$ をアフィン開集合として、$\overline{T} \cap U_{\overline{k}}$ が空でないように選べる。
$U = \Spec(A)$ と書くと、$A$ は有限型 $k$-代数である（Morphisms, Lemma
\ref{morphisms-lemma-locally-finite-type-characterize}）。従って $A_{\overline{k}}$ は有限型
$\overline{k}$-代数であり、特に Noether である。$\mathfrak p = (f_1, \ldots, f_n)$ を
$\overline{T} \cap U_{\overline{k}}$ に対応する素イデアルとする。
$A_{\overline{k}} = A \otimes_k \overline{k}$ なので、有限部分拡大 $\overline{k}/k'/k$ が存在して、
各 $f_i \in A_{k'}$ となる。$\text{Gal}(\overline{k}/k')$ が $\overline{T}$ を固定することは
明らかであり、(1) が証明された。

\medskip\noindent
(2) は、$k'$ 上の状況に Lemma
\ref{varieties-lemma-galois-action-irreducible-components} (1) を適用すれば従う。実際、既約成分
$\overline{T}$ は $T'_{\overline{k}}$ の形であり、ここで $T' \subset X_{k'}$ はある既約成分である。

\medskip\noindent
(3) を示すため、$T \subset X$ を既約成分とする。既約成分
$\overline{T} \subset X_{\overline{k}}$ で $T$ に写るものを選ぶ（Lemma \ref{varieties-lemma-image-irreducible}）。上により、
$\overline{T}$ の軌道は有限であり、$\overline{T}_1, \ldots, \overline{T}_n$ と書ける。このとき
$\overline{T}_1 \cup \ldots \cup \overline{T}_n$ は $\text{Gal}(\overline{k}/k)$-不変な
$X_{\overline{k}}$ の閉部分集合であるから、Lemma
\ref{varieties-lemma-closed-fixed-by-Galois} により $W_{\overline{k}}$ の形であり、ここで
$W \subset X$ はある閉部分集合である。明らかに $W = T$ であり、証明は完了する。
\end{proof}

\begin{lemma}
\label{varieties-lemma-finite-extension-geometrically-irreducible-components}
$k$ を体とし、$X \to \Spec(k)$ を局所有限型とする。$X$ は有限個の既約成分をもつと仮定する。
このとき有限分離拡大 $k'/k$ が存在し、$X_{k'}$ の各既約成分は $k'$ 上
幾何学的既約である。
\end{lemma}

\begin{proof}
$\overline{k}$ を $k$ の分離代数閉包とする。$X$ が有限個の既約成分をもつという仮定と
Lemma \ref{varieties-lemma-galois-action-irreducible-components-locally-finite-type} (3) から、
$X_{\overline{k}}$ は有限個の既約成分 $\overline{T}_1, \ldots, \overline{T}_n$ をもつ。
Lemma \ref{varieties-lemma-galois-action-irreducible-components-locally-finite-type} (2) により、有限拡大
$\overline{k}/k'/k$ と既約成分 $T_i \subset X_{k'}$ が存在し、
$\overline{T}_i = T_{i, \overline{k}}$ となる。これで証明は終わる。
\end{proof}

\begin{lemma}
\label{varieties-lemma-irreducible-components-geometrically-irreducible}
$X$ を体 $k$ 上のスキームとする。$X$ は有限個の既約成分をもち、そのすべてが
幾何学的既約であると仮定する。このとき $X$ は有限個の連結成分をもち、その各々は
幾何学的連結である。
\end{lemma}

\begin{proof}
連結成分は既約成分の合併なので明らかである。詳細は省略する。
\end{proof}







\section{幾何学的整スキーム}
\label{varieties-section-geometrically-integral}

\noindent
体上の整スキーム $X$ は、基礎体を拡大すると $X$ が非被約または可約になることがある。
幾何学的整スキームではこれは起こらない。

\begin{definition}
\label{varieties-definition-geometrically-integral}
$X$ を体 $k$ 上のスキームとする。
\begin{enumerate}
\item $x \in X$ とする。$X$ が
{\it $x$ において幾何学的に点ごとに整}であるとは、任意の体拡大 $k'/k$ と、
任意の点 $x' \in X_{k'}$ で $x$ 上にあるものに対し局所環
$\mathcal{O}_{X_{k'}, x'}$ が整であることをいう。
\item $X$ が各点において幾何学的に点ごとに整であるとき、$X$ は
{\it 幾何学的に点ごとに整}であるという。
\item $X$ が $k$ 上{\it 幾何学的整}であるとは、スキーム $X_{k'}$ が、
$k'$ を $k$ の任意の体拡大とするとき常に整であることをいう。
\end{enumerate}
\end{definition}

\noindent
概念 (2) と (3) を区別する必要がある。例えば $k = \mathbf{R}$ かつ
$X = \Spec(\mathbf{C}[x])$ ならば、$X$ は $\mathbf{R}$ 上幾何学的に点ごとに整であるが、
もちろん幾何学的整ではない。

\begin{lemma}
\label{varieties-lemma-geometrically-integral}
$k$ を体とし、$X$ を $k$ 上のスキームとする。$X$ が $k$ 上幾何学的整であるための
必要十分条件は、$X$ が $k$ 上幾何学的被約かつ幾何学的既約であることである。
\end{lemma}

\begin{proof}
Properties, Lemma \ref{properties-lemma-characterize-integral} を参照せよ。
\end{proof}

\begin{lemma}
\label{varieties-lemma-proper-geometrically-reduced-global-sections}
$k$ を体とし、$X$ を $k$ 上の固有スキームとする。
\begin{enumerate}
\item $A = H^0(X, \mathcal{O}_X)$ は有限次元 $k$-代数である。
\item $A = \prod_{i = 1, \ldots, n} A_i$ は Artin 局所 $k$-代数の積であり、$X$ の
各連結成分に一つの因子が対応する。
\item $X$ が被約ならば、$A = \prod_{i = 1, \ldots, n} k_i$ は体の積であり、各体は
$k$ の有限拡大である。
\item $X$ が幾何学的被約ならば、$k_i$ は $k$ 上有限分離である。
\item $X$ が幾何学的連結ならば、$A$ は $k$ 上幾何学的既約である。
\item $X$ が幾何学的既約ならば、$A$ は $k$ 上幾何学的既約である。
\item $X$ が幾何学的被約かつ幾何学的連結ならば $A = k$ である。
\item $X$ が幾何学的整ならば $A = k$ である。
\end{enumerate}
\end{lemma}

\begin{proof}
Cohomology of Schemes, Lemma
\ref{coherent-lemma-proper-over-affine-cohomology-finite} により、
$A = H^0(X, \mathcal{O}_X)$ は有限次元 $k$-代数である。これで (1) が証明された。

\medskip\noindent
次に $A$ は局所 Artin $k$-代数の積である。Algebra, Lemma
\ref{algebra-lemma-finite-dimensional-algebra} および Proposition
\ref{algebra-proposition-dimension-zero-ring} を参照せよ。$X = Y \amalg Z$ で $Y$ と $Z$ が
$X$ の開部分ならば、冪等元 $e \in A$ を得る。そのためには $\mathcal{O}_X$ の切断であって、
$1$ が $Y$ 上、$0$ が $Z$ 上となるものをとればよい。逆に、$e \in A$ が冪等元ならば、
$X$ の対応する分解を得る。
最後に、$X$ の台位相空間は Noether なので、その連結成分は開である。従って $X$ の
連結成分は $1$ 対 $1$ に $A$ の原始冪等元と対応する。これで (2) が証明された。

\medskip\noindent
$X$ が被約ならば、$A$ は被約である。従って局所環 $A_i = k_i$ は被約であり、ゆえに体である
（例えば Algebra, Lemma \ref{algebra-lemma-minimal-prime-reduced-ring}）。これで (3) が証明された。

\medskip\noindent
$X$ が幾何学的被約ならば、
$A \otimes_k \overline{k} =
H^0(X_{\overline{k}}, \mathcal{O}_{X_{\overline{k}}})$
は被約である（等式は Cohomology of Schemes, Lemma
\ref{coherent-lemma-flat-base-change-cohomology} による）。これは
$k_i \otimes_k \overline{k}$ が体の積であることを意味し、従って例えば Algebra,
Lemmas \ref{algebra-lemma-characterize-separable-field-extensions} および
\ref{algebra-lemma-geometrically-reduced-finite-purely-inseparable-extension} により、
$k_i/k$ は分離拡大である。これで (4) が証明された。

\medskip\noindent
$X$ が幾何学的連結ならば、$A \otimes_k \overline{k} =
H^0(X_{\overline{k}}, \mathcal{O}_{X_{\overline{k}}})$ は (2) により零次元局所環である。
従ってそのスペクトルは一点からなり、特に既約である。ゆえに $A$ は幾何学的既約である。
これで (5) が証明された。もちろん (5) は (6) を導く。

\medskip\noindent
$X$ が幾何学的被約かつ幾何学的連結ならば、$A = k_1$ は体であり、拡大 $k_1/k$ は
有限分離かつ幾何学的既約である。しかしこのとき $k_1 \otimes_k \overline{k}$ は
$[k_1 : k]$ 個の $\overline{k}$ のコピーの積であるから、$k_1 = k$ と結論できる。
これで (7) が証明された。もちろん (7) は (8) を導く。
\end{proof}

\noindent
以下は Stein 因子分解の初歩的な版である。実際の Stein 因子分解は More on Morphisms,
Section \ref{more-morphisms-section-stein-factorization} で論じる。

\begin{lemma}
\label{varieties-lemma-baby-stein}
$X$ を体 $k$ 上の固有スキームとし、$A = H^0(X, \mathcal{O}_X)$ とおく。標準射
$X \to \Spec(A)$ のファイバーは幾何学的連結である。
\end{lemma}

\begin{proof}
$S = \Spec(A)$ とおく。標準射 $X \to S$ は、Schemes, Lemma
\ref{schemes-lemma-morphism-into-affine} によって
$\Gamma(S, \mathcal{O}_S) = A = \Gamma(X, \mathcal{O}_X)$ に対応する射である。
$k$-代数 $A$ は有限積 $A = \prod A_i$ であり、その各因子は $k$ 上有限な局所 Artin
$k$-代数である（Lemma \ref{varieties-lemma-proper-geometrically-reduced-global-sections}）。
点 $s_i \in S$ を、$A_i$ の極大イデアルに対応する点とする。代数閉包
$\overline{k}$ を $k$ 上で選び、$\overline{A} = A \otimes_k \overline{k}$ とおく。
埋め込み $\kappa(s_i) \to \overline{k}$ を $k$ 上で選ぶと、これは $\overline{k}$-代数写像
$$
\sigma_i : \overline{A} = A \otimes_k \overline{k} \to
\kappa(s_i) \otimes_k \overline{k} \to \overline{k}
$$
を定める。基底変換
$$
\xymatrix{
\overline{X} \ar[r] \ar[d] & X \ar[d] \\
\overline{S} \ar[r] & S
}
$$
を考える。これは $X$ の $\overline{S} = \Spec(\overline{A})$ への基底変換である。
Cohomology of Schemes, Lemma
\ref{coherent-lemma-flat-base-change-cohomology} により
$\Gamma(\overline{X}, \mathcal{O}_{\overline{X}}) = \overline{A}$ である。
$\overline{s}_i \in \Spec(\overline{A})$ が $\overline{k}$-有理点で $\sigma_i$ に対応するならば、
$\overline{s}_i$ は $s_i \in S$ に写り、$\overline{X}_{\overline{s}_i}$ は
$X_{s_i}$ の $\Spec(\sigma_i)$ による基底変換である。従って $k$ が代数閉の場合に
補題を証明すれば十分である。

\medskip\noindent
$k$ は代数閉であると仮定する。この場合 $\kappa(s_i)$ は代数閉であり、$X_{s_i}$ が
連結であることを示せばよい。積分解 $A = \prod A_i$ は直和分解
$\Spec(A) = \coprod \Spec(A_i)$ に対応する（Algebra, Lemma
\ref{algebra-lemma-spec-product}）。$X_i$ を $\Spec(A_i)$ の逆像と書く。
Lemma \ref{varieties-lemma-proper-geometrically-reduced-global-sections} の (2) により、
$A_i = \Gamma(X_i, \mathcal{O}_{X_i})$ である。$X_{s_i} \to X_i$ は台位相空間上の同型を
誘導する閉埋め込みであることに注意せよ（$\Spec(A_i)$ は一点集合だからである）。
従って $X_{s_i}$ が連結でなければ $X_i$ も連結でない。ゆえに $X_i$ は空で
$A_i = 0$ であるか、または $X_i$ は $U \amalg V$ と書け、ここで $U$ と $V$ は
開かつ空でない。このとき $A_i$ は非自明な冪等元をもつことになる。
$A_i$ は局所なので、これは矛盾であり、証明が完了する。
\end{proof}

\begin{lemma}
\label{varieties-lemma-geometrically-reduced-stein}
$k$ を体とし、$X$ を $k$ 上の固有かつ幾何学的被約なスキームとする。次は同値である。
\begin{enumerate}
\item $H^0(X, \mathcal{O}_X) = k$ である。
\item $X$ は幾何学的連結である。
\end{enumerate}
\end{lemma}

\begin{proof}
Lemma \ref{varieties-lemma-baby-stein} により (1) $\Rightarrow$ (2) である。
Lemma \ref{varieties-lemma-proper-geometrically-reduced-global-sections} により (2) $\Rightarrow$ (1) である。
\end{proof}




\section{幾何学的正規スキーム}
\label{varieties-section-geometrically-normal}

\noindent
Properties, Definition \ref{properties-definition-normal} において、正規スキームの概念を定義した。
この概念は非 Noether スキームに対しても定義される。従って「幾何学的正則」スキームを
論じた場合とは異なり、ここでは基礎体のすべての体拡大を考える。

\begin{definition}
\label{varieties-definition-geometrically-normal}
$X$ を体 $k$ 上のスキームとする。
\begin{enumerate}
\item $x \in X$ とする。$X$ が
{\it $x$ において幾何学的正規}であるとは、任意の体拡大 $k'/k$ と、任意の点
$x' \in X_{k'}$ で $x$ 上にあるものに対して局所環 $\mathcal{O}_{X_{k'}, x'}$ が
正規であることをいう。
\item $X$ が $k$ 上{\it 幾何学的正規}であるとは、$X$ がすべての
$x \in X$ において幾何学的正規であることをいう。
\end{enumerate}
\end{definition}

\begin{lemma}
\label{varieties-lemma-geometrically-normal-at-point}
$k$ を体とし、$X$ を $k$ 上のスキームとし、$x \in X$ とする。次は同値である。
\begin{enumerate}
\item $X$ は $x$ において幾何学的正規である。
\item 任意の有限純非分離体拡大 $k'$（$k$ の拡大）と、任意の点
$x' \in X_{k'}$ で $x$ 上にあるものに対して、局所環
$\mathcal{O}_{X_{k'}, x'}$ は正規である。
\item 環 $\mathcal{O}_{X, x}$ は $k$ 上幾何学的正規である（Algebra, Definition
\ref{algebra-definition-geometrically-normal} を参照）。
\end{enumerate}
\end{lemma}

\begin{proof}
(1) が (2) を導くことは明らかである。(2) を仮定する。$k'/k$ を有限純非分離体拡大とする
（例えば $k = k'$ でもよい）。環 $\mathcal{O}_{X, x} \otimes_k k'$ を考える。
Algebra, Lemma \ref{algebra-lemma-p-ring-map} により、そのスペクトルは
$\mathcal{O}_{X, x}$ のスペクトルと同じである。従ってこれも局所環である
（Algebra, Lemma \ref{algebra-lemma-characterize-local-ring}）。ゆえに点
$x' \in X_{k'}$ で $x$ 上にあるものは一意であり、
$\mathcal{O}_{X_{k'}, x'} \cong \mathcal{O}_{X, x} \otimes_k k'$ である。
仮定によりこれは正規環である。従って Algebra, Lemma
\ref{algebra-lemma-geometrically-normal} により (3) が従う。

\medskip\noindent
(3) を仮定する。$k'/k$ を体拡大とする。$\Spec(k') \to \Spec(k)$ は全射なので、
$X_{k'} \to X$ も全射である（Morphisms, Lemma
\ref{morphisms-lemma-base-change-surjective}）。任意の点
$x' \in X_{k'}$ で $x$ 上にあるものをとる。局所環 $\mathcal{O}_{X_{k'}, x'}$ は環
$\mathcal{O}_{X, x} \otimes_k k'$ の局所化である。従って仮定により正規であり、
(1) が証明された。
\end{proof}

\begin{lemma}
\label{varieties-lemma-geometrically-normal}
$k$ を体とし、$X$ を $k$ 上のスキームとする。次は同値である。
\begin{enumerate}
\item $X$ は幾何学的正規である。
\item $X_{k'}$ は任意の体拡大 $k'/k$ に対して正規スキームである。
\item $X_{k'}$ は任意の有限生成体拡大 $k'/k$ に対して正規スキームである。
\item $X_{k'}$ は任意の有限純非分離体拡大 $k'/k$ に対して正規スキームである。
\item 任意のアフィン開集合 $U \subset X$ に対し、環 $\mathcal{O}_X(U)$ は
幾何学的正規である（Algebra, Definition
\ref{algebra-definition-geometrically-normal} を参照）。
\item $X_{k^{perf}}$ は正規スキームである。
\end{enumerate}
\end{lemma}

\begin{proof}
(1) を仮定する。このとき任意の体拡大 $k'/k$ と任意の点 $x' \in X_{k'}$ に対し、
$X_{k'}$ の $x'$ における局所環は正規である。定義により、これは $X_{k'}$ が正規であることを
意味する。従って (2) が成り立つ。

\medskip\noindent
(2) が (3) を導き、(3) が (4) を導くことは明らかである。

\medskip\noindent
(4) を仮定し、$U \subset X$ をアフィン開部分スキームとする。このとき $U_{k'}$ は、任意の
有限純非分離拡大 $k'/k$（$k = k'$ を含む）に対して正規スキームである。これは、
$k' \otimes_k \mathcal{O}(U)$ がすべての有限純非分離拡大 $k'/k$ に対して正規環であることを
意味する。従って定義により
$\mathcal{O}(U)$ は幾何学的正規な $k$-代数である。ゆえに (4) は (5) を導く。

\medskip\noindent
(5) を仮定する。任意の体拡大 $k'/k$ に対して、基底変換 $X_{k'}$ は、環
$\mathcal{O}_X(U) \otimes_k k'$ のスペクトルを貼り合わせて得られる。ここで $U$ は $X$ の
アフィン開集合を走る
（Schemes, Section \ref{schemes-section-fibre-products} を参照）。従って $X_{k'}$ は正規である。
ゆえに (1) が成り立つ。

\medskip\noindent
(5) と (6) の同値性は、幾何学的正規代数の定義と、直前に証明した (3) と (4) の同値性から従う。
\end{proof}

\begin{lemma}
\label{varieties-lemma-geometrically-normal-upstairs}
$k$ を体とし、$X$ を $k$ 上のスキームとし、$k'/k$ を体拡大とする。
$x \in X$ を点とし、$x' \in X_{k'}$ を $x$ 上にある点とする。次は同値である。
\begin{enumerate}
\item $X$ は $x$ において幾何学的正規である。
\item $X_{k'}$ は $x'$ において幾何学的正規である。
\end{enumerate}
特に、$X$ が $k$ 上幾何学的正規であるための必要十分条件は、$X_{k'}$ が $k'$ 上
幾何学的正規であることである。
\end{lemma}

\begin{proof}
(1) が (2) を導くことは明らかである。(2) を仮定する。$k''/k$ を有限純非分離体拡大とし、
$x'' \in X_{k''}$ を $x$ 上にある点とする（実際には一意である）。共通体拡大
$k'''/k$ を $k'/k$ と $k''/k$ について選ぶ（Fields, Lemma
\ref{fields-lemma-common-extension-field}）。点 $x''' \in X_{k'''}$ を、$x'$ と $x''$ の
両方の上にあるものとして選ぶ。局所環の写像
$$
\mathcal{O}_{X_{k''}, x''} \longrightarrow \mathcal{O}_{X_{k'''}, x''''}.
$$
を考える。これは平坦局所環準同型であり、従って忠実平坦である。(2) により右辺の局所環は
正規である。ゆえに Algebra, Lemma \ref{algebra-lemma-descent-normal} により
$\mathcal{O}_{X_{k''}, x''}$ は正規である。Lemma
\ref{varieties-lemma-geometrically-normal-at-point} により、$X$ は $x$ において幾何学的正規である。
\end{proof}

\begin{lemma}
\label{varieties-lemma-fibre-product-normal}
$k$ を体とする。$X$ を $k$ 上の幾何学的正規スキームとし、$Y$ を $k$ 上の正規スキームとする。
このとき $X \times_k Y$ は正規スキームである。
\end{lemma}

\begin{proof}
これは Algebra, Lemma \ref{algebra-lemma-geometrically-normal-tensor-normal} に帰着し、
その帰着は Lemma \ref{varieties-lemma-geometrically-normal} による。
\end{proof}

\begin{lemma}
\label{varieties-lemma-base-change-normal-by-separable}
$k$ を体とする。$X$ を $k$ 上の正規スキームとする。$K/k$ を分離体拡大とする。
このとき $X_K$ は正規スキームである。
\end{lemma}

\begin{proof}
Lemma \ref{varieties-lemma-fibre-product-normal} および Algebra, Lemma
\ref{algebra-lemma-separable-field-extension-geometrically-normal} から従う。
\end{proof}

\begin{lemma}
\label{varieties-lemma-geometrically-normal-stein}
$k$ を体とし、$X$ を $k$ 上の固有かつ幾何学的正規なスキームとする。次は同値である。
\begin{enumerate}
\item $H^0(X, \mathcal{O}_X) = k$ である。
\item $X$ は幾何学的連結である。
\item $X$ は幾何学的既約である。
\item $X$ は幾何学的整である。
\end{enumerate}
\end{lemma}

\begin{proof}
Lemma \ref{varieties-lemma-geometrically-reduced-stein} により (1) と (2) は同値である。局所 Noether な
正規スキーム（例えば $X_{\overline{k}}$）は、その既約成分の直和である
（Properties, Lemma \ref{properties-lemma-normal-Noetherian}）。従って (2) と (3) は同値である。
$X_{\overline{k}}$ は被約であると仮定されているので、(3) と (4) も同値である。
\end{proof}






\section{体の変更と局所 Noether スキーム}
\label{varieties-section-locally-Noetherian}

\noindent
$X$ を体 $k$ 上の局所 Noether スキームとする。$X_{k'}$ も局所 Noether になるとは
限らない。例えば $X = \Spec(\overline{\mathbf{Q}})$ かつ $k = \mathbf{Q}$ のとき、
$X_{\overline{\mathbf{Q}}}$ は
$\overline{\mathbf{Q}} \otimes_{\mathbf{Q}} \overline{\mathbf{Q}}$ のスペクトルであり、
これは Noether でない（ヒント：冪等元が多すぎる）。しかし有限生成体拡大によってのみ
基底変換を行うなら、Noether 性は保たれる（または $X$ が $k$ 上局所有限型である場合も、
この性質は基底変換で保たれるため同様である）。

\begin{lemma}
\label{varieties-lemma-locally-Noetherian-base-change}
$k$ を体とし、$X$ を $k$ 上のスキームとし、$k'/k$ を有限生成体拡大とする。
$X$ が局所 Noether であるための必要十分条件は、$X_{k'}$ が局所 Noether であることである。
\end{lemma}

\begin{proof}
Properties, Lemma \ref{properties-lemma-locally-Noetherian} により、$X$ がアフィンの場合へ
帰着できる。例えば $X = \Spec(A)$ と書く。この場合、$A$ が Noether であるための必要十分条件が
$A_{k'}$ が Noether であることであると示せばよい。
$A \to A_{k'} = k' \otimes_k A$ は忠実平坦なので、$A_{k'}$ が Noether ならば $A$ も
Noether である。これは Algebra, Lemma \ref{algebra-lemma-descent-Noetherian} による。
逆に $A$ が Noether ならば、Algebra, Lemma
\ref{algebra-lemma-Noetherian-field-extension} により $A_{k'}$ も Noether である。
\end{proof}






\section{幾何学的正則スキーム}
\label{varieties-section-geometrically-regular}

\noindent
体 $k$ 上の幾何学的正則スキームとは、$k$ 上の局所 Noether スキームであって、適切に
基礎体を変更しても正則であり続けるものである。$k$ 上有限型のスキームが幾何学的正則である
ための必要十分条件は、それが $k$ 上平滑であることである（Lemma
\ref{varieties-lemma-geometrically-regular-smooth}）。幾何学的正則性の概念は、形式的ファイバーのように
平滑性を用いることができない状況で最も重要となる（将来の参照をここに挿入する）。

\medskip\noindent
次の定義では局所 Noether スキームに限る。正則局所環であるという性質は Noether 局所環に対して
のみ定義されるからである。上の Lemma \ref{varieties-lemma-locally-Noetherian-base-change} により、
有限生成体拡大に限れば、この性質は基礎体の変更で保たれる。本節ではこの注意を以後断らずに用いる。
特に次の定義は意味をもつ。

\begin{definition}
\label{varieties-definition-geometrically-regular}
$k$ を体とし、$X$ を $k$ 上の局所 Noether スキームとする。
\begin{enumerate}
\item $x \in X$ とする。$X$ が
{\it $x$ において幾何学的正則}であるとは、$k$ 上で、任意の有限生成体拡大 $k'/k$ と、任意の点
$x' \in X_{k'}$ で $x$ 上にあるものに対し、局所環 $\mathcal{O}_{X_{k'}, x'}$ が
正則であることをいう。
\item $X$ が {\it $k$ 上幾何学的正則}であるとは、$X$ がそのすべての点において
幾何学的正則であることをいう。
\end{enumerate}
\end{definition}

\noindent
同様の定義により、体上の幾何学的 Cohen--Macaulay、$(R_k)$、および $(S_k)$ スキームを
定義できる。必要に応じて、これらについて別に節を加える。

\begin{lemma}
\label{varieties-lemma-geometrically-regular-at-point}
$k$ を体とし、$X$ を $k$ 上の局所 Noether スキームとし、$x \in X$ とする。次は同値である。
\begin{enumerate}
\item $X$ は $x$ において幾何学的正則である。
\item 任意の有限純非分離体拡大 $k'$（$k$ の拡大）と、任意の点
$x' \in X_{k'}$ で $x$ 上にあるものに対し、局所環 $\mathcal{O}_{X_{k'}, x'}$ は正則である。
\item 環 $\mathcal{O}_{X, x}$ は $k$ 上幾何学的正則である（Algebra, Definition
\ref{algebra-definition-geometrically-regular} を参照）。
\end{enumerate}
\end{lemma}

\begin{proof}
(1) が (2) を導くことは明らかである。(2) を仮定する。特にこれは
$\mathcal{O}_{X, x}$ が正則局所環であることを意味する。$k'/k$ を有限純非分離体拡大とする。
環 $\mathcal{O}_{X, x} \otimes_k k'$ を考える。Algebra, Lemma
\ref{algebra-lemma-p-ring-map} により、そのスペクトルは $\mathcal{O}_{X, x}$ の
スペクトルと同じである。従ってこれも局所環である（Algebra, Lemma
\ref{algebra-lemma-characterize-local-ring}）。ゆえに点 $x' \in X_{k'}$ で $x$ 上にあるものは
一意であり、$\mathcal{O}_{X_{k'}, x'} \cong \mathcal{O}_{X, x} \otimes_k k'$ である。
仮定によりこれは正則環である。従って幾何学的正則環の定義から (3) が従う。

\medskip\noindent
(3) を仮定する。$k'/k$ を体拡大とする。$\Spec(k') \to \Spec(k)$ は全射なので、
$X_{k'} \to X$ も全射である（Morphisms, Lemma
\ref{morphisms-lemma-base-change-surjective}）。任意の点 $x' \in X_{k'}$ で $x$ 上にあるものを
とる。局所環 $\mathcal{O}_{X_{k'}, x'}$ は環 $\mathcal{O}_{X, x} \otimes_k k'$ の
局所化である。従って仮定により正則であり、(1) が証明された。
\end{proof}

\begin{lemma}
\label{varieties-lemma-geometrically-regular}
$k$ を体とし、$X$ を $k$ 上の局所 Noether スキームとする。次は同値である。
\begin{enumerate}
\item $X$ は幾何学的正則である。
\item $X_{k'}$ は任意の有限生成体拡大 $k'/k$ に対して正則スキームである。
\item $X_{k'}$ は任意の有限純非分離体拡大 $k'/k$ に対して正則スキームである。
\item 任意のアフィン開集合 $U \subset X$ に対し、環 $\mathcal{O}_X(U)$ は
幾何学的正則である（Algebra, Definition
\ref{algebra-definition-geometrically-regular} を参照）。
\item アフィン開被覆 $X = \bigcup U_i$ であって、各 $\mathcal{O}_X(U_i)$ が $k$ 上
幾何学的正則であるものが存在する。
\end{enumerate}
\end{lemma}

\begin{proof}
(1) を仮定する。このとき任意の有限生成体拡大 $k'/k$ と任意の点
$x' \in X_{k'}$ に対し、$X_{k'}$ の $x'$ における局所環は正則である。Properties, Lemma
\ref{properties-lemma-characterize-regular} により、これは $X_{k'}$ が正則であることを意味する。
従って (2) が成り立つ。

\medskip\noindent
(2) が (3) を導くことは明らかである。

\medskip\noindent
(3) を仮定し、$U \subset X$ をアフィン開部分スキームとする。このとき $U_{k'}$ は、任意の
有限純非分離拡大 $k'/k$（$k = k'$ を含む）に対して正則スキームである。これは
$k' \otimes_k \mathcal{O}(U)$ が、すべての有限純非分離拡大 $k'/k$ に対して正則環であることを
意味する。従って $\mathcal{O}(U)$ は幾何学的正則な $k$-代数であり、(4) が成り立つ。

\medskip\noindent
(4) が (5) を導くことは明らかである。$X = \bigcup U_i$ を (5) のようなアフィン開被覆とする。
任意の体拡大 $k'/k$ に対して、基底変換 $X_{k'}$ は環
$\mathcal{O}_X(U_i) \otimes_k k'$ のスペクトルを貼り合わせて得られる
（Schemes, Section \ref{schemes-section-fibre-products}）。従って $X_{k'}$ は正則であり、
(1) が成り立つ。
\end{proof}

\begin{lemma}
\label{varieties-lemma-geometrically-regular-upstairs}
$k$ を体とし、$X$ を $k$ 上のスキームとし、$k'/k$ を有限生成体拡大とする。
$x \in X$ を点とし、$x' \in X_{k'}$ を $x$ 上にある点とする。次は同値である。
\begin{enumerate}
\item $X$ は $x$ において幾何学的正則である。
\item $X_{k'}$ は $x'$ において幾何学的正則である。
\end{enumerate}
特に、$X$ が $k$ 上幾何学的正則であるための必要十分条件は、$X_{k'}$ が $k'$ 上
幾何学的正則であることである。
\end{lemma}

\begin{proof}
(1) が (2) を導くことは明らかである。(2) を仮定する。$k''/k$ を有限純非分離体拡大とし、
$x'' \in X_{k''}$ を $x$ 上にある点とする（実際には一意である）。共通な有限生成体拡大
$k'''/k$ を $k'/k$ と $k''/k$ についてとることができる。Fields, Lemma
\ref{fields-lemma-common-extension-field} とその証明を参照せよ。そして点
$x''' \in X_{k'''}$ で $x'$ と $x''$ の両方の上にあるものをとる。局所環の写像
$$
\mathcal{O}_{X_{k''}, x''} \longrightarrow \mathcal{O}_{X_{k'''}, x''''}.
$$
を考える。これは Noether 局所環の間の平坦局所環準同型であり、従って忠実平坦である。
(2) により右辺の局所環は正則である。ゆえに Algebra, Lemma
\ref{algebra-lemma-flat-under-regular} により、$\mathcal{O}_{X_{k''}, x''}$ は正則である。
Lemma \ref{varieties-lemma-geometrically-regular-at-point} により、$X$ は $x$ において幾何学的正則である。
\end{proof}

\noindent
次の補題は Algebra, Lemma \ref{algebra-lemma-geometrically-regular-descent} の幾何学的変種である。

\begin{lemma}
\label{varieties-lemma-flat-under-geometrically-regular}
$k$ を体とする。$f : X \to Y$ を $k$ 上の局所 Noether スキームの射とする。
$x \in X$ を点とし、$y = f(x)$ とおく。$X$ が $x$ において幾何学的正則であり、
$f$ が $x$ において平坦ならば、$Y$ は $y$ において幾何学的正則である。
特に、$X$ が $k$ 上幾何学的正則であり、$f$ が平坦かつ全射ならば、$Y$ は $k$ 上
幾何学的正則である。
\end{lemma}

\begin{proof}
$k'$ を $k$ の有限純非分離拡大とする。$f' : X_{k'} \to Y_{k'}$ を $f$ の基底変換とする。
$x' \in X_{k'}$ を $x$ 上にある一意の点とする。$Y_{k'}$ が $y' = f'(x')$ において
正則であることを示せば、Lemma \ref{varieties-lemma-geometrically-regular} により $Y$ は $k$ 上
$y'$ において幾何学的正則である。Morphisms, Lemma
\ref{morphisms-lemma-base-change-module-flat} により、射 $X_{k'} \to Y_{k'}$ は $x'$ において
平坦である。従って環準同型
$$
\mathcal{O}_{Y_{k'}, y'}
\longrightarrow
\mathcal{O}_{X_{k'}, x'}
$$
は局所 Noether 環の間の平坦局所準同型であり、仮定により右辺は正則である。従って Algebra,
Lemma \ref{algebra-lemma-flat-under-regular} により左辺は正則局所環である。
\end{proof}

\begin{lemma}
\label{varieties-lemma-geometrically-regular-smooth}
$k$ を体とし、$X$ を $k$ 上局所有限型のスキームとし、$x \in X$ とする。
$X$ が $x$ において幾何学的正則であるための必要十分条件は、$X \to \Spec(k)$ が
$x$ において平滑であることである（Morphisms, Definition
\ref{morphisms-definition-smooth}）。
\end{lemma}

\begin{proof}
問題は $x$ の近傍で局所的なので、$X = \Spec(A)$ と仮定してよい。ここでこれは有限型
$k$-代数のスペクトルである。
$x$ が素イデアル $\mathfrak p$ に対応するとする。

\medskip\noindent
$A$ が $k$ 上 $\mathfrak p$ において平滑ならば、$A$ を局所化して $A$ が $k$ 上平滑であると
仮定してよい。この場合 $k' \otimes_k A$ は $k'$ 上平滑であり、これはすべての拡大体
$k'/k$ について成り立つ。また、
これらの Noether 環はそれぞれ Algebra, Lemma
\ref{algebra-lemma-characterize-smooth-over-field} により正則である。

\medskip\noindent
$X$ が $x$ において幾何学的正則であると仮定する。剰余体
$K := \kappa(x) = \kappa(\mathfrak p)$ を考える。これは $x$ の剰余体であり、$k$ の有限生成拡大である。
Algebra, Lemma \ref{algebra-lemma-make-separable} により、有限純非分離拡大 $k'/k$ が存在して、
合成体 $k'K$ は $k'$ の分離体拡大となる。素イデアル
$\mathfrak p' \subset A' = k' \otimes_k A$ を $\mathfrak p$ 上にあるものとしてとる。これは $\mathfrak p$ 上にある一意の
素イデアルである（Algebra, Lemma \ref{algebra-lemma-p-ring-map}）。従って剰余体
$K' := \kappa(\mathfrak p')$ は合成体 $k'K$ である。仮定により局所環
$(A')_{\mathfrak p'}$ は正則である。従って Algebra, Lemma
\ref{algebra-lemma-separable-smooth} により、$k' \to A'$ は $\mathfrak p'$ において平滑である。
さらにこれは Algebra, Lemma \ref{algebra-lemma-smooth-field-change-local} により、
$k \to A$ が $\mathfrak p$ において平滑であることを意味する。補題が証明された。
\end{proof}


\begin{example}
\label{varieties-example-geometrically-reduced-not-normal}
$k =\mathbf{F}_p(t)$ とする。正則多様体 $V$ であって $k$ 上で幾何学的被約でないものの例は
容易に与えられる。例えば $\Spec(k[x]/(x^p - t))$ をとればよい。実は、幾何学的被約ではあるが
幾何学的正規ですらない正則多様体 $V$ の例も存在する。すなわち $p > 2$ に対してスキーム
$V = \Spec(k[x, y]/(y^2 - x^p + t))$ をとる。多項式
$y^2 - x^p + t \in k[x, y]$ は既約なので、これは多様体である。
射 $V \to \Spec(k)$ は、点 $v_0 \in V$ であって極大イデアル $(y, x^p - t)$ に対応するものを除くすべての
点において平滑である（$2y$ が可逆だからである）。特に、$V$ は $v_0$ を除くすべての点において
（幾何学的）正則である。局所環
$$
\mathcal{O}_{V, v_0} = \left(k[x, y]/(y^2 - x^p + t)\right)_{(y, x^p - t)}
$$
は次元 $1$ の整域である。その極大イデアルは $1$ 個の元、すなわち $y$ で生成される。従ってこれは
離散付値環であり正則である。$k' = k[t^{1/p}]$ とする。記号
$t' = t^{1/p} \in k'$、$V' = V_{k'}$、および $v'_0 \in V'$ を用い、最後のものを $v_0$ 上にある
一意の点とする。$k'$ 上では $x^p - t = (x - t')^p$ と書けるが、多項式
$y^2 - (x - t')^p$ はなお既約であり、$V'$ はなお多様体である。しかし元
$$
\frac{y}{x - t'} \in (\text{fraction field of }\mathcal{O}_{V', v'_0})
$$
は $\mathcal{O}_{V', v'_0}$ 上整であり（二乗を計算すればよい）、しかもそこに含まれない。
従って $V'$ は $v'_0$ において正規でない。これで例が完了する。
\end{example}







\section{体の変更と Cohen--Macaulay 性}
\label{varieties-section-CM}

\noindent
次の補題は、幾何学的 Cohen--Macaulay スキームを定義しても意味がないことを述べる。
なぜなら、それは Cohen--Macaulay スキームと同じものになるからである。

\begin{lemma}
\label{varieties-lemma-CM-base-change}
$X$ を体 $k$ 上の局所 Noether スキームとする。
$k'/k$ を有限生成体拡大とする。
$x \in X$ を点とし、$x' \in X_{k'}$ を $x$ 上にある点とする。このとき
$$
\mathcal{O}_{X, x}\text{ is Cohen-Macaulay}
\Leftrightarrow
\mathcal{O}_{X_{k'}, x'}\text{ is Cohen-Macaulay}
$$
が成り立つ。$X$ が $k$ 上局所有限型ならば、任意の体拡大 $k'/k$ について同じことが成り立つ。
\end{lemma}

\begin{proof}
補題の第一の場合は Algebra, Lemma \ref{algebra-lemma-CM-geometrically-CM} から従う。
第二の場合は Algebra, Lemma \ref{algebra-lemma-extend-field-CM-locus} と同値である。
\end{proof}







\section{体の変更と Jacobson 性}
\label{varieties-section-overfield}

\noindent
体上局所有限型のスキームには閉点が豊富に存在する。すなわち、それは Jacobson である。
さらに、各剰余体は基礎体の有限拡大である。

\begin{lemma}
\label{varieties-lemma-locally-finite-type-Jacobson}
$X$ を $k$ 上局所有限型のスキームとする。このとき
\begin{enumerate}
\item 任意の閉点 $x \in X$ に対して拡大 $\kappa(x)/k$ は代数的であり、かつ
\item $X$ は Jacobson スキームである
（Properties, Definition \ref{properties-definition-jacobson}）。
\end{enumerate}
\end{lemma}

\begin{proof}
スキームが Jacobson であるための必要十分条件は、それが Jacobson スキームによるアフィン開被覆を
もつことである。Properties, Lemma \ref{properties-lemma-locally-jacobson} を参照せよ。
閉点における剰余体についての性質も $X$ 上局所的である。従って $X$ はアフィンであると仮定してよい。
この場合、結論は Hilbert の零点定理から従う。Algebra, Theorem
\ref{algebra-theorem-nullstellensatz} を参照せよ。またこれは Morphisms, Lemmas
\ref{morphisms-lemma-jacobson-finite-type-points},
\ref{morphisms-lemma-Jacobson-universally-Jacobson}, および
\ref{morphisms-lemma-ubiquity-Jacobson-schemes} を組み合わせても従う。
\end{proof}

\noindent
$X$ が局所有限型でなくても、十分大きな基礎体拡大を行えば同じ結論が得られることが分かる。

\begin{lemma}
\label{varieties-lemma-make-Jacobson}
$X$ を体 $k$ 上のスキームとする。濃度が十分大きい任意の体拡大 $K/k$ に対して
\begin{enumerate}
\item 任意の閉点 $x \in X_K$ に対して拡大 $\kappa(x)/K$ は代数的であり、かつ
\item $X_K$ は Jacobson スキームである
（Properties, Definition \ref{properties-definition-jacobson}）。
\end{enumerate}
\end{lemma}

\begin{proof}
アフィン開被覆 $X = \bigcup U_i$ を選ぶ。Algebra, Lemma
\ref{algebra-lemma-base-change-Jacobson} および Properties, Lemma
\ref{properties-lemma-affine-jacobson} により、濃度 $\kappa_i$ が存在し、
$U_{i, K}$ が $K$ 上で所望の性質をもつための条件として $\#(K) \geq \kappa_i$ で十分である。
$\kappa = \max\{\kappa_i\}$ とおく。このとき $K$ の濃度が $\kappa$ より大きければ、
各 $U_{i, K}$ が補題の結論を満たすことが分かる。従って Properties, Lemma
\ref{properties-lemma-locally-jacobson} により $X_K$ は Jacobson である。
$X_K$ の閉点における剰余体についての主張は、各 $U_{i, K}$ の閉点における剰余体についての
対応する主張から従う。
\end{proof}





\section{体の変更と豊富な可逆層}
\label{varieties-section-change-fields-ample}

\noindent
次の結果は、この節における結果の典型的なものである。

\begin{lemma}
\label{varieties-lemma-ample-after-field-extension}
$k$ を体とする。$X$ を $k$ 上のスキームとする。$X_K$ 上に豊富な可逆層が存在するような
体拡大 $K/k$ があれば、$X$ は豊富な可逆層をもつ。
\end{lemma}

\begin{proof}
$K/k$ を、$X_K$ が豊富な可逆層 $\mathcal{L}$ をもつような体拡大とする。
射 $X_K \to X$ は全射である。従って $X$ は準コンパクトスキームの像として準コンパクトである
（Properties, Definition \ref{properties-definition-ample}）。$X_K$ は準分離的なので
（Properties, Lemma \ref{properties-lemma-affine-s-opens-cover-quasi-separated} による）、
$X$ は準分離的であることが分かる。実際、$U, V \subset X$ がアフィン開ならば、
$(U \cap V)_K = U_K \cap V_K$ は準コンパクトであり、$(U \cap V)_K \to U \cap V$ は全射である。
従って Schemes, Lemma \ref{schemes-lemma-characterize-quasi-separated} が適用できる。

\medskip\noindent
$K = \colim A_i$ と、$K$ を $k$ 上有限型である部分代数の余極限として書く。
$X_i = X \times_{\Spec(k)} \Spec(A_i)$ と記す。$X_K = \lim X_i$ なので、ある $i$ と、
可逆層 $\mathcal{L}_i$ であって $X_i$ 上にあり、その $X_K$ への引き戻しが $\mathcal{L}$ となるものが得られる
（Limits, Lemma \ref{limits-lemma-descend-invertible-modules}；ここおよび以下では、上で示したように
$X$ が準コンパクトかつ準分離的であることを用いる）。Limits, Lemma
\ref{limits-lemma-limit-ample} により、$\mathcal{L}_i$ は、必要なら $i$ を大きくして、豊富であると
仮定してよい。そのような $i$ を固定し、$\mathfrak m \subset A_i$ を極大イデアルとする。
Hilbert の零点定理（Algebra, Theorem \ref{algebra-theorem-nullstellensatz}）により、剰余体
$k' = A_i/\mathfrak m$ は $k$ の有限拡大である。従って $X_{k'} \subset X_i$ は閉部分スキームであり、
よって豊富な可逆層をもつ（Properties, Lemma \ref{properties-lemma-ample-on-closed}）。
$X_{k'} \to X$ は有限局所自由なので、Divisors, Proposition
\ref{divisors-proposition-push-down-ample} により $X$ は豊富な可逆層をもつ。
\end{proof}

\begin{lemma}
\label{varieties-lemma-quasi-affine-after-field-extension}
$k$ を体とする。$X$ を $k$ 上のスキームとする。$X_K$ が準アフィンとなる体拡大 $K/k$ があれば、
$X$ は準アフィンである。
\end{lemma}

\begin{proof}
$K/k$ を、$X_K$ が準アフィンとなる体拡大とする。射 $X_K \to X$ は全射である。
従って $X$ は準コンパクトスキームの像として準コンパクトである
（Properties, Definition \ref{properties-definition-quasi-affine}）。$X_K$ は準分離的なので
（アフィンスキームの開部分スキームだからである）、$X$ は準分離的であることが分かる。
実際、$U, V \subset X$ がアフィン開ならば、$(U \cap V)_K = U_K \cap V_K$ は準コンパクトであり、
$(U \cap V)_K \to U \cap V$ は全射である。従って Schemes, Lemma
\ref{schemes-lemma-characterize-quasi-separated} が適用できる。

\medskip\noindent
$K = \colim A_i$ と、$K$ を $k$ 上有限型である部分代数の余極限として書く。
$X_i = X \times_{\Spec(k)} \Spec(A_i)$ と記す。$X_K = \lim X_i$ なので、ある $i$ であって
$X_i$ が準アフィンとなるものが得られる（Limits, Lemma \ref{limits-lemma-limit-quasi-affine}；ここでは、上で示したように
$X$ が準コンパクトかつ準分離的であることを用いる）。Hilbert の零点定理
（Algebra, Theorem \ref{algebra-theorem-nullstellensatz}）により、剰余体
$k' = A_i/\mathfrak m$ は $k$ の有限拡大である。従って $X_{k'} \subset X_i$ は閉部分スキームであり、
よって準アフィンである（Properties, Lemma
\ref{properties-lemma-quasi-affine-locally-closed}）。$X_{k'} \to X$ は有限局所自由なので、
Divisors, Lemma \ref{divisors-lemma-push-down-quasi-affine} により結論を得る。
\end{proof}

\begin{lemma}
\label{varieties-lemma-quasi-projective-after-field-extension}
$k$ を体とする。$X$ を $k$ 上のスキームとする。$X_K$ が $K$ 上準射影的となる体拡大 $K/k$ があれば、
$X$ は $k$ 上準射影的である。
\end{lemma}

\begin{proof}
定義により、スキームの射 $g : Y \to T$ が準射影的であるとは、それが局所有限型かつ準コンパクトであり、
$g$-豊富な可逆層が $Y$ 上に存在することをいう。$K/k$ を、$X_K$ が $K$ 上準射影的となる体拡大とする。
$\Spec(A) \subset X$ をアフィン開とする。このとき $U_K$ は $X_K$ のアフィン開部分スキームなので、
$A_K$ は有限型 $K$-代数である。すると Algebra, Lemma
\ref{algebra-lemma-finite-type-descends} により $A$ は有限型 $k$-代数である。
従って $X \to \Spec(k)$ は局所有限型である。$X_K \to \Spec(K)$ は準コンパクトなので、
$X_K$ は準コンパクトであり、従って $X$ は準コンパクトであり、ゆえに $X \to \Spec(k)$ は有限型である。
Morphisms, Lemma \ref{morphisms-lemma-finite-type-over-affine-ample-very-ample} により、
$X_K$ は豊富な可逆層をもつ。すると Lemma \ref{varieties-lemma-ample-after-field-extension} により、
$X$ は豊富な可逆層をもつ。従って Morphisms, Lemma
\ref{morphisms-lemma-finite-type-over-affine-ample-very-ample} により、$X \to \Spec(k)$ は準射影的である。
\end{proof}

\noindent
次の補題は Descent, Lemma \ref{descent-lemma-descending-property-proper} の特別な場合である。

\begin{lemma}
\label{varieties-lemma-proper-after-field-extension}
$k$ を体とする。$X$ を $k$ 上のスキームとする。$X_K$ が $K$ 上固有となる体拡大 $K/k$ があれば、
$X$ は $k$ 上固有である。
\end{lemma}

\begin{proof}
$K/k$ を、$X_K$ が $K$ 上固有となる体拡大とする。このことから $X_K$ は分離的かつ準コンパクトである
ことを思い出そう（Morphisms, Definition \ref{morphisms-definition-proper}）。射 $X_K \to X$ は全射である。
従って $X$ は準コンパクトスキームの像として準コンパクトである
（Properties, Definition \ref{properties-definition-ample}）。$X_K$ は分離的なので、$X$ は準分離的である
ことが分かる。実際、$U, V \subset X$ がアフィン開ならば、$(U \cap V)_K = U_K \cap V_K$ は
準コンパクトであり、$(U \cap V)_K \to U \cap V$ は全射である。従って Schemes, Lemma
\ref{schemes-lemma-characterize-quasi-separated} が適用できる。

\medskip\noindent
$K = \colim A_i$ と、$K$ を $k$ 上有限型である部分代数の余極限として書く。
$X_i = X \times_{\Spec(k)} \Spec(A_i)$ と記す。Limits, Lemma
\ref{limits-lemma-eventually-proper} により、ある $i$ が存在して $X_i \to \Spec(A_i)$ が固有となる。
ここでは、上で示したように $X$ が準コンパクトかつ準分離的であることを用いる。
極大イデアル $\mathfrak m \subset A_i$ を選ぶ。Hilbert の零点定理
（Algebra, Theorem \ref{algebra-theorem-nullstellensatz}）により、剰余体
$k' = A_i/\mathfrak m$ は $k$ の有限拡大である。基底変換 $X_{k'} \to \Spec(k')$ は固有である
（Morphisms, Lemma \ref{morphisms-lemma-base-change-proper}）。$k'/k$ は有限なので、
$X_{k'} \to X$ と合成 $X_{k'} \to \Spec(k)$ はともに固有である
（Morphisms, Lemmas \ref{morphisms-lemma-finite-proper},
\ref{morphisms-lemma-base-change-proper}, および
\ref{morphisms-lemma-composition-proper}）。第一のことから $X$ は $k$ 上分離的である。実際、$X_{k'}$ は
分離的である（Morphisms, Lemma
\ref{morphisms-lemma-image-universally-closed-separated}）。第二のことから、Morphisms, Lemma
\ref{morphisms-lemma-image-proper-is-proper} により $X \to \Spec(k)$ は固有である。
\end{proof}

\begin{lemma}
\label{varieties-lemma-projective-after-field-extension}
$k$ を体とする。$X$ を $k$ 上のスキームとする。$X_K$ が $K$ 上射影的となる体拡大 $K/k$ があれば、
$X$ は $k$ 上射影的である。
\end{lemma}

\begin{proof}
$k$ 上のスキームが $k$ 上射影的であるための必要十分条件は、それが $k$ 上準射影的かつ固有であることである。
Morphisms, Lemma \ref{morphisms-lemma-projective-is-quasi-projective-proper} を参照せよ。
従って補題は Lemmas \ref{varieties-lemma-quasi-projective-after-field-extension} および
\ref{varieties-lemma-proper-after-field-extension} から従う。
\end{proof}




\section{接空間}
\label{varieties-section-tangent-spaces}

\noindent
この節では、双対数に値をもつ点を用いて、スキームの射の始域の一点における接空間を定義する。

\begin{definition}
\label{varieties-definition-dual-numbers}
任意の環 $R$ に対し、$R$ 上の {\it 双対数}とは $R$-代数 $R[\epsilon]$ である。
$R$-加群としては基底 $1$, $\epsilon$ をもつ自由加群であり、その $R$-代数構造は
$\epsilon^2 = 0$ とおくことから得られる。
\end{definition}

\noindent
$f : X \to S$ をスキームの射とする。$x \in X$ を点とし、その像を $s = f(x)$ とする。これは $S$ の点である。
実線からなる可換図式
\begin{equation}
\label{varieties-equation-tangent-space}
\vcenter{
\xymatrix{
\Spec(\kappa(x)) \ar[r] \ar[dr] \ar@/^1pc/[rr] &
\Spec(\kappa(x)[\epsilon]) \ar@{.>}[r] \ar[d]&
X \ar[d] \\
&
\Spec(\kappa(s)) \ar[r] &
S
}
}
\end{equation}
を考える。ここで曲がった矢印は $\Spec(\kappa(x))$ から $X$ への標準射である。

\begin{lemma}
\label{varieties-lemma-tangent-space}
(\ref{varieties-equation-tangent-space}) を可換にする点線の矢印全体の集合は、標準的な
$\kappa(x)$-ベクトル空間構造をもつ。
\end{lemma}

\begin{proof}
$\kappa = \kappa(x)$ とおく。スキームの圏において押し出し
$$
\Spec(\kappa[\epsilon]) \amalg_{\Spec(\kappa)} \Spec(\kappa[\epsilon])
= \Spec(\kappa[\epsilon_1, \epsilon_2])
$$
があることに注意せよ。ここで $\kappa[\epsilon_1, \epsilon_2]$ は $\kappa$-代数であり、基底
$1, \epsilon_1, \epsilon_2$ をもち、
$\epsilon_1^2 = \epsilon_1\epsilon_2 = \epsilon_2^2 = 0$ を満たす。
これは、環についての対応する結果と、局所環のスペクトルからスキームへの射の記述
（Schemes, Lemma \ref{schemes-lemma-morphism-from-spec-local-ring}）から直ちに従う。
二つの矢印 $\theta_1, \theta_2 : \Spec(\kappa[\epsilon]) \to X$ が与えられたとき、射
$$
\theta_1 + \theta_2 :
\Spec(\kappa[\epsilon]) \to
\Spec(\kappa[\epsilon_1, \epsilon_2])
\xrightarrow{\theta_1, \theta_2} X
$$
を考えられる。ここで第一の矢印は $\epsilon_i \mapsto \epsilon$ により与えられる。
一方、$\lambda \in \kappa$ が与えられると、$\Spec(\kappa[\epsilon])$ の自己写像がある。
これは $\kappa$-代数 $\kappa[\epsilon]$ の自己準同型であって、$\epsilon$ を $\lambda \epsilon$ に送るものに対応する。
$\theta : \Spec(\kappa[\epsilon]) \to X$ にこの自己写像を前合成して $\lambda \theta$ を得る。
適切なスキームの可換図式の存在を確認すれば、ベクトル空間の公理を検証できる。
詳細は省略する。（別の証明として、すべてを局所環で表し、環準同型の水準で
ベクトル空間の公理を検証することもできる。）
\end{proof}

\begin{definition}
\label{varieties-definition-tangent-space}
$f : X \to S$ をスキームの射とし、$x \in X$ とする。(\ref{varieties-equation-tangent-space}) を可換にする
点線の矢印全体の集合に、その標準的な $\kappa(x)$-ベクトル空間構造を入れたものを
{\it $X$ の $S$ 上 $x$ における接空間}と呼び、$T_{X/S, x}$ と記す。
この空間の元を {\it $X/S$ の $x$ における接ベクトル}と呼ぶ。
\end{definition}

\noindent
$x \in X$ における接ベクトルは、スキーム論的ファイバー $X_s$、すなわち
$f : X \to S$ の $s = f(x)$ 上のファイバーの中にあるので、標準的な同一視
\begin{equation}
\label{varieties-equation-tangent-space-fibre}
T_{X/S, x} = T_{X_s/s, x}
\end{equation}
を得る。点の関手を用いるこの心地よい定義には次の代数的記述がある。この記述から、
{\it $X$ の $S$ 上 $x$ における余接空間}を $\kappa(x)$-ベクトル空間
$$
T^*_{X/S, x} = \Omega_{X/S, x} \otimes_{\mathcal{O}_{X, x}} \kappa(x)
$$
として定義することが示唆される。実際、これは標準的な $\kappa(x)$-双対、すなわち
$X$ の $S$ 上 $x$ における接空間の双対だからである。

\begin{lemma}
\label{varieties-lemma-tangent-space-cotangent-space}
$f : X \to S$ をスキームの射とし、$x \in X$ とする。$\kappa(x)$ 上のベクトル空間の
標準同型
$$
T_{X/S, x} = \Hom_{\mathcal{O}_{X, x}}(\Omega_{X/S, x}, \kappa(x))
$$
が存在する。
\end{lemma}

\begin{proof}
$\kappa = \kappa(x)$ とおく。$\theta \in T_{X/S, x}$ が与えられると、写像
$$
\theta^*\Omega_{X/S} \to
\Omega_{\Spec(\kappa[\epsilon])/\Spec(\kappa(s))} \to
\Omega_{\Spec(\kappa[\epsilon])/\Spec(\kappa)}
$$
を得る。切断をとると、$\mathcal{O}_{X, x}$-線形写像
$\xi_\theta : \Omega_{X/S, x} \to \kappa \text{d}\epsilon$、すなわち補題の式の右辺の元を得る。
$\theta \mapsto \xi_\theta$ が同型であることを示すため、$S$ を $s$ に、$X$ をスキーム論的ファイバー
$X_s$ に置き換えてよい。実際、式の両辺はスキーム論的ファイバーのみに依存する。
$T_{X/S, x}$ については明らかであり、右辺については Morphisms, Lemma
\ref{morphisms-lemma-base-change-differentials} を参照せよ。また、$X$ を
$\mathcal{O}_{X, x}$ のスペクトルに置き換えてもよい。これは $T_{X/S, x}$ を変えず
（Schemes, Lemma \ref{schemes-lemma-morphism-from-spec-local-ring}）、$\Omega_{X/S, x}$ も変えない
（Modules, Lemma \ref{modules-lemma-stalk-module-differentials}）。

\medskip\noindent
$(A, \mathfrak m, \kappa)$ を体 $k$ 上の局所環とする。証明を終えるには、任意の $A$-線形写像
$\xi : \Omega_{A/k} \to \kappa$ が一意な $k$-代数写像 $\varphi : A \to \kappa[\epsilon]$ から生じ、
この写像が標準写像 $c : A \to \kappa$ と $\epsilon$ を法として一致することを示せばよい。
$\varphi(a) = c(a) + D(a) \epsilon$ と書く。すると $a \mapsto D(a)$ は $k$-導分であることが分かる。
$\Omega_{A/k}$ の普遍性を用いると、各 $D$ が一意な $\xi$ に対応し、逆も成り立つ。
これで証明が完了する。
\end{proof}

\begin{lemma}
\label{varieties-lemma-tangent-space-rational-point}
$f : X \to S$ をスキームの射とする。$x \in X$ を点とし、$s = f(x) \in S$ とおく。
$\kappa(x) = \kappa(s)$ と仮定する。このとき標準同型
$$
\mathfrak m_x/(\mathfrak m_x^2 + \mathfrak m_s\mathcal{O}_{X, x})
=
\Omega_{X/S, x} \otimes_{\mathcal{O}_{X, x}} \kappa(x)
$$
および
$$
T_{X/S, x} =
\Hom_{\kappa(x)}(
\mathfrak m_x/(\mathfrak m_x^2 + \mathfrak m_s\mathcal{O}_{X, x}),
\kappa(x))
$$
が存在する。これは、より一般に $\kappa(x)/\kappa(s)$ が分離代数拡大の場合にも成り立つ。
\end{lemma}

\begin{proof}
第二の同型は、Lemma \ref{varieties-lemma-tangent-space-cotangent-space} により第一の同型から従う。
第一の同型については、Morphisms, Lemma \ref{morphisms-lemma-base-change-differentials} により、
$S$ を $s$ に、$X$ を $X_s$ に置き換えてよい。また Modules, Lemma
\ref{modules-lemma-stalk-module-differentials} により、$X$ を $\mathcal{O}_{X, x}$ のスペクトルに
置き換えてよい。従って次の代数的事実を示せばよい。$(A, \mathfrak m, \kappa)$ を体 $k$ 上の局所環とし、
$\kappa/k$ が分離代数拡大であるとする。このとき標準写像
$$
\mathfrak m/\mathfrak m^2
\longrightarrow
\Omega_{A/k} \otimes \kappa
$$
は同型である。$\mathfrak m/\mathfrak m^2 = H_1(\NL_{\kappa/A})$ であることに注意せよ。
Algebra, Lemma \ref{algebra-lemma-exact-sequence-NL} により、$\Omega_{\kappa/k} = 0$ および
$H_1(\NL_{\kappa/k}) = 0$ を示せば十分である。$\kappa$ は $k$ の中にあるその有限分離拡大の合併なので、
$\kappa$ が $k$ の有限分離拡大である場合に示せば十分である
（Algebra, Lemma \ref{algebra-lemma-colimits-NL}）。この場合、環準同型 $k \to \kappa$ はエタールであり、
従って $\NL_{\kappa/k} = 0$ である（ほとんど定義による。Algebra, Section
\ref{algebra-section-etale} を参照せよ）。
\end{proof}

\begin{lemma}
\label{varieties-lemma-map-tangent-spaces}
$f : X \to Y$ を基礎スキーム $S$ 上のスキームの射とする。$x \in X$ を点とし、$y = f(x)$ とおく。
$\kappa(y) = \kappa(x)$ ならば、$f$ は自然な線形写像
$$
\text{d}f : T_{X/S, x} \longrightarrow T_{Y/S, y}
$$
を誘導し、これは Lemma \ref{varieties-lemma-tangent-space-cotangent-space} の同一視を通じて線形写像
$\Omega_{Y/S, y} \otimes \kappa(y) \to \Omega_{X/S, x} \otimes \kappa(x)$ の双対である。
\end{lemma}

\begin{proof}
省略する。
\end{proof}

\begin{lemma}
\label{varieties-lemma-tangent-space-product}
$X$, $Y$ を基礎 $S$ 上のスキームとする。$x \in X$ と $y \in Y$ が同じ像点
$s \in S$ をもち、$\kappa(s) = \kappa(x)$ および $\kappa(s) = \kappa(y)$ を満たすとする。
標準同型
$$
T_{X \times_S Y/S, (x, y)} = T_{X/S, x} \oplus T_{Y/S, y}
$$
が存在する。左辺から右辺への写像は、射影 $X \times_S Y \to X$ と $X \times_S Y \to Y$ から生じる
接空間上の写像により誘導される。右辺から左辺への写像は、ファイバーの接空間との同一視
(\ref{varieties-equation-tangent-space-fibre}) を通じて写像 $1 \times y : X_s \to X_s \times_s Y_s$ および
$x \times 1 : Y_s \to X_s \times_s Y_s$ により誘導される。
\end{lemma}

\begin{proof}
直和分解は Morphisms, Lemma \ref{morphisms-lemma-differential-product} から Lemma
\ref{varieties-lemma-tangent-space-rational-point} を介して従う。写像との両立性は Lemma
\ref{varieties-lemma-map-tangent-spaces} から従う。
\end{proof}

\begin{lemma}
\label{varieties-lemma-injective-tangent-spaces-unramified}
$f : X \to Y$ を基礎スキーム $S$ 上局所有限型のスキームの射とする。$x \in X$ を点とし、
$y = f(x)$ とおき、$\kappa(y) = \kappa(x)$ と仮定する。このとき次は同値である。
\begin{enumerate}
\item $\text{d}f : T_{X/S, x} \longrightarrow T_{Y/S, y}$ は単射である。
\item $f$ は $x$ において非分岐である。
\end{enumerate}
\end{lemma}

\begin{proof}
射 $f$ は Morphisms, Lemma \ref{morphisms-lemma-permanence-finite-type} により局所有限型である。
写像 $\text{d}f$ が単射であるための必要十分条件は、
$\Omega_{Y/S, y} \otimes \kappa(y) \to \Omega_{X/S, x} \otimes \kappa(x)$ が全射であることである
（Lemma \ref{varieties-lemma-map-tangent-spaces}）。完全列
$f^*\Omega_{Y/S} \to \Omega_{X/S} \to \Omega_{X/Y} \to 0$
（Morphisms, Lemma \ref{morphisms-lemma-triangle-differentials}）により、これは
$\Omega_{X/Y, x} \otimes \kappa(x) = 0$ と同値である。従って結論は Morphisms, Lemma
\ref{morphisms-lemma-unramified-at-point} から従う。
\end{proof}







\section{生成的に有限な射}
\label{varieties-section-generically-finite}

\noindent
この節では、Morphisms, Section \ref{morphisms-section-generically-finite} で研究した
スキームの生成的に有限な射という概念を再検討する。

\begin{lemma}
\label{varieties-lemma-quasi-finite-in-codim-1}
$f : X \to Y$ を局所有限型とする。$y \in Y$ を、$\mathcal{O}_{Y, y}$ が次元 $\leq 1$ の
Noether 環となる点とする。さらに次の条件のいずれか一つが満たされると仮定する。
\begin{enumerate}
\item 任意の生成点 $\eta$ であって $X$ の既約成分の生成点となるものに対して、体拡大
$\kappa(\eta)/\kappa(f(\eta))$ は有限（または代数的）である。
\item 任意の生成点 $\eta$ であって $X$ の既約成分の生成点となり $f(\eta) \leadsto y$ を満たすものに対して、
体拡大 $\kappa(\eta)/\kappa(f(\eta))$ は有限（または代数的）である。
\item $f$ は $X$ の既約成分のすべての生成点において準有限である。
\item $Y$ は局所 Noether であり、$f$ は $X$ の稠密な点集合において準有限である。
\item ここにさらに追加する。
\end{enumerate}
このとき $f$ は $X$ の、$y$ 上にあるすべての点において準有限である。
\end{lemma}

\begin{proof}
条件 (4) ならば $X$ は局所 Noether である
（Morphisms, Lemma \ref{morphisms-lemma-finite-type-noetherian}）。射が準有限である点の集合は開である
（Morphisms, Lemma \ref{morphisms-lemma-quasi-finite-points-open}）。局所 Noether スキームの稠密開集合は
既約成分のすべての生成点を含むので、(4) は (3) を導く。条件 (3) は Morphisms, Lemma
\ref{morphisms-lemma-residue-field-quasi-finite} により条件 (1) を導く。条件 (1) は条件 (2) を導く。
従って (2) が成り立つ場合に補題を証明すれば十分である。

\medskip\noindent
(2) が成り立つと仮定する。$\Spec(\mathcal{O}_{Y, y})$ は $Y$ の、$y$ に特殊化する点全体の集合であることを
思い出そう。Schemes, Lemma \ref{schemes-lemma-specialize-points} を参照せよ。これを Morphisms, Lemma
\ref{morphisms-lemma-base-change-quasi-finite} と組み合わせると、$Y$ を $\Spec(\mathcal{O}_{Y, y})$ に
置き換えてよいことが分かる。従って $Y = \Spec(B)$ と仮定してよい。ここで $B$ は次元 $\leq 1$ の
Noether 局所環であり、$y$ は閉点である。

\medskip\noindent
$X = \bigcup X_i$ を、被約閉部分スキームとみなした $X$ の既約成分による表示とする。
各ファイバー $X_{i, y}$ が離散空間であることを示せれば、$X_y = \bigcup X_{i, y}$ も離散であり、
Morphisms, Lemma \ref{morphisms-lemma-quasi-finite-at-point-characterize} により、
$X \to Y$ は $X_y$ のすべての点において準有限であると結論できる。
従って $X$ は整スキームであると仮定してよい。

\medskip\noindent
$X \to Y$ が生成点 $\eta$、すなわち $X$ の生成点を $y$ に写すならば、$X$ は $\kappa(y)$ の有限拡大のスペクトルであり、
結論は成り立つ。$X$ が $\eta$ を、極小素イデアル $\mathfrak q$ であって $B$ に属し
$\mathfrak m_B$ と異なるものに対応する点へ写すと仮定する。分解 $X \to \Spec(B/\mathfrak q) \to \Spec(B)$ を得る。
$x \in X$ を $y$ 上にある点とする。次元公式
（Morphisms, Lemma \ref{morphisms-lemma-dimension-formula}）により
$$
\dim(\mathcal{O}_{X, x}) \leq \dim(B/\mathfrak q) +
\text{trdeg}_{\kappa(\mathfrak q)}(R(X)) - \text{trdeg}_{\kappa(y)} \kappa(x)
$$
を得る。$\dim(B/\mathfrak q) = 1$ であり、$X$ の生成点は $x$ と異なり $x$ に特殊化し、
$R(X)$ は $\kappa(\mathfrak q)$ 上代数的であることが分かっている。従って
$$
1 \leq 1 - \text{trdeg}_{\kappa(y)} \kappa(x)
$$
を得る。ゆえに Morphisms, Lemma
\ref{morphisms-lemma-algebraic-residue-field-extension-closed-point-fibre} により、各点 $x$ であって $X_y$ に属するものは
$X_y$ において閉である。従って Morphisms, Lemma
\ref{morphisms-lemma-quasi-finite-at-point-characterize} により、$X \to Y$ は各点 $x$ であって $X_y$ に属するものにおいて
準有限である（同じ補題から $X_y$ が離散位相空間であることも従う）。
\end{proof}

\begin{lemma}
\label{varieties-lemma-finite-in-codim-1}
$f : X \to Y$ を固有射とする。$y \in Y$ を、$\mathcal{O}_{Y, y}$ が次元 $\leq 1$ の Noether 環となる
点とする。さらに次の条件のいずれか一つが満たされると仮定する。
\begin{enumerate}
\item 任意の生成点 $\eta$ であって $X$ の既約成分の生成点となるものに対して、体拡大
$\kappa(\eta)/\kappa(f(\eta))$ は有限（または代数的）である。
\item 任意の生成点 $\eta$ であって $X$ の既約成分の生成点となり $f(\eta) \leadsto y$ を満たすものに対して、
体拡大 $\kappa(\eta)/\kappa(f(\eta))$ は有限（または代数的）である。
\item $f$ は $X$ のすべての生成点において準有限である。
\item $Y$ は局所 Noether であり、$f$ は $X$ の稠密な点集合において準有限である。
\item ここにさらに追加する。
\end{enumerate}
このとき開近傍 $V \subset Y$ であって $y$ を含み、$f^{-1}(V) \to V$ が有限となるものが存在する。
\end{lemma}

\begin{proof}
Lemma \ref{varieties-lemma-quasi-finite-in-codim-1} により、射 $f$ はファイバー $X_y$ のすべての点において
準有限である。従って $X_y$ は離散位相空間である
（Morphisms, Lemma \ref{morphisms-lemma-quasi-finite-at-point-characterize}）。$f$ は固有なので、
ファイバー $X_y$ は準コンパクト、すなわち有限である。従って Cohomology of Schemes, Lemma
\ref{coherent-lemma-proper-finite-fibre-finite-in-neighbourhood} を適用して結論を得る。
\end{proof}

\begin{lemma}
\label{varieties-lemma-modification-normal-iso-over-codimension-1}
$X$ を Noether スキームとする。$f : Y \to X$ を、$Y$ が被約であるスキームの双有理固有射とする。
$U \subset X$ を、$f$ が同型となる最大の開集合とする。このとき $U$ は次を含む。
\begin{enumerate}
\item 余次元 $0$ の、$X$ におけるすべての点。
\item すべての $x \in X$ であって、余次元 $1$ を $X$ 上もち、$\mathcal{O}_{X, x}$ が離散付値環となるもの。
\item すべての $x \in X$ であって、$Y \to X$ の $x$ 上のファイバーが有限であり、かつ
$\mathcal{O}_{X, x}$ が正規であるもの。
\item すべての $x \in X$ であって、$f$ がある $y \in Y$、すなわち $x$ 上にある点において準有限であり、
$\mathcal{O}_{X, x}$ が正規であるもの。
\end{enumerate}
\end{lemma}

\begin{proof}
(1) は Morphisms, Lemma \ref{morphisms-lemma-birational-isomorphism-over-dense-open} から従う。
(2) は (3) と Lemma \ref{varieties-lemma-finite-in-codim-1} から従う
（また、有限射は有限ファイバーをもつという事実を用いる）。

\medskip\noindent
(3) は (4) と Morphisms, Lemma \ref{morphisms-lemma-finite-fibre} から従うが、直接の証明も与える。
$x \in X$ を (3) のとおりとする。Cohomology of Schemes, Lemma
\ref{coherent-lemma-proper-finite-fibre-finite-in-neighbourhood} により、$f$ は有限であると仮定してよい。
$X$ はアフィンであると仮定してよい。これにより、Noether アフィンスキームの有限双有理射
$Y \to X$ と点 $x \in X$ であって $\mathcal{O}_{X, x}$ が正規整域となる場合に帰着される。
$\mathcal{O}_{X, x}$ は整域であり $X$ は Noether なので、$X$ を $x$ のアフィン開近傍であって整であるものに
置き換えてよい。すると $Y \to X$ は双有理であり $Y$ は被約なので、$Y$ は整であることが分かる。
$X = \Spec(A)$ および $Y = \Spec(B)$ と書けば、$A \subset B$ は同じ分数体をもつ整域の有限包含である。
$\mathfrak p \subset A$ を $x$ に対応する素イデアルとすれば、$A_\mathfrak p$ が正規であることから
$A_\mathfrak p \subset B_\mathfrak p$ は等号となる。$B$ は有限 $A$-加群なので、ある
$a \in A$, $a \not \in \mathfrak p$ が存在して $A_a \to B_a$ は同型となる。

\medskip\noindent
$x \in X$ と $y \in Y$ を (4) のとおりとする。$X$ をアフィン開近傍で置き換えた後、
$X = \Spec(A)$ および $A \subset \mathcal{O}_{X, x}$ と仮定してよい。Properties, Lemma
\ref{properties-lemma-ring-affine-open-injective-local-ring} を参照せよ。すると $A$ は整域であり、従って
$X$ は整である。$f$ は双有理であり $Y$ は被約なので、$Y$ も整である。環準同型
$\mathcal{O}_{X, x} \to \mathcal{O}_{Y, y}$ を考える。これは本質的有限型の環準同型であり、剰余体拡大は
有限であり、$\dim(\mathcal{O}_{Y, y}/\mathfrak m_x\mathcal{O}_{Y, y}) = 0$ である
（これを見るには、Morphisms, Definition \ref{morphisms-definition-quasi-finite} および Algebra, Definition
\ref{algebra-definition-quasi-finite} における準有限写像の定義をたどればよい）。Algebra, Lemma
\ref{algebra-lemma-essentially-finite-type-fibre-dim-zero} により、$\mathcal{O}_{Y, y}$ は有限
$\mathcal{O}_{X, x}$-代数 $B$ の局所化である。もちろん $B$ は、それ自身 $B$ を $\mathcal{O}_{Y, y}$ における像で
置き換えてよく、$B$ は $\mathcal{O}_{Y, y}$ と同じ分数体をもつ整域であると仮定してよい。
このとき $\mathcal{O}_{X, x} \subset B$ は $f$ が双有理なので同じ分数体をもつ。
$\mathcal{O}_{X, x}$ は正規なので、$\mathcal{O}_{X, x} = B$ と結論できる（有限ならば整だからである）。
特に $\mathcal{O}_{X, x} = \mathcal{O}_{Y, y}$ であることが分かる。Morphisms, Lemma
\ref{morphisms-lemma-morphism-defined-local-ring} により、$X$ を縮小した後、切断 $X \to Y$、すなわち $f$ の切断であって、
$x$ を $y$ に写し、局所環上で与えられた同型を誘導するものが存在すると仮定してよい。
$X \to Y$ は閉であり（Schemes, Lemma \ref{schemes-lemma-section-immersion} による）、必然的に $X$ の生成点を
$Y$ の生成点へ写すので、$X \to Y$ の像は $Y$ である。従って $Y = X$ であり、示すべきことを得た。
\end{proof}





\section{Noether 正規化の変種}
\label{varieties-section-noether-normalization}

\noindent
Noether 正規化とは、$k$ が体であり、$A$ が有限型 $k$-代数で、その次元が $d$ ならば、有限かつ単射な
$k$-代数準同型 $k[x_1, \ldots, x_d] \to A$ が存在するという主張である。
Algebra, Lemma \ref{algebra-lemma-Noether-normalization} を参照せよ。幾何学的には、これは
$\Spec(A) \to \mathbf{A}^d_k$ という $\Spec(k)$ 上の有限全射が存在することを意味する。

\begin{lemma}
\label{varieties-lemma-noether-normalization}
$f : X \to S$ をスキームの射とする。$x \in X$ とし、その像を $s \in S$ とする。
$V \subset S$ を $s$ のアフィン開近傍とする。$f$ が局所有限型であり $\dim_x(X_s) = d$ ならば、
アフィン開集合 $U \subset X$ であって $x \in U$ かつ $f(U) \subset V$ を満たすものと、分解
$$
U \xrightarrow{\pi} \mathbf{A}^d_V \to V
$$
であって $f|_U : U \to V$ の分解となり、$\pi$ が準有限であるものが存在する。
\end{lemma}

\begin{proof}
これは Algebra, Lemma \ref{algebra-lemma-quasi-finite-over-polynomial-algebra} から従う。
\end{proof}

\begin{lemma}
\label{varieties-lemma-noether-normalization-affine}
$f : X \to S$ をアフィンスキームの有限型射とする。$s \in S$ とする。$\dim(X_s) = d$ ならば、
分解
$$
X \xrightarrow{\pi} \mathbf{A}^d_S \to S
$$
であって $f$ の分解となり、ファイバーの射 $\pi_s : X_s \to \mathbf{A}^d_{\kappa(s)}$、すなわち
$s$ 上の射が有限であるものが存在する。
\end{lemma}

\begin{proof}
$S = \Spec(A)$ および $X = \Spec(B)$ と書き、$A \to B$ を $f$ に対応する環準同型とする。
$\mathfrak p \subset A$ を $s$ に対応する素イデアルとする。全射 $A[x_1, \ldots, x_r] \to B$ を選べる。
Algebra, Lemma \ref{algebra-lemma-Noether-normalization} により、元
$y_1, \ldots, y_d \in A$ であって、$\mathbf{Z}$-部分代数、すなわち $A$ の
$x_1, \ldots, x_r$ により生成される部分代数に属し、$A$-代数準同型
$A[t_1, \ldots, t_d] \to B$ で $t_i$ を $y_i$ に送るものが有限な $\kappa(\mathfrak p)$-代数準同型
$\kappa(\mathfrak p)[t_1, \ldots, t_d] \to B \otimes_A \kappa(\mathfrak p)$ を誘導するものが存在する。
これで補題が証明された。
\end{proof}

\begin{lemma}
\label{varieties-lemma-geometric-structure-unramified}
$f : X \to S$ をスキームの射とし、$x \in X$ とする。$V = \Spec(A)$ を $f(x)$ の $S$ における
アフィン開近傍とする。$f$ が $x$ において非分岐ならば、アフィン開 $U \subset X$ であって
$x \in U$ かつ $f(U) \subset V$ を満たし、次の可換図式をもつものが存在する。
$$
\xymatrix{
X \ar[d] & U \ar[l] \ar[rd] \ar[r]^-j &
\Spec(A[t]_{g'}/(g)) \ar[d] \ar[r] &
\Spec(A[t]) = \mathbf{A}^1_V \ar[ld] \\
Y & & V \ar[ll]
}
$$
ここで $j$ は埋め込み、$g \in A[t]$ はモニック多項式、$g'$ は $g$ の $t$ に関する導関数である。
$f$ が $x$ においてエタールならば、$j$ が開埋め込みとなるように図式を選べる。
\end{lemma}

\begin{proof}
非分岐の場合は Algebra, Proposition \ref{algebra-proposition-unramified-locally-standard} の翻訳である。
エタールの場合は Algebra, Proposition \ref{algebra-proposition-etale-locally-standard} の翻訳であるか、
同値に、主張にはわずかな相違があるものの Morphisms, Lemma
\ref{morphisms-lemma-etale-locally-standard-etale} から従う。
\end{proof}

\begin{lemma}
\label{varieties-lemma-unramfied-over-affine}
$f : X \to S$ をアフィンスキームの有限型射とする。$x \in X$ とし、その像を $s \in S$ とする。
$$
r =
\dim_{\kappa(x)} \Omega_{X/S, x} \otimes_{\mathcal{O}_{X, x}} \kappa(x) =
\dim_{\kappa(x)} \Omega_{X_s/s, x} \otimes_{\mathcal{O}_{X_s, x}} \kappa(x) =
\dim_{\kappa(x)} T_{X/S, x}
$$
とおく。このとき $f$ の分解
$$
X \xrightarrow{\pi} \mathbf{A}^r_S \to S
$$
であって $\pi$ が $x$ において非分岐となるものが存在する。
\end{lemma}

\begin{proof}
Morphisms, Lemma \ref{morphisms-lemma-finite-type-differentials} により第一の次元は有限である。
第一の等号は、$\Omega_{X/S}$ のファイバーへの制限が微分加群であることから従う。Morphisms, Lemma
\ref{morphisms-lemma-base-change-differentials} を参照せよ。最後の等号は Lemma
\ref{varieties-lemma-tangent-space-cotangent-space} から従う。従って主張が意味をもつことが分かる。

\medskip\noindent
補題を証明するため、$S = \Spec(A)$ および $X = \Spec(B)$ と書き、$A \to B$ を $f$ に対応する
環準同型とする。$\mathfrak q \subset B$ を $x$ に対応する素イデアルとする。$A$-代数の全射
$A[x_1, \ldots, x_t] \to B$ を選ぶ。$\Omega_{B/A}$ は $\text{d}x_1, \ldots, \text{d}x_t$ により
生成されるので、それらの $\Omega_{X/S, x} \otimes_{\mathcal{O}_{X, x}} \kappa(x)$ における像は、
これを $\kappa(x)$-ベクトル空間として生成する。番号を付け直した後、
$\text{d}x_1, \ldots, \text{d}x_r$ が $\Omega_{X/S, x} \otimes_{\mathcal{O}_{X, x}} \kappa(x)$ の基底に
写ると仮定してよい。$P = A[x_1, \ldots, x_r] \to B$ は $\mathfrak q$ において非分岐であると主張する。
これを見るには $\Omega_{B/P, \mathfrak q} = 0$ を示せば十分である
（Algebra, Lemma \ref{algebra-lemma-unramified}）。$\Omega_{B/P}$ は $\Omega_{B/A}$ を
$\text{d}x_1, \ldots, \text{d}x_r$ により生成される部分加群で割った商であることに注意せよ。従って
$\Omega_{B/P, \mathfrak q} \otimes_{B_\mathfrak q} \kappa(\mathfrak q) = 0$ である。
これは $x_1, \ldots, x_r$ の選び方による。
Nakayama の補題、より正確には Algebra, Lemma \ref{algebra-lemma-NAK} の (2) により、
$\Omega_{B/P}$ は有限なので（上の参照を見よ）この補題が適用でき、$\Omega_{B/P, \mathfrak q} = 0$ と結論する。
\end{proof}

\begin{lemma}
\label{varieties-lemma-immersion-into-affine}
$f : X \to S$ をスキームの射とする。$x \in X$ とし、その像を $s \in S$ とする。
$V \subset S$ を $s$ のアフィン開近傍とする。$f$ が局所有限型であり、
$$
r =
\dim_{\kappa(x)} \Omega_{X/S, x} \otimes_{\mathcal{O}_{X, x}} \kappa(x) =
\dim_{\kappa(x)} \Omega_{X_s/s, x} \otimes_{\mathcal{O}_{X_s, x}} \kappa(x) =
\dim_{\kappa(x)} T_{X/S, x}
$$
ならば、次が存在する。
\begin{enumerate}
\item アフィン開 $U \subset X$ であって $x \in U$ かつ $f(U) \subset V$ を満たすものと、分解
$$
U \xrightarrow{j} \mathbf{A}^{r + 1}_V \to V
$$
であって $f|_U$ の分解となり、$j$ が埋め込みであるもの、または
\item アフィン開 $U \subset X$ であって $x \in U$ かつ $f(U) \subset V$ を満たすものと、分解
$$
U \xrightarrow{j} D \to V
$$
であって $f|_U$ の分解となり、$j$ が閉埋め込みで、$D \to V$ が相対次元 $r$ の平滑射であるもの。
\end{enumerate}
\end{lemma}

\begin{proof}
任意のアフィン開 $U \subset X$ であって $x \in U$ かつ $f(U) \subset V$ を満たすものを選ぶ。
Lemma \ref{varieties-lemma-unramfied-over-affine} を $U \to V$ に適用し、その補題の主張にある
$U \to \mathbf{A}^r_V \to V$ を得る。Lemma \ref{varieties-lemma-geometric-structure-unramified} により分解
$$
U \xrightarrow{j} D \xrightarrow{j'} \mathbf{A}^{r + 1}_V
\xrightarrow{p} \mathbf{A}^r_V \to V
$$
を得る。ここで $j$ と $j'$ は埋め込み、$p$ は射影であり、$p \circ j'$ は標準エタールである。
従って特に (1) と (2) が成り立つことが分かる。
\end{proof}


\section{ファイバーの次元}
\label{varieties-section-dimension-fibres}

\noindent
有限型射のファイバーの次元は一般に跳び上がることをすでに見た。この節では、余次元 $1$ では
この現象が起こらないことを論じる。より一般には、ファイバーの次元がどれほど跳び得るかを論じる。
関連する結果を列挙すると次のとおりである。
\begin{enumerate}
\item 有限型射 $X \to S$ に対して、$x \in X$ のうち
$\dim_x(X_{f(x)}) \leq d$ を満たすものの集合は開である。Algebra, Lemma
\ref{algebra-lemma-dimension-fibres-bounded-open-upstairs} および
Morphisms, Lemma \ref{morphisms-lemma-openness-bounded-dimension-fibres} を参照せよ。
\item 次元公式が成り立つ。Algebra, Lemma \ref{algebra-lemma-dimension-formula} および
Morphisms, Lemma \ref{morphisms-lemma-dimension-formula} を参照せよ。
\item 付値環を支配する整有限型スキームではファイバーの次元は一定である。Algebra, Lemma
\ref{algebra-lemma-finite-type-domain-over-valuation-ring-dim-fibres} を参照せよ。
\item $X \to S$ が有限型で、$X$ のすべての生成点において準有限ならば、$X \to S$ は
余次元 $1$ で準有限である。Algebra, Lemma \ref{algebra-lemma-finite-in-codim-1} および
Lemma \ref{varieties-lemma-quasi-finite-in-codim-1} を参照せよ。
\end{enumerate}
上で最後に挙げた結果は次のように一般化される。

\begin{lemma}
\label{varieties-lemma-dimension-fibre-in-codim-1}
$f : X \to Y$ を局所有限型とする。$x \in X$ を点、その像を $y \in Y$ とし、
$\mathcal{O}_{Y, y}$ は次元 $\leq 1$ の Noether 環であるとする。$d \geq 0$ を整数とし、
既約成分の生成点 $\eta$ で、その既約成分が $X$ の $x$ を含むものすべてに対して
$\dim_\eta(X_{f(\eta)}) = d$ が成り立つとする。このとき $\dim_x(X_y) = d$ である。
\end{lemma}

\begin{proof}
$\Spec(\mathcal{O}_{Y, y})$ は $Y$ の点で $y$ へ特殊化するもの全体の集合であることを想起しよう。
Schemes, Lemma \ref{schemes-lemma-specialize-points} を参照せよ。従って $Y$ を
$\Spec(\mathcal{O}_{Y, y})$ に置き換え、$Y = \Spec(B)$、ここで $B$ は次元 $\leq 1$ の
Noether 局所環であり、$y$ は閉点である、と仮定してよい。また $X$ を $x$ のアフィン近傍で
置き換えてよい。

\medskip\noindent
$X = \bigcup X_i$ を、被約閉部分スキームとみなした $X$ の既約成分の合併とする。
各ファイバー $X_{i, y}$ の次元が $d$ であることを示せれば、$X_y = \bigcup X_{i, y}$ の次元も
$d$ である。従って $X$ は整スキームであると仮定してよい。

\medskip\noindent
$X \to Y$ が生成点 $\eta$（$X$ の生成点）を $y$ に写すならば、$X$ は $\kappa(y)$ 上のスキームであり、
主張は仮定により正しい。$X$ が $\eta$ を点 $\xi \in Y$ に写すと仮定する。この点は
極小素イデアル $\mathfrak q$ に対応し、それは $B$ の極小素イデアルであって
$\mathfrak m_B$ とは異なる。分解
$X \to \Spec(B/\mathfrak q) \to \Spec(B)$ を得る。次元公式
（Morphisms, Lemma \ref{morphisms-lemma-dimension-formula}）により
$$
\dim(\mathcal{O}_{X, x}) + \text{trdeg}_{\kappa(y)} \kappa(x) \leq
\dim(B/\mathfrak q) + \text{trdeg}_{\kappa(\mathfrak q)}(R(X))
$$
が成り立つ。$\dim(B/\mathfrak q) = 1$ である。また
$\text{trdeg}_{\kappa(\mathfrak q)}(R(X)) = d$ である。これは仮定
$\dim_\eta(X_\xi) = d$ と Morphisms, Lemma
\ref{morphisms-lemma-dimension-fibre-at-a-point} による。
$\mathcal{O}_{X, x} \to \mathcal{O}_{X_s, x}$ は核をもち
（$\eta \mapsto \xi \not = y$ だからである）、かつ $\mathcal{O}_{X, x}$ は Noether 整域なので、
$\dim(\mathcal{O}_{X, x}) > \dim(\mathcal{O}_{X_y, x})$ であることが分かる。従って
$$
\dim_x(X_s) =
\dim(\mathcal{O}_{X_s, x}) + \text{trdeg}_{\kappa(y)} \kappa(x) \leq d
$$
と結論する（Morphisms, Lemma \ref{morphisms-lemma-dimension-fibre-at-a-point}）。一方、
$\dim_x(X_s) \geq \dim_\eta(X_{f(\eta)}) = d$ が Morphisms, Lemma
\ref{morphisms-lemma-openness-bounded-dimension-fibres} により成り立つ。
\end{proof}

\begin{lemma}
\label{varieties-lemma-dominate-valuation-ring-dimension-fibres}
$f : X \to \Spec(R)$ を、既約スキームから付値環のスペクトルへの射とする。$f$ が局所有限型かつ
全射ならば、特殊ファイバーは等次元であり、その次元は一般ファイバーの次元に等しい。
\end{lemma}

\begin{proof}
$X$ をその被約化で置き換えてよい。これは $X$ の次元も特殊ファイバーの次元も変えないからである。
すると $X$ は整であり、補題は Algebra, Lemma
\ref{algebra-lemma-finite-type-domain-over-valuation-ring-dim-fibres} から従う。
\end{proof}

\noindent
次の補題は Lemma \ref{varieties-lemma-dimension-fibre-in-codim-1} を一般化する。

\begin{lemma}
\label{varieties-lemma-dimension-fibre-in-higher-codimension}
$f : X \to Y$ を局所有限型とする。$x \in X$ を点、その像を $y \in Y$ とし、
$\mathcal{O}_{Y, y}$ は Noether 環であるとする。$d \geq 0$ を整数とし、既約成分の生成点
$\eta$ で、その既約成分が $X$ の $x$ を含むものすべてに対して $f(\eta) \not = y$ かつ
$\dim_\eta(X_{f(\eta)}) = d$ が成り立つとする。このとき
$\dim_x(X_y) \leq d + \dim(\mathcal{O}_{Y, y}) - 1$ である。
\end{lemma}

\begin{proof}
Lemma \ref{varieties-lemma-dimension-fibre-in-codim-1} の証明とまったく同様に、$X = \Spec(A)$ で
$A$ は整域、$Y = \Spec(B)$ で $B$ は極大イデアルが $y$ に対応する Noether 局所環である場合に
帰着する。$B$ を $B/\Ker(B \to A)$ で置き換えた後、$B$ は整域で $B \subset A$ であると
仮定してよい。次元公式（Morphisms, Lemma \ref{morphisms-lemma-dimension-formula}）を用いると
$$
\dim(\mathcal{O}_{X, x}) + \text{trdeg}_{\kappa(y)} \kappa(x) \leq
\dim(B) + \text{trdeg}_B(A)
$$
を得る。$\text{trdeg}_B(A) = d$ である。これは仮定 $\dim_\eta(X_\xi) = d$ と
Morphisms, Lemma \ref{morphisms-lemma-dimension-fibre-at-a-point} による。
$\mathcal{O}_{X, x} \to \mathcal{O}_{X_y, x}$ は核をもち（$f(\eta) \not = y$ だからである）、
かつ $\mathcal{O}_{X, x}$ は Noether 整域なので、
$\dim(\mathcal{O}_{X, x}) > \dim(\mathcal{O}_{X_y, x})$ であることが分かる。従って
$$
\dim_x(X_y) =
\dim(\mathcal{O}_{X_y, x}) + \text{trdeg}_{\kappa(y)} \kappa(x)
< \dim(B) + d
$$
と結論する（等号は Morphisms, Lemma
\ref{morphisms-lemma-dimension-fibre-at-a-point} による）。これで示すべきことが証明された。
\end{proof}

\section{代数的スキーム}
\label{varieties-section-algebraic-schemes}

\noindent
次の定義は \cite[I Definition 6.4.1]{EGA} から採ったものである。

\begin{definition}
\label{varieties-definition-algebraic-scheme}
$k$ を体とする。{\it 代数的 $k$-スキーム}とは、スキーム $X$ であって $k$ 上のもの、かつ構造射
$X \to \Spec(k)$ が有限型であるものをいう。{\it 局所代数的 $k$-スキーム}とは、スキーム $X$ であって
$k$ 上のもの、かつ構造射 $X \to \Spec(k)$ が局所有限型であるものをいう。
\end{definition}

\noindent
任意の（局所）代数的 $k$-スキームは（局所）Noether であることに注意せよ。Morphisms, Lemma
\ref{morphisms-lemma-finite-type-noetherian} を参照せよ。代数的 $k$-スキームの圏は、$k$ 上の多様体の圏とは
異なり、すべての積とファイバー積をもつ。局所代数的 $k$-スキームの圏についても同様である。

\begin{lemma}
\label{varieties-lemma-algebraic-scheme-dim-0}
$k$ を体とする。$X$ を局所代数的 $k$-スキームで次元 $0$ のものとする。このとき $X$ は、
局所 Artin $k$-代数 $A$ で $\dim_k(A) < \infty$ を満たすもののスペクトルの非交和である。
$X$ が代数的 $k$-スキームで次元 $0$ ならば、さらに $X$ はアフィンであり、射
$X \to \Spec(k)$ は有限である。
\end{lemma}

\begin{proof}
$X$ を局所代数的 $k$-スキームで次元 $0$ のものとする。$U = \Spec(A) \subset X$ をアフィン開部分スキームとする。
$\dim(X) = 0$ なので $\dim(A) = 0$ である。Noether 正規化、Algebra, Lemma
\ref{algebra-lemma-Noether-normalization} により、有限な単射 $k \to A$ が存在する、すなわち
$\dim_k(A) < \infty$ である。従って $A$ は Artin 環である。Algebra, Lemma
\ref{algebra-lemma-finite-dimensional-algebra} を参照せよ。これより
$A = A_1 \times \ldots \times A_r$ は有限個の Artin 局所環の積である。Algebra, Lemma
\ref{algebra-lemma-artinian-finite-length} を参照せよ。もちろん $\dim_k(A_i) < \infty$ が各 $i$ に対して
成り立つ。これらの次元の和は $\dim_k(A)$ に等しいからである。

\medskip\noindent
上の議論は、$X$ が有限離散位相空間からなる開被覆をもつことを示す。従って $X$ は離散位相空間である。
ゆえに $X$ は、その連結成分の非交和と同型であり、各連結成分は一点集合である。一点からなるスキームは
アフィンなので、上の段落の結果により、これらの各一点は局所 Artin $k$-代数 $A$ で
$\dim_k(A) < \infty$ を満たすもののスペクトルであると結論する。

\medskip\noindent
最後に、$X$ が代数的 $k$-スキームで次元 $0$ ならば、$X$ は準コンパクトであり、従って有限非交和
$X = \Spec(A_1) \amalg \ldots \amalg \Spec(A_r)$ である。それゆえアフィンであり
（Schemes, Lemma \ref{schemes-lemma-disjoint-union-affines} を参照せよ）、$X \to \Spec(k)$ の有限性は
証明の最初の段落ですでに見た。
\end{proof}

\noindent
次の補題は、局所代数的スキームに関する次元論のいくつかの主張をまとめる。

\begin{lemma}
\label{varieties-lemma-dimension-locally-algebraic}
$k$ を体とし、$X$ を局所代数的 $k$-スキームとする。
\begin{enumerate}
\item
\label{varieties-item-catenary}
$X$ の位相空間は鎖状である
（Topology, Definition \ref{topology-definition-catenary}）。
\item
\label{varieties-item-dimension-at-closed-point}
閉点 $x \in X$ に対して、$\dim_x(X) = \dim(\mathcal{O}_{X, x})$ である。
\item
\label{varieties-item-dimension-irreducible}
$X$ が既約ならば、$\dim(X) = \dim(U)$ が任意の空でない開集合 $U \subset X$ に対して成り立ち、
$\dim(X) = \dim_x(X)$ が任意の $x \in X$ に対して成り立つ。
\item
\label{varieties-item-irreducible-maximal-chains}
$X$ が既約ならば、既約閉部分集合の任意の鎖は極大鎖に延長でき、既約閉部分集合のすべての極大鎖の長さは
$\dim(X)$ に等しい。
\item
\label{varieties-item-dimension-irreducibles-passing-through}
$x \in X$ に対して
$\dim_x(X) = \max \dim(Z) = \min \dim(\mathcal{O}_{X, x'})$
である。ここで最大値は既約成分 $Z \subset X$ で $x$ を含むもの全体にわたり、最小値は
特殊化 $x \leadsto x'$ で $x'$ が $X$ の閉点であるもの全体にわたる。
\item
\label{varieties-item-dimension-irreducible-trdeg}
$X$ が既約で生成点が $x$ ならば、$\dim(X) = \text{trdeg}_k(\kappa(x))$ である。
\item
\label{varieties-item-immediate-specialization}
$x \leadsto x'$ が $X$ の点の直後の特殊化ならば、
$\text{trdeg}_k(\kappa(x)) = \text{trdeg}_k(\kappa(x')) + 1$ である。
\item
\label{varieties-item-dimension-sup-trdeg}
$X$ の次元は数 $\text{trdeg}_k(\kappa(x))$ の上限である。ここで $x$ は $X$ の既約成分の
生成点全体を動く。
\item
\label{varieties-item-specialization}
$x \leadsto x'$ が $X$ の点の非自明な特殊化ならば、
\begin{enumerate}
\item $\dim_x(X) \leq \dim_{x'}(X)$、
\item $\dim(\mathcal{O}_{X, x}) < \dim(\mathcal{O}_{X, x'})$、
\item $\text{trdeg}_k(\kappa(x)) > \text{trdeg}_k(\kappa(x'))$、および
\item 非自明な特殊化の任意の極大鎖
$x = x_0 \leadsto x_1 \leadsto \ldots \leadsto x_n = x'$ の長さは
$n = \text{trdeg}_k(\kappa(x)) - \text{trdeg}_k(\kappa(x'))$ である。
\end{enumerate}
\item
\label{varieties-item-dimension-formula}
$x \in X$ に対して
$\dim_x(X) = \text{trdeg}_k(\kappa(x)) + \dim(\mathcal{O}_{X, x})$ である。
\item
\label{varieties-item-immediate-specialization-local-ring}
$x \leadsto x'$ が $X$ の点の直後の特殊化であり、かつ $X$ が既約または等次元ならば、
$\dim(\mathcal{O}_{X, x'}) = \dim(\mathcal{O}_{X, x}) + 1$ である。
\end{enumerate}
\end{lemma}

\begin{proof}
先に証明したより一般的な結果に依拠する代わりに、これらの主張を有限型 $k$-代数に対する対応する主張へ
帰着し、可換代数の章の結果を引用する。

\medskip\noindent
(\ref{varieties-item-catenary}) の証明。Topology, Lemma \ref{topology-lemma-catenary} により、これは $X$ 上局所的である。
従って $X = \Spec(A)$、ここで $A$ は有限型 $k$-代数である、と仮定してよい。示すべきことは、
$A$ が鎖状であることである（Algebra, Lemma \ref{algebra-lemma-catenary}）。Algebra, Lemma
\ref{algebra-lemma-quotient-catenary} を用いて $k[x_1, \ldots, x_n]$ に帰着し、次いで Algebra, Lemma
\ref{algebra-lemma-dimension-height-polynomial-ring} を適用できる。別の説明として、$k$ は自明に
Cohen--Macaulay であり、Cohen--Macaulay 環は普遍鎖状である
（Algebra, Lemma \ref{algebra-lemma-CM-ring-catenary}）ことからも従う。

\medskip\noindent
(\ref{varieties-item-dimension-at-closed-point}) の証明。アフィン近傍 $U = \Spec(A)$ で $x$ を含むものを選ぶ。このとき
$\dim_x(X) = \dim_x(U)$ である。従ってアフィンの場合に帰着し、これは Algebra, Lemma
\ref{algebra-lemma-dimension-closed-point-finite-type-field} である。

\medskip\noindent
(\ref{varieties-item-dimension-irreducible}) の証明。空でない任意の二つのアフィン開集合 $U, U' \subset X$ が
同じ次元をもつことを示せば十分である（既約部分集合の任意の有限鎖はアフィン開集合と交わる）。
閉点 $x$ で $X$ の点であり、かつ $x \in U \cap U'$ を満たすものを選ぶ。これは可能である。実際、$X$ は既約なので
$U \cap U'$ は空でなく、さらに Lemma \ref{varieties-lemma-locally-finite-type-Jacobson} により $X$ は Jacobson
だから、そのような閉点が存在する。このとき Algebra, Lemma
\ref{algebra-lemma-dimension-spell-it-out} により
$\dim(U) = \dim(\mathcal{O}_{X, x}) = \dim(U')$ である
（厳密には、補題を適用する前に $X$ をその被約化で置き換える必要がある）。

\medskip\noindent
(\ref{varieties-item-irreducible-maximal-chains}) の証明。既約閉部分集合の鎖が与えられたとき、その最小のものと交わる
アフィン開集合 $U \subset X$ を取れる。従って主張は Algebra, Lemma
\ref{algebra-lemma-dimension-spell-it-out} と、(\ref{varieties-item-dimension-irreducible}) で見た
$\dim(U) = \dim(X)$ から従う。

\medskip\noindent
(\ref{varieties-item-dimension-irreducibles-passing-through}) の証明。アフィン近傍
$U = \Spec(A)$ で $x$ を含むものを選ぶ。このとき $\dim_x(X) = \dim_x(U)$ である。対応
$Z \mapsto Z \cap U$ は、$X$ の既約成分で $x$ を通るものと、$U$ の既約成分で $x$ を通るものとの間の
全単射である。また、(\ref{varieties-item-dimension-irreducible}) により
$\dim(Z \cap U) = \dim(Z)$ がそのような $Z$ に対して成り立つ。従って主張は
Algebra, Lemma \ref{algebra-lemma-dimension-at-a-point-finite-type-over-field} から従う。

\medskip\noindent
(\ref{varieties-item-dimension-irreducible-trdeg}) の証明。(\ref{varieties-item-dimension-irreducible}) により、
$X = \Spec(A)$ がアフィンである場合に帰着する。この場合は $A_{red}$ に Algebra, Lemma
\ref{algebra-lemma-dimension-prime-polynomial-ring} を適用すれば従う。

\medskip\noindent
(\ref{varieties-item-immediate-specialization}) の証明。
$Z = \overline{\{x\}} \supset Z' = \overline{\{x'\}}$ とおく。(\ref{varieties-item-irreducible-maximal-chains}) により、
$Z \supset Z'$ は $Z$ の既約閉部分スキームの極大鎖の始まりであり、従って
$\dim(Z) = \dim(Z') + 1$ である。(\ref{varieties-item-dimension-irreducible-trdeg}) により結論を得る。

\medskip\noindent
(\ref{varieties-item-dimension-sup-trdeg}) の証明。単純な位相的議論から
$\dim(X) = \sup \dim(Z)$ が従う。ここで上限は $X$ の既約成分全体にわたる
（ヒント：Topology, Lemma \ref{topology-lemma-irreducible} を用いよ）。従って
(\ref{varieties-item-dimension-irreducible-trdeg}) から主張が従う。

\medskip\noindent
(\ref{varieties-item-specialization}) の証明。(a) は、任意の開集合 $U \subset X$ で $x'$ を含むものが $x$ も含むことから従う。
(b) は、$\mathcal{O}_{X, x}$ が $\mathcal{O}_{X, x'}$ の局所化であり、従って
$\mathcal{O}_{X, x}$ の任意の素イデアル鎖が $\mathcal{O}_{X, x'}$ の素イデアル鎖に対応し、末尾に
$\mathfrak m_{x'}$ を加えて延長できることから従う。(c) と (d) はともに
(\ref{varieties-item-immediate-specialization}) から形式的に従う。

\medskip\noindent
(\ref{varieties-item-dimension-formula}) の証明。アフィン近傍 $U = \Spec(A)$ で $x$ を含むものを選ぶ。このとき
$\dim_x(X) = \dim_x(U)$ である。従ってアフィンの場合に帰着し、これは Algebra, Lemma
\ref{algebra-lemma-dimension-at-a-point-finite-type-field} である。

\medskip\noindent
(\ref{varieties-item-immediate-specialization-local-ring}) の証明。$X$ が等次元ならば
（Topology, Definition \ref{topology-definition-equidimensional}）、$\dim(X)$ は $X$ の各既約成分の次元に
等しい。従って (\ref{varieties-item-dimension-irreducibles-passing-through}) により
$\dim_x(X) = \dim(X) = \dim_{x'}(X)$ である。ゆえに
(\ref{varieties-item-immediate-specialization}) から主張が従う。
\end{proof}

\begin{lemma}
\label{varieties-lemma-dimension-fibres-locally-algebraic}
$k$ を体とする。$f : X \to Y$ を局所代数的 $k$-スキームの射とする。
\begin{enumerate}
\item $y \in Y$ に対して、ファイバー $X_y$ は $\kappa(y)$ 上の局所代数的スキームであり、従って
Lemma \ref{varieties-lemma-dimension-locally-algebraic} のすべての結果を適用できる。
\item $X$ は既約であると仮定する。$Z = \overline{f(X)}$ および
$d = \dim(X) - \dim(Z)$ とおく。このとき
\begin{enumerate}
\item $\dim_x(X_{f(x)}) \geq d$ がすべての $x \in X$ に対して成り立つ、
\item $x \in X$ のうち $\dim_x(X_{f(x)}) = d$ を満たすものの集合は稠密開、
\item $\dim(\mathcal{O}_{Z, f(x)}) \geq 1$ ならば
$\dim_x(X_{f(x)}) \leq d + \dim(\mathcal{O}_{Z, f(x)}) - 1$、
\item $\dim(\mathcal{O}_{Z, f(x)}) = 1$ ならば $\dim_x(X_{f(x)}) = d$。
\end{enumerate}
\item $x \in X$ と $y = f(x)$ に対して
$\dim_x(X_y) \geq \dim_x(X) - \dim_y(Y)$ である。
\end{enumerate}
\end{lemma}

\begin{proof}
Morphisms, Lemma \ref{morphisms-lemma-permanence-finite-type} により射 $f$ は局所有限型である。
従って基底変換 $X_y \to \Spec(\kappa(y))$ は局所有限型である。これで (1) が示された。
証明の残りでは、$X$、$Y$ および $f$ のファイバーに対して Lemma
\ref{varieties-lemma-dimension-locally-algebraic} の結果を自由に用いる。

\medskip\noindent
(2) の証明。$\eta \in X$ を生成点とし、$\xi = f(\eta)$ とおく。このとき
$Z = \overline{\{\xi\}}$ である。従って
$$
d = \dim(X) - \dim(Z) =
\text{trdeg}_k \kappa(\eta) - \text{trdeg}_k \kappa(\xi) =
\text{trdeg}_{\kappa(\xi)} \kappa(\eta) = \dim_\eta(X_\xi)
$$
である。従って (2)(a) と (2)(b) は Morphisms, Lemma
\ref{morphisms-lemma-openness-bounded-dimension-fibres} から従う。(2)(c) と (2)(d) は Lemmas
\ref{varieties-lemma-dimension-fibre-in-higher-codimension} および
\ref{varieties-lemma-dimension-fibre-in-codim-1} から従う。

\medskip\noindent
(3) の証明。$x \in X$ とする。$X' \subset X$ を、$X$ の既約成分で $x$ を通り次元が $\dim_x(X)$ であるものとする。
このとき (2) により $\dim_x(X_y) \geq \dim(X') - \dim(Z')$ である。ここで $Z' \subset Y$ は
$X'$ の像の閉包である。これで (3) が示された。
\end{proof}

\begin{lemma}
\label{varieties-lemma-dimension-product-locally-algebraic}
\begin{slogan}
積の次元は各因子の次元の和である。
\end{slogan}
$k$ を体とし、$X$、$Y$ を局所代数的 $k$-スキームとする。
\begin{enumerate}
\item 点 $z \in X \times Y$ で $(x, y)$ 上にあるものに対して
$\dim_z(X \times Y) = \dim_x(X) + \dim_y(Y)$ である。
\item $\dim(X \times Y) = \dim(X) + \dim(Y)$ である。
\end{enumerate}
\end{lemma}

\begin{proof}
(1) の証明。構造射の分解
$$
X \times Y \longrightarrow Y \longrightarrow \Spec(k)
$$
を考える。第一の射 $p : X \times Y \to Y$ は、平坦射 $X \to \Spec(k)$ の基底変換なので平坦である。
Morphisms, Lemma \ref{morphisms-lemma-base-change-flat} を参照せよ。さらに Morphisms, Lemma
\ref{morphisms-lemma-dimension-fibre-after-base-change} により
$\dim_z(p^{-1}(y)) = \dim_x(X)$ である。従って Morphisms, Lemma
\ref{morphisms-lemma-dimension-fibre-at-a-point-additive} により
$\dim_z(X \times Y) = \dim_x(X) + \dim_y(Y)$ である。(2) は (1) の直接の帰結である。
\end{proof}

\section{完備局所環}
\label{varieties-section-complete-local-rings}

\noindent
体上のスキームの完備局所環に関するいくつかの結果を述べる。

\begin{lemma}
\label{varieties-lemma-complete-local-ring-structure-as-algebra}
$k$ を体とし、$X$ を $k$ 上の局所 Noether スキームとする。$x \in X$ を点とし、その剰余体を
$\kappa$ とする。剰余体上恒等写像を誘導する同型
\begin{equation}
\label{varieties-equation-complete-local-ring}
\kappa[[x_1, \ldots, x_n]]/I \longrightarrow \mathcal{O}_{X, x}^\wedge
\end{equation}
が存在する。一般には (\ref{varieties-equation-complete-local-ring}) を $k$-代数の同型として選ぶことはできない。
しかし拡大 $\kappa/k$ が分離的ならば、(\ref{varieties-equation-complete-local-ring}) を $k$-代数の同型として選べる。
\end{lemma}

\begin{proof}
この同型の存在は Cohen 構造定理の直接の帰結である\footnote{$\kappa$ の標数が $p$ ならば、
この定理が述べるのは、全射
$\Lambda[[x_1, \ldots, x_n]] \to \mathcal{O}_{X, x}^\wedge$ が得られるということにすぎない。
ここで $\Lambda$ は $\kappa$ の Cohen 環である。
しかしもちろん、この場合この写像は $\Lambda/p\Lambda[[x_1, \ldots, x_n]]$ を経由し、
$\Lambda/p\Lambda = \kappa$ である。}
（Algebra, Theorem \ref{algebra-theorem-cohen-structure-theorem}）。

\medskip\noindent
$p$ を奇素数、$k = \mathbf{F}_p(t)$ とし、$A = k[x, y]/(y^2 + x^p - t)$ とおく。
$A^\wedge$ を $A$ の極大イデアル $\mathfrak m = (y)$ における完備化とする。これは環として
$k(t^{1/p})[[z]]$ と同型だが、$k$-代数としては同型でない。その理由は、$A^\wedge$ が
$p$ 乗して $t$ となる元を含まないからである（読者は $y^2$ を法として計算すれば分かる）。
これはまた、(\ref{varieties-equation-complete-local-ring}) のいかなる同型も $k$-代数の同型にはなり得ないことを示す。

\medskip\noindent
$\kappa/k$ が分離的ならば、剰余体上恒等写像を誘導する $k$-代数準同型
$\kappa \to \mathcal{O}_{X, x}^\wedge$ が存在する。More on Algebra, Lemma
\ref{more-algebra-lemma-lift-residue-field} を参照せよ。生成元
$f_1, \ldots, f_n \in \mathfrak m_x$ を選ぶ。写像
$$
\kappa[[x_1, \ldots, x_n]] \longrightarrow \mathcal{O}_{X, x}^\wedge,\quad
x_i \longmapsto f_i
$$
を考える。両辺は $(x_1, \ldots, x_n)$-進完備であり（右辺については Algebra, Lemmas
\ref{algebra-lemma-hathat-finitely-generated} による）、この写像は構成により
$(x_1, \ldots, x_n)$ を法として全射なので、Algebra, Lemma
\ref{algebra-lemma-completion-generalities} により全射である。
\end{proof}

\begin{lemma}
\label{varieties-lemma-base-change-complete-local-ring}
$K/k$ を体の拡大とする。$X$ を局所代数的 $k$-スキームとし、$Y = X_K$ とおく。
$y \in Y$ を点、その像を $x \in X$ とする。
$\dim(\mathcal{O}_{X, x}) = \dim(\mathcal{O}_{Y, y})$ かつ $\kappa(x)/k$ が分離的であると仮定する。
(\ref{varieties-equation-complete-local-ring}) にあるような $k$-代数の同型
$$
\kappa(x)[[x_1, \ldots, x_n]]/(g_1, \ldots, g_m) \longrightarrow
\mathcal{O}_{X, x}^\wedge
$$
を選ぶ。このとき、(\ref{varieties-equation-complete-local-ring}) にあるような $K$-代数の同型
$$
\kappa(y)[[x_1, \ldots, x_n]]/(g_1, \ldots, g_m) \longrightarrow
\mathcal{O}_{Y, y}^\wedge
$$
が得られる。ここでは $\kappa(x) \to \kappa(y)$ により $g_j$ を $\kappa(y)$ 上の形式的べき級数とみなす。
\end{lemma}

\begin{proof}
局所環の写像 $\mathcal{O}_{X, x} \to \mathcal{O}_{Y, y}$ は局所環の写像
$\mathcal{O}_{X, x}^\wedge \to \mathcal{O}_{Y, y}^\wedge$ を誘導する。誘導された写像
$$
\kappa(x) \to \kappa(x)[[x_1, \ldots, x_n]]/(g_1, \ldots, g_m)
\to \mathcal{O}_{X, x}^\wedge \to \mathcal{O}_{Y, y}^\wedge
$$
と $\kappa(y)$ への射影との合成は、標準準同型 $\kappa(x) \to \kappa(y)$ である。
Lemma \ref{varieties-lemma-change-fields-flat} により、剰余体 $\kappa(y)$ は
$\kappa(x) \otimes_k K$ の、素イデアル $\mathfrak p_0$ における局所化である。ここでこの素イデアルは
$\kappa(x) \otimes_k K \to \kappa(y)$ の核である。一方、Lemma
\ref{varieties-lemma-change-fields-algebraic-unramified} により局所環
$(\kappa(x) \otimes_k K)_{\mathfrak p_0}$ は $\kappa(y)$ に等しい。従って写像
$$
\kappa(x) \otimes_k K \to \mathcal{O}_{Y, y}^\wedge
$$
は標準的に $\kappa(y)$ を経由する。可換図式
$$
\xymatrix{
\kappa(y) \ar[rr] & & \mathcal{O}_{Y, y}^\wedge \\
\kappa(x) \ar[r] \ar[u] &
\kappa(x)[[x_1, \ldots, x_n]]/(g_1, \ldots, g_m) \ar[r] &
\mathcal{O}_{X, x}^\wedge \ar[u]
}
$$
を得る。$f_i \in \mathfrak m_x^\wedge \subset \mathcal{O}_{X, x}^\wedge$ を $x_i$ の像とする。
写像が全射なので $\mathfrak  m_x^\wedge = (f_1, \ldots, f_n)$ であることに注意せよ。写像
$$
\kappa(y)[[x_1, \ldots, x_n]] \longrightarrow \mathcal{O}_{Y, y}^\wedge,\quad
x_i \longmapsto f_i
$$
を考える。ここで $f_i$ は実際には $f_i$ の $\mathfrak m_y^\wedge$ における像を意味する。
Lemma \ref{varieties-lemma-change-fields-algebraic-unramified} により
$\mathfrak m_x \mathcal{O}_{Y, y} = \mathfrak m_y$ なので、右辺は $(x_1, \ldots, x_n)$ に関して完備である
（それが完備局所環であることには Algebra, Lemma
\ref{algebra-lemma-hathat-finitely-generated} を用いよ）。両辺は $(x_1, \ldots, x_n)$-進完備なので、
写像は $(x_1, \ldots, x_n)$ を法として全射であることから Algebra, Lemma
\ref{algebra-lemma-completion-generalities} により全射である。もちろん形式的べき級数
$g_1, \ldots, g_m$ はこの写像で零に写る。実際、これらはすでに $\mathcal{O}_{X, x}^\wedge$ で零に写る。
従って可換図式
$$
\xymatrix{
\kappa(y)[[x_1, \ldots, x_n]]/(g_1, \ldots, g_m) \ar[r] &
\mathcal{O}_{Y, y}^\wedge \\
\kappa(x)[[x_1, \ldots, x_n]]/(g_1, \ldots, g_m) \ar[r] \ar[u] &
\mathcal{O}_{X, x}^\wedge \ar[u]
}
$$
を得る。なお上の水平矢印が同型であることを示す必要がある。これが全射であることはすでに分かっている。
$\mathcal{O}_{X, x} \to \mathcal{O}_{Y, y}$ は平坦である
（Lemma \ref{varieties-lemma-change-fields-flat}）。従って
$\mathcal{O}_{X, x}^\wedge \to \mathcal{O}_{Y, y}^\wedge$ は平坦である
（More on Algebra, Lemma \ref{more-algebra-lemma-flat-completion}）。ゆえに Algebra, Lemma
\ref{algebra-lemma-mod-injective} を、
$R = \kappa(x)[[x_1, \ldots, x_n]]/(g_1, \ldots, g_m)$、
$S = \kappa(y)[[x_1, \ldots, x_n]]/(g_1, \ldots, g_m)$、
$M = \mathcal{O}_{Y, y}^\wedge$、$N = S$ として適用し、写像が単射であると結論する。
\end{proof}

\section{大域生成}
\label{varieties-section-global-generation}

\noindent
準連接加群の大域生成に関するいくつかの補題を述べる。

\begin{lemma}
\label{varieties-lemma-globally-generated-base-change}
$X \to \Spec(A)$ をスキームの射とする。$A \subset A'$ を忠実平坦な環準同型とする。
$\mathcal{F}$ を準連接 $\mathcal{O}_X$-加群とする。このとき $\mathcal{F}$ が大域生成されることと、
基底変換 $\mathcal{F}_{A'}$ が大域生成されることとは同値である。
\end{lemma}

\begin{proof}
より正確には $X_{A'} = X \times_{\Spec(A)} \Spec(A')$ とおく。
$\mathcal{F}_{A'} = p^*\mathcal{F}$ とおく。ここで $p : X_{A'} \to X$ は射影である。Cohomology of Schemes, Lemma
\ref{coherent-lemma-flat-base-change-cohomology} により
$H^0(X_{k'}, \mathcal{F}_{A'}) = H^0(X, \mathcal{F}) \otimes_A A'$ である。従って $s_i$（$i \in I$）が
$H^0(X, \mathcal{F})$ の $A$-加群としての生成元ならば、$H^0(X_{A'}, \mathcal{F}_{A'})$ におけるそれらの像は
$H^0(X_{A'}, \mathcal{F}_{A'})$ の $A'$-加群としての生成元である。従って示すべきことは、写像
$\alpha : \bigoplus_{i \in I} \mathcal{O}_X \to \mathcal{F}$、
$(f_i) \mapsto \sum f_i s_i$ が全射であることと $p^*\alpha$ が全射であることとの同値性である。
これはアフィン開集合 $U = \Spec(B)$ で $X$ のものの上で調べればよい。このとき $\mathcal{F}|_U$ は
$B$-加群 $M$ に対応し、$s_i|_U$ は元 $x_i \in M$ に対応する。従って示すべきことは、
$\bigoplus_{i \in I} B \to M$ が全射であることと、基底変換
$\bigoplus_{i \in I} B \otimes_A A' \to M \otimes_A A'$ が全射であることとの同値性である。
これは $A \to A'$ が忠実平坦なので正しい。
\end{proof}

\begin{lemma}
\label{varieties-lemma-very-ample-vanish-at-point}
$k$ を無限体とする。$X$ を $k$ 上有限型のスキームとし、$\mathcal{L}$ を $X$ 上の非常に豊富な
可逆層とする。$n \geq 0$ とし、点 $x, x_1, \ldots, x_n \in X$ は、$x$ が $k$-有理点、すなわち
$\kappa(x) = k$ であり、$x \not = x_i$ が $i = 1, \ldots, n$ に対して成り立つとする。このとき
$s \in H^0(X, \mathcal{L})$ であって、$x$ で消えるが $x_i$ では消えないものが存在する。
\end{lemma}

\begin{proof}
$n = 0$ ならば結果は自明なので、$n > 0$ と仮定する。非常に豊富な可逆層の定義により、この補題は直ちに
$X = \mathbf{P}^r_k$ となるある $r > 0$ について、かつ $\mathcal{L} = \mathcal{O}_X(1)$ の場合に帰着する。
$\mathbf{P}^r_k = \text{Proj}(k[T_0, \ldots, T_r])$ と書き、
$V = H^0(X, \mathcal{L}) = kT_0 \oplus \ldots \oplus kT_r$ とおく。$x$ は $k$-有理点なので、
$s \in V$ で $x$ において消えるものの集合は余次元 $1$ の部分空間 $W \subset V$ であり、$W$ は $x$ に対応する
斉次素イデアルを生成する。$x_i \not = x$ なので、対応する斉次素イデアル
$\mathfrak p_i \subset k[T_0, \ldots, T_r]$ は $W$ を含まない。$k$ は無限なので
$W \not = \bigcup W \cap \mathfrak q_i$ となり、証明は完了する。
\end{proof}

\begin{lemma}
\label{varieties-lemma-generated-by-dim-plus-1-sections}
$k$ を無限体とする。$X$ を代数的 $k$-スキームとし、$\mathcal{L}$ を可逆 $\mathcal{O}_X$-加群とする。
$V \to \Gamma(X, \mathcal{L})$ を $k$-ベクトル空間の線形写像とし、その像は $\mathcal{L}$ を生成するとする。
このとき、部分空間 $W \subset V$ で $\dim_k(W) \leq \dim(X) + 1$ を満たし、$\mathcal{L}$ を生成するものが存在する。
\end{lemma}

\begin{proof}
証明を通じて、各 $x \in X$ に対する線形写像
$$
\psi_x : V \to \Gamma(X, \mathcal{L}) \to \mathcal{L}_x \to
\mathcal{L}_x \otimes_{\mathcal{O}_{X, x}} \kappa(x)
$$
が零でないことを用いる。$\dim(X)$ に関する帰納法で証明する。

\medskip\noindent
基底の場合は $\dim(X) = 0$ である。この場合 $X$ は有限個の点をもち、
$X = \{x_1, \ldots, x_n\}$ である（例えば Lemma \ref{varieties-lemma-algebraic-scheme-dim-0} を参照せよ）。
$k$ は無限なので、ベクトル $v \in V$ で $\psi_{x_i}(v) \not = 0$ がすべての $i$ に対して成り立つものが存在する。
このとき $W = k\cdot v$ が求めるものである。

\medskip\noindent
$\dim(X) > 0$ と仮定する。$X_i \subset X$ を、次元が $\dim(X)$ に等しい既約成分とする。
$X$ は Noether なので、そのようなものは有限個しかない。各 $i$ に対して点 $x_i \in X_i$ を選ぶ。
上と同様に、$v \in V$ を、$\psi_{x_i}(v) \not = 0$ がすべての $i$ に対して成り立つように選ぶ。
$Z \subset X$ を、$v$ の $\Gamma(X, \mathcal{L})$ における像の零点スキームとする。Divisors, Definition
\ref{divisors-definition-zero-scheme-s} を参照せよ。構成により $\dim(Z) < \dim(X)$ である。
帰納法により、$W \subset V$ で $\dim(W) \leq \dim(X)$ を満たし、$W$ が
$\mathcal{L}|_Z$ を生成するものを取れる。このとき $W + k\cdot v$ は $\mathcal{L}$ を生成する。
\end{proof}

\section{点と接ベクトルの分離}
\label{varieties-section-separating-points-tangent-vectors}

\noindent
これは次の結果にほかならない。

\begin{lemma}
\label{varieties-lemma-separate-points-tangent-vectors}
$k$ を代数閉体とする。$X$ を固有 $k$-スキームとし、$\mathcal{L}$ を可逆 $\mathcal{O}_X$-加群とする。
$V \subset H^0(X, \mathcal{L})$ を $k$-部分ベクトル空間とする。次を仮定する。
\begin{enumerate}
\item 相異なる閉点の各対 $x, y \in X$ に対して、切断 $s \in V$ で $x$ で消えるが $y$ では消えないものが存在し、
\item 各閉点 $x \in X$ と零でない接ベクトル $\theta \in T_{X/k, x}$ に対して、切断 $s \in V$ で
$x$ で消えるが $\theta$ による引き戻しが零でないものが存在する。
\end{enumerate}
このとき $\mathcal{L}$ は非常に豊富であり、標準射
$\varphi_{\mathcal{L}, V} : X \to \mathbf{P}(V)$ は閉埋め込みである。
\end{lemma}

\begin{proof}
条件 (1) は特に $V$ の元が $\mathcal{L}$ を $X$ 上で生成することを含意する。従って
Constructions, Example \ref{constructions-example-projective-space} により標準射
$$
\varphi = \varphi_{\mathcal{L}, V} :
X
\longrightarrow
\mathbf{P}(V)
$$
を得る。Morphisms, Lemma \ref{morphisms-lemma-image-proper-scheme-closed} により射 $\varphi$ は固有である。
(1) により写像 $\varphi$ は閉点上単射である（計算は省略する）。特に $\mathbf{P}(V)$ の任意の閉点上の
ファイバーは一点集合である（細部は省略する）。従って、例えば Cohomology of Schemes, Lemma
\ref{coherent-lemma-proper-finite-fibre-finite-in-neighbourhood} を用いると、$\varphi$ は有限である。
証明を終えるには、写像
$$
\varphi^\sharp :
\mathcal{O}_{\mathbf{P}(V)}
\longrightarrow
\varphi_*\mathcal{O}_X
$$
が全射であることを示せば十分である。これは閉点における茎上で調べればよい。$x \in X$ を閉点とし、
その像を閉点 $p = \varphi(x) \in \mathbf{P}(V)$ とする。(1) により
$\varphi^{-1}(\{p\}) = \{x\}$ であり、かつ $\varphi$ は固有、従って閉なので、
$\varphi^{-1}(U)$ は $x$ の開近傍の基本系を動くことが分かる。ここで $U$ は $p$ の開近傍の
基本系を動く。
従って $p$ における茎上で写像
$$
\varphi^\sharp_x :
\mathcal{O}_{\mathbf{P}(V), p}
\longrightarrow
\mathcal{O}_{X, x}
$$
を得る。特に $\mathcal{O}_{X, x}$ は有限 $\mathcal{O}_{\mathbf{P}(V), p}$-加群である。さらに、$x$ と
$p$ の剰余体はいずれも $k$ に等しい（$k$ は代数閉なので、Hilbert の零点定理を用いよ）。最後に条件 (2) は
写像
$$
T_{X/k, x} \longrightarrow T_{\mathbf{P}(V)/k, p}
$$
が単射であることを含意する。実際、この写像の核にある零でない $\theta$ は (2) の結論を満たし得ない。
上の局所環の写像で言い換えると、これは
$$
\mathfrak m_p/\mathfrak m_p^2 \longrightarrow
\mathfrak m_x/\mathfrak m_x^2
$$
が全射であることを意味する。Lemma \ref{varieties-lemma-tangent-space-rational-point} を参照せよ。Algebra, Lemma
\ref{algebra-lemma-when-surjective-local} を適用すれば証明が完了する。
\end{proof}

\begin{lemma}
\label{varieties-lemma-variant-separate-points-tangent-vectors}
$k$ を代数閉体とする。$X$ を固有 $k$-スキームとし、$\mathcal{L}$ を可逆 $\mathcal{O}_X$-加群とする。
各閉部分スキーム $Z \subset X$ で次元 $0$、次数 $2$、かつ $k$ 上のものに対して写像
$$
H^0(X, \mathcal{L}) \longrightarrow H^0(Z, \mathcal{L}|_Z)
$$
が全射であると仮定する。このとき $\mathcal{L}$ は $X$ 上で、$k$ に関して非常に豊富である。
\end{lemma}

\begin{proof}
これは Lemma \ref{varieties-lemma-separate-points-tangent-vectors} の言い換えである。実際、相異なる閉点
$x, y \in X$ が与えられたとき、$Z = x \cup y$ とおけば（閉部分スキームとみなす）、補題の条件 (1) を得る。
また零でない接ベクトル $\theta \in T_{X/k, x}$ が与えられたとき、射
$\theta : \Spec(k[\epsilon]) \to X$ は閉埋め込みである。$Z = \Im(\theta)$ とおけば補題の条件 (2) を得る。
\end{proof}

\section{積の閉包}
\label{varieties-section-closure-of-products}

\noindent
閉包と積との関係に関するいくつかの結果を述べる。

\begin{lemma}
\label{varieties-lemma-closure-of-product}
$k$ を体とする。$X$、$Y$ を $k$ 上のスキームとし、$A \subset X$、$B \subset Y$ を部分集合とする。
$$
AB =
\{z \in X \times_k Y \mid \text{pr}_X(z) \in A, \ \text{pr}_Y(z) \in B\}
\subset X \times_k Y
$$
とおく。このとき集合論的に
$$
\overline{A} \times_k \overline{B} = \overline{AB}
$$
が成り立つ。
\end{lemma}

\begin{proof}
包含 $\overline{AB} \subset \overline{A} \times_k \overline{B}$ は直ちに分かる。$X$ と $Y$ をそれぞれ
被約閉部分スキーム $\overline{A}$ と $\overline{B}$ で置き換えてよい。
$W \subset X \times_k Y$ を空でない開部分集合とする。Morphisms, Lemma
\ref{morphisms-lemma-scheme-over-field-universally-open} により、部分集合
$U = \text{pr}_X(W)$ は $X$ の空でない開部分集合である。従って $A \cap U$ は空でない。
$a \in A \cap U$ を選ぶ。ファイバーを $Y_a = \{a\} \times_k Y \cong Y_{\kappa(a)}$ と書く。
これは $\text{pr}_X : X \times_k Y \to X$ の $a$ 上のファイバーである。再び Morphisms, Lemma
\ref{morphisms-lemma-scheme-over-field-universally-open} により、射 $Y_a \to Y$ は開である。実際、
$\Spec(\kappa(a)) \to \Spec(k)$ は普遍開である。従って空でない開部分集合
$W_a = W \times_{X \times_k Y} Y_a$ は $Y$ の空でない開部分集合に写る。この像に入る
$b \in B$ が存在すると結論する。ゆえに望みどおり $AB \cap W \not = \emptyset$ である。
\end{proof}

\begin{lemma}
\label{varieties-lemma-closure-image-product-map}
$k$ を体とする。$f : A \to X$、$g : B \to Y$ を $k$ 上のスキームの射とする。このとき集合論的に
$$
\overline{f(A)} \times_k \overline{g(B)} =
\overline{(f \times g)(A \times_k B)}
$$
が成り立つ。
\end{lemma}

\begin{proof}
Lemma \ref{varieties-lemma-closure-of-product} の記法では $f \times g$ の像が $f(A)g(B)$ なので、同補題から従う。
\end{proof}

\begin{lemma}
\label{varieties-lemma-scheme-theoretic-image-product-map}
$k$ を体とする。$f : A \to X$、$g : B \to Y$ を $k$ 上のスキームの準コンパクトな射とする。
$Z \subset X$ を $f$ のスキーム論的像とする。Morphisms, Definition
\ref{morphisms-definition-scheme-theoretic-image} を参照せよ。同様に、$Z' \subset Y$ を $g$ の
スキーム論的像とする。このとき $Z \times_k Z'$ は $f \times g$ のスキーム論的像である。
\end{lemma}

\begin{proof}
$Z$ は $X$ の閉部分スキームのうち $f$ が経由する最小のものであることを想起せよ。$Z'$ についても同様である。
$W \subset X \times_k Y$ を $f \times g$ のスキーム論的像とする。$f \times g$ は
$Z \times_k Z'$ を経由するので、$W \subset Z \times_k Z'$ である。

\medskip\noindent
逆の包含を示すため、$U \subset X$ と $V \subset Y$ をアフィン開集合とする。Morphisms, Lemma
\ref{morphisms-lemma-quasi-compact-scheme-theoretic-image} により、スキーム $Z \cap U$ は
$f|_{f^{-1}(U)} : f^{-1}(U) \to U$ のスキーム論的像であり、$Z' \cap V$ および
$W \cap U \times_k V$ についても同様である。従って $X$ と $Y$ はアフィンであると仮定してよい。
$f$ と $g$ は準コンパクトなので、これは $A = \bigcup U_i$ が有限個のアフィンの合併であり、
$B = \bigcup V_j$ も有限個のアフィンの合併であることを含意する。そこで $A$ を $\coprod U_i$ で、
$B$ を $\coprod V_j$ で置き換えてよい。すなわち $A$ と $B$ もアフィンであると仮定してよい。
この場合、$Z$ は
$\Ker(\Gamma(X, \mathcal{O}_X) \to \Gamma(A, \mathcal{O}_A))$ で切り出され、$Z'$ と $W$ についても
同様である。従って、$A$ と $B$ はアフィンなので成り立つ等式
$$
\Gamma(A \times_k B, \mathcal{O}_{A \times_k B})
=
\Gamma(A, \mathcal{O}_A) \otimes_k \Gamma(B, \mathcal{O}_B)
$$
から結果が従う。細部は省略する。
\end{proof}

\section{体上平滑なスキーム}
\label{varieties-section-smooth}

\noindent
体上平滑なスキームを特徴づける二つの補題を述べる。

\begin{lemma}
\label{varieties-lemma-char-zero-differentials-free-smooth}
$k$ を体とし、$X$ を $k$ 上のスキームとする。次を仮定する。
\begin{enumerate}
\item $X$ は $k$ 上局所有限型である。
\item $\Omega_{X/k}$ は局所自由である。
\item $k$ の標数は零である。
\end{enumerate}
このとき構造射 $X \to \Spec(k)$ は平滑である。
\end{lemma}

\begin{proof}
Algebra, Lemma \ref{algebra-lemma-characteristic-zero-local-smooth} から従う。
\end{proof}

\noindent
正標数では、微分層が自由である被約でない有限型スキームが存在する。例えば
$\Spec(\mathbf{F}_p[t]/(t^p))$ over $\Spec(\mathbf{F}_p)$ である。基礎体 $k$ が標数 $p$ の
非完全体ならば、被約スキーム $X/k$ で $\Omega_{X/k}$ が自由だが平滑でないものが存在する。例えば
$\Spec(k[t]/(t^p-a)$ であり、ここで $a \in k$ は $p$ 乗ではない。

\begin{lemma}
\label{varieties-lemma-char-p-differentials-free-smooth}
$k$ を体とし、$X$ を $k$ 上のスキームとする。次を仮定する。
\begin{enumerate}
\item $X$ は $k$ 上局所有限型である。
\item $\Omega_{X/k}$ は局所自由である。
\item $X$ は被約である。
\item $k$ は完全である。
\end{enumerate}
このとき構造射 $X \to \Spec(k)$ は平滑である。
\end{lemma}

\begin{proof}
$x \in X$ を点とする。$X$ は局所 Noether なので（Morphisms, Lemma
\ref{morphisms-lemma-finite-type-noetherian} を参照せよ）、既約成分
$X_1, \ldots, X_n$ で $x$ を通るものは有限個である（Properties, Lemma
\ref{properties-lemma-Noetherian-topology} および Topology, Lemma
\ref{topology-lemma-Noetherian} を参照せよ）。$\eta_i \in X_i$ を生成点とする。$X$ は被約なので
$\mathcal{O}_{X, \eta_i} = \kappa(\eta_i)$ である。Algebra, Lemma
\ref{algebra-lemma-minimal-prime-reduced-ring} を参照せよ。さらに $\kappa(\eta_i)$ は完全体 $k$ の
有限生成体拡大なので、$k$ 上分離生成である（Algebra, Section
\ref{algebra-section-separability} を参照せよ）。従って
$\Omega_{X/k, \eta_i} = \Omega_{\kappa(\eta_i)/k}$ は、$\kappa(\eta_i)$ の $k$ 上の超越次数を階数とする
自由加群である。Morphisms, Lemma \ref{morphisms-lemma-dimension-fibre-at-a-point} により
$\dim_{\eta_i}(X_i) = \text{rank}_{\eta_i}(\Omega_{X/k})$ と結論する。
$x \in X_1 \cap \ldots \cap X_n$ なので
$$
\text{rank}_x(\Omega_{X/k}) = \text{rank}_{\eta_i}(\Omega_{X/k}) = \dim(X_i).
$$
従って Algebra, Lemma \ref{algebra-lemma-dimension-at-a-point-finite-type-over-field} により
$\dim_x(X) = \text{rank}_x(\Omega_{X/k})$ である。ゆえに、例えば Algebra, Lemma
\ref{algebra-lemma-characterize-smooth-over-field} により $X \to \Spec(k)$ は $x$ において平滑である。
\end{proof}

\begin{lemma}
\label{varieties-lemma-smooth-regular}
\begin{slogan}
体上平滑ならば正則である
\end{slogan}
$X \to \Spec(k)$ を平滑射とし、$k$ は体であるとする。このとき $X$ は正則スキームである。
\end{lemma}

\begin{proof}
（Lemma \ref{varieties-lemma-geometrically-regular-smooth} も参照せよ。）Algebra, Lemma
\ref{algebra-lemma-characterize-smooth-over-field} により、各局所環 $\mathcal{O}_{X, x}$ は正則である。
また $X$ は $k$ 上局所有限型なので局所 Noether である。従って Properties, Lemma
\ref{properties-lemma-characterize-regular} により $X$ は正則である。
\end{proof}

\begin{lemma}
\label{varieties-lemma-smooth-geometrically-normal}
$X \to \Spec(k)$ を平滑射とし、$k$ は体であるとする。このとき $X$ は $k$ 上、幾何学的正則、
幾何学的正規、かつ幾何学的被約である。
\end{lemma}

\begin{proof}
（Lemma \ref{varieties-lemma-geometrically-regular-smooth} も参照せよ。）$k'$ を $k$ の有限純非分離拡大とする。
$X_{k'}$ が正則、正規、被約であることを示せば十分である。Lemmas
\ref{varieties-lemma-geometrically-regular}、\ref{varieties-lemma-geometrically-normal}、および
\ref{varieties-lemma-check-only-finite-inseparable-extensions} を参照せよ。Morphisms, Lemma
\ref{morphisms-lemma-base-change-smooth} により射 $X_{k'} \to \Spec(k')$ も平滑である。従って、体上平滑な
スキーム $X$ が正則、正規、被約であることを示せば十分である。Lemma
\ref{varieties-lemma-smooth-regular} により $X$ は正則である。従って Properties, Lemma
\ref{properties-lemma-regular-normal} により $X$ は正規である。
\end{proof}

\begin{lemma}
\label{varieties-lemma-affine-space-over-field}
$k$ を体とし、$d \geq 0$ とする。$W \subset \mathbf{A}^d_k$ を空でない開集合とする。このとき閉点
$w \in W$ であって $k \subset \kappa(w)$ が有限分離的であるものが存在する。
\end{lemma}

\begin{proof}
必要なら $W$ を縮小して、$W = \mathbf{A}^d_k \setminus V(f)$ と仮定してよい。ここで
$f \in k[x_1, \ldots, x_d]$ である。補題が偽ならば、$f(a_1, \ldots, a_d) = 0$ がすべての
$(a_1, \ldots, a_d) \in (k^{sep})^d$ に対して成り立つ。
$k^{sep}$ は無限体なので、これは不合理である。
\end{proof}

\begin{lemma}
\label{varieties-lemma-smooth-separable-closed-points-dense}
$k$ を体とする。$X$ が $\Spec(k)$ 上平滑ならば、集合
$$
\{x \in X\text{ closed such that }k \subset \kappa(x)
\text{ is finite separable}\}
$$
は $X$ で稠密である。
\end{lemma}

\begin{proof}
空でない平滑な $X$ が $k$ 上与えられたとき、剰余体が $k$ 上有限分離的である閉点が少なくとも一つ
存在することを示せば十分である。このため、図式
$$
\xymatrix{
X & U \ar[l] \ar[r]^-\pi & \mathbf{A}^d_k
}
$$
で $\pi$ がエタールであるものを選ぶ。Morphisms, Lemma
\ref{morphisms-lemma-smooth-etale-over-affine-space} を参照せよ。射 $\pi : U \to \mathbf{A}^d_k$ は開である。
Morphisms, Lemma \ref{morphisms-lemma-etale-open} を参照せよ。Lemma
\ref{varieties-lemma-affine-space-over-field} により、閉点 $w \in \pi(U)$ で剰余体が $k$ 上有限分離的であるものを選べる。
任意の $x \in U$ で $\pi(x) = w$ を満たすものを選ぶ。Morphisms, Lemma
\ref{morphisms-lemma-etale-over-field} により体拡大 $\kappa(x)/\kappa(w)$ は有限分離的である。
従って $\kappa(x)/k$ は有限分離的である。点 $x$ は $X$ の閉点である。Morphisms, Lemma
\ref{morphisms-lemma-algebraic-residue-field-extension-closed-point-fibre} を参照せよ。
\end{proof}

\begin{lemma}
\label{varieties-lemma-geometrically-reduced-dense-smooth-open}
$X$ を、体 $k$ 上局所有限型な被約スキームとする。このとき、$X$ が $k$ 上幾何学的被約であることと、
$X$ が $k$ 上平滑な稠密開集合を含むこととは同値である。
\end{lemma}

\begin{proof}
問題は $X$ 上局所的なので、$X$ は準コンパクトであると仮定してよい。
$X = X_1 \cup \ldots \cup X_n$ を $X$ の既約成分とする。まず $X$ が $k$ 上平滑な稠密開集合を含むと
仮定する。このとき $X$ は各 $X_i$ の生成点において幾何学的被約である
（例えば Lemma \ref{varieties-lemma-smooth-geometrically-normal} による）。Lemma
\ref{varieties-lemma-generic-points-geometrically-reduced} により $X$ は幾何学的被約である。

\medskip\noindent
逆に $X$ が幾何学的被約であると仮定する。
$Z = \bigcup_{i \not = j} X_i \cap X_j$ は $X$ で至る所疎なので、$X$ を $X \setminus Z$ で置き換えてよい。
$X \setminus Z$ は既約スキームの非交和なので、$X$ が既約である場合に帰着する。$X$ は既約かつ被約なので
整である。Properties, Lemma \ref{properties-lemma-characterize-integral} を参照せよ。その生成点を
$\eta \in X$ とする。函数体
$K = k(X) = \kappa(\eta) = \mathcal{O}_{X, \eta}$ は $k$ 上幾何学的被約であり、従って $k$ 上分離的である。
Algebra, Lemma \ref{algebra-lemma-characterize-separable-field-extensions} を参照せよ。
$U = \Spec(A) \subset X$ を任意の空でないアフィン開集合とし、$K = A_{(0)}$ を $A$ の分数体とする。
Algebra, Lemma \ref{algebra-lemma-separable-smooth} を適用すると、$A$ は $(0)$ において $k$ 上平滑である。
定義により、これは $A$ のある主局所化が $k$ 上平滑であることを意味し、主張を得る。
\end{proof}

\begin{lemma}
\label{varieties-lemma-dense-smooth-open-variety-over-perfect-field}
$k$ を完全体とする。$X$ を局所代数的な被約 $k$-スキーム、例えば $k$ 上の多様体とする。このとき
$$
\{x \in X \mid X \to \Spec(k)\text{ is smooth at }x\} =
\{x \in X \mid \mathcal{O}_{X, x}\text{ is regular}\}
$$
であり、これは $X$ の稠密開部分スキームである。
\end{lemma}

\begin{proof}
二つの集合の等しさは Algebra, Lemma \ref{algebra-lemma-separable-smooth} と定義から直ちに従う
（完全体の定義については Algebra, Definition \ref{algebra-definition-perfect} を参照せよ）。この集合は開である。
スキームの射が平滑である点の集合は開だからである。Morphisms, Definition
\ref{morphisms-definition-smooth} を参照せよ。最後に、これが稠密であることを示す二つの議論を与える。
(1) $X$ の生成点はこの集合に入る。生成点における局所環は体であり（Algebra, Lemma
\ref{algebra-lemma-minimal-prime-reduced-ring}）、従って正則だからである。
(2) Lemma \ref{varieties-lemma-perfect-reduced} により $X$ は幾何学的被約であるから、Lemma
\ref{varieties-lemma-geometrically-reduced-dense-smooth-open} を適用する。
\end{proof}

\begin{lemma}
\label{varieties-lemma-flat-under-smooth}
$k$ を体とする。$f : X \to Y$ を、$k$ 上局所有限型なスキームの射とする。$x \in X$ を点とし、
$y = f(x)$ とおく。$X \to \Spec(k)$ が $x$ において平滑で、$f$ が $x$ において平坦ならば、
$Y \to \Spec(k)$ は $y$ において平滑である。特に $X$ が $k$ 上平滑で、$f$ が平坦かつ全射ならば、
$Y$ は $k$ 上平滑である。
\end{lemma}

\begin{proof}
$Y$ が $y$ において幾何学的正則であることを示せば十分である。Lemma
\ref{varieties-lemma-geometrically-regular-smooth} を参照せよ。これは Lemma
\ref{varieties-lemma-flat-under-geometrically-regular}（および $(X, x)$ に適用した Lemma
\ref{varieties-lemma-geometrically-regular-smooth}）から従う。
\end{proof}

\begin{lemma}
\label{varieties-lemma-variety-with-smooth-rational-point}
$k$ を体とする。$X$ を $k$ 上の多様体とし、$k$-有理点 $x$ をもち、$X$ は $x$ において平滑であるとする。
このとき $X$ は $k$ 上幾何学的整である。
\end{lemma}

\begin{proof}
$U \subset X$ を $X$ の平滑軌跡とする。仮定により $U$ は空でなく、従って稠密かつスキーム論的に稠密である。
すると $U_{\overline{k}} \subset X_{\overline{k}}$ も稠密かつスキーム論的に稠密である
（細部は省略する）。従って $U$ が幾何学的整であることを示せば十分である。$U$ は $k$-有理点をもつので、
Lemma \ref{varieties-lemma-geometrically-connected-if-connected-and-point} により幾何学的連結である。一方、
$U_{\overline{k}}$ は被約かつ正規である（Lemma \ref{varieties-lemma-smooth-geometrically-normal}）。
連結な正規 Noether スキームは整なので（Properties, Lemma
\ref{properties-lemma-normal-Noetherian}）、証明は完了する。
\end{proof}

\begin{lemma}
\label{varieties-lemma-smooth-stratification}
$X$ を体 $k$ 上有限型のスキームとする。有限純非分離拡大 $k'/k$、整数 $t \geq 0$、および閉部分スキーム
$$
X_{k'} \supset Z_0 \supset Z_1 \supset \ldots \supset Z_t = \emptyset
$$
であって、$Z_0 = (X_{k'})_{red}$ であり、$Z_i \setminus Z_{i + 1}$ が $k'$ 上平滑となることが
すべての $i$ に対して成り立つものが存在する。
\end{lemma}

\begin{proof}
$\dim(X)$ に関する帰納法を用いてよい。Lemma \ref{varieties-lemma-finite-extension-geometrically-reduced} により、
有限純非分離拡大 $k'/k$ であって $(X_{k'})_{red}$ が $k'$ 上幾何学的被約となるものを取れる。Lemma
\ref{varieties-lemma-geometrically-reduced-dense-smooth-open} により、至る所疎な閉部分スキーム
$X' \subset (X_{k'})_{red}$ であって $(X_{k'})_{red} \setminus X'$ が $k'$ 上平滑となるものが存在する。
このとき $\dim(X') < \dim(X)$ である。帰納法の仮定により、有限純非分離拡大 $k''/k'$、整数
$t' \geq 0$、および閉部分スキーム
$$
X'_{k''} \supset Y_0 \supset Y_1 \supset \ldots \supset Y_{t'} = \emptyset
$$
であって、$Y_0 = (X'_{k''})_{red}$ であり、$Y_i \setminus Y_{i + 1}$ が $k''$ 上平滑となることが
すべての $i$ に対して成り立つものが存在する。そこで $t = t' + 1$ とし、
$$
X_{k''} \supset Z_0 \supset Z_1 \supset \ldots \supset Z_t = \emptyset
$$
を考える。これは $Z_0 = (X_{k''})_{red}$ および $Z_i = Y_{i - 1}$（$i > 0$）で与えられる。
$X'_{k''}$ は $X_{k''}$ の閉部分スキームなので、これは意味をなす。述べたすべての性質の確認は省略する。
\end{proof}

\section{多様体の型}
\label{varieties-section-types}

\noindent
この短い節では、多様体のいくつかの初等的な大域的性質を論じる。

\begin{definition}
\label{varieties-definition-variety-type}
$k$ を体とし、$X$ を $k$ 上の多様体とする。
\begin{enumerate}
\item $X$ がアフィンスキームであるとき、$X$ を{\it アフィン多様体}という。これは、$X$ が
$\mathbf{A}^n_k$ の閉部分スキームと同型となるような $n$ が存在することと同値である。
\item $X$ を{\it 射影多様体}というのは、構造射 $X \to \Spec(k)$ が射影的であるときである。Morphisms, Lemma
\ref{morphisms-lemma-characterize-locally-projective} により、これは $X$ が $\mathbf{P}^n_k$ の
閉部分スキームと同型となるような $n$ が存在することと同値である。
\item $X$ を{\it 準射影多様体}というのは、構造射 $X \to \Spec(k)$ が準射影的であるときである。Morphisms, Lemma
\ref{morphisms-lemma-characterize-locally-quasi-projective} により、これは $X$ が $\mathbf{P}^n_k$ の
局所閉部分スキームと同型となるような $n$ が存在することと同値である。
\item 射 $X \to \Spec(k)$ が固有であるような多様体を{\it 固有多様体}という。
\item 射 $X \to \Spec(k)$ が平滑であるような多様体を{\it 平滑多様体}という。
\end{enumerate}
\end{definition}

\noindent
射影多様体は固有多様体であることに注意せよ。Morphisms, Lemma
\ref{morphisms-lemma-locally-projective-proper} を参照せよ。また $\mathbf{A}^n_k$ は
$\mathbf{P}^n_k$ の開部分スキームと同型なので、アフィン多様体は準射影的である。Constructions, Lemma
\ref{constructions-lemma-standard-covering-projective-space} を参照せよ。

\begin{lemma}
\label{varieties-lemma-regular-functions-proper-variety}
$X$ を $k$ 上の固有多様体とする。このとき
\begin{enumerate}
\item $K = H^0(X, \mathcal{O}_X)$ は体であり、体 $k$ の有限拡大である。
\item $X$ が幾何学的被約ならば $K/k$ は分離的である。
\item $X$ が幾何学的既約ならば $K/k$ は純非分離的である。
\item $X$ が幾何学的整ならば $K = k$ である。
\end{enumerate}
\end{lemma}

\begin{proof}
Lemma \ref{varieties-lemma-proper-geometrically-reduced-global-sections} の特別な場合である。
\end{proof}

\section{正規化}
\label{varieties-section-normalization}

\noindent
正規化に伴ういくつかの問題を扱う。

\begin{lemma}
\label{varieties-lemma-normalization-locally-algebraic}
$k$ を体とし、$X$ を $k$ 上の局所代数的スキームとする。$\nu : X^\nu \to X$ を正規化射とする。
Morphisms, Definition \ref{morphisms-definition-normalization} を参照せよ。このとき
\begin{enumerate}
\item $\nu$ は有限かつ支配的であり、$X^\nu$ は $k$ 上の正規既約局所代数的スキームの非交和である。
\item $\nu$ は $X^\nu \to X_{red} \to X$ と分解し、第一の射は $X_{red}$ の正規化射である。
\item $X$ が被約代数的スキームならば $\nu$ は双有理である。
\item $X$ が多様体ならば $X^\nu$ は多様体であり、$\nu$ は多様体の有限双有理射である。
\end{enumerate}
\end{lemma}

\begin{proof}
$X$ は体上局所有限型なので、$X$ は局所 Noether であり
（Morphisms, Lemma \ref{morphisms-lemma-finite-type-noetherian}）、従って各準コンパクト開集合は有限個の
既約成分しかもたない（Properties, Lemma
\ref{properties-lemma-Noetherian-irreducible-components}）。ゆえに Morphisms, Definition
\ref{morphisms-definition-normalization} を適用できる。正規化 $X^\nu$ は常に正規整スキームの非交和であり、
正規化射 $\nu$ は常に支配的である。Morphisms, Lemma
\ref{morphisms-lemma-normalization-normal} を参照せよ。$X$ は普遍 Nagata なので
（Morphisms, Lemma \ref{morphisms-lemma-ubiquity-nagata}）、$\nu$ は有限である
（Morphisms, Lemma \ref{morphisms-lemma-nagata-normalization}）。従って $X^\nu$ も局所代数的である。
これで (1) が証明された。

\medskip\noindent
(2) は Morphisms, Lemma \ref{morphisms-lemma-normalization-reduced} である。

\medskip\noindent
(3) は Morphisms, Lemma \ref{morphisms-lemma-normalization-birational} である。

\medskip\noindent
(4) は (1)、(2)、(3)、および $X^\nu$ が分離スキーム上有限なスキームとして分離的であることから従う。
\end{proof}

\begin{lemma}
\label{varieties-lemma-normalize-projective}
$k$ を体とし、$X$ を $k$ 上の固有スキームとする。$\nu : X^\nu \to X$ を正規化射とする。
Morphisms, Definition \ref{morphisms-definition-normalization} を参照せよ。このとき $X^\nu$ は $k$ 上固有である。
$X$ が $k$ 上射影的ならば $X^\nu$ は $k$ 上射影的である。
\end{lemma}

\begin{proof}
Lemma \ref{varieties-lemma-normalization-locally-algebraic} により射 $\nu$ は有限である。従って Morphisms, Lemmas
\ref{morphisms-lemma-finite-proper} および \ref{morphisms-lemma-composition-proper} により $X^\nu$ は $k$ 上固有である。
Morphisms, Lemma \ref{morphisms-lemma-finite-projective} により射 $\nu$ は射影的である。従って $X$ が $k$ 上
射影的ならば、Morphisms, Lemma \ref{morphisms-lemma-composition-projective} により $X^\nu$ は $k$ 上射影的である。
\end{proof}

\begin{lemma}
\label{varieties-lemma-relative-normalization-finite}
$k$ を体とする。$f : Y \to X$ を $k$ 上の局所代数的スキームの準コンパクトな射とする。
$X'$ を $X$ の $Y$ における正規化とする。$Y$ が被約ならば $X' \to X$ は有限である。
\end{lemma}

\begin{proof}
$Y$ は準分離的なので（Properties, Lemma \ref{properties-lemma-locally-Noetherian-quasi-separated} および
Morphisms, Lemma \ref{morphisms-lemma-finite-type-noetherian} により）、射 $f$ は準分離的である
（Schemes, Lemma \ref{schemes-lemma-compose-after-separated}）。従って Morphisms, Definition
\ref{morphisms-definition-normalization-X-in-Y} を適用できる。結果は Morphisms, Lemma
\ref{morphisms-lemma-nagata-normalization-finite-general} から従う。ここでは、局所代数的スキームが局所 Noether
（従って局所有限個の既約成分をもつ）であり、かつ Nagata であること
（Morphisms, Lemma \ref{morphisms-lemma-ubiquity-nagata}）を用いる。いくつかの細部は省略する。
\end{proof}

\begin{lemma}
\label{varieties-lemma-finite-extension-geometrically-normal}
$k$ を体とし、$X$ を代数的 $k$-スキームとする。このとき有限純非分離拡大 $k'/k$ であって、
正規化 $Y$ が $X_{k'}$ に対するもので、かつ $k'$ 上幾何学的正規となるものが存在する。
\end{lemma}

\begin{proof}
$K = k^{perf}$ を完全閉包とする。$Y_K$ を $X_K$ の正規化とする。Lemma
\ref{varieties-lemma-normalization-locally-algebraic} を参照せよ。Limits, Lemma
\ref{limits-lemma-descend-finite-presentation} により、有限部分拡大 $K/k'/k$ と有限表示の射
$\nu : Y \to X_{k'}$ が存在し、その $K$ への基底変換は正規化射 $\nu_K : Y_K \to X_K$ である。
$Y$ は $k'$ 上幾何学的正規であることに注意せよ（Lemma \ref{varieties-lemma-geometrically-normal}）。
$k'$ を大きくした後、$Y \to X_{k'}$ は有限であると仮定してよい
（Limits, Lemma \ref{limits-lemma-descend-finite-finite-presentation}）。
$\nu_K : Y_K \to X_K$ は正規化射なので、双有理射 $Y_K \to (X_K)_{red}$ を誘導する。従って稠密開集合
$V_K \subset X_K$ であって $\nu_K^{-1}(V_K) \to V_K$ が閉埋め込みとなるものが存在する
（これは $\nu_K^{-1}(V_K)$ と $V_{K, red}$ との同型を誘導する。例えば Morphisms, Lemma
\ref{morphisms-lemma-birational-isomorphism-over-dense-open} を参照せよ）。$k'$ を大きくした後、$V_K$ は
稠密開集合 $V \subset Y$ の基底変換であり、射 $\nu^{-1}(V) \to V$ は閉埋め込みであると仮定できる
（Limits, Lemmas \ref{limits-lemma-descend-opens} および
\ref{limits-lemma-descend-closed-immersion-finite-presentation}）。これより容易に $\nu$ が正規化射であることが従い、
証明は完了する。
\end{proof}

\begin{lemma}
\label{varieties-lemma-normalization-and-change-of-fields}
$k$ を体とし、$X$ を局所代数的 $k$-スキームとする。$K/k$ を体の拡大とする。
$\nu : X^\nu \to X$ を $X$ の正規化とし、$Y^\nu \to X_K$ を基底変換の正規化とする。このとき標準射
$$
Y^\nu \longrightarrow X^\nu \times_{\Spec(k)} \Spec(K)
$$
は $K/k$ が分離的ならば同型であり、一般には普遍同相である。
\end{lemma}

\begin{proof}
$Y = X_K$ とおく。$X^{(0)}$ および $Y^{(0)}$ を、それぞれ $X$ および $Y$ の既約成分の生成点の集合とする。
このとき射影射 $\pi : Y \to X$ は $\pi(Y^{(0)}) = X^{(0)}$ を満たす。これは $\pi$ が全射、開、かつ一般化を
保つからである。Morphisms, Lemmas \ref{morphisms-lemma-scheme-over-field-universally-open} および
\ref{morphisms-lemma-open-generizing} を参照せよ。$X^{(0)}$ および $Y^{(0)}$ をそれぞれ（被約）スキームと
みなすと、$X^\nu$ および $Y^\nu$ はそれぞれ $X$ および $Y$ の、$X^{(0)}$ および $Y^{(0})$ における
正規化である。従って Morphisms, Lemma \ref{morphisms-lemma-functoriality-normalization} は標準射
$Y^\nu \to X^\nu$ を与える。この射は $Y \to X$ 上にあり、さらにファイバー積の普遍性により補題の標準射を与える。

\medskip\noindent
この射が補題で述べた性質をもつことを示すには、$X = \Spec(A)$ がアフィンであると仮定してよい。
$Q(A_{red})$ を $A_{red}$ の全商環とする。このとき $X^\nu$ は、整閉包 $A'$、すなわち $A$ の
$Q(A_{red})$ における整閉包のスペクトルである。Morphisms, Lemmas \ref{morphisms-lemma-normalization-reduced} および
\ref{morphisms-lemma-description-normalization} を参照せよ。同様に $Y^\nu$ は、整閉包 $B'$、すなわち
$A \otimes_k K$ の $Q((A \otimes_k K)_{red})$ における整閉包のスペクトルである。標準写像
$Q(A_{red}) \to Q((A \otimes_k K)_{red})$ と標準写像 $A' \to B'$ があり、補題の射は誘導写像
$$
A' \otimes_k K \longrightarrow B'
$$
という $K$-代数の写像に対応する。その核は冪零元からなる。実際、
$Q(A_{red}) \otimes_k K \to Q((A \otimes_k K)_{red})$ の核は冪零元全体の集合である。

\medskip\noindent
$K/k$ が分離的ならば、Lemma \ref{varieties-lemma-base-change-normal-by-separable} により $A' \otimes_k K$ は正規である。
特にこれは被約であり、従って $Q((A \otimes_k K)_{red}) = Q(A' \otimes_k K)$ かつ
$B' = A' \otimes_k K$ である。Algebra, Lemma
\ref{algebra-lemma-characterize-reduced-ring-normal} を参照せよ。

\medskip\noindent
$K/k$ が分離的でないと仮定する。このとき $k$ の標数は $p > 0$ である。各 $b \in B'$ に対して、
ある $q$ が存在し、それが $p$ の冪であって $b^q$ が $A' \otimes_k K$ の像に入ることを示す。これにより表示した写像が
普遍同相であることが Algebra, Lemma \ref{algebra-lemma-p-ring-map} から従う。与えられた $b$ に対して、
部分体 $F \subset K$ であって $F/k$ が有限生成であり、$b$ が $Q((A \otimes_k F)_{red})$ に含まれ、
$A \otimes_k F$ 上整となるものが存在する。モニック多項式
$P = T^d + \alpha_1 T^{d - 1} + \ldots + \alpha _d$ であって $P(b) = 0$ かつ
$\alpha_i \in A \otimes_k F$ を満たすものを選ぶ。超越基底 $t_1, \ldots, t_r$ を $F$ の $k$ 上で選ぶ。
$F/F'/k(t_1, \ldots, t_r)$ を最大分離部分拡大とする
（Fields,  Lemma \ref{fields-lemma-separable-first}）。$F/F'$ は有限純非分離なので、ある $q$ が存在し、
$\lambda^q \in F'$ がすべての $\lambda \in F$ に対して成り立つ。このとき $b^q$ は
$Q((A \otimes_k F')_{red})$ に入り、多項式
$T^d + \alpha_1^q T^{d - 1} + \ldots + \alpha _d^q$ を満たし、その係数は
$\alpha_i^q \in A \otimes_k F'$ である。分離的な場合により
$b^q \in A' \otimes_k F'$ であり、証明は完了する。
\end{proof}

\begin{lemma}
\label{varieties-lemma-geometrically-normal-in-codim-1}
$k$ を体とし、$X$ を局所代数的 $k$-スキームとする。$\nu : X^\nu \to X$ を $X$ の正規化とする。
$x \in X$ を点とし、(a) $\mathcal{O}_{X, x}$ は被約、(b) $\dim(\mathcal{O}_{X, x}) = 1$、かつ (c)
すべての $x' \in X^\nu$ で $\nu(x') = x$ を満たすものに対して拡大 $\kappa(x')/k$ は分離的であるとする。
このとき $X$ は $x$ において幾何学的被約であり、$X^\nu$ は $x'$ において幾何学的正則である。ただし
$\nu(x') = x$ とする。
\end{lemma}

\begin{proof}
Lemma \ref{varieties-lemma-normalization-locally-algebraic} の結果を断りなく用いる。$x' \in X^\nu$ を $x$ 上の点とする。
次元論（Section \ref{varieties-section-algebraic-schemes}）により $\dim(\mathcal{O}_{X^\nu, x'}) = 1$ である。
$X^\nu$ は正規なので、$\mathcal{O}_{X^\nu, x'}$ は離散付値環である
（Properties, Lemma \ref{properties-lemma-criterion-normal}）。従って $\mathcal{O}_{X^\nu, x'}$ は正則局所
$k$-代数であり、その剰余体は $k$ 上分離的である。ゆえに $k \to \mathcal{O}_{X^\nu, x'}$ は
$\mathfrak m_{x'}$-進位相に関して形式的平滑である。More on Algebra, Lemma
\ref{more-algebra-lemma-regular-implies-fs} を参照せよ。すると More on Algebra, Theorem
\ref{more-algebra-theorem-regular-fs} により $\mathcal{O}_{X^\nu, x'}$ は $k$ 上幾何学的正則である。
従って Lemma \ref{varieties-lemma-geometrically-regular-at-point} により $X^\nu$ は $x'$ において幾何学的正則である。

\medskip\noindent
$\mathcal{O}_{X, x}$ は被約なので、写像の族 $\mathcal{O}_{X, x} \to \mathcal{O}_{X^\nu, x'}$ は単射である。
$\mathcal{O}_{X^\nu, x'}$ は幾何学的被約 $k$-代数なので、直ちに $\mathcal{O}_{X, x}$ は幾何学的被約
$k$-代数であることが従う。従って Lemma \ref{varieties-lemma-geometrically-reduced-at-point} により $X$ は $x$ において
幾何学的被約である。
\end{proof}

\section{可逆関数の群}
\label{varieties-section-units}

\noindent
$\mathcal{O}^*(X)/k^*$ は、$X$ が $k$ 上の多様体である場合、しばしば（ただし常にではない）
有限生成アーベル群となる。これを一連の補題によって示す。すべては次の特別な場合に基づく。

\begin{lemma}
\label{varieties-lemma-open-in-normal-proper}
$k$ を代数閉体とする。$\overline{X}$ を $k$ 上の固有多様体とし、
$X \subset \overline{X}$ を開部分スキームとする。$X$ は正規であると仮定する。
このとき $\mathcal{O}^*(X)/k^*$ は有限生成アーベル群である。
\end{lemma}

\begin{proof}
主張は $X$ のみに関するので、$\overline{X}$ を $k$ 上の別の固有多様体に置き換えてよい。
$\nu : \overline{X}^\nu \to
\overline{X}$ を正規化射とする。Lemma
\ref{varieties-lemma-normalization-locally-algebraic} により、$\nu$ は有限であり、$\overline{X}^\nu$ は多様体である。
$X$ は正規なので、$\nu^{-1}(X) \to X$ は同型である（ごく小さな細部は省略する）。最後に、
$\overline{X}^\nu$ は $k$ 上固有である。実際、有限射は固有であり
（Morphisms, Lemma \ref{morphisms-lemma-finite-proper}）、固有射の合成は固有である
（Morphisms, Lemma \ref{morphisms-lemma-composition-proper}）。従って、$\overline{X}$ は正規であると仮定してよく、
以下そう仮定する。

\medskip\noindent
任意のアフィン開集合 $U$（$\overline{X}$ の開集合）に対して、環 $\mathcal{O}(U)$ は有限生成
$k$-代数であり、Noether 環、整域、かつ正規であることを、以後断りなく用いる。Algebra, Lemma
\ref{algebra-lemma-Noetherian-permanence}、Properties, Definition
\ref{properties-definition-integral}、Properties, Lemmas
\ref{properties-lemma-locally-Noetherian} および
\ref{properties-lemma-locally-normal}、ならびに Morphisms, Lemma
\ref{morphisms-lemma-locally-finite-type-characterize} を参照せよ。

\medskip\noindent
$\xi_1, \ldots, \xi_r$ を、$X$ の $\overline{X}$ における補集合の生成点とする。
$\overline{X}$ の台位相は Noether なので、これらは有限個である（Morphisms, Lemma
\ref{morphisms-lemma-finite-type-noetherian}、Properties, Lemma
\ref{properties-lemma-Noetherian-topology}、および Topology, Lemma
\ref{topology-lemma-Noetherian} を参照せよ）。各 $i$ に対し、局所環
$\mathcal{O}_i = \mathcal{O}_{X, \xi_i}$ は Noether 正規整域の局所化として正規 Noether 局所整域である。
$J \subset \{1, \ldots, r\}$ を、添字 $i$ であって $\dim(\mathcal{O}_i) = 1$ となるものの集合とする。
$j \in J$ に対して局所環 $\mathcal{O}_j$ は離散付値環である。Algebra, Lemma
\ref{algebra-lemma-characterize-dvr} を参照せよ。従って付値
$$
v_j : k(\overline{X})^* \longrightarrow \mathbf{Z}
$$
を得て、これは $v_j(f) \geq 0 \Leftrightarrow f \in \mathcal{O}_j$ を満たす。

\medskip\noindent
$\mathcal{O}(X)$ を $k$-代数 $k(X) = k(\overline{X})$ の部分代数とみなす。写像
$$
\mathcal{O}(X)^* \longrightarrow
\prod\nolimits_{j \in J} \mathbf{Z},
\quad
f \longmapsto \prod v_j(f)
$$
の核は $k^*$ であると主張する。この主張が補題を示すことは明らかである。実際、
$f \in \mathcal{O}(X)$ がこの写像の核に属すると仮定する。
$U = \Spec(B) \subset \overline{X}$ を任意のアフィン開集合とする。このとき $B$ は Noether 正規整域である。
高さ一の任意の素イデアル $\mathfrak q \subset B$ と、それに対応する点 $\xi \in X$ に対し、
$\xi = \xi_j$ となる $j \in J$ があるか、または $\xi \in X$ である。というのも、Properties, Lemma
\ref{properties-lemma-codimension-local-ring} により
$\text{codim}(\overline{\{\xi\}}, \overline{X}) = 1$ であり、従って
$\xi \in \overline{X} \setminus X$ ならば、それは $\overline{X} \setminus X$ の生成点、すなわちある
$\xi_j$（$j \in J$）に等しいからである。いずれの場合にも
$f \in \mathcal{O}_{X, \xi} = B_{\mathfrak q}$ である。実際、$f$ は写像の核に属する。従って
$f \in \bigcap_{\text{ht}(\mathfrak q) = 1} B_{\mathfrak q} = B$ である。Algebra, Lemma
\ref{algebra-lemma-normal-domain-intersection-localizations-height-1} を参照せよ。言い換えれば、
$f \in \Gamma(\overline{X}, \mathcal{O}_{\overline{X}})$ である。しかし $k$ は代数閉なので、Lemma
\ref{varieties-lemma-regular-functions-proper-variety} により $f \in k$ と結論できる。
\end{proof}

\noindent
次に、体は引き続き代数閉としたまま、初等的な議論によって上の場合を一般化する。

\begin{lemma}
\label{varieties-lemma-units-integral-finite-type-algebraically-closed}
$k$ を代数閉体とする。$X$ を $k$ 上局所有限型な整スキームとする。
このとき $\mathcal{O}^*(X)/k^*$ は有限生成アーベル群である。
\end{lemma}

\begin{proof}
$X$ は整なので、制限写像 $\mathcal{O}(X) \to \mathcal{O}(U)$ は、任意の空でない開部分スキーム
$U \subset X$ に対して単射である。従って $X$ はアフィンであると仮定してよい。
閉埋め込み $X \to \mathbf{A}^n_k$ を選び、$\overline{X}$ を $X$ の $\mathbf{P}^n_k$ における閉包と書く。
ここでは通常の埋め込み $\mathbf{A}^n_k \to \mathbf{P}^n_k$ を用いる。従って $X$ は射影多様体
$\overline{X}$ のアフィン開集合であると仮定してよい。

\medskip\noindent
$\nu : \overline{X}^\nu \to \overline{X}$ を正規化射とする。Morphisms, Definition
\ref{morphisms-definition-normalization} を参照せよ。Morphisms, Lemmas
\ref{morphisms-lemma-normalization-normal}、
\ref{morphisms-lemma-Japanese-normalization}、および
\ref{morphisms-lemma-ubiquity-nagata} により、$\nu$ は有限かつ支配的であり、
$\overline{X}^\nu$ は正規既約スキームである。さらに $\overline{X}^\nu$ は固有多様体である。
実際、$\overline{X} \to \Spec(k)$ は有限射と固有射の合成として固有である
（Morphisms, Sections \ref{morphisms-section-proper} および
\ref{morphisms-section-integral} の結果を参照せよ）。また $\nu$ は全射である。
なぜなら、$\nu$ の像は閉であり、$\overline{X}$ の生成点を含むからである。従って
$X^\nu = \nu^{-1}(X)$ とおけば、$X^\nu$ に対して結果を示せば十分である。言い換えれば、
$X$ は正規固有多様体 $\overline{X}$ の空でない開集合であると仮定してよい。この場合は Lemma
\ref{varieties-lemma-open-in-normal-proper} で扱われている。
\end{proof}

\noindent
直前の補題から、次のわずかな一般化が従う。

\begin{lemma}
\label{varieties-lemma-units-general-algebraically-closed}
$k$ を代数閉体とする。$X$ を $k$ 上局所有限型な連結被約スキームで、有限個の既約成分をもつものとする。
このとき $\mathcal{O}^*(X)/k^*$ は有限生成アーベル群である。
\end{lemma}

\begin{proof}
$X = \bigcup X_i$ を既約成分分解とする。Lemma
\ref{varieties-lemma-units-integral-finite-type-algebraically-closed} により、
$\mathcal{O}(X_i)^*/k^*$ は有限生成アーベル群である。写像
$$
\mathcal{O}(X)^* \longrightarrow \prod \mathcal{O}(X_i)^*/k^*.
$$
の核に $f \in \mathcal{O}(X)^*$ が属するとする。各 $i$ に対し、ある $\lambda_i \in k$ が存在して
$f|_{X_i} = \lambda_i$ となる。$X_i \cap X_j$ に制限することにより、
$\lambda_i = \lambda_j$ が $X_i \cap X_j \not = \emptyset$ の場合に成り立つと分かる。$X$ は連結なので、
すべての $\lambda_i$ は一致し、従って $f \in k^*$ である。これにより
$$
\mathcal{O}(X)^*/k^* \subset \prod \mathcal{O}(X_i)^*/k^*
$$
が示された。右辺は有限個の有限生成アーベル群の積なので、補題が従う。
\end{proof}

\begin{lemma}
\label{varieties-lemma-integral-closure-ground-field}
$k$ を体とする。$X$ を連結かつ被約な $k$ 上のスキームとする。このとき
$k$ の $\Gamma(X, \mathcal{O}_X)$ における整閉包は体である。
\end{lemma}

\begin{proof}
$k' \subset \Gamma(X, \mathcal{O}_X)$ を $k$ の整閉包とする。このとき
$X \to \Spec(k)$ は $\Spec(k')$ を経由して分解する。Schemes, Lemma
\ref{schemes-lemma-morphism-into-affine} を参照せよ。$X$ は被約なので、$k'$ は非零冪零元をもたない。
$k \to k'$ は整なので、$k'$ の各素イデアルは極大イデアルであると同時に極小素イデアルであり、
$\Spec(k')$ は全不連結である。Algebra, Lemmas \ref{algebra-lemma-integral-no-inclusion} および
\ref{algebra-lemma-ring-with-only-minimal-primes} を参照せよ。$X$ は連結なので、射
$X \to \Spec(k')$ は定値であり、その像を $\mathfrak p \subset k'$ に対応する点とする。
このとき任意の $f \in k'$ で $f \not \in \mathfrak p$ を満たすものは、$\mathcal{O}_X$ の可逆元に写る。
$k'$ の定義により、これは $f$ が $k'$ の単元であることを強制する。従って $k'$ は極大イデアル
$\mathfrak p$ をもつ局所環である。Algebra, Lemma \ref{algebra-lemma-characterize-local-ring} を参照せよ。
すでに $k'$ が被約であることを見たので、これは $k'$ が体であることを意味する。Algebra, Lemma
\ref{algebra-lemma-minimal-prime-reduced-ring} を参照せよ。
\end{proof}

\begin{proposition}
\label{varieties-proposition-units-general}
$k$ を体とし、$X$ を $k$ 上のスキームとする。$X$ は $k$ 上局所有限型、連結、被約であり、
有限個の既約成分をもつと仮定する。上の条件に加えて次の条件の少なくとも一つが満たされるならば、
$\mathcal{O}(X)^*/k^*$ は有限生成アーベル群である。
\begin{enumerate}
\item $k$ の $\Gamma(X, \mathcal{O}_X)$ における整閉包は $k$ である。
\item $X$ は $k$-有理点をもつ。
\item $X$ は幾何学的整である。
\end{enumerate}
\end{proposition}

\begin{proof}
$\overline{k}$ を $k$ の代数閉包とする。$Y$ を $(X_{\overline{k}})_{red}$ の連結成分とする。
標準射 $p : Y \to X$ は開（Morphisms, Lemma
\ref{morphisms-lemma-scheme-over-field-universally-open} による）かつ閉
（Morphisms, Lemma \ref{morphisms-lemma-integral-universally-closed} による）であることに注意せよ。
従って $p(Y) = X$ である。これは $X$ が連結と仮定されているためである。特に、$X$ は被約なので、これは
$\mathcal{O}(X) \subset \mathcal{O}(Y)$ を意味する。Lemma
\ref{varieties-lemma-galois-action-irreducible-components-locally-finite-type} により、$Y$ は有限個の既約成分をもつ
（補題中の体 $\overline{k}$ は分離代数閉包であるが、Lemma
\ref{varieties-lemma-separably-closed-field-irreducible-components} により二つの基底変換 $X_{\overline{k}}$ は
同じ既約成分をもつので、この違いは問題にならない）。従って Lemma
\ref{varieties-lemma-units-general-algebraically-closed} を $Y$ に適用できる。このことから、
$\mathcal{O}(X)^*/k^*$ が有限生成アーベル群でないならば、元
$f \in \mathcal{O}(X)$、$f \not \in k$ で、$\overline{k}$ の元に写るものが存在する。この像は写像
$\mathcal{O}(X) \to \mathcal{O}(Y)$ による。この場合 $f$ は $k$ 上代数的であり、従って $k$ 上整である。
ゆえに条件 (1) が成り立つならば、このことは起こりえない。証明を終えるため、条件 (2) と (3) が
(1) を含意することを示す。

\medskip\noindent
$k \subset k' \subset \Gamma(X, \mathcal{O}_X)$ を、$k$ の $\Gamma(X, \mathcal{O}_X)$ における整閉包とする。
Lemma \ref{varieties-lemma-integral-closure-ground-field} により $k'$ は体である。
$e : \Spec(k) \to X$ が $k$-有理点ならば、
$e^\sharp : \Gamma(X, \mathcal{O}_X) \to k$ は包含写像
$k \to \Gamma(X, \mathcal{O}_X)$ の切断である。特に $e^\sharp$ の $k'$ への制限は体準同型
$k' \to k$ であり、これは $k$ 上の写像である。このことは明らかに (2) が (1) を含意することを示す。

\medskip\noindent
整閉包 $k'$、すなわち $k$ の $\Gamma(X, \mathcal{O}_X)$ における整閉包が自明でないならば、
$X$ は、$k'/k$ が純非分離でない場合には幾何学的連結でないか、または $X$ は、
$k'/k$ が非自明な純非分離拡大である場合には幾何学的被約でない。細部は省略する。従って (3) は (1) を含意する。
\end{proof}

\begin{lemma}
\label{varieties-lemma-units-variety}
$k$ を体とし、$X$ を $k$ 上の多様体とする。次の条件の少なくとも一つが成り立つならば、群
$\mathcal{O}(X)^*/k^*$ は有限生成アーベル群である。
\begin{enumerate}
\item $k$ は $\Gamma(X, \mathcal{O}_X)$ において整閉である。
\item $k$ は $k(X)$ において代数閉である。
\item $X$ は $k$ 上幾何学的整である。
\item $k$ は体拡大 $\kappa(x)/k$ の「共通部分」である。ここで $x$ は $x$ の閉点全体を動く。
\end{enumerate}
\end{lemma}

\begin{proof}
(1) で十分であることは Proposition \ref{varieties-proposition-units-general} から分かる。
(2)、(3)、(4) の各々が (1) を含意することの確認は省略する。
\end{proof}

\section{K"unneth 公式、I}
\label{varieties-section-kunneth}

\noindent
本節では、基礎が体であり、準連接加群のコホモロジーを考える場合の K"unneth 公式を証明する。
より一般的な形については Derived Categories of Schemes, Section
\ref{perfect-section-kunneth} を参照せよ。

\begin{lemma}
\label{varieties-lemma-kunneth}
$k$ を体とする。$X$ と $Y$ を $k$ 上のスキームとし、$\mathcal{F}$ および $\mathcal{G}$ を、それぞれ
$\mathcal{O}_X$-加群および $\mathcal{O}_Y$-加群である準連接加群とする。このとき標準同型
$$
H^n(X \times_{\Spec(k)} Y, \text{pr}_1^*\mathcal{F}
\otimes_{\mathcal{O}_{X \times_{\Spec(k)} Y}} \text{pr}_2^*\mathcal{G}) =
\bigoplus\nolimits_{p + q = n}
H^p(X, \mathcal{F}) \otimes_k H^q(Y, \mathcal{G})
$$
が存在する。ただし $X$ と $Y$ は準コンパクトで、アフィン対角射をもつものとする\footnote{$X$ と $Y$ が
準分離的である場合は、下の Lemma \ref{varieties-lemma-kunneth-general} で論じる。}
（例えば $X$ と $Y$ が分離的ならばよい）。
\end{lemma}

\begin{proof}
この証明では、修飾のない積およびテンソル積は $k$ 上で取る。写像
$$
H^p(X, \mathcal{F}) \otimes H^q(Y, \mathcal{G})
\longrightarrow
H^n(X \times Y, \text{pr}_1^*\mathcal{F}
\otimes_{\mathcal{O}_{X \times Y}} \text{pr}_2^*\mathcal{G})
$$
として、まずコホモロジーの関手性から得られる写像
$H^p(X, \mathcal{F}) \to H^p(X \times Y, \text{pr}_1^*\mathcal{F})$ および
$H^p(Y, \mathcal{G}) \to H^p(X \times Y, \text{pr}_2^*\mathcal{G})$ を用い、次いでカップ積
$$
\cup :
H^p(X \times Y, \text{pr}_1^*\mathcal{F})
\otimes H^q(X \times Y, \text{pr}_2^*\mathcal{G})
\longrightarrow
H^n(X \times Y, \text{pr}_1^*\mathcal{F}
\otimes_{\mathcal{O}_{X \times Y}} \text{pr}_2^*\mathcal{G})
$$
を用いる。$X$ と $Y$ がアフィンである場合、アフィン上の高次コホモロジー群の消滅
（Cohomology of Schemes, Lemma
\ref{coherent-lemma-quasi-coherent-affine-cohomology-zero}）および定義
（準連接加群の引き戻しと準連接加群のテンソル積の定義）により結果は正しい。

\medskip\noindent
有限アフィン開被覆
$\mathcal{U} : X = \bigcup_{i \in I} U_i$ および
$\mathcal{V} : Y = \bigcup_{j \in J} V_j$ を選ぶ。これらはアフィン開被覆
$\mathcal{W} : X \times Y = \bigcup_{(i, j) \in I \times J} U_i \times V_j$ を定める。
$\mathcal{W}$ は $\text{pr}_1^{-1}\mathcal{U}$ および $\text{pr}_2^{-1}\mathcal{V}$ の細分であることに注意せよ。
従って Cohomology, Lemma \ref{cohomology-lemma-functoriality-cech} により写像
$$
\check{\mathcal{C}}^\bullet(\mathcal{U}, \mathcal{F}) \to
\check{\mathcal{C}}^\bullet(\mathcal{W}, \text{pr}_1^*\mathcal{F})
\quad\text{and}\quad
\check{\mathcal{C}}^\bullet(\mathcal{V}, \mathcal{G}) \to
\check{\mathcal{C}}^\bullet(\mathcal{W}, \text{pr}_2^*\mathcal{G})
$$
を得る。これらはコホモロジー上の引き戻し写像と両立する。Cohomology, Equation
(\ref{cohomology-equation-needs-signs}) では、コホモロジー上のカップ積を定める複体の写像
$$
\text{Tot}(
\check{\mathcal{C}}^\bullet(\mathcal{W}, \text{pr}_1^*\mathcal{F})
\otimes
\check{\mathcal{C}}^\bullet(\mathcal{W}, \text{pr}_2^*\mathcal{G}))
\longrightarrow
\check{\mathcal{C}}^\bullet(\mathcal{W},
\text{pr}_1^*\mathcal{F} \otimes_{\mathcal{O}_{X \times Y}}
\text{pr}_2^*\mathcal{G})
$$
を構成した。以上を組み合わせると、複体の写像
\begin{equation}
\label{varieties-equation-kunneth-on-cech}
\text{Tot}(
\check{\mathcal{C}}^\bullet(\mathcal{U}, \mathcal{F})
\otimes
\check{\mathcal{C}}^\bullet(\mathcal{V}, \mathcal{G}))
\longrightarrow
\check{\mathcal{C}}^\bullet(\mathcal{W},
\text{pr}_1^*\mathcal{F} \otimes_{\mathcal{O}_{X \times Y}}
\text{pr}_2^*\mathcal{G})
\end{equation}
を得る。この写像は複体の同型ではないことを注意しておく。我々の被覆を用いて準連接層の
コホモロジーを計算できることを思い出そう（Cohomology of Schemes, Lemmas
\ref{coherent-lemma-affine-diagonal} および
\ref{coherent-lemma-cech-cohomology-quasi-coherent}）。従ってコホモロジー上で
(\ref{varieties-equation-kunneth-on-cech}) は補題の写像を再現する。

\medskip\noindent
準連接加群の短完全列 $0 \to \mathcal{F} \to \mathcal{F}' \to \mathcal{F}'' \to 0$ を考える。
(\ref{varieties-equation-kunneth-on-cech}) の構成は $\mathcal{F}$ に関して関手的であり、関連する {\v C}ech 複体の形成は
変数 $\mathcal{F}$ に関して完全である（アフィン開集合上の切断を取っているため）。従って短完全複体列の間の写像
$$
\xymatrix{
\text{Tot}(
\check{\mathcal{C}}^\bullet(\mathcal{U}, \mathcal{F})
\otimes
\check{\mathcal{C}}^\bullet(\mathcal{V}, \mathcal{G})) \ar[r] \ar[d] &
\text{Tot}(
\check{\mathcal{C}}^\bullet(\mathcal{U}, \mathcal{F}')
\otimes
\check{\mathcal{C}}^\bullet(\mathcal{V}, \mathcal{G})) \ar[r] \ar[d] &
\text{Tot}(
\check{\mathcal{C}}^\bullet(\mathcal{U}, \mathcal{F}'')
\otimes
\check{\mathcal{C}}^\bullet(\mathcal{V}, \mathcal{G})) \ar[d] \\
\check{\mathcal{C}}^\bullet(\mathcal{W},
\text{pr}_1^*\mathcal{F} \otimes_{\mathcal{O}_{X \times Y}}
\text{pr}_2^*\mathcal{G}) \ar[r] &
\check{\mathcal{C}}^\bullet(\mathcal{W},
\text{pr}_1^*\mathcal{F}' \otimes_{\mathcal{O}_{X \times Y}}
\text{pr}_2^*\mathcal{G}) \ar[r] &
\check{\mathcal{C}}^\bullet(\mathcal{W},
\text{pr}_1^*\mathcal{F}'' \otimes_{\mathcal{O}_{X \times Y}}
\text{pr}_2^*\mathcal{G})
}
$$
を得る（外側の零は省略した）。長完全コホモロジー列を見ると、補題の結果が $2$-out-of-$3$ の意味で
$\mathcal{F}, \mathcal{F}', \mathcal{F}''$ のうち二つに対して成り立てば、残る一つにも成り立つ。

\medskip\noindent
$X$ は準連接加群に関して有限コホモロジー次元をもつ。Cohomology of Schemes, Lemma
\ref{coherent-lemma-vanishing-nr-affines} を参照せよ。
$d(\mathcal{F}) = \max \{d \mid H^d(X, \mathcal{F}) \not = 0\}$ に関する帰納法により、
$d(\mathcal{F}) = 0$ の場合へ帰着する。$d(\mathcal{F}) > 0$ と仮定する。Cohomology of Schemes, Lemma
\ref{coherent-lemma-affine-diagonal-universal-delta-functor} により、埋め込み
$\mathcal{F} \to \mathcal{F}'$ であって、$H^p(X, \mathcal{F}') = 0$ がすべての
$p \geq 1$ に対して成り立つものが存在することを既に見た。
$\mathcal{F}'' = \Coker(\mathcal{F} \to \mathcal{F}')$ とおくと
$d(\mathcal{F}'') < d(\mathcal{F})$ である。そこで前段落の結果を適用すれば、補題を $\mathcal{F}'$ および
$\mathcal{F}''$ に対して示せば十分であり、帰納段階が証明される。

\medskip\noindent
$\mathcal{G}$ に対して同様に論じることにより、$\mathcal{F}$ と $\mathcal{G}$ はいずれも次数 $0$ でのみ
非零コホモロジーをもつと仮定してよい。$V \subset Y$ をアフィン開集合とする。アフィン開被覆
$\mathcal{U}_V : X \times V = \bigcup_{i \in I} U_i \times V$ を考える。直ちに
$$
\check{\mathcal{C}}^\bullet(\mathcal{U}, \mathcal{F}) \otimes \mathcal{G}(V) =
\check{\mathcal{C}}^\bullet(\mathcal{U}_V,
\text{pr}_1^*\mathcal{F} \otimes_{\mathcal{O}_{X \times Y}}
\text{pr}_2^*\mathcal{G})
$$
を得る（複体としての等号である）。従って
$$
R\text{pr}_{2, *}(\text{pr}_1^*\mathcal{F} \otimes_{\mathcal{O}_{X \times Y}}
\text{pr}_2^*\mathcal{G})
\cong
\Gamma(X, \mathcal{F}) \otimes_k \mathcal{G} \cong
\bigoplus\nolimits_{\alpha \in A} \mathcal{G}
$$
が $Y$ 上で成り立つ。ここで $A$ は $k$-ベクトル空間 $\Gamma(X, \mathcal{F})$ の基底である。
$Y$ 上のコホモロジーは直和と可換である（Cohomology, Lemma
\ref{cohomology-lemma-quasi-separated-cohomology-colimit}）。$\text{pr}_2$ に対する Leray スペクトル系列
（Cohomology, Lemma \ref{cohomology-lemma-apply-Leray} を介する）を用いると、
$H^n(X \times Y, \text{pr}_1^*\mathcal{F} \otimes_{\mathcal{O}_{X \times Y}}
\text{pr}_2^*\mathcal{G})$ は $n > 0$ に対して零であり、
$H^0(X, \mathcal{F}) \otimes H^0(Y, \mathcal{G})$ と同型となるのは $n = 0$ の場合である。これで証明は終わる
（ただし、この同型が次数 $0$ のカップ積によって実際に与えられることを確認すべきであるが、その確認は省略する）。
\end{proof}

\begin{lemma}
\label{varieties-lemma-kunneth-general}
$k$ を体とする。$X$ と $Y$ を $k$ 上のスキームとし、$\mathcal{F}$ および $\mathcal{G}$ を、それぞれ
$\mathcal{O}_X$-加群および $\mathcal{O}_Y$-加群である準連接加群とする。このとき標準同型
$$
H^n(X \times_{\Spec(k)} Y, \text{pr}_1^*\mathcal{F}
\otimes_{\mathcal{O}_{X \times_{\Spec(k)} Y}} \text{pr}_2^*\mathcal{G}) =
\bigoplus\nolimits_{p + q = n}
H^p(X, \mathcal{F}) \otimes_k H^q(Y, \mathcal{G})
$$
が存在する。ただし $X$ と $Y$ は準コンパクトかつ準分離的であるとする。
\end{lemma}

\begin{proof}
$X$ と $Y$ が分離的である場合、またはより一般にアフィン対角射をもつ場合には、「よりよい」証明として
Lemma \ref{varieties-lemma-kunneth} を参照せよ（この証明に対するその証明の長所は、写像を引き戻しとそれに続くカップ積として
同定する点にある）。$X'$ および $Y'$ を、それぞれ $X$ および $Y$ の無限小肥厚で、構造層が
$\mathcal{O}_{X'} = \mathcal{O}_X \oplus \mathcal{F}$ および
$\mathcal{O}_{Y'} = \mathcal{O}_Y \oplus \mathcal{G}$ となるものとする。ここで $\mathcal{F}$ および
$\mathcal{G}$ はそれぞれ二乗零イデアルである。このとき
$$
\mathcal{O}_{X' \times Y'} =
\mathcal{O}_{X \times Y}
\oplus \text{pr}_1^*\mathcal{F}
\oplus \text{pr}_2^*\mathcal{G}
\oplus \text{pr}_1^*\mathcal{F}
\otimes_{\mathcal{O}_{X \times Y}} \text{pr}_2^*\mathcal{G}
$$
が $X \times Y$ 上の層として成り立つ。このようにして
$$
H^n(X \times Y, \mathcal{O}_{X \times Y}) =
\bigoplus\nolimits_{p + q = n}
H^p(X, \mathcal{O}_X) \otimes_k H^q(Y, \mathcal{O}_Y)
$$
を $k$ 上の任意の準コンパクトかつ準分離的なスキームの対について証明すれば十分である。いくつかの細部は省略する。

\medskip\noindent
この主張を証明するため、Cohomology of Schemes, Lemma
\ref{coherent-lemma-hypercoverings} の形のコホモロジーと基底変換を用いる。この補題によれば、下に有界な
$k$-ベクトル空間の複体、すなわち複体 $\mathcal{K}^\bullet$ であって、$\Spec(k)$ 上の準連接加群からなり、
$Y$ の $\Spec(k)$ 上のコホモロジーを普遍的に計算するものが存在する。特に
$$
R\text{pr}_{1, *}(\mathcal{O}_{X \times Y}) \cong
(X \to \Spec(k))^*\mathcal{K}^\bullet
$$
が $D(\mathcal{O}_X)$ において成り立つ。複体 $\mathcal{K}^\bullet$ はホモトピーを除いて
$\bigoplus_{q \geq 0} H^q(Y, \mathcal{O}_Y)[-q]$ と同型である。これは体上のベクトル空間の任意の複体について
成り立つからである。従って
$$
R\text{pr}_{1, *}(\mathcal{O}_{X \times Y}) \cong
\bigoplus\nolimits_{q \geq 0}
H^q(Y, \mathcal{O}_Y)[-q] \otimes_k \mathcal{O}_X
$$
が $D(\mathcal{O}_X)$ において成り立つ。すると
\begin{align*}
R\Gamma(X \times Y, \mathcal{O}_{X \times Y})
& =
R\Gamma(X, R\text{pr}_{1, *}(\mathcal{O}_{X \times Y})) \\
& =
R\Gamma(X, \bigoplus\nolimits_{q \geq 0}
H^q(Y, \mathcal{O}_Y)[-q] \otimes_k \mathcal{O}_X) \\
& =
\bigoplus\nolimits_{q \geq 0}
R\Gamma(X,  H^q(Y, \mathcal{O}_Y) \otimes \mathcal{O}_X)[-q] \\
& =
\bigoplus\nolimits_{q \geq 0}
R\Gamma(X, \mathcal{O}_X) \otimes_k H^q(Y, \mathcal{O}_Y)[-q] \\
& =
\bigoplus\nolimits_{p, q \geq 0}
H^p(X, \mathcal{O}_X)[-p] \otimes_k H^q(Y, \mathcal{O}_Y)[-q]
\end{align*}
となり、望む結果を得る。第一の等号は $\text{pr}_1$ に対する Leray による
（Cohomology, Lemma \ref{cohomology-lemma-before-Leray}）。第二の等号は上で与えた全順像の分解による。
第三の等号は、コホモロジーが常に有限直和と可換であることによる
（また $Y$ のコホモロジーは十分大きい次数で消滅する。Cohomology of Schemes, Lemma
\ref{coherent-lemma-vanishing-nr-affines-quasi-separated} を参照せよ）。第四の等号は、$X$ 上のコホモロジーが無限直和と
可換であることによる。Cohomology, Lemma
\ref{cohomology-lemma-quasi-separated-cohomology-colimit} を参照せよ。最後の等号は、上で述べた体の導来圏に関する
注意による。
\end{proof}

\section{多様体の Picard 群}
\label{varieties-section-picard-groups}

\noindent
本節では代数多様体の Picard 群に関するいくつかの初等的な結果を集める。

\begin{lemma}
\label{varieties-lemma-change-rings-pic-pre}
$A \to B$ を忠実平坦な環準同型とする。$X$ を $A$ 上の準コンパクトかつ準分離的なスキームとする。
$\mathcal{L}$ を可逆 $\mathcal{O}_X$-加群とし、その $X_B$ への引き戻しは自明であるとする。このとき
$H^0(X, \mathcal{L})$ および $H^0(X, \mathcal{L}^{\otimes -1})$ は可逆
$H^0(X, \mathcal{O}_X)$-加群であり、乗法写像は同型
$$
H^0(X, \mathcal{L}) \otimes_{H^0(X, \mathcal{O}_X)}
H^0(X, \mathcal{L}^{\otimes -1}) \longrightarrow
H^0(X, \mathcal{O}_X)
$$
を誘導する。
\end{lemma}

\begin{proof}
$\mathcal{L}_B$ を $\mathcal{L}$ の $X_B$ への引き戻しと書き、同型
$\mathcal{L}_B \to \mathcal{O}_{X_B}$ を選ぶ。$R = H^0(X, \mathcal{O}_X)$、
$M = H^0(X, \mathcal{L})$ とおき、$M$ を $R$-加群とみなす。任意の準連接
$\mathcal{O}_X$-加群 $\mathcal{F}$ と、その引き戻し $\mathcal{F}_B$（$X_B$ 上）に対して標準同型
$H^0(X_B, \mathcal{F}_B) = H^0(X, \mathcal{F}) \otimes_A B$ がある。Cohomology of Schemes, Lemma
\ref{coherent-lemma-flat-base-change-cohomology} を参照せよ。従って
$$
M \otimes_R (R \otimes_A B) =
M \otimes_A B = H^0(X_B, \mathcal{L}_B) \cong
H^0(X_B, \mathcal{O}_{X_B}) = R \otimes_A B
$$
である。$R \to R \otimes_A B$ は忠実平坦写像 $A \to B$ の基底変換として忠実平坦なので、Algebra,
Proposition \ref{algebra-proposition-ffdescent-finite-projectivity} により $M$ は可逆 $R$-加群である。
同様に $N = H^0(X, \mathcal{L}^{\otimes -1})$ は可逆 $R$-加群である。テンソル積に関する主張を見るには、
$X_B$ へ引き戻した後にそれが成り立つことと、$R \to R \otimes_A B$ の忠実平坦性を用いればよい。
いくつかの細部は省略する。
\end{proof}

\begin{lemma}
\label{varieties-lemma-change-rings-pic}
$A \to B$ を忠実平坦な環準同型とする。$X$ を次を満たす $A$ 上のスキームとする。
\begin{enumerate}
\item $X$ は準コンパクトかつ準分離的である。
\item $R = H^0(X, \mathcal{O}_X)$ は半局所環である。
\end{enumerate}
このとき引き戻し写像 $\Pic(X) \to \Pic(X_B)$ は単射である。
\end{lemma}

\begin{proof}
$\mathcal{L}$ を可逆 $\mathcal{O}_X$-加群とし、その引き戻し $\mathcal{L}'$（$X_B$ への）は自明であるとする。
$M = H^0(X, \mathcal{L})$ および $N = H^0(X, \mathcal{L}^{\otimes - 1})$ とおく。
Lemma \ref{varieties-lemma-change-rings-pic-pre} により $R$-加群 $M$ と $N$ は可逆である。$R$ は半局所なので
$M \cong R$ かつ $N \cong R$ である。Algebra, Lemma
\ref{algebra-lemma-locally-free-semi-local-free} を参照せよ。生成元 $s \in M$ および $t \in N$ を選ぶ。
このとき Lemma \ref{varieties-lemma-change-rings-pic-pre} の最後の部分により
$st \in R = H^0(X, \mathcal{O}_X)$ は単元である。従って $s$ と $t$ は $\mathcal{L}$ および
$\mathcal{L}^{\otimes -1}$ の $X$ 上の自明化を定める。
\end{proof}

\begin{lemma}
\label{varieties-lemma-change-fields-pic}
$k'/k$ を体拡大とする。$X$ を次を満たす $k$ 上のスキームとする。
\begin{enumerate}
\item $X$ は準コンパクトかつ準分離的である。
\item $R = H^0(X, \mathcal{O}_X)$ は半局所である。例えば $\dim_k R < \infty$ ならばよい。
\end{enumerate}
このとき引き戻し写像 $\Pic(X) \to \Pic(X_{k'})$ は単射である。
\end{lemma}

\begin{proof}
Lemma \ref{varieties-lemma-change-rings-pic} の特別な場合である。$\dim_k R < \infty$ ならば
$R$ は Artin 環であり、従って半局所である（Algebra, Lemmas
\ref{algebra-lemma-finite-dimensional-algebra} および
\ref{algebra-lemma-artinian-finite-nr-max}）。
\end{proof}

\begin{example}
\label{varieties-example-need-condition}
Lemma \ref{varieties-lemma-change-fields-pic} は、スキーム $X$ が体 $k$ 上にあるというだけで、ほかに何らの条件も
課さなければ成り立たない。
次が一例である。$k$ を体とし、$t \in \mathbf{P}^1_k$ を閉点とする。
$X = \mathbf{P}^1 \setminus \{t\}$ とおく。このとき全射
$$
\mathbf{Z} = \Pic(\mathbf{P}^1_k) \longrightarrow \Pic(X)
$$
がある。第一の等号は Divisors, Lemma \ref{divisors-lemma-Pic-projective-space-UFD} により、全射性は
Divisors, Lemma \ref{divisors-lemma-extend-invertible-module} による
（$\mathbf{P}^1_k$ は次元 $1$ で $k$ 上平滑であり、従ってそのすべての局所環は離散付値環である）。
$\mathcal{L}$ が表示した写像の核に属するならば、
$\mathcal{L} \cong \mathcal{O}_{\mathbf{P}^1_k}(nt)$ がある $n \in \mathbf{Z}$ に対して成り立つ。読者に
$\mathcal{O}_{\mathbf{P}^1_k}(t) \cong \mathcal{O}_{\mathbf{P}^1_k}(d)$ を示すことを委ねる。ここで
$d = [\kappa(t) : k]$ である。従って
$$
\Pic(X) = \mathbf{Z}/d\mathbf{Z}
$$
である。ゆえに $t$ が $k$-有理点でなければ $d > 1$ であり、この Picard 群は非零である。
一方、基礎体 $k$ を任意の体拡大 $k'$ に拡大し、そこに $k$-埋め込み
$\kappa(t) \to k'$ が存在すると、
$\mathbf{P}^1_{k'} \setminus X_{k'}$ は $k'$-有理点 $t'$ をもつ。従って
$\mathcal{O}_{\mathbf{P}^1_{k'}}(1) = \mathcal{O}_{\mathbf{P}^1_{k'}}(t')$ は写像
$\mathbf{Z} \to \Pic(X_{k'})$ の核に属し、上と同様に $\Pic(X_{k'}) = 0$ が従う。
\end{example}

\noindent
次の補題によれば、正規多様体上では「有理同値な可逆加群」は同型である。

\begin{lemma}
\label{varieties-lemma-rational-equivalence-for-Pic}
$k$ を体とし、$X$ を $k$ 上の正規多様体とする。$U \subset \mathbf{A}^n_k$ を開部分スキームとし、
$k$-有理点 $p, q \in U(k)$ をもつとする。任意の可逆加群 $\mathcal{L}$（$X \times_{\Spec(k)} U$ 上）に対し、
その制限 $\mathcal{L}|_{X \times p}$ と $\mathcal{L}|_{X \times q}$ は同型である。
\end{lemma}

\begin{proof}
$X \times_{\Spec(k)} U \to X$ のファイバーは、体上のアフィン $n$-空間の開部分スキームである。
従って Divisors, Lemma \ref{divisors-lemma-open-subscheme-UFD} により、これらのファイバーの Picard 群は自明である。
Divisors, Lemma \ref{divisors-lemma-in-image-pullback} を適用すると、$\mathcal{L}$ は可逆加群
$\mathcal{N}$（$X$ 上）の引き戻しであることが分かる。
\end{proof}

\section{基礎体の一意性}
\label{varieties-section-base-field}

\noindent
「$X$ を $k$ 上のスキームとする」という句は、$X$ が射 $X \to \Spec(k)$ を備えたスキームであることを意味する。
そこで体 $k$ がスキーム $X$ によって一意に定まるかを問うことができる。もちろん一般にはそうではない。
例えば、$\mathbf{A}^1_{\mathbf{C}}$ は通常は複素数体 $\mathbf{C}$ 上のスキームとみなされるが、
$\mathbf{Q}$ 上のスキームとみなすこともできる。しかし、射 $X \to \Spec(k)$ が
$X \to \Spec(k') \to \Spec(k)$ と、いかなる非自明な体拡大 $k'/k$ に対しても分解しないことを要求したらどうだろうか。
言い換えれば、$k$ が、$X$ がその上に存在する体 $k$ のうち何らかの意味で極大であることを要求する。

\medskip\noindent
これでも $k$ の一意性が保証されないことを示す例は、スキーム
$$
X =
\Spec\left(
\mathbf{Q}(x)[y]\left[\frac{1}{P(y)}, P \in \mathbf{Q}[y], P \not = 0\right]
\right)
$$
である。一見するとこれは $\mathbf{Q}(x)$ 上のスキームに見えるが、見直せば
$\mathbf{Q}(y)$ 上のスキームでもあることは明らかである。さらに、体 $\mathbf{Q}(x)$ と
$\mathbf{Q}(y)$ は $R = \Gamma(X, \mathcal{O}_X)$ の部分体であり、$R$ の部分体のうち極大である
（細部は省略する）。特に $\mathbf{Q}(x)$ と $\mathbf{Q}(y)$ はともに上の意味で極大である。
二つの射 $X \to \Spec(\mathbf{Q}(x))$ および
$X \to \Spec(\mathbf{Q}(y))$ はともに「本質的に有限型」であることに注意せよ
（すなわち対応する環準同型が本質的に有限型である）。従って $X$ は有限次元 Noether スキームであり、
完全に病的な例ではない。

\medskip\noindent
一意性を妨げうる別の問題は、スキーム $X$ が非被約でありうることである。この場合、$X$ から与えられた体の
スペクトルへの射が多数存在しうる。明示的な例として双対数
$D = \mathbf{C}[y]/(y^2) = \mathbf{C} \oplus \epsilon \mathbf{C}$ を考える。
任意の導分 $\theta : \mathbf{C} \to \mathbf{C}$ で $\mathbf{Q}$ 上のものに対して環準同型
$$
\mathbf{C} \longrightarrow D, \quad
c \longmapsto c + \epsilon \theta(c).
$$
を得る。これらすべての写像が一致する $\mathbf{C}$ の部分体は、$\mathbf{Q}$ の代数閉包である。
これは、$X$ が存在しうるすべての体の共通部分を取ると、$X$ が非被約な場合には非常に小さな体になりうることを
意味する。

\medskip\noindent
これに関連する一つの観察は次のとおりである。体 $k$ と二つの部分体 $k_1, k_2$ が与えられ、後二者が $k$ に含まれるとする。
$k$ が $k_1$ 上および $k_2$ 上有限であっても、一般には $k$ が $k_1 \cap k_2$ 上有限であるとは
{\it 限らない}。例として体 $k = \mathbf{Q}(t)$ と、その部分体 $k_1 = \mathbf{Q}(t^2)$ および
$\mathbf{Q}((t + 1)^2)$ がある。実際、この場合 $k_1 \cap k_2 = \mathbf{Q}$ である。
従って以下で体の共通部分を取る際には注意が必要である。

\medskip\noindent
以上を踏まえ、$X$ が体上局所有限型であるならば、ある一意性が成り立つことを示す。正確な結果は次である。

\begin{proposition}
\label{varieties-proposition-unique-base-field}
$X$ をスキームとする。$a : X \to \Spec(k_1)$ および $b : X \to \Spec(k_2)$ を、
$X$ から体のスペクトルへの射とする。$a, b$ は局所有限型であり、$X$ は被約かつ連結であると仮定する。
このとき $k_1' = k_2'$ である。ここで $k_i' \subset \Gamma(X, \mathcal{O}_X)$ は、
$k_i$ の $\Gamma(X, \mathcal{O}_X)$ における整閉包である。
\end{proposition}

\begin{proof}
まず、$X$ が準コンパクトである場合に補題が成り立つと仮定する（準コンパクトな場合は後で証明する）。
$X$ は体上局所有限型なので局所 Noether である。Morphisms, Lemma
\ref{morphisms-lemma-finite-type-noetherian} を参照せよ。特に、これは局所連結であり、開部分集合の連結成分は開で、
準コンパクト開集合の交わりは準コンパクトである。Properties, Lemma
\ref{properties-lemma-Noetherian-topology}、Topology, Lemma
\ref{topology-lemma-locally-connected}、Topology, Section
\ref{topology-section-noetherian}、および Topology, Lemma
\ref{topology-lemma-constructible-Noetherian-space} を参照せよ。開被覆
$X = \bigcup_{i \in I} U_i$ を、各 $U_i$ が準コンパクトかつ連結となるよう選ぶ。各 $i$ に対し、$K_i \subset \mathcal{O}_X(U_i)$ を
$k_1$ と $k_2$ の整閉包とする。各対 $i, j \in I$ に対し
$$
U_i \cap U_j = \coprod U_{i, j, l}
$$
と、その有限個の連結成分へ分解する。$K_{i, j, l} \subset \mathcal{O}(U_{i, j, l})$ を
$k_1$ と $k_2$ の整閉包と書く。Lemma \ref{varieties-lemma-integral-closure-ground-field} により、環
$K_i$ と $K_{i, j, l}$ は体である。ここで $k_1'$ と $k_2'$ はともに写像
$$
\prod K_i \longrightarrow \prod K_{i, j, l}, \quad
(x_i)_i \longmapsto x_i|_{U_{i, j, l}} - x_j|_{U_{i, j, l}}
$$
の核に等しいと主張する。これで望むことが示される。実際、$k_1'$ がこの核に含まれることは明らかである。
逆に $(x_i)_i$ が核に属すると仮定する。層条件により $(x_i)_i$ は $f \in \mathcal{O}(X)$ に対応する。
$i_0 \in I$ を一つ選び、$P(T) \in k_1[T]$ を $P(x_{i_0}) = 0$ を満たすモニック多項式とする。
このとき $P(f) = 0$ と主張する。これは $f \in k_1$ を示す。これを証明するには、$P(x_i) = 0$ を
すべての $i$ に対して示さねばならない。$i \in I$ を選ぶ。$X$ は連結なので、列
$i_0, i_1, \ldots, i_n = i \in I$ であって
$U_{i_t} \cap U_{i_{t + 1}} \not = \emptyset$ を満たすものが存在する。従って各 $t$ に対してある $l_t$ が存在し、
$x_{i_t}$ と $x_{i_{t + 1}}$ は体 $K_{i_t, i_{t + 1}, l_t}$ の同じ元へ写る。
ゆえに $P(x_{i_t}) = 0$ ならば $P(x_{i_{t + 1}}) = 0$ である。
$P(x_{i_0}) = 0$ から始めて帰納すると、望むとおり $P(x_i) = 0$ を得る。

\medskip\noindent
補題の証明を終えるため、$X$ が準コンパクトであるという追加仮定の下で補題を証明する。Lemma
\ref{varieties-lemma-integral-closure-ground-field} により、$k_i$ を $k_i'$ に置き換えた後、$k_i$ は
$\Gamma(X, \mathcal{O}_X)$ において整閉であると仮定してよい。このことから
$\mathcal{O}(X)^*/k_i^*$ は有限生成アーベル群である。Proposition
\ref{varieties-proposition-units-general} を参照せよ。$k_{12} = k_1 \cap k_2$ を $\mathcal{O}(X)$ の部分環として定める。
$k_{12}$ は体であることに注意せよ。
$$
k_1^*/k_{12}^* \longrightarrow \mathcal{O}(X)^*/k_2^*
$$
なので、$k_1^*/k_{12}^*$ も有限生成アーベル群である。従って、$\alpha_1, \ldots, \alpha_n \in k_1^*$ が存在し、
すべての元 $\lambda \in k_1$ が
$$
\lambda = c \alpha_1^{e_1} \ldots \alpha_n^{e_n}
$$
という形をもつ。ここで
$e_i \in \mathbf{Z}$ かつ $c \in k_{12}$ である。特に環準同型
$$
k_{12}[x_1, \ldots, x_n, \frac{1}{x_1 \ldots x_n}] \longrightarrow k_1, \quad
x_i \longmapsto \alpha_i
$$
は全射である。Hilbert の零点定理 Algebra, Theorem \ref{algebra-theorem-nullstellensatz} により、
$k_1$ は $k_{12}$ の有限拡大である。同様に $k_2$ は $k_{12}$ の有限拡大であると結論する。
特に $k_1$ と $k_2$ はともに整閉包 $k_{12}'$ に含まれる。ここでこれは、$k_{12}$ の
$\Gamma(X, \mathcal{O}_X)$ における整閉包である。
しかし Lemma \ref{varieties-lemma-integral-closure-ground-field} により $k_{12}'$ は体であり、$k_i$ が
$\Gamma(X, \mathcal{O}_X)$ において整閉となるよう選んだので、望むとおり
$k_1 = k_{12} = k_2$ と結論する。
\end{proof}

\section{自己同型}
\label{varieties-section-automorphisms}

\noindent
本節では体上のスキームの自己同型を扱う。曲線の（無限小）自己同型に関する情報については、
Algebraic Curves, Section \ref{curves-section-vector-fields} および
Moduli of Curves, Section \ref{moduli-curves-section-finite-aut} を参照せよ。

\begin{lemma}
\label{varieties-lemma-automorphism}
$X$ を体 $k$ 上有限型の被約スキームとする。$f : X \to X$ を、$k$ 上の自己同型であって
$X$ の台位相空間上に恒等写像を誘導するものとする。このとき
\begin{enumerate}
\item $f^*\mathcal{F} \cong \mathcal{F}$ が任意のコヒーレント $\mathcal{O}_X$-加群に対して成り立ち、
\item $\dim(Z) > 0$ が任意の既約成分 $Z \subset X$ に対して成り立つならば、$f$ は恒等射である。
\end{enumerate}
\end{lemma}

\begin{proof}
(1) は (2) と、$X$ の $0$ 次元連結成分が体のスペクトルであることから従う。

\medskip\noindent
$Z \subset X$ を、整閉部分スキームとみなした既約成分とする。明らかに $f(Z) \subset Z$ であり、
$f|_Z : Z \to Z$ は $k$ 上の自己同型であって、$Z$ の台位相空間上に恒等写像を誘導する。
$X$ は被約なので、射 $f|_Z : Z \to Z$ が恒等射であることを示せば十分である。これで次段落で論じる場合に帰着する。

\medskip\noindent
$X$ は次元 $> 0$ の既約スキームであると仮定する。空でないアフィン開集合 $U \subset X$ を選ぶ。
$f(U) \subset U$ であり、$U \subset X$ はスキーム論的に稠密なので、$f|_U : U \to U$ が恒等射であることを
証明すれば十分である。

\medskip\noindent
$X = \Spec(A)$ はアフィンかつ既約で次元 $> 0$ とし、$k$ は無限体とする。非定数元 $g \in A$ を取る。集合
$$
S = \bigcup\nolimits_{\lambda \in k} V(g - \lambda)
$$
は $X$ で稠密である。実際、これは稠密部分集合 $\mathbf{A}^1_k(k)$ の、非定数射
$g : X \to \mathbf{A}^1_k$ による逆像である。$x \in S$ ならば、$g(x)$、すなわち $g$ の $\kappa(x)$ における像は
$k \to \kappa(x)$ の像に属する。従って $f^\sharp : \kappa(x) \to \kappa(x)$ は $g(x)$ を固定する。
ゆえに $f^\sharp(g)$ の $\kappa(x)$ における像は $g(x)$ に等しい。従って
$$
S \subset V(g - f^\sharp(g))
$$
であり、$X$ は被約かつ $S$ は稠密なので $g=f^\sharp(g)$ を得る。$f^\sharp = \text{id}_A$ が従う。
実際 $A$ は $k$-代数として上のような元 $g$ によって生成される（細部は省略する。ヒント：定数関数全体は
$k$ 上有限次元の $A$ のベクトル部分空間である）。従って $f = \text{id}_X$ である。

\medskip\noindent
$X = \Spec(A)$ はアフィンかつ既約で次元 $> 0$ とし、$k$ は有限体とする。任意の $1$ 次元整閉部分スキーム
$C \subset X$ に対して制限 $f|_C : C \to C$ が恒等射ならば、$f$ は恒等射である。従って $X$ が曲線である場合に
帰着する。有限体上の曲線の自己同型群は有限である（細部は省略する）。従って $f$ は有限位数、例えば $n$ をもつ。
そこで上と同様に非定数射 $g : X \to \mathbf{A}^1_k$ を選び、
$$
S = \{x \in X\text{ closed such that }[\kappa(g(x)) : k]
\text{ is prime to }n\}
$$
を考える。先と同様に論じると $S$ は $X$ で稠密である。閉点 $x \in X$ に対し、写像
$f^\sharp : \kappa(x) \to \kappa(x)$ は位数が $n$ を割る自己同型である。従って $x \in S$ に対して、
この自己同型は $g$ の $\kappa(x)$ における像が生成する部分体上、自明に作用する。ゆえに
$S \subset V(g - f^\sharp(g))$ であり、先と同様に結論を得る。
\end{proof}

\section{オイラー標数}
\label{varieties-section-euler}

\noindent
本節では、体上固有なスキーム上のコヒーレント層のオイラー標数について、いくつかの初等的性質を証明する。

\begin{definition}
\label{varieties-definition-euler-characteristic}
$k$ を体とする。$X$ を $k$ 上固有なスキームとする。$\mathcal{F}$ をコヒーレント $\mathcal{O}_X$-加群とする。
このとき、{\it $\mathcal{F}$ のオイラー標数}とは整数
$$
\chi(X, \mathcal{F}) = \sum\nolimits_i (-1)^i \dim_k H^i(X, \mathcal{F}).
$$
のことである。この式が正当であることは後述する。
\end{definition}

\noindent
定義の状況では、ベクトル空間 $H^i(X, \mathcal{F})$ のうち非零なものは有限個しかない
（Cohomology of Schemes, Lemma \ref{coherent-lemma-quasi-coherence-higher-direct-images}）。また各空間は有限次元である
（Cohomology of Schemes, Lemma \ref{coherent-lemma-proper-over-affine-cohomology-finite}）。従って
$\chi(X, \mathcal{F}) \in \mathbf{Z}$ は良定義である。この定義は体 $k$ に依存し、対
$(X, \mathcal{F})$ だけでは定まらないことに注意せよ。

\begin{lemma}
\label{varieties-lemma-euler-characteristic-additive}
$k$ を体とし、$X$ を $k$ 上固有なスキームとする。
$0 \to \mathcal{F}_1 \to \mathcal{F}_2 \to \mathcal{F}_3 \to 0$ を $X$ 上のコヒーレント加群の短完全列とする。このとき
$$
\chi(X, \mathcal{F}_2) = \chi(X, \mathcal{F}_1) + \chi(X, \mathcal{F}_3)
$$
である。
\end{lemma}

\begin{proof}
この短完全列に付随するコホモロジー長完全列
$$
0 \to H^0(X, \mathcal{F}_1) \to H^0(X, \mathcal{F}_2) \to
H^0(X, \mathcal{F}_3) \to H^1(X, \mathcal{F}_1) \to \ldots
$$
を考える。線型代数の階数・退化次数定理から
$$
0 = \dim H^0(X, \mathcal{F}_1) - \dim H^0(X, \mathcal{F}_2)
+ \dim H^0(X, \mathcal{F}_3) - \dim H^1(X, \mathcal{F}_1) + \ldots
$$
を得る。これは直ちに補題を含意する。
\end{proof}

\begin{lemma}
\label{varieties-lemma-chi-tensor-finite}
$k$ を体とし、$X$ を $k$ 上固有なスキームとする。$\mathcal{F}$ を
$\dim(\text{Supp}(\mathcal{F})) \leq 0$ を満たすコヒーレント層とする。このとき
\begin{enumerate}
\item $\mathcal{F}$ は大域切断で生成され、
\item $H^0(X, \mathcal{F}) =
\bigoplus_{x \in \text{Supp}(\mathcal{F})} \mathcal{F}_x$ であり、
\item $H^i(X, \mathcal{F}) = 0$ が $i > 0$ に対して成り立ち、
\item $\chi(X, \mathcal{F}) = \dim_k H^0(X, \mathcal{F})$ であり、さらに
\item $\chi(X, \mathcal{F} \otimes \mathcal{E}) = n\chi(X, \mathcal{F})$ が、任意の局所自由加群
$\mathcal{E}$ で階数 $n$ のものに対して成り立つ。
\end{enumerate}
\end{lemma}

\begin{proof}
Cohomology of Schemes, Lemma \ref{coherent-lemma-coherent-support-closed} により
$\mathcal{F} = i_*\mathcal{G}$ である。ここで $i : Z \to X$ は $\mathcal{F}$ のスキーム論的台の包含であり、
$\mathcal{G}$ はコヒーレント $\mathcal{O}_Z$-加群である。スキーム論的台の定義により、$Z$ の台位相空間は
$\text{Supp}(\mathcal{F})$ である。$Z$ の次元は $0$ なので、$Z$ はアフィンである
（Properties, Lemma \ref{properties-lemma-locally-Noetherian-dimension-0}）。従って $\mathcal{G}$ は大域生成され、
$\mathcal{G}$ の高次コホモロジー群は零である
（Cohomology of Schemes, Lemma \ref{coherent-lemma-quasi-coherent-affine-cohomology-zero}）。実際、Lemma
\ref{varieties-lemma-algebraic-scheme-dim-0} により $Z$ は局所 Artin 環のスペクトルの有限非交和である。従って対応して
$H^0(Z, \mathcal{G}) = \bigoplus_{z \in Z} \mathcal{G}_z$ となる。$\mathcal{F}$ と $\mathcal{G}$ のコホモロジーは
Cohomology of Schemes, Lemma \ref{coherent-lemma-relative-affine-cohomology} により一致する。従って
$H^i(X, \mathcal{F}) = 0$ が $i > 0$ に対して成り立ち、$H^0(X, \mathcal{F}) = H^0(Z, \mathcal{G})$ である。
特に (3) が成り立つ。$z \in Z$ を $x \in \text{Supp}(\mathcal{F})$ に対応する点とすると、
$\mathcal{G}_z = (i_*\mathcal{G})_x = \mathcal{F}_x$ である。従って (2) が成り立つ。もちろん (2) は (1) を含意する。
オイラー標数 $\chi(X, \mathcal{F})$ の定義と (3) から (4) を得る。射影公式
（Cohomology, Lemma \ref{cohomology-lemma-projection-formula}）により
$$
i_*(\mathcal{G} \otimes i^*\mathcal{E}) = \mathcal{F} \otimes \mathcal{E}.
$$
である。$Z$ は次元 $0$ なので、局所自由層 $i^*\mathcal{E}$ は $\mathcal{O}_Z^{\oplus n}$ と同型である。
従って上と同様に論じれば (5) が成り立つ。
\end{proof}

\begin{lemma}
\label{varieties-lemma-euler-characteristic-extend-base-field}
$k'/k$ を体拡大とする。$X$ を $k$ 上固有なスキームとし、$\mathcal{F}$ を $X$ 上のコヒーレント層とする。
$\mathcal{F}'$ を $\mathcal{F}$ の $X_{k'}$ への引き戻しとする。このとき
$\chi(X, \mathcal{F}) = \chi(X', \mathcal{F}')$ である。
\end{lemma}

\begin{proof}
これは平坦基底変換
$$
H^i(X_{k'}, \mathcal{F}') = H^i(X, \mathcal{F}) \otimes_k k'
$$
による。Cohomology of Schemes, Lemma \ref{coherent-lemma-flat-base-change-cohomology} を参照せよ。
\end{proof}

\begin{lemma}
\label{varieties-lemma-euler-characteristic-morphism}
$k$ を体とする。$f : Y \to X$ を $k$ 上固有なスキームの射とする。$\mathcal{G}$ をコヒーレント
$\mathcal{O}_Y$-加群とする。このとき
$$
\chi(Y, \mathcal{G}) = \sum (-1)^i \chi(X, R^if_*\mathcal{G})
$$
である。
\end{lemma}

\begin{proof}
この式には意味がある。すなわち層 $R^if_*\mathcal{G}$ はコヒーレントであり、非零なものは有限個しかない。
Cohomology of Schemes, Proposition \ref{coherent-proposition-proper-pushforward-coherent} および
Lemma \ref{coherent-lemma-quasi-coherence-higher-direct-images} を参照せよ。Cohomology, Lemma
\ref{cohomology-lemma-Leray} により、
$$
E_2^{p, q} = H^p(X, R^qf_*\mathcal{G})
$$
をもち $H^{p + q}(Y, \mathcal{G})$ に収束するスペクトル系列がある。$X$ 上のコホモロジーの有限性により、
非零な $E_2^{p, q}$ は有限個しかなく、各 $E_2^{p, q}$ は有限次元ベクトル空間である。従って
$E_r^{p, q}$ についても $r \geq 2$ ならば同様であり、
$$
\sum (-1)^{p + q} \dim_k E_r^{p, q}
$$
は $r$ に依存しない。十分大きい $r$ に対して $E_r^{p, q} = E_\infty^{p, q}$ であり、収束とは
$H^n(Y, \mathcal{G})$ 上に、$E_\infty^{p, q}$ を $p + q = n$ のとき次数付き部分とするフィルトレーションが
存在することを意味する（これがスペクトル系列の収束の意味である）ので、結論を得る。より一般的な
Homology, Lemma \ref{homology-lemma-biregular-ss-relation-in-K0} とも比較せよ。
\end{proof}

\section{射影空間}
\label{varieties-section-projective-space}

\noindent
体上の射影空間に関するいくつかの結果を示す。

\begin{lemma}
\label{varieties-lemma-projective-space-smooth}
\begin{slogan}
射影空間は平滑である。
\end{slogan}
$k$ を体とし、$n \geq 0$ とする。このとき $\mathbf{P}^n_k$ は次元 $n$ の $k$ 上平滑射影多様体である。
\end{lemma}

\begin{proof}
省略する。
\end{proof}

\begin{lemma}
\label{varieties-lemma-intersection-in-affine-space}
$k$ を体とし、$n \geq 0$ とする。$X, Y \subset \mathbf{A}^n_k$ を閉部分集合とする。
$X$ と $Y$ は等次元で、$\dim(X) = r$ および $\dim(Y) = s$ と仮定する。このとき
$X \cap Y$ の各既約成分の次元は $\geq r + s - n$ である。
\end{lemma}

\begin{proof}
閉部分スキーム $X \times Y \subset \mathbf{A}^{2n}_k$ を考え、座標
$x_1, \ldots, x_n, y_1, \ldots, y_n$ を用いる。このとき
$X \cap Y = X \times Y \cap V(x_1 - y_1, \ldots, x_n - y_n)$ である。
閉点 $t \in X \cap Y \subset X \times Y$ を取る。Lemma
\ref{varieties-lemma-dimension-product-locally-algebraic} により
$\dim_t(X \times Y) = \dim(X) + \dim(Y)$ である。従って Lemma
\ref{varieties-lemma-dimension-locally-algebraic} により $\dim(\mathcal{O}_{X \times Y, t}) = r + s$ である。
Algebra, Lemma \ref{algebra-lemma-one-equation} から
$$
\dim(\mathcal{O}_{X \cap Y, t}) =
\dim(\mathcal{O}_{X \times Y, t}/(x_1 - y_1, \ldots, x_n - y_n)) \geq
r + s - n
$$
を得る。Lemma \ref{varieties-lemma-dimension-locally-algebraic} により、これは結論を含意する。
\end{proof}

\begin{lemma}
\label{varieties-lemma-intersection-in-projective-space}
$k$ を体とし、$n \geq 0$ とする。$X, Y \subset \mathbf{P}^n_k$ を空でない閉部分集合とする。
$\dim(X) = r$、$\dim(Y) = s$、および $r + s \geq n$ ならば、$X \cap Y$ は空でなく、
$\dim(X \cap Y) \geq r + s - n$ である。
\end{lemma}

\begin{proof}
$\mathbf{A}^n = \Spec(k[x_0, \ldots, x_n])$ および
$\mathbf{P}^n = \text{Proj}(k[T_0, \ldots, T_n])$ と書く。射
$\pi : \mathbf{A}^{n + 1} \setminus \{0\} \to \mathbf{P}^n$ で、
$(x_0, \ldots, x_n)$ を点 $[x_0 : \ldots : x_n]$ へ送るものを考える。より正確には、これは対
$(\mathcal{O}_{\mathbf{A}^{n + 1} \setminus \{0\}}, (x_0, \ldots, x_n))$ に付随する射である。
Constructions, Lemma \ref{constructions-lemma-projective-space} を参照せよ。標準アフィン開集合 $D_+(T_i)$ 上では、環準同型
$$
k\left[\frac{T_0}{T_i}, \ldots, \frac{T_n}{T_i}\right]
\longrightarrow
k\left[T_0, \ldots, T_n, \frac{1}{T_i}\right] \cong
k\left[\frac{T_0}{T_i}, \ldots, \frac{T_n}{T_i}\right]
\left[T_i, \frac{1}{T_i}\right]
$$
に付随する射を得る。この射は全射かつ相対次元 $1$ で平滑であり、そのファイバーは既約である（細部は省略する）。
従って $\pi^{-1}(X)$ と $\pi^{-1}(Y)$ は、それぞれ次元 $r + 1$ と $s + 1$ の空でない閉部分集合である。
$V \subset \pi^{-1}(X)$ を次元 $r + 1$ の既約成分として選び、
$W \subset \pi^{-1}(Y)$ を次元 $s + 1$ の既約成分として選ぶ。このことから、$V$ と $W$ は交わる $\pi$ の各ファイバーをすべて含む
（$\pi$ のファイバーは次元 $1$ で既約であり、Lemma \ref{varieties-lemma-dimension-fibres-locally-algebraic} により
$V \to \pi(V)$ と $W \to \pi(W)$ のファイバーは次元 $\geq 1$ だからである）。
$\overline{V}$ と $\overline{W}$ を、$V$ と $W$ の $\mathbf{A}^{n + 1}$ における閉包とする。
$0 \in \mathbf{A}^{n + 1}$ は $\pi$ の各ファイバーの閉包に属するので、
$0 \in \overline{V} \cap \overline{W}$ である。Lemma \ref{varieties-lemma-intersection-in-affine-space} により
$\dim(\overline{V} \cap \overline{W}) \geq r + s - n + 1$ である。再び Lemma
\ref{varieties-lemma-dimension-fibres-locally-algebraic} を用いて上と同様に論じれば、
$\pi(V \cap W) \subset X \cap Y$ の次元は少なくとも $r + s - n$ である。これが望む結論である。
\end{proof}

\begin{lemma}
\label{varieties-lemma-equation-codim-1-in-projective-space}
$k$ を体とする。$Z \subset \mathbf{P}^n_k$ を埋め込まれた点をもたない閉部分スキームで、$Z$ の各既約成分が
次元 $n - 1$ をもつものとする。このときイデアル $I(Z) \subset k[T_0, \ldots, T_n]$ で $Z$ に対応するものは単項である。
\end{lemma}

\begin{proof}
これは Divisors, Lemma \ref{divisors-lemma-equation-codim-1-in-projective-space} の特別な場合である。
\end{proof}

\section{射影空間上のコヒーレント層}
\label{varieties-section-coherent-Pn}

\noindent
本節では、\cite{Mum} に見られる、体上の $\mathbf{P}^n$ 上のコヒーレント層のコホモロジーに関する
いくつかの結果を証明する。これらは後で Quot スキームと Hilbert スキームを論じる際に有用となる。

\subsection{準備}
\label{varieties-subsection-preliminaries}

\noindent
$k$ を体、$n \geq 1$、$d \geq 1$ とし、
$s \in \Gamma(\mathbf{P}_k^n, \mathcal{O}(d))$ を非零切断とする。本節では射影空間上の構造層の
記号 $\mathcal{O}(d)$ で、射影空間上の構造層の第 $d$ 捩れを表す（Constructions, Definitions
\ref{constructions-definition-twist} および \ref{constructions-definition-projective-space}）。
$\mathbf{P}^n_k$ は多様体なのでこの切断は正則である。従って $s$ は $\mathcal{O}(d)$ の正則切断であり、
有効 Cartier 因子 $H = Z(s) \subset \mathbf{P}^n_k$ を定める。Divisors, Section
\ref{divisors-section-effective-Cartier-divisors} を参照せよ。このような因子 $H$ を {\it 超曲面}といい、
$d = 1$ ならば {\it 超平面}という。

\begin{lemma}
\label{varieties-lemma-hyperplane}
$k$ を体とし、$n \geq 1$ とする。$i : H \to \mathbf{P}^n_k$ を超平面とする。このとき同型
$$
\varphi : \mathbf{P}^{n - 1}_k \longrightarrow H
$$
であって、その下で $i^*\mathcal{O}(1)$ が $\mathcal{O}(1)$ に引き戻されるものが存在する。
\end{lemma}

\begin{proof}
$\mathbf{P}^n_k = \text{Proj}(k[T_0, \ldots, T_n])$ である。切断 $s$ は $T_0, \ldots, T_n$ の次数 $1$ の
斉次形式に対応する。Cohomology of Schemes, Section
\ref{coherent-section-cohomology-projective-space} を参照せよ。$s = \sum a_i T_i$ としよう。
Constructions, Lemma \ref{constructions-lemma-closed-in-projective-space} により、
$H = \text{Proj}(k[T_0, \ldots, T_n]/I)$ である。ここで次数付きイデアル $I$ は、$I_d$ を写像
$\Gamma(\mathbf{P}^n_k, \mathcal{O}(d)) \to \Gamma(H, i^*\mathcal{O}(d))$ の核とすることで定義される。
Divisors, Definition \ref{divisors-definition-zero-scheme-s} における $Z(s)$ の構成から、$D_{+}(T_j)$ 上で
$H$ のイデアルは元 $\sum a_i T_i/T_j$ により、多項式環 $k[T_0/T_j, \ldots, T_n/T_j]$ の中で生成される。
従って $I$ が $\sum a_i T_i$ で生成されるイデアルであることは明らかである。次に
$$
k[T_0, \ldots, T_n]/I = k[T_0, \ldots, T_n]/(\sum a_i T_i) \cong
k[S_0, \ldots, S_{n - 1}]
$$
は次数付き環として同型であることに注意せよ。例えば $a_n \not = 0$ ならば、$S_i$ を $T_i$ の剰余類へ写せばよい。
$\text{Proj}$ の関手性により望む同型を得る。構造層の捩れの等しさは、例えば Constructions, Lemma
\ref{constructions-lemma-surjective-graded-rings-generated-degree-1-map-proj} から従う。
\end{proof}

\begin{lemma}
\label{varieties-lemma-exact-sequence-induction}
$k$ を無限体とし、$n \geq 1$ とする。$\mathcal{F}$ を $\mathbf{P}^n_k$ 上のコヒーレント加群とする。
このとき非零切断 $s \in \Gamma(\mathbf{P}^n_k, \mathcal{O}(1))$ と短完全列
$$
0 \to \mathcal{F}(-1) \to \mathcal{F} \to i_*\mathcal{G} \to 0
$$
が存在する。ここで $i : H \to \mathbf{P}^n_k$ は超平面 $H$ の包含であり、この超平面は $s$ に付随し、
$\mathcal{G} = i^*\mathcal{F}$ である。
\end{lemma}

\begin{proof}
写像 $\mathcal{F}(-1) \to \mathcal{F}$ は、$n = 1$、$m = -1$ および切断 $s$（$\mathcal{O}(1)$ の切断）に対して
Constructions, Equation (\ref{constructions-equation-multiply-on-sheaf}) から得られる。
$D_{+}(T_0)$ に制限したときのこの写像を具体的に調べよう。
$D_{+}(T_0) = \Spec(k[x_1, \ldots, x_n])$ と書き、$x_i = T_i/T_0$ とする。
$\mathcal{O}(1)|_{D_{+}(T_0)}$ と $\mathcal{O}$ の同一視には切断 $T_0$ を用いる。従って $s = \sum a_iT_i$ ならば、上で選んだ
同一視の下で $s|_{D_{+}(T_0)} = a_0 + \sum a_ix_i$ である。さらに、
$\mathcal{F}|_{D_{+}(T_0)}$ が有限 $k[x_1, \ldots, x_n]$-加群 $M$ に対応すると仮定する。同一視
$\mathcal{F}(-1) = \mathcal{F} \otimes \mathcal{O}(-1)$ と、選んだ $\mathcal{O}(1)$ の自明化により、
$\mathcal{F}(-1)$ も $M$ に対応することが分かる。従って $\mathcal{F}(-1) \to \mathcal{F}$ を
$D_{+}(T_0)$ に制限すると写像
$$
M \xrightarrow{a_0 + \sum a_ix_i} M
$$
を得る。この射が単射であることを示すには、$a_0 + \sum a_ix_i$ を $M$ のどの随伴素イデアルにも属さないよう
選べば十分である。Algebra, Lemma \ref{algebra-lemma-ass-zero-divisors} を参照せよ。Algebra, Lemma
\ref{algebra-lemma-finite-ass} により、集合 $\text{Ass}(M)$、すなわち $M$ の随伴素イデアル全体は有限である。
$\mathfrak p \in \text{Ass}(M)$ に対して、共通部分
$\mathfrak p \cap \{a_0 + \sum a_i x_i\}$ は真の $k$-ベクトル部分空間であることに注意せよ。
従って真のベクトル部分空間の有限族
$V_1, \ldots, V_m \subset \Gamma(\mathbf{P}^n_k, \mathcal{O}(1))$ が存在し、$s$ を $\bigcup V_i$ の外に取れば、
$s$ による乗法は $D_{+}(T_0)$ 上で単射になる。同様に $D_{+}(T_j)$ への制限を
$j = 1, \ldots, n$ に対して考える。$k$ は無限なので、真のベクトル部分空間の有限和は全空間に決して等しくない。
従って写像が単射になるよう $s$ を選べる。
$\mathcal{F}(-1) \to \mathcal{F}$ の余核は
$\Im(s : \mathcal{O}(-1) \to \mathcal{O})$ により零化され、これは $H$ のイデアル層である。Divisors, Definition
\ref{divisors-definition-zero-scheme-s} を参照せよ。従って Cohomology of Schemes, Lemma
\ref{coherent-lemma-i-star-equivalence} を用いて $\mathcal{G}$ を $H$ 上に得る。
\end{proof}

\begin{remark}
\label{varieties-remark-exact-sequence-induction}
$k$ を無限体とし、$n \geq 1$ とする。有限個のコヒーレント加群 $\mathcal{F}_i$ が $\mathbf{P}^n_k$ 上に
与えられたとき、Lemma \ref{varieties-lemma-exact-sequence-induction} の主張がそれぞれに対して成り立つような一つの
$s \in \Gamma(\mathbf{P}^n_k, \mathcal{O}(1))$ を選べる。これを証明するには、補題を
$\bigoplus \mathcal{F}_i$ に適用すればよい。
\end{remark}

\begin{remark}
\label{varieties-remark-exact-sequence-induction-cohomology}
Lemmas \ref{varieties-lemma-hyperplane} および \ref{varieties-lemma-exact-sequence-induction} の状況では、
$H \cong \mathbf{P}^{n - 1}_k$ があり、その Serre 捩れは
$\mathcal{O}_H(d) = i^*\mathcal{O}_{\mathbf{P}^n_k}(d)$ である。任意の $d \in \mathbf{Z}$ に対して短完全列
$$
0 \to \mathcal{F}(d - 1) \to \mathcal{F}(d) \to i_*(\mathcal{G}(d)) \to 0
$$
がある。実際、$\mathcal{O}_{\mathbf{P}^n_k}(d)$ によるテンソル積は完全関手であり、射影公式
（Cohomology, Lemma \ref{cohomology-lemma-projection-formula}）により
$i_*(\mathcal{G}(d)) = i_*\mathcal{G} \otimes \mathcal{O}_{\mathbf{P}^n_k}(d)$ である。
対応する長完全列
$$
H^i(\mathbf{P}^n_k, \mathcal{F}(d - 1)) \to
H^i(\mathbf{P}^n_k, \mathcal{F}(d)) \to
H^i(H, \mathcal{G}(d)) \to
H^{i + 1}(\mathbf{P}^n_k, \mathcal{F}(d - 1))
$$
を得る。これは上の議論と、Cohomology of Schemes, Lemma
\ref{coherent-lemma-relative-affine-cohomology} による
$H^i(\mathbf{P}^n_k, i_*\mathcal{G}(d)) = H^i(H, \mathcal{G}(d))$ という事実から従う
（閉埋め込みはアフィンである）。
\end{remark}





\subsection{正則性}
\label{varieties-subsection-regularity}

\noindent
定義は次のとおりである。

\begin{definition}
\label{varieties-definition-regularity}
$k$ を体、$n \geq 0$ とし、$\mathcal{F}$ を $\mathbf{P}^n_k$ 上のコヒーレント層とする。
$\mathcal{F}$ が {\it $m$-正則}であるとは、
$$
H^i(\mathbf{P}^n_k, \mathcal{F}(m - i)) = 0
$$
が $i = 1, \ldots, n$ に対して成り立つことをいう。
\end{definition}

\noindent
$\mathcal{F} = \mathcal{O}(d)$ が $m$-正則であるための必要十分条件は $d \geq -m$ であることに注意せよ。
これは Cohomology of Schemes, Equation (\ref{coherent-equation-identify}) におけるコホモロジー群の計算から従う。
実際、$H^n(\mathbf{P}^n_k, \mathcal{O}(d)) = 0$ であるための必要十分条件は $d \geq -n$ である。

\begin{lemma}
\label{varieties-lemma-m-regular-extend-base-field}
$k'/k$ を体拡大とし、$n \geq 0$ とする。$\mathcal{F}$ を $\mathbf{P}^n_k$ 上のコヒーレント層とする。
$\mathcal{F}'$ を $\mathcal{F}$ の $\mathbf{P}^n_{k'}$ への引き戻しとする。このとき $\mathcal{F}$ が
$m$-正則であるための必要十分条件は、$\mathcal{F}'$ が $m$-正則であることである。
\end{lemma}

\begin{proof}
これは平坦基底変換
$$
H^i(\mathbf{P}^n_{k'}, \mathcal{F}') =
H^i(\mathbf{P}^n_k, \mathcal{F}) \otimes_k k'
$$
による。Cohomology of Schemes, Lemma \ref{coherent-lemma-flat-base-change-cohomology} を参照せよ。
\end{proof}

\begin{lemma}
\label{varieties-lemma-m-regular}
Lemma \ref{varieties-lemma-exact-sequence-induction} の状況で、$\mathcal{F}$ が $m$-正則ならば、
$\mathcal{G}$ は $m$-正則であり、その基礎空間は $H \cong \mathbf{P}^{n - 1}_k$ である。
\end{lemma}

\begin{proof}
Cohomology of Schemes, Lemma \ref{coherent-lemma-relative-affine-cohomology} により
$H^i(\mathbf{P}^n_k, i_*\mathcal{G}) = H^i(H, \mathcal{G})$ であることを想起せよ。従って $i \geq 1$ に対して、
Remark \ref{varieties-remark-exact-sequence-induction-cohomology} により
$$
H^i(\mathbf{P}^n_k, \mathcal{F}(m - i)) \to
H^i(H, \mathcal{G}(m - i)) \to
H^{i + 1}(\mathbf{P}^n_k, \mathcal{F}(m - 1 - i))
$$
を得る。これから補題が従う。
\end{proof}

\begin{lemma}
\label{varieties-lemma-m-regular-up}
$k$ を体、$n \geq 0$ とし、$\mathcal{F}$ を $\mathbf{P}^n_k$ 上のコヒーレント層とする。
$\mathcal{F}$ が $m$-正則ならば、$\mathcal{F}$ は $(m + 1)$-正則である。
\end{lemma}

\begin{proof}
$n$ に関する帰納法で証明する。$n = 0$ ならばすべての層は任意の $m$ に対して $m$-正則なので、証明すべきことはない。
Lemma \ref{varieties-lemma-m-regular-extend-base-field} により $k$ を無限上体で置き換え、$k$ が無限であると仮定してよい。
従って Lemma \ref{varieties-lemma-exact-sequence-induction} を適用できる。Lemma \ref{varieties-lemma-m-regular} により
$\mathcal{G}$ は $m$-正則である。$n$ に関する帰納法により $\mathcal{G}$ は $(m + 1)$-正則である。Remark
\ref{varieties-remark-exact-sequence-induction-cohomology} と同様に、列
$$
0 \to \mathcal{F}(m - i) \to \mathcal{F}(m + 1 - i)
\to i_*\mathcal{G}(m + 1 - i) \to 0
$$
に付随するコホモロジー長完全列を考える。読者は $i \geq 1$ に対して、
$H^i(\mathbf{P}^n_k, \mathcal{F}(m + 1 - i))$ の消滅を、既知の
$H^i(\mathbf{P}^n_k, \mathcal{F}(m - i))$ および
$H^i(\mathbf{P}^n_k, \mathcal{G}(m + 1 - i))$ の消滅から容易に導ける。
\end{proof}

\begin{lemma}
\label{varieties-lemma-m-regular-multiply}
$k$ を体、$n \geq 0$ とし、$\mathcal{F}$ を $\mathbf{P}^n_k$ 上のコヒーレント層とする。
$\mathcal{F}$ が $m$-正則ならば、乗法写像
$$
H^0(\mathbf{P}^n_k, \mathcal{F}(m)) \otimes_k
H^0(\mathbf{P}^n_k, \mathcal{O}(1))
\longrightarrow
H^0(\mathbf{P}^n_k, \mathcal{F}(m + 1))
$$
は全射である。
\end{lemma}

\begin{proof}
$k'/k$ を体拡大とする。$\mathcal{F}'$ を Lemma \ref{varieties-lemma-m-regular-extend-base-field} と同様に定める。
Cohomology of Schemes, Lemma \ref{coherent-lemma-flat-base-change-cohomology} により、補題の線型写像を $k'$ へ
基底変換したものは、層 $\mathcal{F}'$ に対する同じ線型写像である。$k \to k'$ は忠実平坦なので、補題を
$k'$ 上で証明すれば十分である。すなわち $k$ は無限であると仮定してよい。

\medskip\noindent
$k$ は無限であると仮定する。$n$ に関する帰納法で補題を証明する。$n = 0$ の場合は
$\mathcal{O}(1) \cong \mathcal{O}$ が $T_0$ により生成されるので自明である。$n > 0$ に対して Lemma
\ref{varieties-lemma-exact-sequence-induction} を適用し、その列を $\mathcal{O}(m + 1)$ とテンソルして
$$
0 \to \mathcal{F}(m) \xrightarrow{s} \mathcal{F}(m + 1) \to
i_*\mathcal{G}(m + 1) \to 0
$$
を得る。Remark \ref{varieties-remark-exact-sequence-induction-cohomology} を参照せよ。
$t \in H^0(\mathbf{P}^n_k, \mathcal{F}(m + 1))$ とする。帰納法により、像
$\overline{t} \in H^0(H, \mathcal{G}(m + 1))$ は $\sum g_i \otimes \overline{s}_i$ の像である。ここで
$\overline{s}_i \in \Gamma(H, \mathcal{O}(1))$ かつ $g_i \in H^0(H, \mathcal{G}(m))$ である。
$\mathcal{F}$ は $m$-正則なので $H^1(\mathbf{P}^n_k, \mathcal{F}(m - 1)) = 0$ である。従って短完全列
$$
0 \to \mathcal{F}(m - 1) \xrightarrow{s} \mathcal{F}(m) \to
i_*\mathcal{G}(m) \to 0
$$
に付随するコホモロジー長完全列から、$g_i$ を
$f_i \in H^0(\mathbf{P}^n_k, \mathcal{F}(m))$ へ持ち上げられる。また $\overline{s}_i$ を
$s_i \in H^0(\mathbf{P}^n_k, \mathcal{O}(1))$ へ持ち上げられる（例えば Lemma \ref{varieties-lemma-hyperplane} の
証明を参照せよ）。$\sum f_i \otimes s_i$ の像を $t$ から引けば、$\overline{t} = 0$ と仮定してよい。
しかしこれは、$t$ が $f \otimes s$ の像であり、ここである
$f \in H^0(\mathbf{P}^n_k, \mathcal{F}(m))$ を取れることをまさに意味する。これで証明が終わる。
\end{proof}

\begin{lemma}
\label{varieties-lemma-m-regular-globally-generated}
$k$ を体、$n \geq 0$ とし、$\mathcal{F}$ を $\mathbf{P}^n_k$ 上のコヒーレント層とする。
$\mathcal{F}$ が $m$-正則ならば、$\mathcal{F}(m)$ は大域生成される。
\end{lemma}

\begin{proof}
すべての $d \gg 0$ に対して層 $\mathcal{F}(d)$ は大域生成される。これは例えば Cohomology of Schemes,
Lemma \ref{coherent-lemma-coherent-projective} の前半から従う。$d \geq m$ を、
$\mathcal{F}(d)$ が大域生成されるように選ぶ。基底 $f_1, \ldots, f_r \in H^0(\mathbf{P}^n_k, \mathcal{F})$ を選ぶ。Lemma
\ref{varieties-lemma-m-regular-multiply} により、任意の元 $f \in H^0(\mathbf{P}^n_k, \mathcal{F}(d))$ は
$f = \sum P_if_i$ と書ける。ここで $P_i \in k[T_0, \ldots, T_n]$ は次数 $d - m$ の斉次多項式である。
切断 $f$ は $\mathcal{F}(d)$ を生成するので、切断 $f_i$ は $\mathcal{F}(m)$ を生成する。
\end{proof}

\subsection{Hilbert 多項式}
\label{varieties-subsection-hilbert}

\noindent
次の補題は、より一般的な Lemma \ref{varieties-lemma-numerical-polynomial-from-euler}
によって後に不要となる。

\begin{lemma}
\label{varieties-lemma-hilbert-polynomial}
$k$ を体、$n \geq 0$ とし、$\mathcal{F}$ を $\mathbf{P}^n_k$ 上のコヒーレント層とする。関数
$$
d \longmapsto \chi(\mathbf{P}^n_k, \mathcal{F}(d))
$$
は多項式である。
\end{lemma}

\begin{proof}
$n$ に関する帰納法で証明する。$n = 0$ ならば
$\mathbf{P}^n_k = \Spec(k)$ かつ $\mathcal{F}(d) = \mathcal{F}$ である。
従ってこの場合、関数は定数、すなわち次数 $0$ の多項式である。$n > 0$ と仮定する。
Lemma \ref{varieties-lemma-euler-characteristic-extend-base-field} により、$k$ は無限であると仮定してよい。
Lemma \ref{varieties-lemma-exact-sequence-induction} を適用する。捩った完全列
$0 \to \mathcal{F}(d - 1) \to \mathcal{F}(d) \to i_*\mathcal{G}(d) \to 0$
に Lemma \ref{varieties-lemma-euler-characteristic-additive} を適用して
$$
\chi(\mathbf{P}^n_k, \mathcal{F}(d)) -
\chi(\mathbf{P}^n_k, \mathcal{F}(d - 1)) =
\chi(H, \mathcal{G}(d))
$$
を得る。Remark \ref{varieties-remark-exact-sequence-induction-cohomology} を参照せよ。
$H \cong \mathbf{P}^{n - 1}_k$ なので、帰納法により右辺は多項式である。
最後に、任意の関数 $f : \mathbf{Z} \to \mathbf{Z}$ について、写像
$d \mapsto f(d) - f(d - 1)$ が多項式ならばその関数自身も多項式であることに注意すればよい。
この事実の証明は省略する（ヒント：Algebra, Lemma
\ref{algebra-lemma-numerical-polynomial} と比較せよ）。
\end{proof}

\begin{definition}
\label{varieties-definition-hilbert-polynomial}
$k$ を体、$n \geq 0$ とし、$\mathcal{F}$ を $\mathbf{P}^n_k$ 上のコヒーレント層とする。関数
$d \mapsto \chi(\mathbf{P}^n_k, \mathcal{F}(d))$ を $\mathcal{F}$ の
{\it Hilbert 多項式}という。
\end{definition}

\noindent
Hilbert 多項式の係数は $\mathbf{Q}$ に属し、一般には $\mathbf{Z}$ には属さない。例えば
$\mathcal{O}_{\mathbf{P}^n_k}$ の Hilbert 多項式は
$$
d \longmapsto {d + n \choose n} = \frac{d^n}{n!} + \ldots
$$
である。これは次の補題と
$$
H^0(\mathbf{P}^n_k, \mathcal{O}_{\mathbf{P}^n_k}(d)) = k[T_0, \ldots, T_n]_d
$$
から従う。この空間は次数 $d$ の部分であり、その $k$ 上の次元は ${d + n \choose n}$ である。

\begin{lemma}
\label{varieties-lemma-hilbert-polynomial-H0}
$k$ を体、$n \geq 0$ とし、$\mathcal{F}$ を $\mathbf{P}^n_k$ 上の、Hilbert 多項式
$P \in \mathbf{Q}[t]$ をもつコヒーレント層とする。このとき
$$
P(d) = \dim_k H^0(\mathbf{P}^n_k, \mathcal{F}(d))
$$
がすべての $d \gg 0$ に対して成り立つ。
\end{lemma}

\begin{proof}
これは $\mathcal{F}$ の十分高い捩れのコホモロジーが消滅することから従う。
Cohomology of Schemes, Lemma
\ref{coherent-lemma-coherent-projective} を参照せよ。
\end{proof}



\subsection{商の有界性}
\label{varieties-subsection-boundedness}

\noindent
本小節では、$\mathbf{P}^n$ 上の与えられたコヒーレント層の商の正則性を
Hilbert 多項式によって評価する。

\begin{lemma}
\label{varieties-lemma-bound-quotients-free}
$k$ を体、$n \geq 0$、$r \geq 1$ とし、$P \in \mathbf{Q}[t]$ とする。
整数 $m$ が存在し、これは $n$、$r$、$P$ のみに依存して、次の性質をもつ。すなわち
$$
0 \to \mathcal{K} \to \mathcal{O}^{\oplus r} \to \mathcal{F} \to 0
$$
が $\mathbf{P}^n_k$ 上のコヒーレント層の短完全列であり、$\mathcal{F}$ の Hilbert 多項式が $P$ ならば、
$\mathcal{K}$ は $m$-正則である。
\end{lemma}

\begin{proof}
$n$ に関する帰納法で証明する。$n = 0$ ならば
$\mathbf{P}^n_k = \Spec(k)$ であり、任意のコヒーレント加群は $0$-正則で、任意の全射は大域切断上でも全射である。
$n > 0$ と仮定し、補題にある完全列を考える。多項式 $P' \in \mathbf{Q}[t]$ を
$P'(t) = P(t) - P(t - 1)$ と置く。整数 $m'$ を、$n - 1$、$r$、$P'$ に対して成り立つものとする。
Lemmas \ref{varieties-lemma-m-regular-extend-base-field} および
\ref{varieties-lemma-euler-characteristic-extend-base-field} により、$k$ を体拡大で置き換えられるので、
$k$ は無限であると仮定してよい。コヒーレント層 $\mathcal{F}$ に
Lemma \ref{varieties-lemma-exact-sequence-induction} を適用する。
$\mathcal{F}' = i^*\mathcal{F}$ の Hilbert 多項式は $P'$ である
（Lemma \ref{varieties-lemma-hilbert-polynomial} の証明を参照せよ）。
$i^*$ は右完全なので、$\mathcal{F}'$ は
$\mathcal{O}_H^{\oplus r} = i^*\mathcal{O}^{\oplus r}$ の商である。
従って帰納法の仮定を $\mathcal{F}'$ に、$H \cong \mathbf{P}^{n - 1}_k$ 上で適用できる
（Lemma \ref{varieties-lemma-hyperplane}）。$\mathcal{K}(-1) \to \mathcal{K}$ は、
$\mathcal{K} \subset \mathcal{O}^{\oplus r}$ なので単射であり、その余核は
$i_*\mathcal{H}$ である。ここで $\mathcal{H} = i^*\mathcal{K}$ とする。
蛇の補題（Homology, Lemma \ref{homology-lemma-snake}）により、列と行が完全な可換図式
$$
\xymatrix{
& 0 \ar[d] & 0 \ar[d] & 0 \ar[d] \\
0 \ar[r] &
\mathcal{K}(-1) \ar[r] \ar[d] &
\mathcal{O}^{\oplus r}(-1) \ar[r] \ar[d] &
\mathcal{F}(-1) \ar[d] \ar[r] & 0 \\
0 \ar[r] &
\mathcal{K} \ar[r] \ar[d] &
\mathcal{O}^{\oplus r} \ar[r] \ar[d] &
\mathcal{F} \ar[d] \ar[r] & 0\\
0 \ar[r] &
i_*\mathcal{H} \ar[r] \ar[d] &
i_*\mathcal{O}_H^{\oplus r} \ar[r] \ar[d] &
i_*\mathcal{F}' \ar[r] \ar[d] & 0 \\
& 0 & 0 & 0
}
$$
を得る。従って帰納法の仮定を完全列
$0 \to \mathcal{H} \to \mathcal{O}_H^{\oplus r} \to \mathcal{F}' \to 0$
に、$H \cong \mathbf{P}^{n - 1}_k$ 上で適用でき（Lemma \ref{varieties-lemma-hyperplane}）、$\mathcal{H}$ は $m'$-正則である。
これは、$\mathcal{H}$ の $d$-正則性がすべての $d \geq m'$ に対して成り立つことを意味する
（Lemma \ref{varieties-lemma-m-regular-up}）。

\medskip\noindent
$i \geq 2$ かつ $d \geq m'$ とする。上の図式の左列に付随するコホモロジー長完全列と
$H^{i - 1}(H, \mathcal{H}(d))$ の消滅から、写像
$$
H^i(\mathbf{P}^n_k, \mathcal{K}(d - 1))
\longrightarrow
H^i(\mathbf{P}^n_k, \mathcal{K}(d))
$$
が単射であることが従う。これらの群は $d \gg 0$ で零になるので
（Cohomology of Schemes, Lemma
\ref{coherent-lemma-coherent-projective}）、
$H^i(\mathbf{P}^n_k, \mathcal{K}(d))$ はすべての $d \geq m'$ および $i \geq 2$ に対して零である。

\medskip\noindent
なお $H^1$ を制御しなければならない。まず、すべての写像
$$
H^1(\mathbf{P}^n_k, \mathcal{K}(m' - 1)) \to
H^1(\mathbf{P}^n_k, \mathcal{K}(m')) \to
H^1(\mathbf{P}^n_k, \mathcal{K}(m' + 1)) \to \ldots
$$
は、$H^1(H, \mathcal{H}(d))$ が $d \geq m'$ に対して消滅することにより全射である。
$d > m'$ が
$$
H^1(\mathbf{P}^n_k, \mathcal{K}(d - 1))
\longrightarrow
H^1(\mathbf{P}^n_k, \mathcal{K}(d))
$$
を単射にすると仮定する。このとき
$H^0(\mathbf{P}^n_k, \mathcal{K}(d)) \to H^0(H, \mathcal{H}(d))$
は全射である。可換図式
$$
\xymatrix{
H^0(\mathbf{P}^n_k, \mathcal{K}(d)) \otimes_k
H^0(\mathbf{P}^n_k, \mathcal{O}(1))
\ar[r] \ar[d] &
H^0(\mathbf{P}^n_k, \mathcal{K}(d + 1)) \ar[d] \\
H^0(H, \mathcal{H}(d)) \otimes_k
H^0(H, \mathcal{O}_H(1))
\ar[r] &
H^0(H, \mathcal{H}(d + 1))
}
$$
を考える。Lemma \ref{varieties-lemma-m-regular-multiply} により、下の水平射は全射である。
従って右の垂直射は全射である。従って
$$
H^1(\mathbf{P}^n_k, \mathcal{K}(d))
\longrightarrow
H^1(\mathbf{P}^n_k, \mathcal{K}(d + 1))
$$
は単射である。帰納法により
$$
H^1(\mathbf{P}^n_k, \mathcal{K}(d - 1)) \to
H^1(\mathbf{P}^n_k, \mathcal{K}(d)) \to
H^1(\mathbf{P}^n_k, \mathcal{K}(d + 1)) \to \ldots
$$
はすべて単射である。これらの群は最終的に消滅するので
$H^1(\mathbf{P}^n_k, \mathcal{K}(d - 1)) = 0$ と結論する。
従って群 $H^1(\mathbf{P}^n_k, \mathcal{K}(d))$ の次元は、$d \geq m'$ に対して零になるまで狭義単調減少する。
ゆえに $\mathcal{K}$ の正則性は
$m' + \dim_k H^1(\mathbf{P}^n_k, \mathcal{K}(m'))$ で上から抑えられる。
一方、高次コホモロジー群の消滅により
$$
\dim_k H^1(\mathbf{P}^n_k, \mathcal{K}(m')) = 
- \chi(\mathbf{P}^n_k, \mathcal{K}(m')) +
\dim_k H^0(\mathbf{P}^n_k, \mathcal{K}(m'))
$$
である。ここで $H^0$ の次元は
$H^0(\mathbf{P}^n_k, \mathcal{O}^{\oplus r}(m'))$ の次元で上から抑えられ、これは
これは高々 $r{n + m' \choose n}$（$m' > 0$ の場合）であり、そうでなければ零である。
最後に、項 $\chi(\mathbf{P}^n_k, \mathcal{K}(m'))$ は
$r{n + m' \choose n} - P(m')$ に等しい。これで所望の型の上界が得られ、補題の証明が終わる。
\end{proof}

\section{Frobenius 射}
\label{varieties-section-frobenius}

\noindent
$p$ を素数とする。$X$ がスキームならば、$X$ は「標数 $p$ をもつ」、または
「$X$ は標数 $p$ である」、または「$X$ は標数 $p$ の中にある」という。ただしこれは、
$p$ が $\mathcal{O}_X$ の中で零であることを意味する。

\begin{definition}
\label{varieties-definition-absolute-frobenius}
$p$ を素数とし、$X$ を標数 $p$ のスキームとする。{\it $X$ の絶対 Frobenius 射}とは、射
$F_X : X \to X$ であって、台となる位相空間上では恒等写像により与えられ、
$F_X^\sharp : \mathcal{O}_X \to \mathcal{O}_X$ が $g \mapsto g^p$ で与えられるものをいう。
\end{definition}

\noindent
この定義が意味をもつのは、任意の環 $A$ で標数が $p$ であるものに対し、写像
$F_A : A \to A$、$a \mapsto a^p$ が $\Spec(A)$ 上に恒等写像を誘導する環自己準同型だからである。
さらに $A$ が局所環ならば $F_A$ は局所準同型である。このようにして、$X$ の絶対 Frobenius 射は
スキームの圏における $X$ の自己準同型であることが分かる。実際、絶対 Frobenius 射は
標数 $p$ のスキームの圏上の恒等関手の自己写像を定める。

\begin{lemma}
\label{varieties-lemma-frobenius-endomorphism-identity}
$p > 0$ を素数とする。$f : X \to Y$ を標数 $p$ のスキームの射とする。このとき図式
$$
\xymatrix{
X \ar[d]_f \ar[r]_{F_X} & X \ar[d]^f \\
Y \ar[r]^{F_Y} & Y
}
$$
は可換である。
\end{lemma}

\begin{proof}
これは次の自明な代数的事実から従う。$\varphi : A \to B$ が標数 $p$ の環の準同型ならば、
$\varphi(a^p) = \varphi(a)^p$ である。
\end{proof}

\begin{lemma}
\label{varieties-lemma-frobenius}
$p > 0$ を素数とし、$X$ を標数 $p$ のスキームとする。このとき絶対 Frobenius 射
$F_X : X \to X$ は普遍同相であり、整であり、純非分離な剰余体拡大を誘導する。
\end{lemma}

\begin{proof}
これは Frobenius 自己準同型 $F_A : A \to A$ に対する対応する結果から従う。ここで $A$ は標数 $p > 0$ の環である。
例えば Algebra, Section \ref{algebra-section-universal-homeomorphism}、特に
Lemma \ref{algebra-lemma-p-ring-map} を参照せよ。
\end{proof}

\noindent
固定した基底上のスキームを扱う場合には、Frobenius 射の相対版がある。

\begin{definition}
\label{varieties-definition-relative-frobenius}
$p > 0$ を素数とし、$S$ を標数 $p$ のスキームとする。$X$ を $S$ 上のスキームとする。
$$
X^{(p)} = X^{(p/S)} = X \times_{S, F_S} S
$$
を $S$ 上のスキームとして定義する。Lemma \ref{varieties-lemma-frobenius-endomorphism-identity} を適用すると、
次の可換図式に入る一意な射 $F_{X/S} : X \longrightarrow X^{(p)}$ が $S$ 上に存在する。
$$
\xymatrix{
X \ar[rr]_{F_{X/S}} \ar[rrd] \ar@/^1em/[rrrr]^{F_X}
& & X^{(p)} \ar[rr] \ar[d] & & X \ar[d] \\
& & S \ar[rr]^{F_S} & & S
}
$$
ここで右の正方形は Cartesian である。射 $F_{X/S}$ を
{\it $X/S$ の相対 Frobenius 射}という。
\end{definition}

\noindent
$X \mapsto X^{(p)}$ は関手であることに注意せよ。これは絶対 Frobenius 射
$F_S : S \to S$ による基底変換関手である。相対 Frobenius 射についても、先ほどと同じ補題が成り立つ。

\begin{lemma}
\label{varieties-lemma-relative-frobenius-endomorphism-identity}
$p > 0$ を素数とし、$S$ を標数 $p$ のスキームとする。
$f : X \to Y$ を $S$ 上のスキームの射とする。このとき図式
$$
\xymatrix{
X \ar[d]_f \ar[r]_{F_{X/S}} & X^{(p)} \ar[d]^{f^{(p)}} \\
Y \ar[r]^{F_{Y/S}} & Y^{(p)}
}
$$
は可換である。
\end{lemma}

\begin{proof}
これは Lemma \ref{varieties-lemma-frobenius-endomorphism-identity} と定義から従う。
\end{proof}

\begin{lemma}
\label{varieties-lemma-relative-frobenius}
$p > 0$ を素数とし、$S$ を標数 $p$ のスキームとする。$X$ を $S$ 上のスキームとする。
このとき相対 Frobenius 射 $F_{X/S} : X \to X^{(p)}$ は普遍同相であり、整であり、
純非分離な剰余体拡大を誘導する。
\end{lemma}

\begin{proof}
Lemma \ref{varieties-lemma-frobenius} により、射 $F_X : X \to X$ と射
$h : X^{(p)} \to X$（これは $F_S$ の基底変換）は普遍同相である。
$h \circ F_{X/S} = F_X$ なので、$F_{X/S}$ は普遍同相である
（Morphisms, Lemma \ref{morphisms-lemma-permanence-universal-homeomorphism}）。
Morphisms, Lemmas \ref{morphisms-lemma-universal-homeomorphism} および
\ref{morphisms-lemma-universally-injective} により、$F_{X/S}$ は他の性質ももつ。
\end{proof}

\begin{lemma}
\label{varieties-lemma-relative-frobenius-omega}
$p > 0$ を素数とし、$S$ を標数 $p$ のスキーム、$X$ を $S$ 上のスキームとする。このとき
$\Omega_{X/S} = \Omega_{X/X^{(p)}}$ である。
\end{lemma}

\begin{proof}
これは次の代数的事実に翻訳される。$A \to B$ を標数 $p$ の環の準同型とする。
$B' = B \otimes_{A, F_A} A$ と置き、環写像
$F_{B/A} : B' \to B$、$b \otimes a \mapsto b^pa$ を考える。このとき主張は
$\Omega_{B/A} = \Omega_{B/B'}$ である。実際、$\text{d}(b^pa) = 0$ である。ただし
$\text{d} : B \to \Omega_{B/A}$ は普遍導分であり、従って $\text{d}$ は $B'$-導分である。
\end{proof}

\begin{lemma}
\label{varieties-lemma-relative-frobenius-finite}
$p > 0$ を素数とし、$S$ を標数 $p$ のスキーム、$X$ を $S$ 上のスキームとする。
$X \to S$ が局所有限型ならば、$F_{X/S}$ は有限である。
\end{lemma}

\begin{proof}
これは次の代数的事実に翻訳される。$A \to B$ を標数 $p$ の環の有限型準同型とする。
$B' = B \otimes_{A, F_A} A$ と置き、環写像
$F_{B/A} : B' \to B$、$b \otimes a \mapsto b^pa$ を考える。このとき主張は
$F_{B/A}$ が有限であることである。実際、$x_1, \ldots, x_n \in B$ が $A$ 上の生成元ならば、
$x_i$ は $B'$ 上整である。なぜなら $x_i^p = F_{B/A}(x_i \otimes 1)$ だからである。
従って Algebra, Lemma \ref{algebra-lemma-characterize-finite-in-terms-of-integral} により
$F_{B/A} : B' \to B$ は有限である。
\end{proof}

\begin{lemma}
\label{varieties-lemma-geometrically-reduced-p}
$k$ を標数 $p > 0$ の体とし、$X$ を $k$ 上のスキームとする。このとき
$X$ が幾何学的被約であることと $X^{(p)}$ が被約であることは同値である。
\end{lemma}

\begin{proof}
絶対 Frobenius 射 $F_k : k \to k$ を考える。このとき $F_k(k) = k^p$ である。言い換えれば、
$F_k : k \to k$ は $k$ から $k^{1/p}$ への埋め込みと同型である。従って補題は
Lemma \ref{varieties-lemma-geometrically-reduced} から従う。
\end{proof}

\begin{lemma}
\label{varieties-lemma-inseparable-deg-p-smooth}
$k$ を標数 $p > 0$ の体とし、$X$ を $k$ 上の多様体とする。次は同値である。
\begin{enumerate}
\item $X^{(p)}$ は被約である。
\item $X$ は幾何学的被約である。
\item 空でない開集合 $U \subset X$ が存在し、これは $k$ 上平滑である。
\end{enumerate}
この場合 $X^{(p)}$ は $k$ 上の多様体であり、
$F_{X/k} : X \to X^{(p)}$ は次数 $p^{\dim(X)}$ の有限支配射である。
\end{lemma}

\begin{proof}
(1) と (2) の同値性は Lemma \ref{varieties-lemma-geometrically-reduced-p} で示した。
(2) が (3) を含意することは Lemma \ref{varieties-lemma-geometrically-reduced-dense-smooth-open} で示した。
(3) が成り立つならば $U$ は幾何学的被約であり
（例えば Lemma \ref{varieties-lemma-geometrically-regular-smooth} を参照せよ）、従って
Lemma \ref{varieties-lemma-generic-points-geometrically-reduced} により $X$ は幾何学的被約である。
こうして (1)、(2)、(3) は同値である。

\medskip\noindent
(1)、(2)、(3) が成り立つと仮定する。$F_{X/k}$ は同相写像なので
（Lemma \ref{varieties-lemma-relative-frobenius}）、$X^{(p)}$ は多様体である。
次に Lemma \ref{varieties-lemma-relative-frobenius-finite} により $F_{X/k}$ は有限である。
これは全射なので支配的である。次数を計算するには
（Morphisms, Definition \ref{morphisms-definition-degree}）、
$F_{U/k} : U \to U^{(p)}$ の次数を計算すれば十分である
（Lemma \ref{varieties-lemma-relative-frobenius-endomorphism-identity} により
$F_{U/k} = F_{X/k}|_U$ である）。$U$ を少し縮小すれば、
エタール射 $h : U \to \mathbf{A}^n_k$ が存在すると仮定してよい。
Morphisms, Lemma \ref{morphisms-lemma-smooth-etale-over-affine-space} を参照せよ。
もちろん $n = \dim(U)$ である。実際、
$\mathbf{A}^n_k \to \Spec(k)$ は相対次元 $n$ の平滑射であり、エタール射 $h$ は相対次元 $0$ の平滑射であり、
$U \to \Spec(k)$ は相対次元 $\dim(U)$ の平滑射であって、相対次元は正しく加算される
（Morphisms, Lemma \ref{morphisms-lemma-composition-relative-dimension-d}）。
$h$ は多様体の生成的に有限な支配射なので、$\deg(h)$ が定義されることに注意せよ。
Lemma \ref{varieties-lemma-relative-frobenius-endomorphism-identity} により、可換図式
$$
\xymatrix{
U \ar[rr]_{F_{X/k}} \ar[d]_h & & U^{(p)} \ar[d]^{h^{(p)}} \\
\mathbf{A}^n_k \ar[rr]^{F_{\mathbf{A}^n_k/k}} & &
(\mathbf{A}^n_k)^{(p)}
}
$$
を得る。$h^{(p)}$ は $h$ の基底変換なので、これもエタールであり、従って
$h^{(p)}$ も多様体の生成的に有限な支配射である。$h^{(p)}$ の次数は体拡大
$k(X^{(p)})/k((\mathbf{A}^n_k)^{(p)})$ の次数であり、これは体拡大
$k(X)/k(\mathbf{A}^n_k)$ の次数と同じである。なぜなら $h^{(p)}$ は $h$ の基底変換だからである
（細部は省略する）。次数の乗法性
（Morphisms, Lemma \ref{morphisms-lemma-degree-composition}）により、
$F_{\mathbf{A}^n_k/k}$ の次数が $p^n$ であることを示せば十分である。このため
$(\mathbf{A}^n_k)^{(p)} = \mathbf{A}^n_k$ であり、
$F_{\mathbf{A}^n_k/k}$ が座標をその $p$ 乗へ送る写像によって与えられることに注意せよ。
\end{proof}

\begin{remark}
\label{varieties-remark-n-fold-relative-frobenius}
$p > 0$ を素数とし、$S$ を標数 $p$ のスキーム、$X$ を $S$ 上のスキームとする。$n \geq 1$ に対して
$$
X^{(p^n)} = X^{(p^n/S)} = X \times_{S, F_S^n} S
$$
を $S$ 上のスキームとみなす。$X \mapsto X^{(p^n)}$ は関手であることに注意せよ。
Lemma \ref{varieties-lemma-frobenius-endomorphism-identity} を適用すると、
$F_{X/S, n} = (F_X^n, \text{id}_S) : X \longrightarrow X^{(p^n)}$
は、次の可換図式に入る $S$ 上の射である。
$$
\xymatrix{
X \ar[rr]_{F_{X/S, n}} \ar[rrd] \ar@/^1em/[rrrr]^{F_X^n}
& & X^{(p^n)} \ar[rr] \ar[d] & & X \ar[d] \\
& & S \ar[rr]^{F_S^n} & & S
}
$$
ここで右の正方形は Cartesian である。射 $F_{X/S, n}$ は
{\it $n$ 重相対 Frobenius 射（$X/S$ に対する）}と呼ばれることがある。この呼び方が意味をもつのは、公式
$$
F_{X/S, n} =
F_{X^{(p^{n - 1})}/S} \circ \ldots \circ F_{X^{(p)}/S} \circ F_{X/S}
$$
があり、$F_{X/S, n}$ が $n$ 個の相対 Frobenius 射の合成であることを示すからである。さらに
$$
F_{X^{(p^m)}/S} = F_{X^{(p^{m - 1})}/S}^{(p)} = \ldots = F_{X/S}^{(p^m)}
$$
であるから（細部は省略する）、次も得る。
$$
F_{X/S, n} =
F_{X/S}^{(p^{n - 1})} \circ \ldots \circ F_{X/S}^{(p)} \circ F_{X/S}
$$
\end{remark}

\section{1 次元環の貼り合わせ}
\label{varieties-section-glueing-dim-1}

\noindent
本節では、分離スキームの余次元 $1$ の点の有限集合がアフィン開集合に含まれることを証明するための
代数的準備を行う。

\begin{situation}
\label{varieties-situation-glue}
ここでは環の可換図式
$$
\xymatrix{
A \ar[r] & K \\
R \ar[u] \ar[r] & B \ar[u]
}
$$
が与えられているとする。ここで $K$ は体であり、$A$、$B$ は分数体 $K$ をもつ $K$ の部分環である。
最後に $R = A \times_K B = A \cap B$ とする。
\end{situation}

\begin{lemma}
\label{varieties-lemma-glue-valuation-ring}
Situation \ref{varieties-situation-glue} において $B$ が付値環であると仮定する。このとき単元 $u$ を $A$ の任意の単元とすれば、
$u \in R$ または $u^{-1} \in R$ のいずれかが成り立つ。
\end{lemma}

\begin{proof}
実際、像 $c$、すなわち $u$ の $K$ における像が $B$ に属すれば $u \in R$ である。そうでなければ
$c^{-1} \in B$ であり
（Algebra, Lemma \ref{algebra-lemma-valuation-ring-x-or-x-inverse}）、
$u^{-1} \in R$ である。
\end{proof}

\noindent
次の補題は、以下でしばしば現れる条件「$A \otimes B \to K$ が全射である」の意味を説明する。

\begin{lemma}
\label{varieties-lemma-glue-separated}
Situation \ref{varieties-situation-glue} において、$A$ が次元 $1$ の Noether 環であると仮定する。次は同値である。
\begin{enumerate}
\item $A \otimes B \to K$ は全射でない。
\item 離散付値環 $\mathcal{O} \subset K$ が存在し、これは $A$ と $B$ の両方を含む。
\end{enumerate}
\end{lemma}

\begin{proof}
(2) が (1) を含意することは明らかである。逆に $A \otimes B \to K$ が全射でないとすると、その像
$C \subset K$ は体でないので、$C$ は非零極大イデアル $\mathfrak m$ をもつ。
付値環 $\mathcal{O} \subset K$ で $C_\mathfrak m$ を支配するものを選ぶ。
$A \subset \mathcal{O}$ に Algebra, Lemma \ref{algebra-lemma-krull-akizuki} を適用すると、
環 $\mathcal{O}$ は Noether である。従って Algebra, Lemma
\ref{algebra-lemma-valuation-ring-Noetherian-discrete} により $\mathcal{O}$ は離散付値環である。
\end{proof}

\begin{lemma}
\label{varieties-lemma-semi-local}
Situation \ref{varieties-situation-glue} において、次を仮定する。
\begin{enumerate}
\item $A$ は次元 $1$ の Noether 半局所整域である。
\item $B$ は離散付値環である。
\end{enumerate}
このとき次の二つの場合がある。
\begin{enumerate}
\item[(a)] $A^*$ が $R$ に含まれないならば、
$\Spec(A) \to \Spec(R)$ と $\Spec(B) \to \Spec(R)$ は $\Spec(R)$ を被覆する開埋め込みであり、
$K = A \otimes_R B$ である。
\item[(b)] $A^*$ が $R$ に含まれるならば、$B$ は $A$ のある極大イデアルにおける局所環を支配し、
$A \otimes B \to K$ は全射でない。
\end{enumerate}
\end{lemma}

\begin{proof}
仮定 (a) により、単元 $u$ を $A$ から選んで、その $K$ における像を $B$ の極大イデアルに属するようにできる。
このとき $u$ は $R$ の零因子でない元であり、任意の $a \in A$ に対して、ある $n$ が存在して
$u^n a \in R$ となる。従って $A = R_u$ である。

\medskip\noindent
$\mathfrak m_A$ を $A$ の Jacobson 根基とし、$x \in \mathfrak m_A$ を非零元とする。
$\dim(A) = 1$ なので $K = A_x$ である。$x$ を $x^n u^m$ で置き換えることにより
（ここで $n \geq 1$ かつ $m \in \mathbf{Z}$）、$x$ の $B$ における像が単元であると仮定してよい。
任意の $b \in B$ に対して、$x^nb$ が $R$ の像に属するような $n$ が存在する。従って $B = R_x$ である。

\medskip\noindent
$z \in R$ とする。$z \not \in \mathfrak m_A$ であり、かつ $z$ の像が $\mathfrak m_B$ に属さなければ、
$z$ は可逆である。従って $x + u$ は $R$ で可逆である。ゆえに
$\Spec(R) = D(x) \cup D(u)$ である。上で $D(u) = \Spec(A)$ および $D(x) = \Spec(B)$ を示した。

\medskip\noindent
場合 (b) を考える。$x \in \mathfrak m_A$ ならば $1 + x$ は単元であり、
従って $1 + x \in R$、すなわち $x \in R$ である。ゆえに
$\mathfrak m_A \subset R \subset A$ である。実際、この場合 $A$ は $R$ 上整である。すなわち
$A/\mathfrak m_A = \kappa_1 \times \ldots \times \kappa_n$
を体の積として書く。元
$x = (c_1, \ldots, c_r, 0, \ldots, 0)$ を、$c_i \not = 0$ となるように取る。このとき
$$
x^2 - x(c_1, \ldots, c_r, 1, \ldots, 1)  = 0
$$
である。$R$ はすべての単元を含むので、$A/\mathfrak m_A$ はその中の $R$ の像の上で整であり、
従って $A$ は $R$ 上整である。従って $R \subset A \subset B$ であり、これは $B$ が整閉だからである。
さらに $x \in \mathfrak m_A$ が非零ならば
$K = A_x = \bigcup x^{-n}A = \bigcup x^{-n}R$ である。
従って $x^{-1} \not \in B$、すなわち $x \in \mathfrak m_B$ である。
$\mathfrak m_A \subset \mathfrak m_B$ と結論する。従って
$A \cap \mathfrak m_B$ は $A$ の極大イデアルであり、証明が終わる。
\end{proof}

\begin{lemma}
\label{varieties-lemma-semi-local-dimension-one-conductor}
$B$ を次元 $1$ の Noether 半局所整域とし、$B'$ をその分数体における $B$ の整閉包とする。
このとき $B'$ は半局所 Dedekind 整域である。$x$ を $B'$ の Jacobson 根基の非零元とする。
このとき任意の $y \in B'$ に対して、ある $n$ が存在して $x^n y \in B$ となる。
\end{lemma}

\begin{proof}
$\mathfrak m_B$ を $B$ の Jacobson 根基とする。$B'$ の構造は
Algebra, Lemma \ref{algebra-lemma-integral-closure-Dedekind} から従う。
補題の主張にある $x, y \in B'$ が与えられたとし、部分環
$B \subset A \subset B'$、すなわち $x$ と $y$ で生成される部分環を考える。このとき $A$ は $B$ 上有限である
（Algebra, Lemma \ref{algebra-lemma-characterize-finite-in-terms-of-integral}）。
$B$ と $A$ の分数体は等しいので、有限加群 $A/B$ は $B$ の閉点集合上に台をもつ。
従って $\mathfrak m_B^n A \subset B$ となる適当な $n$ がある。さらに
$\Spec(B') \to \Spec(A)$ は全射なので
（Algebra, Lemma \ref{algebra-lemma-integral-overring-surjective}）、$A$ も半局所である。
また $x$ は Jacobson 根基 $\mathfrak m_A$、すなわち $A$ の Jacobson 根基に属する。ここで
$\mathfrak m_A = \sqrt{\mathfrak m_B A}$ であることに注意せよ。
従って $x^m y \in \mathfrak m_B A$ となる $m$ があり、
$x^{nm} y \in B$ となる。
\end{proof}

\begin{lemma}
\label{varieties-lemma-semi-local-both-side}
Situation \ref{varieties-situation-glue} において、次を仮定する。
\begin{enumerate}
\item $A$ は次元 $1$ の Noether 半局所整域である。
\item $B$ は次元 $1$ の Noether 半局所整域である。
\item $A \otimes B \to K$ は全射である。
\end{enumerate}
このとき $\Spec(A) \to \Spec(R)$ と $\Spec(B) \to \Spec(R)$ は
$\Spec(R)$ を被覆する開埋め込みであり、$K = A \otimes_R B$ である。
\end{lemma}

\begin{proof}
まず $B$ が $K$ において整閉である特別な場合を考える。これは
$B$ が Dedekind 整域であることを意味する
（Algebra, Lemma \ref{algebra-lemma-characterize-Dedekind}）。
従って、その極大イデアルにおける局所化はすべて離散付値環である。
$\mathfrak m_1, \ldots, \mathfrak m_r$ を $B$ の極大イデアルとする。次のように置く。
$$
R_1 = A \times_K B_{\mathfrak m_1}
$$
$A \otimes_{R_1} B_{\mathfrak m_1} \to K$ が全射であることに注意すると、
Lemma \ref{varieties-lemma-semi-local} より、$A$ と
$B_{\mathfrak m_1}$ は $\Spec(R_1)$ を被覆する開部分スキームを定め、かつ
$K = A \otimes_{R_1} B_{\mathfrak m_1}$ である。特に $R_1$ は
次元 $1$ の Noether 半局所環である。帰納的に
$$
R_{i + 1} = R_i \times_K B_{\mathfrak m_{i + 1}}
$$
を $i = 1, \ldots, r - 1$ に対して定める。
$R = R_r$ であることに注意せよ。これは
$B = B_{\mathfrak m_1} \cap \ldots \cap B_{\mathfrak m_r}$ だからである（Algebra, Lemma
\ref{algebra-lemma-normal-domain-intersection-localizations-height-1} 参照）。
この帰納的手順から、$R \to A$ は開埋め込み
$\Spec(A) \to \Spec(R)$ を定める。一方、この開部分に属さない
極大イデアル $\mathfrak n_i$（$R$ のもの）は
極大イデアル $\mathfrak m_i$（$B$ のもの）に対応し、実際、環準同型
$R \to B$ は同型
$R_{\mathfrak n_i} \to B_{\mathfrak m_i}$ を定める
（細部は省略する。各段階で $\Spec(R_i)$ にちょうど一つの極大イデアルを加えたことを用いる）。
従って $\Spec(B) \to \Spec(R)$ は望むように開埋め込みである。

\medskip\noindent
一般の場合を考える。$B' \subset K$ を $B$ の整閉包とする。
Lemma \ref{varieties-lemma-semi-local-dimension-one-conductor} を参照せよ。
このとき特別な場合を $R' = A \times_K B'$ に適用できる。
$x \in R'$ を、$A$ のどの極大イデアルにも属さず、
$B'$ のすべての極大イデアルに属するように選ぶ
（Algebra, Lemma \ref{algebra-lemma-chinese-remainder} 参照）。
Lemma \ref{varieties-lemma-semi-local-dimension-one-conductor} により、
整数 $n$ が存在して $x^n \in R = A \times_K B$ となる。
$x$ を $x^n$ で置き換え、$x \in R$ とする。任意の $y \in R'$ に対し、
整数 $n$ が存在して $x^n y \in R$ となる。一方、
$R'_x = A$ は明らかである。従って $R_x = A$ である。
$A$ と $B$ の役割を交換すれば、$y \in R$ であって $B = R_y$ となるものも得られる。
$x$ と $y$ の両方を逆元化すると $(0)$ 以外の素イデアルは残らない。従って
$K = R_{xy} = R_x \otimes_R R_y$ である。これで証明が終わる。
\end{proof}

\begin{lemma}
\label{varieties-lemma-glue-a-bunch-of-local-rings}
$K$ を体とする。$A_1, \ldots, A_r \subset K$ を、次元 $1$ の Noether 半局所環であって、
分数体が $K$ であるものとする。$A_i \otimes A_j \to K$ が
すべての $i \not = j$ に対して全射ならば、Noether 半局所整域
$A \subset K$ であって、次元 $1$ かつ $A_1, \ldots, A_r$ に含まれるものが存在し、次を満たす。
\begin{enumerate}
\item $A \to A_i$ は開埋め込み $j_i : \Spec(A_i) \to \Spec(A)$ を誘導する。
\item $\Spec(A)$ は開集合 $j_i(\Spec(A_i))$ の合併である。
\item $\Spec(A)$ の各閉点は、これらの開集合のうちちょうど一つに属する。
\end{enumerate}
\end{lemma}

\begin{proof}
実際、$A = A_1 \cap \ldots \cap A_r$ と取ればよい。まず、(1) と (2) が証明されたならば、
(3) は $A_i \otimes A_j \to K$ が全射であるという仮定から従うことに注意する。というのも、
$\mathfrak m \in j_i(\Spec(A_i)) \cap j_j(\Spec(A_j))$ ならば、
$A_i \otimes A_j \to K$ の像は $A_\mathfrak m$ に入るからである。
(1) と (2) を示すため、$r$ に関する帰納法を用いる。
$r > 1$ ならば、帰納法により
$B = A_2 \cap \ldots \cap A_r$ に対して (1) と (2) が成り立つ。
そこで Lemma \ref{varieties-lemma-semi-local-both-side} を適用すると、
$A = A_1 \cap B$ に対しても成り立つことが分かる。
\end{proof}

\begin{lemma}
\label{varieties-lemma-create-globally-generated}
$A$ を分数体 $K$ をもつ整域とする。$B_1, \ldots, B_r \subset K$ を、
Noether な $1$ 次元半局所整域であって、分数体が $K$ であるものとする。
$A \otimes B_i \to K$ が $i = 1, \ldots, r$ に対して全射ならば、
$x \in A$ であって、$x^{-1}$ が $B_i$ の Jacobson 根基に
$i = 1, \ldots, r$ のそれぞれについて属するものが存在する。
\end{lemma}

\begin{proof}
$B_i'$ を $B_i$ の $K$ における整閉包とする。非零元
$x \in A$ であって、$x^{-1}$ が $B'_i$ の Jacobson 根基に
$i = 1, \ldots, r$ のそれぞれについて属するものが見つかったと仮定する。このとき
Lemma \ref{varieties-lemma-semi-local-dimension-one-conductor} により、
$x$ をある冪で置き換えれば $x^{-1} \in B_i$ となる。
$\Spec(B'_i) \to \Spec(B_i)$ は全射なので、$x^{-1}$ は
$B_i$ の Jacobson 根基にも属する。従って各 $B_i$ が半局所 Dedekind 整域であると仮定してよい。

\medskip\noindent
$B_i$ が局所環でなければ、$B_i$ をリストから除き、有限個の局所環
$(B_i)_\mathfrak m$ を代わりに加える。従って
$B_i$ が $i = 1, \ldots, r$ に対して離散付値環であると仮定してよい。

\medskip\noindent
$v_i : K \to \mathbf{Z}$（$i = 1, \ldots, r$）を対応する離散付値とする
（Algebra, Lemma \ref{algebra-lemma-characterize-Dedekind} 参照）。
求めるのは、非零元 $x \in A$ であって、$v_i(x) < 0$ が
$i = 1, \ldots, r$ に対して成り立つものである。これを $r$ に関する帰納法で証明する。

\medskip\noindent
$r = 1$ で結論が偽ならば、$A \subset B$ となり、写像
$A \otimes B \to K$ は全射でない。これは矛盾である。

\medskip\noindent
$r > 1$ ならば、帰納法により、非零元 $x \in A$ であって
$v_i(x) < 0$ が $i = 1, \ldots, r - 1$ に対して成り立つものを取れる。
$v_r(x) < 0$ ならば終わりなので、$v_r(x) \geq 0$ と仮定してよい。
基底の場合により、非零元 $y \in A$ であって $v_r(y) < 0$ となるものを取れる。
$x$ をある冪で置き換えれば、$v_i(x) < v_i(y)$ が
$i = 1, \ldots, r - 1$ に対して成り立つと仮定してよい。このとき $x + y$ が求める元である。
\end{proof}

\begin{lemma}
\label{varieties-lemma-power-equal}
$A$ を次元 $1$ の Noether 局所環とする。$L = \prod A_\mathfrak p$ と置き、
積は $A$ の極小素イデアル全体にわたるものとする。
$a_1, a_2 \in \mathfrak m_A$ が $L$ の同じ元に写るならば、
$a_1^n = a_2^n$ となる $n > 0$ が存在する。
\end{lemma}

\begin{proof}
$a_1 = a_2 + x$ と書く。このとき $x$ は $L$ で零に写る。
従って $x$ は $A$ の冪零元である。これは $\bigcap \mathfrak p$ が $(0)$ の根基だからである。
また、零化イデアル $I$（$x$ を零化するもの）は極大イデアルのある冪を含む。実際、すべての極小素イデアルについて
$\mathfrak p \not \in V(I)$ である。$x^k = 0$ かつ $\mathfrak m^n \subset I$ とする。このとき
$$
a_1^{k + n} = a_2^{k + n} + {n + k \choose 1} a_2^{n + k - 1} x +
{n + k \choose 2} a_2^{n + k - 2} x^2 + \ldots +
{n + k \choose k - 1} a_2^{n + 1} x^{k - 1} = a_2^{n + k}
$$
である。これは $a_2 \in \mathfrak m_A$ による。
\end{proof}

\begin{lemma}
\label{varieties-lemma-power-works}
$A$ を次元 $1$ の Noether 局所環とする。$L = \prod A_\mathfrak p$ および
$I = \bigcap \mathfrak p$ と置き、積と共通部分はいずれも
$A$ の極小素イデアル全体にわたるものとする。$f \in L$ が
$f = i + a$ の形であり、$a \in \mathfrak m_A$ かつ $i \in IL$ とする。
このとき $f$ のある冪は $A \to L$ の像に属する。
\end{lemma}

\begin{proof}
$A$ は Noether なので、$I^t = 0$ となる $t > 0$ が存在する。
$f = a + i$ かつ $i \in I^kL$ であることが既知だと仮定する。
このとき $f^n = a^n + na^{n - 1}i \bmod I^{k + 1}L$ である。
従って、$na^{n - 1}i$ が $I^k \to I^kL$ の像に属するような
$n \gg 0$ が存在することを示せば十分である。
これを見るため、$g \in A$ を、$\mathfrak m_A = \sqrt{(g)}$ となるように選ぶ
（Algebra, Lemma \ref{algebra-lemma-height-1}）。
すると、例えば Algebra, Proposition
\ref{algebra-proposition-dimension-zero-ring} により $L = A_g$ である。
一方、$n$ が存在して $a^n \in (g)$ となる。従って
$L$ の元の分母は、$a$ の十分高い冪を掛けることで払える。
\end{proof}

\begin{lemma}
\label{varieties-lemma-good-intersection}
$A$ を次元 $1$ の Noether 局所環とする。$L = \prod A_\mathfrak p$ と置き、
積は $A$ の極小素イデアル全体にわたるものとする。
$K \to L$ を整な環準同型とする。このとき、
$a \in \mathfrak m_A$ と $x \in K$ であって、$L$ の同じ元に写り、
$\mathfrak m_A = \sqrt{(a)}$ を満たすものが存在する。
\end{lemma}

\begin{proof}
Lemma \ref{varieties-lemma-power-works} により、$A$ を
$A/(\bigcap \mathfrak p)$ で置き換え、$A$ が被約環であると仮定してよい
（細部は省略する）。また $K$ を $K \to L$ の像で置き換えてよい。
すると $K$ は被約環である。写像 $\Spec(L) \to \Spec(K)$ は
全射かつ閉である（細部は省略する）。従って $\Spec(K)$ は有限離散空間である。
このことから $K$ は体の有限直積である。

\medskip\noindent
$\mathfrak p_j$（$j = 1, \ldots, m$）を $A$ の極小素イデアルとする。
$L_j$ を $A_j$ の分数体と置き、
$L = \prod_{j = 1, \ldots, m} L_j$ とする。
$A_j$ を $A/\mathfrak p_j$ の正規化とする。このとき
$A_j$ は少なくとも一つの極大イデアルをもつ半局所 Dedekind 整域である。Algebra, Lemma
\ref{algebra-lemma-integral-closure-Dedekind} を参照せよ。
$n$ を $A_1, \ldots, A_m$ に含まれる極大イデアルの個数の総和とする。
そのような極大イデアル $\mathfrak m \subset A_j$ に対して、写像
$$
v_{\mathfrak m} : L \to L_j \to \mathbf{Z} \cup \{\infty\}
$$
を考える。ここで第 2 の矢印は、離散付値環 $(A_j)_{\mathfrak m}$ に対応する離散付値を
$0$ を $\infty$ に写すことで拡張したものである。このようにして $n$ 個の写像
$v_1, \ldots, v_n : L \to \mathbf{Z} \cup \{\infty\}$ を得る。
元 $x \in K$ であって、$v_i(x) < 0$ が
$i = 1, \ldots, n$ に対して成り立つものを見つける。

\medskip\noindent
まず、各 $i$ に対して元 $x \in K$ が存在し、$v_i(x) < 0$ となると主張する。
実際、$v_i$ が $\mathfrak m \subset A_j$ に対応するとする。
条件 $v_i(x) \geq 0$ がすべての $x \in K$ について成り立つならば、
$K$ は $(A_j)_{\mathfrak m}$ に写り、これは分数体 $L_j$（$A_j$ のもの）の中にある。
$K$ の $L_j$ における像は体であり、Algebra, Lemma
\ref{algebra-lemma-integral-under-field} により、その上 $L_j$ は代数的である。
これらを合わせると Algebra, Lemma
\ref{algebra-lemma-valuation-ring-cap-field-finite} に矛盾する。

\medskip\noindent
元 $x \in K$ が見つかり、$v_1(x) < 0, \ldots, v_r(x) < 0$ が
ある $r < n$ に対して成り立つとする。$v_{r + 1}(x) < 0$ ならば、
$x$ は $r + 1$ に対しても使える。そうでなければ、前段落の結果により可能であるように、
$y \in K$ を $v_{r + 1}(y) < 0$ となるように選ぶ。
$x$ を $x^n$ で置き換える。ここで $n > 0$ は適当に選ぶ。
すると $v_i(x) < v_i(y)$ が $i = 1, \ldots, r$ に対して成り立つと仮定してよい。
付値の性質により $v_j(x + y) = v_j(x) < 0$ が
$j = 1, \ldots, r$ に対して成り立ち、同様に
$v_{r + 1}(x + y) = v_{r + 1}(y) < 0$ である。
帰納法により、$x \in K$ であって $v_i(x) < 0$ が
$i = 1, \ldots, n$ に対して成り立つものが得られる。

\medskip\noindent
特に元 $x \in K$ は $K$ の各因子への非零な射影をもつ
（$K$ は体の有限直積であり、$x$ のある成分が零ならば、
$v_i(x)$ の値の一つが $\infty$ となることを思い出そう）。
従って $x$ は可逆であり、$x^{-1} \in K$ は
$\infty > v_i(x^{-1}) > 0$ をすべての $i$ に対して満たす元である。
Lemma \ref{varieties-lemma-semi-local-dimension-one-conductor} により、ある $e < 0$ に対し、
元 $x^e \in K$ は $\mathfrak m_A/\mathfrak p_j \subset A/\mathfrak p_j$ の元に
すべての $j = 1, \ldots, m$ について写る。
写像 $\mathfrak m_A \to \prod \mathfrak m_A/\mathfrak p_j$ の余核は
$\mathfrak m_A$ のある冪で零化されることに注意せよ。従って
$e$ をより小さい $e$ で置き換えれば、元 $a \in \mathfrak m_A$ であって、
$\mathfrak m_A/\mathfrak p_j$ における像が $x^e$ の像に等しいものが得られる。
組 $(a, x^e)$ は補題の結論を満たす。
\end{proof}

\begin{lemma}
\label{varieties-lemma-localization-semi-local}
$A$ を環とする。$\mathfrak p_1, \ldots, \mathfrak p_r$ を
$A$ の素イデアルの有限集合とする。
$S = A \setminus \bigcup \mathfrak p_i$ と置く。このとき $S$ は乗法系であり、
$S^{-1}A$ は、その極大イデアルが集合 $\{\mathfrak p_i\}$ の極大元に対応する半局所環である。
\end{lemma}

\begin{proof}
$a, b \in A$ かつ $a, b \in S$ ならば $a, b \not \in \mathfrak p_i$ であり、
従って $ab \not \in \mathfrak p_i$、従って $ab \in S$ である。また $1 \in S$ である。
従って $S$ は $A$ の乗法的部分集合である。
Algebra, Lemma \ref{algebra-lemma-spec-localization} における
$\Spec(S^{-1}A)$ の記述と
Algebra, Lemma \ref{algebra-lemma-silly} により、
$S^{-1}A$ の素イデアルは、$A$ の素イデアルであって、いずれかの
$\mathfrak p_i$ に含まれるものに対応する。従って $S^{-1}A$ の極大イデアルは、集合
$\{\mathfrak p_1, \ldots, \mathfrak p_r\}$ の（包含関係に関する）極大元と一対一に対応する。
\end{proof}









\section{1 次元 Noether スキーム}
\label{varieties-section-dimension-one}

\noindent
本節の主結果は、次元 $1$ の Noether 分離スキームが
豊富な可逆層をもつということである。Proposition
\ref{varieties-proposition-dim-1-noetherian-separated-has-ample} を参照せよ。

\begin{lemma}
\label{varieties-lemma-affine}
$X$ を、そのすべての局所環が次元 $\leq 1$ の Noether 環であるスキームとする。
$U \subset X$ を逆コンパクトな開集合とし、
$j : U \to X$ で包含射を表す。このとき
$R^pj_*\mathcal{F} = 0$、$p > 0$ は、任意の準連接
$\mathcal{O}_U$-加群 $\mathcal{F}$ に対して成り立つ。
\end{lemma}

\begin{proof}
$R^pj_*\mathcal{F}$ の消滅は茎で確かめてよい。
$R^qj_*$ の形成は平坦基底変換と可換である。Cohomology of Schemes, Lemma
\ref{coherent-lemma-flat-base-change-cohomology} を参照せよ。
従って $X$ は次元 $\leq 1$ の Noether 局所環のスペクトルであると仮定してよい。
この場合 $X$ は閉点 $x$ と、$x_1, \ldots, x_n$ という有限個の他の点をもち、
これらは $x$ に特殊化するが互いには特殊化しない（Algebra, Lemma
\ref{algebra-lemma-Noetherian-irreducible-components} 参照）。
$x \in U$ ならば $U = X$ であり、結論は明らかである。そうでなければ、点を番号づけ直すことにより
$U = \{x_1, \ldots, x_r\}$ とある $r$ に対して書ける。このとき $U$ はアフィンである
（Schemes, Lemma \ref{schemes-lemma-scheme-finite-discrete-affine}）。
従って結論は Cohomology of Schemes, Lemma
\ref{coherent-lemma-relative-affine-vanishing} から従う。
\end{proof}

\begin{lemma}
\label{varieties-lemma-open-in-affine-curve-affine}
$X$ を、そのすべての局所環が次元 $\leq 1$ の Noether 環であるアフィンスキームとする。
このとき任意の準コンパクト開集合 $U \subset X$ はアフィンである。
\end{lemma}

\begin{proof}
$j : U \to X$ で包含射を表す。$\mathcal{F}$ を準連接
$\mathcal{O}_U$-加群とする。
Lemma \ref{varieties-lemma-affine} により高次順像 $R^pj_*\mathcal{F}$ は零である。
$\mathcal{O}_X$-加群 $j_*\mathcal{F}$ は準連接である
（Schemes, Lemma \ref{schemes-lemma-push-forward-quasi-coherent}）。
従って Cohomology of Schemes, Lemma
\ref{coherent-lemma-quasi-coherent-affine-cohomology-zero} により、その高次コホモロジー群は消滅する。
Leray スペクトル系列 Cohomology, Lemma
\ref{cohomology-lemma-apply-Leray} により、
$H^p(U, \mathcal{F}) = 0$ はすべての $p > 0$ に対して成り立つ。
従って、例えば Cohomology of Schemes, Lemma
\ref{coherent-lemma-quasi-compact-h1-zero-covering} により $U$ はアフィンである。
\end{proof}

\begin{lemma}
\label{varieties-lemma-complement-codim-1-closed-points}
$X$ をスキーム、$U \subset X$ を開集合とする。次を仮定する。
\begin{enumerate}
\item $U$ は $X$ の逆コンパクトな開集合である。
\item $X \setminus U$ は離散的である。
\item $x \in X \setminus U$ に対して局所環
$\mathcal{O}_{X, x}$ は次元 $\leq 1$ の Noether 環である。
\end{enumerate}
このとき、(1) 可逆 $\mathcal{O}_X$-加群 $\mathcal{L}$ と切断 $s$ であって
$U = X_s$ となるものが存在し、(2) 写像
$\Pic(X) \to \Pic(U)$ は全射である。
\end{lemma}

\begin{proof}
$X \setminus U = \{x_i; i \in I\}$ とする。
アフィン開集合 $U_i \subset X$ を、$x_i \in U_i$ かつ
$x_j \not \in U_i$（$j \not = i$）となるように選ぶ。これは条件 (2) により可能である。
$U_i = \Spec(A_i)$ と書く。$\mathfrak m_i \subset A_i$ を
$x_i$ に対応する極大イデアルとする。局所環に関する仮定により、
$\mathfrak q \subset \mathfrak m_i$、$\mathfrak q \not = \mathfrak m_i$ を満たす素イデアルは有限個しかない
（Algebra, Lemma \ref{algebra-lemma-Noetherian-irreducible-components} 参照）。
従って素イデアル回避（Algebra, Lemma
\ref{algebra-lemma-silly}）により、それらのどの素イデアルにも含まれない
$f_i \in \mathfrak m_i$ を取れる。このとき
$V(f_i) = \{\mathfrak m_i\} \amalg Z_i$ となる閉部分集合
$Z_i \subset U_i$ が存在する。これは、$Z_i$ が $V(f_i)$ の逆コンパクトな開部分集合であって
特殊化の下で閉じていることによる。Algebra, Lemma
\ref{algebra-lemma-constructible-stable-specialization-closed} を参照せよ。
$U_i$ を縮めることにより $V(f_i) = \{x_i\}$ と仮定してよい。このとき
$$
\mathcal{U} : X = U \cup \bigcup U_i
$$
は $X$ の開被覆である。$2$-コサイクルであって、$\mathcal{O}_X^*$ に値をもち、
$f_i$ を $U \cap U_i$ 上に、$f_i/f_j$ を $U_i \cap U_j$ 上に割り当てるものを考える。
これは直線束 $\mathcal{L}$ を定め、切断 $s$ は $1$ により $U$ 上で、
$f_i$ により $U_i$ 上で定められ、補題の主張を満たす。

\medskip\noindent
$\mathcal{N}$ を可逆 $\mathcal{O}_U$-加群とする。
$N_i$ を可逆 $(A_i)_{f_i}$-加群であって
$\mathcal{N}|_{U \cap U_i}$ が $\tilde N_i$ に等しくなるものとする。
環 $(A_{\mathfrak m_i})_{f_i}$ は Artin 環である
（これは次元零の Noether 環だからである。Algebra, Lemma
\ref{algebra-lemma-Noetherian-dimension-0} 参照）。
従ってそれは局所環の直積であり
（Algebra, Lemma \ref{algebra-lemma-artinian-finite-length}）、
ゆえに Picard 群は自明である。従って $U_i$ を縮めることにより
（すなわち、$A_i$ を $(A_i)_g$ で置き換えることにより、ここで
$g \in A_i$ かつ $g \not \in \mathfrak m_i$ である）、
$N_i = (A_i)_{f_i}$、すなわち $\mathcal{N}|_{U \cap U_i}$ は自明であると仮定してよい。
この場合、$\mathcal{N}$ を $X$ 上の可逆層へ延長する方法は明らかである
（各 $U_i$ 上で自明な可逆加群により延長すればよい）。
\end{proof}

\begin{lemma}
\label{varieties-lemma-find-globally-generated}
$X$ を整分離スキームとする。$U \subset X$ を空でないアフィン開集合とし、
$X \setminus U$ は有限個の点 $x_1, \ldots, x_r$ からなり、
$\mathcal{O}_{X, x_i}$ は次元 $1$ の Noether 環であるとする。
このとき、大域生成された可逆 $\mathcal{O}_X$-加群 $\mathcal{L}$ と切断 $s$ であって
$U = X_s$ となるものが存在する。
\end{lemma}

\begin{proof}
$U = \Spec(A)$ と書き、$K$ を $X$ の函数体とする。
$B_i = \mathcal{O}_{X, x_i}$ および $\mathfrak m_i = \mathfrak m_{x_i}$ と書く。
$x_i \not \in U$ なので、開集合 $U \times_X \Spec(B_i)$ は
$\Spec(B_i)$ の中で一点しかもたない。すなわち
$U \times_X \Spec(B_i) = \Spec(K)$ である。
$X$ は分離的なので、$\Spec(K)$ は $U \times \Spec(B_i)$ の閉部分スキームである。
すなわち写像 $A \otimes B_i \to K$ は全射である。
Lemma \ref{varieties-lemma-create-globally-generated} により、非零元
$f \in A$ であって、$f^{-1} \in \mathfrak m_i$ が
$i = 1, \ldots, r$ に対して成り立つものを取れる。
開集合 $x_i \in U_i \subset X$ を、$f^{-1} \in \mathcal{O}(U_i)$ となるように選ぶ。このとき
$$
\mathcal{U} : X = U \cup \bigcup U_i
$$
は $X$ の開被覆である。$2$-コサイクルであって、$\mathcal{O}_X^*$ に値をもち、
$f$ を $U \cap U_i$ 上に、$1$ を $U_i \cap U_j$ 上に割り当てるものを考える。
これは直線束 $\mathcal{L}$ を定め、二つの切断をもつ。
\begin{enumerate}
\item 切断 $s$ は $1$ により $U$ 上で、$f^{-1}$ により $U_i$ 上で定められ、補題の主張を満たす。
\item 切断 $t$ は $f$ により $U$ 上で、$1$ により $U_i$ 上で定められる。
\end{enumerate}
$X_t \supset U_1 \cup \ldots \cup U_r$ であることに注意せよ。
従って $s, t$ は $\mathcal{L}$ を生成し、補題が証明された。
\end{proof}

\begin{lemma}
\label{varieties-lemma-enough-globally-generated-ample}
$X$ を準コンパクトスキームとする。すべての $x \in X$ に対し、
組 $(\mathcal{L}, s)$ であって、大域生成された可逆層 $\mathcal{L}$ と
大域切断 $s$ からなり、$x \in X_s$ かつ $X_s$ がアフィンとなるものが存在するならば、
$X$ は豊富な可逆層をもつ。
\end{lemma}

\begin{proof}
$X$ は準コンパクトなので、組 $(\mathcal{L}_i, s_i)$、
$i = 1, \ldots, n$ の有限族であって、$\mathcal{L}_i$ が大域生成され、
$X_{s_i}$ がアフィンであり、$X = \bigcup X_{s_i}$ となるものを取れる。
再び $X$ の準コンパクト性により、各 $i$ に対して切断の有限族
$t_{i, j}$ であって、$\mathcal{L}_i$ の切断からなり、
$j = 1, \ldots, m_i$ かつ
$X = \bigcup X_{t_{i, j}}$ となるものを取れる。$t_{i, 0} = s_i$ と置く。
可逆層
$$
\mathcal{L} = \mathcal{L}_1 
\otimes_{\mathcal{O}_X} \ldots
\otimes_{\mathcal{O}_X} \mathcal{L}_n
$$
および大域切断
$$
\tau_J = t_{1, j_1} \otimes \ldots \otimes t_{n, j_n}
$$
を考える。Properties, Lemma \ref{properties-lemma-affine-cap-s-open} により、
開集合 $X_{\tau_J}$ は、$j_i = 0$ がある $i$ に対して成り立てばアフィンである。
これらの開集合が $X$ を被覆することは容易に分かる。
従って定義により $\mathcal{L}$ は豊富である。
\end{proof}

\begin{lemma}
\label{varieties-lemma-dim-1-noetherian-integral-separated-has-ample}
$X$ を次元 $1$ の Noether 整分離スキームとする。
このとき $X$ は豊富な可逆層をもつ。
\end{lemma}

\begin{proof}
アフィン開被覆 $X = U_1 \cup \ldots \cup U_n$ を選ぶ。
$X$ は Noether なので、各集合 $X \setminus U_i$ は有限である。
従って Lemma \ref{varieties-lemma-find-globally-generated} により、
組 $(\mathcal{L}_i, s_i)$ であって、大域生成された可逆層 $\mathcal{L}_i$ と
大域切断 $s_i$ からなり、$U_i = X_{s_i}$ となるものを取れる。
Lemma \ref{varieties-lemma-enough-globally-generated-ample} により、
$X$ は豊富な可逆層をもつと結論する。
\end{proof}

\begin{lemma}
\label{varieties-lemma-surjective-pic-birational-finite}
$f : X \to Y$ をスキームの有限射とする。開集合 $V \subset Y$ が存在し、
$f^{-1}(V) \to V$ は同型、$Y \setminus V$ は離散空間であると仮定する。
このとき任意の可逆 $\mathcal{O}_X$-加群は、ある可逆
$\mathcal{O}_Y$-加群の引き戻しである。
\end{lemma}

\begin{proof}
$\Pic(X) = H^1(X, \mathcal{O}_X^*)$ を用いる。Cohomology, Lemma
\ref{cohomology-lemma-h1-invertible} を参照せよ。
Abel 層 $\mathcal{O}_X^*$ と $f$ に対する Leray スペクトル系列を考える。
Cohomology, Lemma \ref{cohomology-lemma-Leray} を参照せよ。
誘導される写像
$$
H^1(X, \mathcal{O}_X^*) \longrightarrow H^0(Y, R^1f_*\mathcal{O}_X^*)
$$
を考える。Divisors, Lemma
\ref{divisors-lemma-finite-trivialize-invertible-upstairs} は、この写像が零であるとまさに述べている。
従って Leray により
$H^1(X, \mathcal{O}_X^*) = H^1(Y, f_*\mathcal{O}_X^*)$ である。
次に写像
$$
f^\sharp : \mathcal{O}_Y^* \longrightarrow f_*\mathcal{O}_X^*
$$
を考える。仮定により、この写像の核と余核は閉部分集合
$T = Y \setminus V$（$Y$ の）上に台をもつ。
$T$ は仮定により離散位相空間なので、$Y$ 上の任意の Abel 層であって
$T$ 上に台をもつものの高次コホモロジー群は零である
（Cohomology, Lemma \ref{cohomology-lemma-cohomology-and-closed-immersions}、
Modules, Lemma \ref{modules-lemma-i-star-exact}、および
$H^i(T, \mathcal{F}) = 0$ が任意の $i > 0$ と任意の Abel 層
$\mathcal{F}$（$T$ 上）に対して成り立つことから従う）。
上の写像を短完全列
$$
0 \to \Ker(f^\sharp) \to \mathcal{O}_Y^* \to \Im(f^\sharp) \to 0,\quad
0 \to \Im(f^\sharp) \to f_*\mathcal{O}_X^* \to \Coker(f^\sharp) \to 0
$$
に分解すると、まず
$H^1(Y, \mathcal{O}_Y^*) \to H^1(Y, \Im(f^\sharp))$ が全射であり、次いで
$H^1(Y, \Im(f^\sharp)) \to H^1(Y, f_*\mathcal{O}_X^*)$ が全射であると分かる。
以上を合わせると、望むように
$H^1(Y, \mathcal{O}_Y^*) \to
H^1(X, \mathcal{O}_X^*)$ は全射である。
\end{proof}

\begin{lemma}
\label{varieties-lemma-glue-invertible-sheaves}
$X$ をスキームとする。$Z_1, \ldots, Z_n \subset X$ を閉部分スキームとし、
$\mathcal{L}_i$ を $Z_i$ 上の可逆層とする。次を仮定する。
\begin{enumerate}
\item $X$ は被約である。
\item $X = \bigcup Z_i$ が集合論的に成り立つ。
\item $Z_i \cap Z_j$ は $i \not = j$ に対して離散位相空間である。
\end{enumerate}
このとき、可逆層 $\mathcal{L}$ が $X$ 上に存在し、その
$Z_i$ への制限は $\mathcal{L}_i$ である。さらに、
$s_i \in \Gamma(Z_i, \mathcal{L}_i)$ という切断が与えられ、これらが
$Z_i \cap Z_j$ の点で消えないならば、$\mathcal{L}$ を、ある
$s \in \Gamma(X, \mathcal{L})$ が存在して
$s|_{Z_i} = s_i$ がすべての $i$ に対して成り立つように選べる。
\end{lemma}

\begin{proof}
$\mathcal{L}$ の存在は Lemma \ref{varieties-lemma-surjective-pic-birational-finite}
から導けるが、ここでは直接証明も与え、その証明を用いて切断に関する主張を確かめる。
$T = \bigcup_{i \not = j} Z_i \cap Z_j$ と置く。$X$ は被約なので、スキームとして
$$
X \setminus T = \bigcup (Z_i \setminus T)
$$
である。仮定 (3) により $T$ は $X$ の離散部分集合である。
従って各 $t \in T$ に対し、開集合 $U_t \subset X$ であって
$t \in U_t$ であり、$t' \not \in U_t$ が
$t' \in T$、$t' \not = t$ に対して成り立つものを取れる。
必要なら $U_t$ を縮めることにより、同型
$\varphi_{t, i} : \mathcal{L}_i|_{U_t \cap Z_i} \to
\mathcal{O}_{U_t \cap Z_i}$ が存在すると仮定してよい。さらに各 $i$ に対して開被覆
$$
Z_i \setminus T = \bigcup\nolimits_j U_{ij}
$$
を、同型
$\varphi_{i, j} : \mathcal{L}_i|_{U_{ij}} \cong \mathcal{O}_{U_{ij}}$
が存在するように選ぶ。ここで
$$
\mathcal{U} : X = \bigcup U_t \cup \bigcup U_{ij}
$$
は $X$ の開被覆である。同型 $\varphi_{t, i}$ および $\varphi_{i, j}$
を用いて、この被覆に関する $2$-コサイクルであって $\mathcal{O}_X^*$ に値をもち、
補題の主張にある $\mathcal{L}$ を定めるものを構成できると主張する。

\medskip\noindent
実際、$i \not = i'$ ならば $U_{i, j} \cap U_{i', j'} = \emptyset$
なので何もすることはない。$U_{i, j} \cap U_{i, j'}$ に対しては、証明冒頭の注意により
$\mathcal{O}_X(U_{i, j} \cap U_{i, j'}) =
\mathcal{O}_{Z_i}(U_{i, j} \cap U_{i, j'})$ である。
従って $\mathcal{L}_i$ の遷移函数（より正確には
$\varphi_{i, j} \circ \varphi_{i, j'}^{-1}$）が、この交わり上のコサイクルの値を定める。
$U_t \cap U_{i, j}$ に対しても同様にできる。
最後に、$t \not = t'$ に対して
$$
U_t \cap U_{t'} = \coprod (U_t \cap U_{t'}) \cap Z_i
$$
であり、さらに交わり $U_t \cap U_{t'} \cap Z_i$ は
$Z_i \setminus T$ に含まれる。従って先と同じ議論により
$$
\mathcal{O}_X(U_t \cap U_{t'}) =
\prod \mathcal{O}_{Z_i}(U_t \cap U_{t'} \cap Z_i)
$$
である。そこで $\mathcal{L}_i$ の遷移函数（より正確には
$\varphi_{t, i} \circ \varphi_{t', i}^{-1}$）を用いて、コサイクルの
$U_t \cap U_{t'}$ 上の値を定められる。これで $\mathcal{L}$ の存在証明が完了する。

\medskip\noindent
補題の最後の主張にある切断 $s_i$ が与えられたとする。上の議論で、
$U_t$ を $s_i|_{U_t \cap Z_i}$ が消えないように選び、
$\varphi_{t, i}$ を $\varphi_{t, i}(s_i|_{U_t \cap Z_i}) = 1$
となるように選ぶ。このとき、$1$ を $U_t$ 上で用い、
$\varphi_{i, j}(s_i|_{U_{i, j}})$ を $U_{i, j}$ 上で用いると、
$\mathcal{L}$ の切断であって $s_i$ に $Z_i$ 上で制限されるものが定まる。
\end{proof}

\begin{remark}
\label{varieties-remark-conductor}
$A$ を被約環とする。$I, J$ を $A$ のイデアルとし、
$V(I) \cup V(J) = \Spec(A)$ と仮定する。$B = A/J$ と置く。
このとき $I \to IB$ は $A$-加群の同型である。実際、
$IB = I + J/J = I/(I \cap J)$ であり、$I \cap J$ は零である。
なぜなら $A$ は被約であり、
$\Spec(A) = V(I) \cup V(J) = V(I \cap J)$ だからである。
従って任意の射影的 $A$-加群 $P$ に対しても $IP = I(P/JP)$ である。
\end{remark}

\begin{lemma}
\label{varieties-lemma-dim-1-noetherian-reduced-separated-has-ample}
$X$ を次元 $1$ の Noether 被約分離スキームとする。
このとき $X$ は豊富な可逆層をもつ。
\end{lemma}

\begin{proof}
$Z_i$、$i = 1, \ldots, n$ を $X$ の既約成分とする。
これらを $X$ の被約閉部分スキームとみなす。
Lemma \ref{varieties-lemma-dim-1-noetherian-integral-separated-has-ample} により、
豊富な可逆層 $\mathcal{L}_i$ が $Z_i$ 上に存在する。
$T = \bigcup_{i \not = j} Z_i \cap Z_j$ と置く。$X$ は次元 $1$ の Noether スキームなので、
集合 $T$ は有限であり、$X$ の閉点からなる。各 $i$ に対し、必要なら
$\mathcal{L}_i$ をその冪で置き換えることにより、
$s_i \in \Gamma(Z_i, \mathcal{L}_i)$ を、$(Z_i)_{s_i}$ がアフィンで
$T \cap Z_i$ を含むように選べる。Properties, Lemma
\ref{properties-lemma-ample-finite-set-in-principal-affine} を参照せよ。

\medskip\noindent
Lemma \ref{varieties-lemma-glue-invertible-sheaves} により、可逆層
$\mathcal{L}$ を $X$ 上に、また $s \in \Gamma(X, \mathcal{L})$ であって
$(\mathcal{L}, s)|_{Z_i} = (\mathcal{L}_i, s_i)$ となるものを取れる。
$X_s$ は $T$ を含み、集合論的にはアフィン閉部分スキーム
$(Z_i)_{s_i}$ に等しいことに注意せよ。従って Limits, Lemma
\ref{limits-lemma-affines-glued-in-closed-affine} によりアフィンである。
証明を終えるには、各 $x \in X$、$x \not \in T$ に対して、整数
$m > 0$ と切断 $t \in \Gamma(X, \mathcal{L}^{\otimes m})$ であって
$X_t$ がアフィンかつ $x \in X_t$ となるものを見つければ十分である。
$x \not \in T$ なので、$x \in Z_i$ となる $i$ は一意に存在する。
$i = 1$ としてよい。$Z \subset X$ を、台となる位相空間が
$Z_2 \cup \ldots \cup Z_n$ である被約閉部分スキームとする。
$\mathcal{I} \subset \mathcal{O}_X$ を $Z$ のイデアル層とする。
$\mathcal{I}_1 \subset \mathcal{O}_{Z_1}$ を、包含射
$Z_1 \to X$ によるこのイデアル層の逆像とする。このとき
$$
\Gamma(X, \mathcal{I}\mathcal{L}^{\otimes m}) =
\Gamma(Z_1, \mathcal{I}_1 \mathcal{L}_1^{\otimes m})
$$
である。Remark \ref{varieties-remark-conductor} を参照せよ。従って、
$m > 0$ と $t \in \Gamma(Z_1, \mathcal{I}_1 \mathcal{L}_1^{\otimes m})$
であって、$x \in (Z_1)_t$ がアフィンとなるものを見つければ十分である。
$\mathcal{L}_1$ は豊富であり、$x$ は
$Z_1 \cap T = V(\mathcal{I}_1)$ に属さないので、切断
$t_1 \in \Gamma(Z_1, \mathcal{I}_1 \mathcal{L}_1^{\otimes m_1})$
であって $x \in (Z_1)_{t_1}$ となるものを取れる。Properties, Proposition
\ref{properties-proposition-characterize-ample} を参照せよ。
$\mathcal{L}_1$ は豊富なので、切断
$t_2 \in \Gamma(Z_1, \mathcal{L}_1^{\otimes m_2})$
であって、$x \in (Z_1)_{t_2}$ かつ $(Z_1)_{t_2}$ がアフィンとなるものを取れる。
Properties, Definition \ref{properties-definition-ample} を参照せよ。
$m = m_1 + m_2$ および $t = t_1 t_2$ と置く。このとき構成により
$t \in \Gamma(Z_1, \mathcal{I}_1 \mathcal{L}_1^{\otimes m})$ かつ
$x \in (Z_1)_t$ であり、$(Z_1)_t$ は Properties, Lemma
\ref{properties-lemma-affine-cap-s-open} によりアフィンである。
\end{proof}

\begin{lemma}
\label{varieties-lemma-lift-line-bundle-from-reduction-dimension-1}
$i : Z \to X$ をスキームの閉埋め込みとする。
$X$ の台となる位相空間が Noether であり、
$\dim(X) \leq 1$ ならば、$\Pic(X) \to \Pic(Z)$ は全射である。
\end{lemma}

\begin{proof}
$X$ 上の Abel 群の層の短完全列
$$
0 \to (1 + \mathcal{I}) \cap \mathcal{O}_X^* \to
\mathcal{O}^*_X \to i_*\mathcal{O}^*_Z \to 0
$$
を考える。ここで $\mathcal{I}$ は $Z$ に対応する準連接イデアル層である。
$\dim(X) \leq 1$ なので、$H^2(X, \mathcal{F}) = 0$ は
任意の Abel 層 $\mathcal{F}$ に対して成り立つ。Cohomology, Proposition
\ref{cohomology-proposition-vanishing-Noetherian} を参照せよ。
従って写像 $H^1(X, \mathcal{O}^*_X) \to H^1(X, i_*\mathcal{O}_Z^*)$
は全射である。Cohomology, Lemma
\ref{cohomology-lemma-cohomology-and-closed-immersions} により
$H^1(X, i_*\mathcal{O}_Z^*) = H^1(Z, \mathcal{O}_Z^*)$ である。
Cohomology, Lemma \ref{cohomology-lemma-h1-invertible} により補題が従う。
\end{proof}

\begin{proposition}
\label{varieties-proposition-dim-1-noetherian-separated-has-ample}
$X$ を次元 $1$ の Noether 分離スキームとする。
このとき $X$ は豊富な可逆層をもつ。
\end{proposition}

\begin{proof}
$Z \subset X$ を $X$ の被約化とする。Lemma
\ref{varieties-lemma-dim-1-noetherian-reduced-separated-has-ample} により、
スキーム $Z$ は豊富な可逆層をもつ。
従って Lemma \ref{varieties-lemma-lift-line-bundle-from-reduction-dimension-1} により、
可逆 $\mathcal{O}_X$-加群 $\mathcal{L}$ が $X$ 上に存在し、その
$Z$ への制限は豊富である。このとき $\mathcal{L}$ は Cohomology of Schemes, Lemma
\ref{coherent-lemma-ample-on-reduction} の適用により豊富である。
\end{proof}

\begin{remark}
\label{varieties-remark-useless-generalization}
実際、$X$ を、その被約化が次元 $1$ の Noether 分離スキームであるスキームとすると、
$X$ は豊富な可逆層をもつ。この証明は Proposition
\ref{varieties-proposition-dim-1-noetherian-separated-has-ample} の証明と同じであるが、
Limits, Lemma \ref{limits-lemma-ample-on-reduction} を
Cohomology of Schemes, Lemma \ref{coherent-lemma-ample-on-reduction} の代わりに用いる。
\end{remark}

\noindent
次の補題は、読者が Zariski の主定理（More on Morphisms, Lemma
\ref{more-morphisms-lemma-quasi-finite-separated-pass-through-finite}）と
Lemma \ref{varieties-lemma-complement-codim-1-closed-points} を用いて確かめられるように、
実際には準有限分離射に対しても成り立つ。

\begin{lemma}
\label{varieties-lemma-surjection-on-pic-quasi-finite}
$f : X \to Y$ をスキームの射とする。$Y$ は次元 $\leq 1$ の Noether スキーム、
$f$ は有限射であり、稠密開集合 $V \subset Y$ が存在して
$f^{-1}(V) \to V$ が閉埋め込みになると仮定する。このとき任意の可逆
$\mathcal{O}_X$-加群は可逆 $\mathcal{O}_Y$-加群の引き戻しである。
\end{lemma}

\begin{proof}
$f$ を $X \to Z \to Y$ と分解する。ここで $Z$ は $f$ のスキーム論的像である。
このとき $X \to Z$ は $V \cap Z$ 上で同型であり、
Lemma \ref{varieties-lemma-surjective-pic-birational-finite} が適用できる。
一方、Lemma \ref{varieties-lemma-lift-line-bundle-from-reduction-dimension-1} は
$Z \to Y$ に適用できる。いくつかの細部は省略する。
\end{proof}








\section{デルタ不変量}
\label{varieties-section-delta-invariant}

\noindent
本節では、$\delta$-不変量を、被約な $1$ 次元 Nagata スキームの
特異点に対して定義する。

\begin{lemma}
\label{varieties-lemma-pre-pre-delta-invariant}
$(A, \mathfrak m)$ を Noether $1$ 次元局所環とし、$f \in \mathfrak m$ とする。
次は同値である。
\begin{enumerate}
\item $\mathfrak m = \sqrt{(f)}$。
\item $f$ は $A$ のどの極小素イデアルにも含まれない。
\item $A_f = \prod_{\mathfrak p\text{ minimal}} A_\mathfrak p$ が $A$-代数として成り立つ。
\end{enumerate}
このような $f \in \mathfrak m$ は存在する。$\text{depth}(A) = 1$ ならば
（例えば $A$ が被約ならば）、(1)--(3) は次とも同値である。
\begin{enumerate}
\item[(4)] $f$ は非零因子である。
\item[(5)] $A_f$ は $A$ の全商環である。
\end{enumerate}
$A$ が被約ならば、(1)--(5) はさらに次とも同値である。
\begin{enumerate}
\item[(6)] $A_f$ は $A$ の極小素イデアルにおける剰余体の直積である。
\end{enumerate}
\end{lemma}

\begin{proof}
$A$ のスペクトルは、有限個の素イデアル
$\mathfrak p_1, \ldots, \mathfrak p_n$ を $\mathfrak m$ 以外にもち、
これらはすべて極小である。
Algebra, Lemma \ref{algebra-lemma-Noetherian-irreducible-components} を参照せよ。
(1) と (2) の同値性は Algebra, Lemma
\ref{algebra-lemma-Zariski-topology} から従う。
(3) が (2) を含意することは明らかである。逆に (2) が成り立つならば、
$A_f$ のスペクトルは、部分集合
$\{\mathfrak p_1, \ldots, \mathfrak p_n\}$（$\Spec(A)$ の）に
誘導位相を入れたものである。
Algebra, Lemma \ref{algebra-lemma-spec-localization} を参照せよ。
これは有限離散位相空間である。従って Algebra, Proposition
\ref{algebra-proposition-dimension-zero-ring} により
$A_f = \prod_{\mathfrak p\text{ minimal}} A_\mathfrak p$ である。
$f$ の存在は Algebra, Lemma \ref{algebra-lemma-height-1} で述べられている。

\medskip\noindent
$A$ の深さが $1$ であると仮定する（これは Algebra, Lemma
\ref{algebra-lemma-bound-depth} により最大値であり、$A$ が被約ならば Algebra, Lemma
\ref{algebra-lemma-criterion-reduced} により成り立つ）。
このとき $\mathfrak m$ は $A$ の随伴素イデアルではない。
$A$ の任意の極小素イデアルは随伴素イデアルである
（Algebra, Proposition
\ref{algebra-proposition-minimal-primes-associated-primes}）。
従って Algebra, Lemma \ref{algebra-lemma-ass-zero-divisors} により、
$A$ の非零因子全体は、極小素イデアルのいずれにも含まれない元全体とちょうど一致する。
ゆえに (4) は (2) と同値である。(5) は Algebra, Lemma
\ref{algebra-lemma-total-ring-fractions-no-embedded-points} により (3) と同値である。

\medskip\noindent
このとき $A_\mathfrak p$ は、極小な $\mathfrak p \subset A$ に対して体である。
Algebra, Lemma \ref{algebra-lemma-minimal-prime-reduced-ring} を参照せよ。
従って (3) は (6) と同値である。
\end{proof}

\begin{lemma}
\label{varieties-lemma-pre-delta-invariant}
$(A, \mathfrak m)$ を被約 Nagata $1$ 次元局所環とする。
$A'$ を、$A$ の全商環における $A$ の整閉包とする。
このとき $A'$ は正規 Nagata 環であり、$A \to A'$ は有限で、
$A'/A$ は $A$-加群として有限の長さをもつ。
\end{lemma}

\begin{proof}
全商環は $A$ 上本質的有限型なので、$A \to A'$ は有限である。
これは $A$ が Nagata であることと Algebra, Lemma
\ref{algebra-lemma-nagata-in-reduced-finite-type-finite-integral-closure} による。
環 $A'$ は、例えば Algebra, Lemma
\ref{algebra-lemma-characterize-reduced-ring-normal} および
\ref{algebra-lemma-Noetherian-irreducible-components} により正規である。
環 $A'$ は、例えば Algebra, Lemma
\ref{algebra-lemma-quasi-finite-over-nagata} により Nagata である。
Lemma \ref{varieties-lemma-pre-pre-delta-invariant} にあるような
$f \in \mathfrak m$ を選ぶ。$A' \subset A_f$ なので
$A_f = A'_f$ は明らかである。従って有限 $A$-加群 $A'/A$ の台は
$\{\mathfrak m\}$ に含まれる。Algebra, Lemma
\ref{algebra-lemma-support-point} により、その長さは有限である。
\end{proof}

\begin{definition}
\label{varieties-definition-delta-invariant-algebra}
$A$ を次元 $1$ の被約 Nagata 局所環とする。
{\it $\delta$-不変量（環 $A$ の）}とは $\text{length}_A(A'/A)$ をいう。
ここで $A'$ は Lemma \ref{varieties-lemma-pre-delta-invariant} におけるものとする。
\end{definition}

\noindent
この不変量の振る舞いに関するいくつかの補題を証明する。

\begin{lemma}
\label{varieties-lemma-delta-invariant-is-zero}
$A$ を次元 $1$ の被約 Nagata 局所環とする。
$\delta$-不変量（$A$ の）が $0$ であるための必要十分条件は、
$A$ が離散付値環であることである。
\end{lemma}

\begin{proof}
$A$ が離散付値環ならば $A$ は正規であり、環 $A'$ は $A$ に等しい。
逆に、$\delta$-不変量（$A$ の）が $0$ ならば、$A$ はその全商環の中で整閉である。
これは $A$ が正規であることを含意し
（Algebra, Lemma \ref{algebra-lemma-characterize-reduced-ring-normal}）、
さらに Algebra, Lemma \ref{algebra-lemma-characterize-dvr} により
$A$ は離散付値環でなければならない。
\end{proof}

\begin{lemma}
\label{varieties-lemma-normalization-same-after-completion}
$A$ を次元 $1$ の被約 Nagata 局所環とする。
$A \to A'$ を Lemma \ref{varieties-lemma-pre-delta-invariant} におけるものとする。
$A^h$、$A^{sh}$、resp.\ $A^\wedge$ を、それぞれ $A$ の Hensel 化、強 Hensel 化、
resp.\ 完備化とする。このとき $A^h$、$A^{sh}$、resp. $A^\wedge$ は
被約 Nagata $1$ 次元局所環であり、
$A' \otimes_A A^h$、$A' \otimes_A A^{sh}$、resp. $A' \otimes_A A^\wedge$ は、
それぞれ $A^h$、$A^{sh}$、resp.\ $A^\wedge$ の全商環における整閉包である。
\end{lemma}

\begin{proof}
$A^\wedge$ は被約である。More on Algebra, Lemma
\ref{more-algebra-lemma-completion-reduced} を参照せよ。
環 $A^h$ と $A^{sh}$ は More on Algebra, Lemma
\ref{more-algebra-lemma-henselization-reduced} により被約である。
$A$、$A^h$、$A^{sh}$、および $A^\wedge$ の次元は、More on Algebra, Lemmas
\ref{more-algebra-lemma-completion-dimension} および
\ref{more-algebra-lemma-henselization-dimension} により等しい。

\medskip\noindent
Noether 局所環が Nagata であるための必要十分条件は、
$A$ の形式的ファイバーが幾何学的被約であることだった。More on Algebra, Lemma
\ref{more-algebra-lemma-Nagata-local-ring} を参照せよ。
この性質は $A^h$ および $A^{sh}$ に継承される。More on Algebra, Section
\ref{more-algebra-section-properties-formal-fibres}、特に Lemma
\ref{more-algebra-lemma-henselization-P-ring} の内容を参照せよ。
完備化は Algebra, Lemma
\ref{algebra-lemma-Noetherian-complete-local-Nagata} により Nagata である。

\medskip\noindent
次に整閉包に関する主張を示す。先へ進む前に、Lemma
\ref{varieties-lemma-pre-pre-delta-invariant} にあるような
$f \in \mathfrak m$ を選ぶ。このとき $f$ の $A^h$、$A^{sh}$、および
$A^\wedge$ における像は、明らかにその環で性質 (1)--(6) を満たす。

\medskip\noindent
$A \to A'$ は有限なので、$A' \otimes_A A^h$ と
$A' \otimes_A A^{sh}$ は、それぞれ $A^h$ と $A^{sh}$ 上有限な
Hensel 局所環の直積である。Algebra, Lemma
\ref{algebra-lemma-finite-over-henselian} を参照せよ。
これらの各局所環は、$A'$ の、極大イデアル
$\mathfrak m' \subset A'$（$\mathfrak m$ の上にある）における Hensel 化である。
Algebra, Lemma
\ref{algebra-lemma-quasi-finite-henselization} または
\ref{algebra-lemma-quasi-finite-strict-henselization} を参照せよ。
従ってこれらの局所環は More on Algebra, Lemma
\ref{more-algebra-lemma-henselization-normal} により正規整域である。
従って $A' \otimes_A A^h$ と $A' \otimes_A A^{sh}$ は正規環である。
$A^h \to A' \otimes_A A^h$ および
$A^{sh} \to A' \otimes_A A^{sh}$ は有限、従って整であり、また
$A' \otimes_A A^h \subset (A^h)_f = Q(A^h)$ および
$A' \otimes_A A^{sh} \subset (A^{sh})_f = Q(A^{sh})$ なので、
$A' \otimes_A A^h$ と $A' \otimes_A A^{sh}$ は求める整閉包である。

\medskip\noindent
完備化についてもまったく同様に論じる。まず Algebra, Lemma
\ref{algebra-lemma-completion-finite-extension} により
$$
A' \otimes_A A^\wedge = (A')^\wedge =
\prod\nolimits (A'_{\mathfrak m'})^\wedge
$$
である。局所環 $A'_{\mathfrak m'}$ は正規で次元 $1$ である
（例えば Algebra, Lemma \ref{algebra-lemma-finite-in-codim-1}、
または Algebra, Section
\ref{algebra-section-homomorphism-dimension} の議論による）。
従って $A'_{\mathfrak m'}$ は離散付値環である。Algebra, Lemma
\ref{algebra-lemma-characterize-dvr} を参照せよ。
ゆえに $(A'_{\mathfrak m'})^\wedge$ は More on Algebra, Lemma
\ref{more-algebra-lemma-completion-dvr} により離散付値環である。
従って $A' \otimes_A A^\wedge$ は正規環であり、先とまったく同様に結論できる。
\end{proof}

\begin{lemma}
\label{varieties-lemma-delta-same-after-completion}
$A$ を次元 $1$ の被約 Nagata 局所環とする。
$\delta$-不変量（$A$ の）は、その Hensel 化、強 Hensel 化、または完備化の
$\delta$-不変量（$A$ に由来する）と同じである。
\end{lemma}

\begin{proof}
完備化 $B = A^\wedge$ の場合を示す。他の場合もまったく同様に証明される。
$A'$、resp.\ $B'$ を、それぞれ $A$、resp.\ $B$ の全商環における整閉包とする。
このとき Lemma \ref{varieties-lemma-normalization-same-after-completion} により
$B' = A' \otimes_A B$ である。従って $B'/B = (A'/A) \otimes_A B$ である。
ここから、Algebra, Lemma \ref{algebra-lemma-pullback-module} と
$B \otimes_A \kappa_A = \kappa_B$ であることにより等式が従う。
\end{proof}

\begin{definition}
\label{varieties-definition-delta-invariant}
$k$ を体とし、$X$ を局所代数的 $k$-スキームとする。
$x \in X$ を、$\mathcal{O}_{X, x}$ が被約かつ
$\dim(\mathcal{O}_{X, x}) = 1$ となる点とする。
{\it $\delta$-不変量（$X$ の $x$ における）}とは、
Definition \ref{varieties-definition-delta-invariant-algebra} で定義された
$\delta$-不変量（$\mathcal{O}_{X, x}$ の）をいう。
\end{definition}

\noindent
これは、局所代数的スキームの局所環が Algebra, Proposition
\ref{algebra-proposition-ubiquity-nagata} により Nagata であるため意味をもつ。
もちろん、より一般に、$x \in X$ がスキームの点であり、その局所環
$\mathcal{O}_{X, x}$ が被約な次元 $1$ の Nagata 環であるときは常に、この定義を使える。
Lemma \ref{varieties-lemma-delta-same-after-completion} から、$\delta$-不変量（$X$ の
$x$ における）は
$$
\delta\text{-invariant of }X\text{ at }x =
\delta\text{-invariant of }\mathcal{O}_{X, x}^h =
\delta\text{-invariant of }\mathcal{O}_{X, x}^\wedge
$$
である。従って $\delta$-不変量は、その点の完備局所環の不変量である。

\begin{lemma}
\label{varieties-lemma-delta-invariant-and-change-of-fields}
$k$ を体とし、$X$ を局所代数的 $k$-スキームとする。
$K/k$ を体拡大とし、$Y = X_K$ と置く。
$y \in Y$ の像を $x \in X$ とする。$X$ が $x$ において幾何学的被約であり、
$\dim(\mathcal{O}_{X, x}) = \dim(\mathcal{O}_{Y, y}) = 1$ と仮定する。このとき
$$
\delta\text{-invariant of }X\text{ at }x \leq
\delta\text{-invariant of }Y\text{ at }y
$$
である。
\end{lemma}

\begin{proof}
$A = \mathcal{O}_{X, x}$ および $B = \mathcal{O}_{Y, y}$ と置く。
Lemma \ref{varieties-lemma-geometrically-reduced-at-point} により $A$ は幾何学的被約である。
従って $B$ は $A \otimes_k K$ の局所化である。
$A \to A'$ を Lemma \ref{varieties-lemma-pre-delta-invariant} におけるものとする。このとき
$$
B' = B \otimes_{(A \otimes_k K)} (A' \otimes_k K)
$$
は $B$ 上有限であり、$B \to B'$ は全商環の同型を誘導する。
実際、Lemma \ref{varieties-lemma-pre-pre-delta-invariant} の (1)--(6) を満たす
$f \in \mathfrak m_A$ を選ぶ。$\dim(B) = 1$ なので、
$f \in \mathfrak m_B$ は $B$ に対して同じ役割を果たし、
$B_f = B'_f$ である。これは $A_f = A'_f$ による。
$B''$ を、Lemma \ref{varieties-lemma-pre-delta-invariant} にあるような
$B$ の全商環における整閉包とする。このとき $B' \subset B''$ である。
従って $\delta$-不変量（$Y$ の $y$ における）は
$\text{length}_B(B''/B)$ であり、
\begin{align*}
\text{length}_B(B''/B)
& \geq
\text{length}_B(B'/B) \\
& =
\text{length}_B((A'/A) \otimes_A B) \\
& =
\text{length}_B(B/\mathfrak m_A B) \text{length}_A(A'/A)
\end{align*}
である。これは Algebra, Lemma \ref{algebra-lemma-pullback-module} による。
実際、$A \to B$ は平坦である（$A \to A \otimes_k K$ の局所化だからである）。
$\text{length}_A(A'/A)$ は $\delta$-不変量（$X$ の $x$ における）であり、
$\text{length}_B(B/\mathfrak m_A B) \geq 1$ なので、補題が証明された。
\end{proof}

\begin{lemma}
\label{varieties-lemma-delta-invariant-and-change-of-fields-better}
$k$ を体とし、$X$ を局所代数的 $k$-スキームとする。
$K/k$ を体拡大とし、$Y = X_K$ と置く。
$y \in Y$ の像を $x \in X$ とする。Lemma
\ref{varieties-lemma-geometrically-normal-in-codim-1} の仮定 (a), (b), (c) が
$x \in X$ に対して成り立ち、$\dim(\mathcal{O}_{Y, y}) = 1$ であると仮定する。
このとき $\delta$-不変量（$X$ の $x$ における）は、
$\delta$-不変量（$Y$ の $y$ における）に等しい。
\end{lemma}

\begin{proof}
$A = \mathcal{O}_{X, x}$ および $B = \mathcal{O}_{Y, y}$ と置く。
Lemma \ref{varieties-lemma-geometrically-normal-in-codim-1} により
$A$ は幾何学的被約である。従って $B$ は $A \otimes_k K$ の局所化である。
$A \to A'$ を Lemma \ref{varieties-lemma-pre-delta-invariant} におけるものとする。
Lemma \ref{varieties-lemma-geometrically-normal-in-codim-1} により
$A' \otimes_k K$ は正規である。従って
$$
B' = B \otimes_{(A \otimes_k K)} (A' \otimes_k K)
$$
は正規で $B$ 上有限であり、$B \to B'$ は全商環の同型を誘導する。
実際、Lemma \ref{varieties-lemma-pre-pre-delta-invariant} の (1)--(6) を満たす
$f \in \mathfrak m_A$ を選ぶ。$\dim(B) = 1$ なので、
$f \in \mathfrak m_B$ は $B$ に対して同じ役割を果たし、
$B_f = B'_f$ である。これは $A_f = A'_f$ による。
従って $B \to B'$ は $B$ に対して Lemma
\ref{varieties-lemma-pre-delta-invariant} におけるものとなる。ゆえに示すべきことは
$\text{length}_A(A'/A) =
\text{length}_B(B'/B) = \text{length}_B((A'/A) \otimes_A B)$
である。$A \to B$ は平坦であり（$A \to A \otimes_k K$ の局所化だからである）、
また $\mathfrak m_B = \mathfrak m_A B$ である。実際、
$B/\mathfrak m_A B$ は上の注意により零次元であり、
$K \otimes_k \kappa(x)$ の局所化である。これは $\kappa(x)$ が
$k$ 上分離的なので被約である。従って Algebra, Lemma
\ref{algebra-lemma-pullback-module} により結論が従う。
\end{proof}





\section{分枝数}
\label{varieties-section-number-of-branches}

\noindent
スキームの点における分枝数は Properties, Section
\ref{properties-section-unibranch} で定義した。

\begin{lemma}
\label{varieties-lemma-number-of-branches}
$X$ をスキームとする。$X$ の任意の準コンパクト開集合が有限個の
既約成分をもつと仮定する。$\nu : X^\nu \to X$ を $X$ の正規化とし、
$x \in X$ とする。
\begin{enumerate}
\item $X$ の $x$ における分枝数は、$x$ の $X^\nu$ における逆像の個数である。
\item $X$ の $x$ における幾何学的分枝数は
$\sum_{\nu(x^\nu) = x} [\kappa(x^\nu) : \kappa(x)]_s$ である。
\end{enumerate}
\end{lemma}

\begin{proof}
まず、$X$ に関する仮定は正規化が定義されることとちょうど同値である。
Morphisms, Definition \ref{morphisms-definition-normalization} を参照せよ。
このとき茎 $A' = (\nu_*\mathcal{O}_{X^\nu})_x$ は、$A = \mathcal{O}_{X, x}$
の $A_{red}$ の全商環における整閉包である。Morphisms, Lemma
\ref{morphisms-lemma-stalk-normalization} を参照せよ。
$\nu$ は整射なので、$X^\nu$ の $x$ の上にある点は、$A'$ の素イデアルのうち、
極大イデアル $\mathfrak m$（$A$ の）の上にあるものに対応する。
$A \to A'$ は整なので、これは $A'$ の極大イデアルと同じものである
（Algebra, Lemmas \ref{algebra-lemma-integral-no-inclusion} および
\ref{algebra-lemma-integral-going-up}）。
従って補題は、その代数版である More on Algebra, Lemma
\ref{more-algebra-lemma-number-of-branches-1} から従う。
\end{proof}

\begin{lemma}
\label{varieties-lemma-geometric-branches-and-change-of-fields}
$k$ を体とし、$X$ を局所代数的 $k$-スキームとする。
$K/k$ を体拡大とする。$y \in X_K$ を、その像 $x$ が $X$ にある点とする。
このとき $X$ の $x$ における幾何学的分枝数は、
$X_K$ の $y$ における幾何学的分枝数に等しい。
\end{lemma}

\begin{proof}
$Y = X_K$ と書き、$X^\nu$、resp.\ $Y^\nu$ をそれぞれ $X$、resp.\ $Y$
の正規化とする。可換図式
$$
\xymatrix{
Y^\nu \ar[r] \ar[d] & X^\nu_K \ar[r] \ar[d]_{\nu_K} & X^\nu \ar[d]_\nu \\
Y \ar@{=}[r] & Y \ar[r] & X
}
$$
を考える。Lemma \ref{varieties-lemma-normalization-and-change-of-fields} により、
上段左の水平射は普遍同相射である。従ってこれは純非分離剰余体拡大を誘導する。
Morphisms, Lemmas \ref{morphisms-lemma-universal-homeomorphism} および
\ref{morphisms-lemma-universally-injective} を参照せよ。
従って Lemma \ref{varieties-lemma-number-of-branches} により、$Y$ の $y$ における
幾何学的分枝数は
$\sum_{\nu_K(y') = y} [\kappa(y') : \kappa(y)]_s$ である。同様に、
$\sum_{\nu(x') = x} [\kappa(x') : \kappa(x)]_s$ は $X$ の $x$ における
幾何学的分枝数である。Schemes, Lemma
\ref{schemes-lemma-points-fibre-product} を用いると、主張は次の代数的事実から従う。
体拡大 $l/\kappa$ と代数的体拡大 $m/\kappa$ が与えられたとき、
$$
\sum\nolimits_{m \otimes_\kappa l \to m'} [m' : l']_s = [m : \kappa]_s
$$
である。ここで和は $m \otimes_\kappa l$ の商体全体にわたる。
これは初等的に証明できる。あるいは
$$
\Spec(m \otimes_\kappa l) \times_{\Spec(l)} \Spec(\overline{l}) =
\Spec(m) \otimes_{\Spec(\kappa)} \Spec(\overline{l})
\longrightarrow
\Spec(m) \times_{\Spec(\kappa)} \Spec(\overline{\kappa})
$$
に Lemma \ref{varieties-lemma-separably-closed-field-connected-components} を適用してもよい。
実際、$[m : \kappa]_s$ は右辺の連結成分の個数と解釈でき、和
$\sum_{m \otimes_\kappa l \to m'} [m' : l']_s$ は左辺の連結成分の個数と解釈できる。
\end{proof}

\begin{lemma}
\label{varieties-lemma-geometrically-unibranch-and-change-of-fields}
$k$ を体とし、$X$ を局所代数的 $k$-スキームとする。
$K/k$ を体拡大とし、$y \in X_K$ を、その像 $x$ が $X$ にある点とする。
このとき $X$ が $x$ において幾何学的単枝であるための必要十分条件は、
$X_K$ が $y$ において幾何学的単枝であることである。
\end{lemma}

\begin{proof}
Lemma \ref{varieties-lemma-geometric-branches-and-change-of-fields} および More on Algebra,
Lemma \ref{more-algebra-lemma-number-of-branches-1} から直ちに従う。
\end{proof}

\begin{definition}
\label{varieties-definition-wedge}
$A$ および $A_i$（$1 \leq i \leq n$）を局所環とする。
{\it $A$ は $A_1, \ldots, A_n$ の楔和である}というのは、同型
$$
\kappa_{A_1} \to \kappa_{A_2} \to \ldots \to \kappa_{A_n}
$$
が存在し、$A$ が、$n$-組
$(a_1, \ldots, a_n) \in A_1 \times \ldots \times A_n$ のうち
$\kappa_{A_n}$ の同じ元へ写るもの全体からなる環と同型であるときである。
\end{definition}

\noindent
基礎環 $\Lambda$ が与えられ、$A$ および $A_i$ が $\Lambda$-代数である場合には、
$\kappa_{A_i} \to \kappa_{A_{i + 1}}$ が $\Lambda$-代数同型であること、および
$A$ が $\Lambda$-代数として、$\Lambda$-代数、すなわち
$n$-組 $(a_1, \ldots, a_n) \in A_1 \times \ldots \times A_n$ のうち
$\kappa_{A_n}$ の同じ元へ写るもの全体からなる環と同型であることを要求する。
特に、$\Lambda = k$ が体で、写像 $k \to \kappa_{A_i}$ が同型ならば、
同型 $\kappa_{A_i} \to \kappa_{A_{i + 1}}$ の選び方は一意であり、しばしば
{\it $A_1, \ldots, A_n$ の楔和}という。

\begin{lemma}
\label{varieties-lemma-delta-number-branches-inequality-sh}
$(A, \mathfrak m)$ を強 Hensel な $1$ 次元被約 Nagata 局所環とする。このとき
$$
\delta\text{-invariant of }A \geq \text{number of geometric branches of }A - 1
$$
である。等号が成り立つならば、$A$ は $n \geq 1$ 個の強 Hensel 離散付値環の
楔和である。
\end{lemma}

\begin{proof}
幾何学的分枝数は $A$ の分枝数に等しい
（More on Algebra, Definition \ref{more-algebra-definition-number-of-branches}
から直ちに従う）。$A \to A'$ を Lemma
\ref{varieties-lemma-pre-delta-invariant} におけるものとする。
$A$ の分枝数は $A'$ の極大イデアルの個数であることに注意する。
More on Algebra, Lemma \ref{more-algebra-lemma-number-of-branches-1} を参照せよ。
全射
$$
A'/A \longrightarrow
\left(\prod\nolimits_{\mathfrak m'} \kappa(\mathfrak m')\right)/
\kappa(\mathfrak m)
$$
がある。$\dim_{\kappa(\mathfrak m)} \prod \kappa(\mathfrak m')$ は
$\geq$ 分枝数なので、不等式は明らかである。

\medskip\noindent
等号が成り立つならば、$\kappa(\mathfrak m') = \kappa(\mathfrak m)$ が
すべての $\mathfrak m' \subset A'$ に対して成り立ち、上の表示された射は同型である。
$A$ は Hensel で $A \to A'$ は有限なので、$A'$ は局所 Hensel 環の直積である。
Algebra, Lemma \ref{algebra-lemma-finite-over-henselian} を参照せよ。
各因子は局所環 $A'_{\mathfrak m'}$ であり、$A'$ は正規なので、これらの因子は
離散付値環である（Algebra, Lemma \ref{algebra-lemma-characterize-dvr}）。
表示された射が同型なので、$A$ は確かにこれらの局所環の楔和である。
\end{proof}

\begin{lemma}
\label{varieties-lemma-delta-number-branches-inequality}
$(A, \mathfrak m)$ を $1$ 次元被約 Nagata 局所環とする。このとき
$$
\delta\text{-invariant of }A \geq \text{number of geometric branches of }A - 1
$$
である。
\end{lemma}

\begin{proof}
$A$ を $A$ の強 Hensel 化で置き換えてよい。この置換は
$\delta$-不変量を変えず（Lemma
\ref{varieties-lemma-delta-same-after-completion}）、また $A$ の幾何学的分枝数も変えない
（これは定義から直ちに従う。More on Algebra, Definition
\ref{more-algebra-definition-number-of-branches} を参照せよ）。
従って $A$ は強 Hensel であると仮定でき、Lemma
\ref{varieties-lemma-delta-number-branches-inequality-sh} を適用できる。
\end{proof}





\section{1 次元スキームの正規化}
\label{varieties-section-normalization-one-dimensional}

\noindent
Krull--Akizuki の結果により、次元 $1$ の Noether スキームの正規化射は、
予想外に良い性質をもつ。

\begin{lemma}
\label{varieties-lemma-normalize-noetherian-dim-1}
$X$ を次元 $1$ の局所 Noether スキームとする。
$\nu : X^\nu \to X$ を正規化とする。このとき次が成り立つ。
\begin{enumerate}
\item $\nu$ は整かつ全射であり、既約成分の間に全単射を誘導する。
\item 分解 $X^\nu \to X_{red} \to X$ が存在し、射
$X^\nu \to X_{red}$ は $X_{red}$ の正規化である。
\item $X^\nu \to X_{red}$ は双有理である。
\item 任意の閉点 $x \in X$ に対し、茎
$(\nu_*\mathcal{O}_{X^\nu})_x$ は、$\mathcal{O}_{X, x}$ の整閉包である。
ここで全商環は $(\mathcal{O}_{X, x})_{red} = \mathcal{O}_{X_{red}, x}$ のものである。
\item $\nu$ のファイバーは有限であり、剰余体拡大は有限である。
\item $X^\nu$ は整正規局所 Noether スキームの非交和であり、各アフィン開集合は
Dedekind 整域の有限直積のスペクトルである。
\end{enumerate}
\end{lemma}

\begin{proof}
これらの結果の多くは、実際には正規化射の一般的性質である。
Morphisms, Lemmas \ref{morphisms-lemma-normalization-reduced},
\ref{morphisms-lemma-stalk-normalization},
\ref{morphisms-lemma-normalization-normal}, および
\ref{morphisms-lemma-normalization-birational} を参照せよ。
明らかでないのは、ファイバーが有限であること、誘導される剰余体拡大が有限であること、
および $X^\nu$ が局所的に Dedekind 整域のスペクトルのように見える
（従って Noether である）ことである。
これを見るため、$X = \Spec(A)$ はアフィン、Noether、次元 $1$ であり、
$A$ は被約であると仮定してよい。次いで Morphisms, Lemma
\ref{morphisms-lemma-description-normalization} の記述を用いて、$A$ が
次元 $1$ の Noether 整域である場合に帰着できる。この場合、求める性質は
Algebra, Lemma \ref{algebra-lemma-integral-closure-Dedekind} に述べられた形の
Krull--Akizuki から従う。
\end{proof}

\noindent
もちろん、$X$ が被約でない場合にも次の補題の変種がある。

\begin{lemma}
\label{varieties-lemma-prepare-delta-invariant}
$X$ を次元 $1$ の被約 Nagata スキームとする。$\nu : X^\nu \to X$
を正規化とする。$x \in X$ を閉点とする。このとき次が成り立つ。
\begin{enumerate}
\item $\nu : X^\nu \to X$ は有限、全射、双有理である。
\item $\mathcal{O}_X \subset \nu_*\mathcal{O}_{X^\nu}$ であり、
$\nu_*\mathcal{O}_{X^\nu}/\mathcal{O}_X$ は、値 $\mathcal{Q}_x$ を
$x$ にもつ摩天楼層の直和である。また、これが非零であるための必要十分条件は、
$x$ が $X$ の特異点であることである。
\item $A' = (\nu_*\mathcal{O}_{X^\nu})_x$ は、$A = \mathcal{O}_{X, x}$
の全商環における整閉包である。
\item $\mathcal{Q}_x = A'/A$ は有限の長さをもち、その長さは
$\delta$-不変量（$X$ の $x$ における）に等しい。
\item $A'$ は Dedekind 整域の有限直積である半局所環である。
\item $A^\wedge$ は次元 $1$ の被約 Noether 完備局所環である。
\item $(A')^\wedge$ は、$A^\wedge$ の全商環における整閉包である。
\item $(A')^\wedge$ は完備離散付値環の有限直積である。
\item $A'/A \cong (A')^\wedge/A^\wedge$ である。
\end{enumerate}
\end{lemma}

\begin{proof}
Lemma \ref{varieties-lemma-normalize-noetherian-dim-1} のすべての結果を用いてよく、
実際そうする。$\nu$ の有限性は Morphisms, Lemma
\ref{morphisms-lemma-nagata-normalization} から従う。
$X$ は被約、Nagata、次元 $1$ なので、More on Algebra, Proposition
\ref{more-algebra-proposition-ubiquity-J-2} により正則点集合は稠密開集合
$U \subset X$ である。正則スキームは正規なので、$\nu$ は $U$ 上同型である。
$\dim(X) \leq 1$ なので、これは $\nu$ が閉点 $x \in X$ の離散集合の上で
同型でないことを含意する。特に短完全列
$$
0 \to \mathcal{O}_X \to \nu_*\mathcal{O}_{X^\nu} \to
\bigoplus\nolimits_{x \in X \setminus U} \mathcal{Q}_x \to 0
$$
を得る。Lemma \ref{varieties-lemma-normalize-noetherian-dim-1} により
$\nu_*\mathcal{O}_{X^\nu}$ の茎の記述があるので、$Q_x = A'/A$ は確かに
有限の長さをもち、その長さは $\delta$-不変量（$X$ の $x$ における）に等しい。
例えば Lemma \ref{varieties-lemma-delta-invariant-is-zero} により、
$Q_x \not = 0$ であることと $x$ が特異点であることはちょうど同値である。
$A'$ が半局所 Dedekind 整域の直積であるという記述も
Lemma \ref{varieties-lemma-normalize-noetherian-dim-1} から従う。
$A$、$A'$、および $(A')^\wedge$ の関係は Lemma
\ref{varieties-lemma-normalization-same-after-completion}（およびその証明）で見た。
\end{proof}




\section{アフィン開集合を見つけること}
\label{varieties-section-finding-affine-opens}

\noindent
Properties, Section \ref{properties-section-finding-affine-opens} で始めた議論を続ける。
分離スキーム上では、余次元 $1$ の任意に与えられた有限個の点を含む
アフィンを見つけられることが分かる。Proposition
\ref{varieties-proposition-finite-set-of-points-of-codim-1-in-affine} を参照せよ。

\medskip\noindent
次の補題は Descent, Lemma
\ref{descent-lemma-characterize-open-immersion} で改良する。

\begin{lemma}
\label{varieties-lemma-characterize-open-immersion}
$f : X \to Y$ をスキームの射とする。$X^0$ を $X$ の既約成分の生成点全体とする。
次を仮定する。
\begin{enumerate}
\item $f$ は分離的である。
\item 開被覆 $X = \bigcup U_i$ が存在し、
$f|_{U_i} : U_i \to Y$ は開埋め込みである。
\item $\xi, \xi' \in X^0$ かつ $\xi \not = \xi'$ ならば、
$f(\xi) \not = f(\xi')$ である。
\end{enumerate}
このとき $f$ は開埋め込みである。
\end{lemma}

\begin{proof}
$y = f(x) = f(x')$ と仮定する。特殊化 $y_0 \leadsto y$ を選び、ここで
$y_0$ は $Y$ の既約成分の生成点とする。$f$ は始域上局所的に同型なので、
特殊化 $x_0 \leadsto x$ および $x'_0 \leadsto x'$ であって、
$y_0 \leadsto y$ へ写るものを選べる。$x_0, x'_0 \in X^0$ であることに注意する。
従って仮定 (3) により $x_0 = x'_0$ である。
$f$ は分離的なので $x = x'$ と結論できる。従って $f$ は開埋め込みである。
\end{proof}

\begin{lemma}
\label{varieties-lemma-local-isomorphism}
$X \to S$ をスキームの射とする。$x \in X$ を、その像が $s \in S$ である点とする。
次を仮定する。
\begin{enumerate}
\item $\mathcal{O}_{X, x} = \mathcal{O}_{S, s}$。
\item $S$ は被約である。
\item $X \to S$ は有限型である。
\item $S$ は有限個の既約成分をもつ。
\end{enumerate}
このとき開近傍 $U$ が存在し、これは $x$ を含み、$f|_U$ は開埋め込みである。
\end{lemma}

\begin{proof}
$S$ の（有限個の）既約成分のうち $s$ を含まないものを除いてよい。
$S$ を $s$ のアフィン開近傍で置き換えてよい。
$X$ を $x$ のアフィン開近傍で置き換えてよい。
$S = \Spec(A)$ および $X = \Spec(B)$ と書く。
$\mathfrak q \subset B$、resp.\ $\mathfrak p \subset A$ を
$x$、resp.\ $s$ に対応する素イデアルとする。
$A$ は被約であり、$A$ のすべての極小素イデアルは $\mathfrak p$ に含まれるので、
$A \subset A_\mathfrak p$ である。$X \to S$ は有限型なので、
$B$ は $A$ 上有限型である。$b_1, \ldots, b_n \in B$ を、
$B$ を $A$ 上生成する元とする。$A_\mathfrak p = B_\mathfrak q$ なので、
$f \in A$、$f \not \in \mathfrak p$ および $a_i \in A$ であって、
$b_i$ と $a_i/f$ が $B_\mathfrak q$ で同じ像をもつものを見つけられる。
従って $g \in B$、$g \not \in \mathfrak q$ であって
$g(fb_i - a_i) = 0$ が $B$ で成り立つものを見つけられる。
このことから、$A_f \to B_{fg}$ の像は $b_1, \ldots, b_n$ の像を含み、
特に $g$ の像も含む。
$n \geq 0$ および $f' \in A$ を、$f'/f^n$ が $g$ の像
（$B_{fg}$ における）へ写るように選ぶ。$A_\mathfrak p = B_\mathfrak q$ なので、
$f' \not \in \mathfrak p$ である。従って $A_{ff'} \to B_{fg}$ は全射である。
最後に、$A_{ff'} \subset A_\mathfrak p = B_\mathfrak q$ であるから（上を参照）、
写像 $A_{ff'} \to B_{fg}$ は単射であり、従って同型である。
\end{proof}

\begin{lemma}
\label{varieties-lemma-points-in-affine}
$f : T \to X$ をスキームの射とする。$X^0$、resp.\ $T^0$ を
既約成分の生成点全体とする。$t_1, \ldots, t_m \in T$ を有限個の点とし、
その像を $x_j = f(t_j)$ とする。次を仮定する。
\begin{enumerate}
\item $T$ はアフィンである。
\item $X$ は準分離的である。
\item $X^0$ は有限である。
\item $f(T^0) \subset X^0$ であり、$f : T^0 \to X^0$ は単射である。
\item $\mathcal{O}_{X, x_j} = \mathcal{O}_{T, t_j}$。
\end{enumerate}
このとき $X$ のアフィン開集合で $x_1, \ldots, x_r$ を含むものが存在する。
\end{lemma}

\begin{proof}
Limits, Proposition \ref{limits-proposition-affine} を用いれば、
$X$ と $T$ が被約である場合へ直ちに帰着する。詳細は省略する。

\medskip\noindent
$X$ と $T$ は被約であると仮定する。Limits, Lemma
\ref{limits-lemma-relative-approximation} により、有向極限
$T = \lim_{i \in I} T_i$ として書け、ここで各スキームは $X$ 上有限表示で、
遷移射はアフィンである。Limits, Lemma \ref{limits-lemma-limit-affine} により、
$i \in I$ を、$T_i$ がアフィンとなるように選ぶ。
$T_i = \Spec(R_i)$ および $T = \Spec(R)$ と書く。
$R' \subset R$ を $R_i \to R$ の像とする。このとき $T' = \Spec(R')$ は
アフィン、被約、$X$ 上有限型であり、$T \to T'$ は支配的である。
$j = 1, \ldots, r$ に対し、$t'_j \in T'$ を $t_j$ の像とする。
局所環の写像
$$
\mathcal{O}_{X, x_j} \to
\mathcal{O}_{T', t'_j} \to
\mathcal{O}_{T, t_j}
$$
を考える。$(T')^0$ を $T'$ の既約成分の生成点全体とする。
特殊化 $\xi \leadsto t'_j$ を $\xi \in (T')^0$ としてとる。
$T \to T'$ は支配的なので、$\eta \in T^0$ で $\xi$ へ写るものを選べる
（注意：先験的には $\eta$ が $t_j$ へ特殊化するとは分からない）。
$\eta$ に仮定 (3) を適用すると、像 $\theta$（$\xi$ の $X$ におけるもの）は
$\mathcal{O}_{X, x_j}$ の極小素イデアルに対応することが分かる。
(5) の同型を介して $\xi$ を持ち上げると、特殊化
$\eta' \leadsto t_j$ を $\eta' \in T^0$ として得て、これは
$\theta \leadsto x_j$ へ写る。
(4) の単射性により $\eta = \eta'$ である。従って
$\mathcal{O}_{T', t'_j}$ の任意の極小素イデアルは
$\mathcal{O}_{T, t_j}$ の極小素イデアルの下にある。ゆえに
$\mathcal{O}_{T', t'_j} \to \mathcal{O}_{T, t_j}$ は単射であり、
上の二つの写像はいずれも同型である。

\medskip\noindent
Lemma \ref{varieties-lemma-local-isomorphism} により、開集合 $U \subset T'$ であって
すべての点 $t'_j$ を含み、$U \to X$ が Lemma
\ref{varieties-lemma-characterize-open-immersion} における局所同型となるものが存在する。
その補題により $U \to X$ は開埋め込みである。最後に Properties, Lemma
\ref{properties-lemma-ample-finite-set-in-affine} により、
開集合 $W \subset U \subset T'$ ですべての $t'_j$ を含むものを見つけられる。
$W$ の $X$ における像が、求めるアフィン開集合である。
\end{proof}




\begin{lemma}
\label{varieties-lemma-finite-set-codim-1-points-in-affine}
$X$ を整分離スキームとする。$x_1, \ldots, x_r \in X$ を有限個の点とし、
$\mathcal{O}_{X, x_i}$ が Noether で次元 $\leq 1$ であるとする。
このとき $X$ のアフィン開部分スキームで、$x_1, \ldots, x_r$ のすべてを
含むものが存在する。
\end{lemma}

\begin{proof}
$K$ を $X$ の有理関数体とする。$A_i = \mathcal{O}_{X, x_i}$ と置く。
このとき $A_i \subset K$ であり、$K$ は $A_i$ の分数体である。
$X$ は分離的なので、$x_i \not = x_j$ ならば、付値環
$\mathcal{O} \subset K$ で $A_i$ と $A_j$ の両方を支配するものは存在しえない。
実際、図式
$$
\xymatrix{
\Spec(\mathcal{O}) \ar[r] \ar[d] & \Spec(A_1) \ar[d] \\
\Spec(A_2) \ar[r] & X
}
$$
を考え、分離性の付値判定法（Schemes, Lemma
\ref{schemes-lemma-separated-implies-valuative}）を適用すると、
$x_i = x_j$ を得てしまう。従って Lemma \ref{varieties-lemma-glue-separated} により、
$A_i \otimes A_j \to K$ はすべての $i \not = j$ に対して全射である。
Lemma \ref{varieties-lemma-glue-a-bunch-of-local-rings} により、
$A = A_1 \cap \ldots \cap A_r$ は極大イデアルをちょうど $r$ 個もつ Noether
半局所環であり、それらを $\mathfrak m_1, \ldots, \mathfrak m_r$ とすれば
$A_i = A_{\mathfrak m_i}$ である。さらに
$$
\Spec(A) = \Spec(A_1) \cup \ldots \cup \Spec(A_r)
$$
は開被覆であり、この被覆の任意の二つの部分の共通部分は $\Spec(K)$ である。
従って与えられた射 $\Spec(A_i) \to X$ は貼り合わさってスキームの射
$$
\Spec(A) \longrightarrow X
$$
となり、$\mathfrak m_i$ を $x_i$ へ写し、局所環の同型を誘導する。
ゆえに Lemma \ref{varieties-lemma-points-in-affine} から結果が従う。
\end{proof}

\begin{lemma}
\label{varieties-lemma-extra-silly}
$A$ を環、$I \subset A$ をイデアル、
$\mathfrak p_1, \ldots, \mathfrak p_r$ を $A$ の素イデアル、
$\overline{f} \in A/I$ を元とする。$I \not \subset \mathfrak p_i$ が
すべての $i$ に対して成り立つならば、$f \in A$、$f \not \in \mathfrak p_i$
であって、$\overline{f}$ へ $A/I$ で写るものが存在する。
\end{lemma}

\begin{proof}
$\mathfrak p_i$ の間に包含関係がないと仮定してよい
（より小さい素イデアルを除けばよい）。まず任意の $f \in A$ を、
$\overline{f}$ を持ち上げるものとして選ぶ。$S$ を、
$s \in \{1, \ldots, r\}$ のうち $f \in \mathfrak p_s$ となるもの全体とする。
$S$ が空ならば終わりである。
そうでなければ、イデアル $J = I \prod_{i \not \in S} \mathfrak p_i$ を考える。
$J$ は $\mathfrak p_s$ に含まれない（$s \in S$）ことに注意する。
実際、$\mathfrak p_i$ の間に包含がなく、$I$ はどの $\mathfrak p_i$ にも含まれない。
従って Algebra, Lemma \ref{algebra-lemma-silly} により、
$g \in J$、$g \not \in \mathfrak p_s$ が $s \in S$ に対して成り立つように選べる。
このとき $f + g$ が補題で求める解である。
\end{proof}

\begin{lemma}
\label{varieties-lemma-finite-set-codim-1-points-in-affine-per-component}
$X$ をスキームとする。$T \subset X$ を有限個の点の集合とする。次を仮定する。
\begin{enumerate}
\item $X$ は有限個の既約成分 $Z_1, \ldots, Z_t$ をもつ。
\item $Z_i \cap T$ は、$Z_i$ に対応する被約誘導部分スキームの
あるアフィン開集合に含まれる。
\end{enumerate}
このとき $X$ のアフィン開部分スキームで $T$ を含むものが存在する。
\end{lemma}

\begin{proof}
Limits, Proposition \ref{limits-proposition-affine} を用いれば、
$X$ が被約である場合へ直ちに帰着する。詳細は省略する。
以下の証明では、$X$ の任意の閉部分集合に、誘導された被約閉部分スキーム構造を入れる。

\medskip\noindent
帰納法により、$U \subset Z_1 \cup \ldots \cup Z_r$ というアフィン開集合で
$T \cap (Z_1 \cup \ldots \cup Z_r)$ を含むものを見つけられることを示す。
$r = 1$ ならば仮定により成り立つ。$r > 1$ とし、
$U \subset Z_1 \cup \ldots \cup Z_{r -  1}$ を
$T \cap (Z_1 \cup \ldots \cup Z_{r - 1})$ を含むアフィン開集合とする。
$V \subset X_r$ を $T \cap Z_r$ を含むアフィン開集合とする
（仮定により存在する）。このとき $U \cap V$ は
$T \cap ( Z_1 \cup \ldots \cup Z_{r - 1} ) \cap Z_r$ を含む。従って
$$
\Delta = (U \cap Z_r) \setminus (U \cap V)
$$
は $T$ のどの元も含まない。$\Delta$ は $U$ の閉部分集合であることに注意する。
素イデアル回避（Algebra, Lemma \ref{algebra-lemma-silly}）により、
標準開集合 $U'$（$U$ の）で、$T \cap U$ を含み $\Delta$ を避けるものを見つけられる。
すなわち $U' \cap Z_r \subset U \cap V$ である。
$U$ を $U'$ で置き換えた後、$U \cap V$ は $U$ で閉であると仮定できる。

\medskip\noindent
上と同じ議論により集合
$\Delta' = (U \cap (Z_1 \cup \ldots \cup Z_{r - 1})) \setminus (U \cap V)$
も $T$ のどの元も含まない。従って $h \in \mathcal{O}(V)$ であって、
$D(h) \subset V$ が $T \cap V$ を含み、かつ
$U \cap D(h) \subset U \cap V$ となるものを見つけられる。
$U \cap V$ は $U$ で閉なので、Lemma \ref{varieties-lemma-extra-silly} を用いて、
$g \in \mathcal{O}(U)$ で、その $U \cap V$ への制限が $h$ の
$U \cap V$ への制限に等しく、$T \cap U \subset D(g)$ となるものを見つけられる。
そこで $U$ を $D(g)$ で、$V$ を $D(h)$ で置き換えると、
$U \cap V$ が $U$ と $V$ の両方で閉である状況に達する。
この場合、スキーム $U \cup V$ は Limits, Lemma
\ref{limits-lemma-affines-glued-in-closed-affine} によりアフィンである。
これで帰納段階、従って補題が証明された。
\end{proof}

\noindent
以上の内容から次の結論を得る。

\begin{proposition}
\label{varieties-proposition-finite-set-of-points-of-codim-1-in-affine}
$X$ を、任意の準コンパクト開集合が有限個の既約成分をもつ分離スキームとする。
$x_1, \ldots, x_r \in X$ を、$\mathcal{O}_{X, x_i}$ が Noether で
次元 $\leq 1$ となる点とする。このとき $X$ のアフィン開部分スキームで、
$x_1, \ldots, x_r$ のすべてを含むものが存在する。
\end{proposition}

\begin{proof}
$X$ を $x_1, \ldots, x_r$ を含む準コンパクト開集合で置き換えられるので、
$X$ は有限個の既約成分をもつと仮定してよい。Lemma
\ref{varieties-lemma-finite-set-codim-1-points-in-affine-per-component} により、
$X$ が整である場合へ帰着する。この場合は Lemma
\ref{varieties-lemma-finite-set-codim-1-points-in-affine} である。
\end{proof}




\section{曲線}
\label{varieties-section-curves}

\noindent
Stacks project では、次を曲線の定義として用いる。

\begin{definition}
\label{varieties-definition-curve}
$k$ を体とする。{\it 曲線}とは、次元 $1$ の $k$ 上の多様体である。
\end{definition}

\noindent
$k$ 上の曲線の標準的な二例は、アフィン直線 $\mathbf{A}^1_k$ と
射影直線 $\mathbf{P}^1_k$ である。スキーム $X = \Spec(k[x, y]/(f))$
が曲線であるための必要十分条件は、$f \in k[x, y]$ が既約であることである。

\medskip\noindent
曲線の定義には、多様体の定義と同じ問題がある。Definition
\ref{varieties-definition-variety} の後の議論を参照せよ。さらに、この定義では、
各曲線に定義体が指定されている。例えば $X = \Spec(\mathbf{C}[x])$ は
$\mathbf{C}$ 上の曲線だが、これを $\mathbf{R}$ 上の曲線と見ることもできる。
スキーム $\Spec(\mathbf{Z})$ は曲線ではないが、スキーム
$\Spec(\mathbf{Z})$ と $\mathbf{A}^1_{\mathbf{F}_p}$ は多くの点で同様に振る舞う。

\begin{lemma}
\label{varieties-lemma-proper-minus-point}
$X$ を、次元 $> 0$ の、体 $k$ 上の分離既約スキームとする。
$x \in X$ を閉点とする。開部分スキーム $X \setminus \{x\}$ は
$k$ 上固有ではない。
\end{lemma}

\begin{proof}
$X$ は既約なので、$U = X \setminus \{x\}$ は $X$ で閉ではない。
特に埋め込み $U \to X$ は固有ではない。Morphisms, Lemma
\ref{morphisms-lemma-image-proper-scheme-closed} により（ここで
$X$ が分離的であることを用いる）、$U \to \Spec(k)$ も固有ではない。
\end{proof}

\begin{lemma}
\label{varieties-lemma-dim-1-quasi-projective}
$X$ を体 $k$ 上の分離有限型スキームとする。$\dim(X) \leq 1$ ならば、
$X$ は $k$ 上 H-準射影的である。
\end{lemma}

\begin{proof}
Proposition \ref{varieties-proposition-dim-1-noetherian-separated-has-ample} により、
スキーム $X$ は豊富な可逆層 $\mathcal{L}$ をもつ。
Morphisms, Lemma \ref{morphisms-lemma-quasi-projective-finite-type-over-S} により、
$X$ は $\mathbf{P}^n_k$ の、$\Spec(k)$ 上の局所閉部分スキームと同型である。
これは $k$ 上 H-準射影的であることの定義である。Morphisms, Definition
\ref{morphisms-definition-quasi-projective} を参照せよ。
\end{proof}

\begin{lemma}
\label{varieties-lemma-dim-1-proper-projective}
$X$ を体 $k$ 上の固有スキームとする。$\dim(X) \leq 1$ ならば、
$X$ は $k$ 上 H-射影的である。
\end{lemma}

\begin{proof}
Lemma \ref{varieties-lemma-dim-1-quasi-projective} により、$X$ は
$\mathbf{P}^n_k$ の局所閉部分スキームであり、これはある体 $k$ に対するものである。
$X$ は $k$ 上固有なので、$X$ は $\mathbf{P}^n_k$ の閉部分スキームである
（Morphisms, Lemma \ref{morphisms-lemma-image-proper-scheme-closed}）。
\end{proof}

\begin{lemma}
\label{varieties-lemma-dim-1-projective-completion}
$X$ を $k$ 上有限型な分離スキームとする。$\dim(X) \leq 1$ ならば、
次の性質をもつ開埋め込み $j : X \to \overline{X}$ が存在する。
\begin{enumerate}
\item $\overline{X}$ は $k$ 上 H-射影的である。すなわち
$\overline{X}$ は $\mathbf{P}^d_k$（ある $d$ に対する）の閉部分スキームである。
\item $j(X) \subset \overline{X}$ は稠密かつスキーム論的に稠密である。
\item $\overline{X} \setminus X = \{x_1, \ldots, x_n\}$ が、ある閉点
$x_i \in \overline{X}$ に対して成り立つ。
\end{enumerate}
\end{lemma}

\begin{proof}
Lemma \ref{varieties-lemma-dim-1-quasi-projective} により、$X$ は
$\mathbf{P}^d_k$（ある $d$ に対する）の局所閉部分スキームであると仮定してよい。
$\overline{X} \subset \mathbf{P}^d_k$ を $X \to \mathbf{P}^d_k$ の
スキーム論的像とする。Morphisms, Definition
\ref{morphisms-definition-scheme-theoretic-image} を参照せよ。
Morphisms, Lemma \ref{morphisms-lemma-quasi-compact-immersion} の記述から
性質 (1) と (2) を得る。例えば生成点を見れば
$\dim(X) = 1 \Rightarrow \dim(\overline{X}) = 1$ である。
Lemma \ref{varieties-lemma-dimension-locally-algebraic} を参照せよ。
$\overline{X}$ は Noether なので、
$\overline{X} \setminus X = \{x_1, \ldots, x_n\}$ は有限個の閉点の集合である。
\end{proof}

\begin{lemma}
\label{varieties-lemma-reduced-dim-1-projective-completion}
$X$ を $k$ 上有限型な分離スキームとする。$X$ が被約で
$\dim(X) \leq 1$ ならば、次を満たす開埋め込み
$j : X \to \overline{X}$ が存在する。
\begin{enumerate}
\item $\overline{X}$ は $k$ 上 H-射影的である。すなわち
$\overline{X}$ は $\mathbf{P}^d_k$（ある $d$ に対する）の閉部分スキームである。
\item $j(X) \subset \overline{X}$ は稠密かつスキーム論的に稠密である。
\item $\overline{X} \setminus X = \{x_1, \ldots, x_n\}$ が、ある閉点
$x_i \in \overline{X}$ に対して成り立つ。
\item 局所環 $\mathcal{O}_{\overline{X}, x_i}$ は
$i = 1, \ldots, n$ に対して離散付値環である。
\end{enumerate}
\end{lemma}

\begin{proof}
$j : X \to \overline{X}$ を Lemma
\ref{varieties-lemma-dim-1-projective-completion} におけるものとする。
$X'$ を、$\overline{X}$ の $X$ における正規化とする。
Lemma \ref{varieties-lemma-relative-normalization-finite} により、射
$X' \to \overline{X}$ は有限である。Morphisms, Lemma
\ref{morphisms-lemma-finite-projective} により $X' \to \overline{X}$ は射影的である。
Morphisms, Lemma
\ref{morphisms-lemma-projective-over-quasi-projective-is-H-projective} により
$X' \to \overline{X}$ は H-射影的である。Morphisms, Lemma
\ref{morphisms-lemma-H-projective-composition} により
$X' \to \Spec(k)$ は H-射影的である。
$\{x'_1, \ldots, x'_m\} \subset X'$ を
$\{x_1, \ldots, x_n\} = \overline{X} \setminus X$ の逆像とする。
このとき $\dim(\mathcal{O}_{X', x'_i}) = 1$ は、すべての
$1 \leq i \leq m$ に対して成り立つ。従って局所環
$\mathcal{O}_{X', x'}$ は Morphisms, Lemma
\ref{morphisms-lemma-relative-normalization-normal-codim-1} により
離散付値環である。そこで $X \to X'$ および
$\{x'_1, \ldots, x'_m\}$ が求めるものである。
\end{proof}

\begin{lemma}
\label{varieties-lemma-dim-1-projective-completion-after-insep}
$X$ を $k$ 上有限型な分離スキームで、$\dim(X) \leq 1$ とする。
このとき、次の可換図式
$$
\xymatrix{
\overline{Y}_1 \amalg \ldots \amalg \overline{Y}_n \ar[rd] &
Y_1 \amalg \ldots \amalg Y_n \ar[r]_-\nu \ar[d] \ar[l]^j &
X_{k'} \ar[r] \ar[d] &
X \ar[d]^f \\
& \Spec(k'_1) \amalg \ldots \amalg \Spec(k'_n) \ar[r] &
\Spec(k') \ar[r] &
\Spec(k)
}
$$
が存在し、次の性質をもつ。
\begin{enumerate}
\item $k'/k$ は有限純非分離体拡大である。
\item $\nu$ は $X_{k'}$ の正規化である。
\item $j$ は稠密な像をもつ開埋め込みである。
\item $k'_i/k'$ は $i = 1, \ldots, n$ に対して有限分離拡大である。
\item $\overline{Y}_i$ は滑らか、射影的、幾何学的既約、かつ次元
$\leq 1$ の $k'_i$ 上のスキームである。
\end{enumerate}
\end{lemma}

\begin{proof}
$X$ をその被約化で置き換えてよいので、$X$ は被約であると仮定し、以後そうする。
$X \to \overline{X}$ を Lemma
\ref{varieties-lemma-reduced-dim-1-projective-completion} におけるものとして選ぶ。
$\overline{X}$ に対して補題を示せれば、$X$ に対しても従う（詳細は省略する）。
従って $X$ は射影的であると仮定し、以後そうする。

\medskip\noindent
$k'/k$ を有限純非分離拡大であって、$X_{k'}$ の正規化が
$k'$ 上幾何学的正規となるものとして選ぶ。Lemma
\ref{varieties-lemma-finite-extension-geometrically-normal} を参照せよ。
$Y = (X_{k'})^\nu$ を正規化と記す。正規化の性質については Section
\ref{varieties-section-normalization} を参照せよ。$Y$ は、次元 $\leq 1$ では
正規と正則が同じなので、幾何学的正則である。Properties, Lemma
\ref{properties-lemma-normal-dimension-1-regular} を参照せよ。
従って Lemma \ref{varieties-lemma-geometrically-regular-smooth} により
$Y$ は $k'$ 上滑らかである。$Y = Y_1 \amalg \ldots \amalg Y_n$ を
$Y$ の既約成分への分解とする。$k'_i = \Gamma(Y_i, \mathcal{O}_{Y_i})$
と置く。これらは Lemma
\ref{varieties-lemma-proper-geometrically-reduced-global-sections} により
$k'$ の有限分離拡大である。Lemma \ref{varieties-lemma-baby-stein} により証明が完了する。
\end{proof}

\begin{lemma}
\label{varieties-lemma-regular-point-on-curve}
$k$ を体とする。$X$ を $k$ 上の曲線とする。$x \in X$ を閉点とする。
$x$ を、$X$ の（被約）閉部分スキームであってイデアル層 $\mathcal{I}$ をもつものとみなす。
次の条件は同値である。
\begin{enumerate}
\item $\mathcal{O}_{X, x}$ は正則である。
\item $\mathcal{O}_{X, x}$ は正規である。
\item $\mathcal{O}_{X, x}$ は離散付値環である。
\item $\mathcal{I}$ は可逆 $\mathcal{O}_X$ 加群である。
\item $x$ は $X$ 上の有効 Cartier 因子である。
\end{enumerate}
$k$ が完全であるか、または $\kappa(x)$ が $k$ 上分離的ならば、これらはさらに
\begin{enumerate}
\item[(6)] $X \to \Spec(k)$ は $x$ において滑らかである。
\end{enumerate}
とも同値である。
\end{lemma}

\begin{proof}
$X$ は曲線なので、局所環 $\mathcal{O}_{X, x}$ は次元 $1$ の Noether 局所整域である
（Lemma \ref{varieties-lemma-dimension-locally-algebraic}）。(4) と (5) は定義により同値であり、
$\mathcal{I}_x = \mathfrak m_x \subset \mathcal{O}_{X, x}$ が一つの生成元をもつこととも同値である
（Divisors, Lemma \ref{divisors-lemma-effective-Cartier-in-points}）。従って (1), (2),
(3), (4), (5) の同値性は Algebra, Lemma \ref{algebra-lemma-characterize-dvr} から従う。
最後の主張は、$k$ が完全な場合には Lemma
\ref{varieties-lemma-dense-smooth-open-variety-over-perfect-field} から従う。
$\kappa(x)/k$ が分離的な場合には、同値性は Algebra, Lemma
\ref{algebra-lemma-separable-smooth} から従う。
\end{proof}

\begin{remark}
\label{varieties-remark-smooth-projective-compactification}
$k$ を体とする。$X$ を $k$ 上の正則曲線とする。Lemmas
\ref{varieties-lemma-regular-point-on-curve} および
\ref{varieties-lemma-reduced-dim-1-projective-completion} により、非特異射影曲線
$\overline{X}$ が存在し、これは $X$ のコンパクト化である。すなわち、補集合が有限個の閉点からなる
開埋め込み $j : X \to \overline{X}$ が存在する。$k$ が完全ならば、
$X$ と $\overline{X}$ は $k$ 上滑らかであり、$\overline{X}$ は
$X$ の滑らかな射影的コンパクト化である。
\end{remark}

\noindent
$X$ が $k$ 上のアフィンスキームであり、かつ $k$ 上固有ならば、$X$ は $k$ 上有限であり
（Morphisms, Lemma \ref{morphisms-lemma-finite-proper}）、従って次元 $0$ である
（Algebra, Lemma \ref{algebra-lemma-finite-dimensional-algebra} および
Proposition \ref{algebra-proposition-dimension-zero-ring}）ことに注意せよ。
従って、次元 $> 0$ の $k$ 上のスキームがアフィンであると同時に $k$ 上固有であることはない。
ゆえに、次の補題における二つの可能性は互いに排他的である。

\begin{lemma}
\label{varieties-lemma-curve-affine-projective}
$X$ を $k$ 上の曲線とする。このとき、$X$ はアフィンスキームであるか、または
$X$ は $k$ 上 H-射影的である。
\end{lemma}

\begin{proof}
Lemma \ref{varieties-lemma-reduced-dim-1-projective-completion} におけるように、
$X \to \overline{X}$ を、$\overline{X} \setminus X = \{x_1, \ldots, x_r\}$
を満たすものとして選ぶ。このとき $\overline{X}$ も曲線である。
$r = 0$ ならば $X = \overline{X}$ は $k$ 上 H-射影的である。
従って $r \geq 1$ と仮定してよく、目標は $X$ がアフィンであることを示すことである。
Lemma \ref{varieties-lemma-open-in-affine-curve-affine} により、
$\overline{X} \setminus \{x_1\}$ がアフィンであることを示せば十分である。
これは、次の段落で述べる主張へ帰着させる。

\medskip\noindent
$X$ を $k$ 上 H-射影的な曲線とする。$x \in X$ を、
$\mathcal{O}_{X, x}$ が離散付値環となる閉点とする。主張：
$U = X \setminus \{x\}$ はアフィンである。Lemma
\ref{varieties-lemma-regular-point-on-curve} により、点 $x$ は $X$ の有効 Cartier 因子を定める。
$n \geq 1$ に対して、$nx = x + \ldots + x$ により $n$ 重和を表す。Divisors, Definition
\ref{divisors-definition-sum-effective-Cartier-divisors} を参照せよ。
$\mathcal{O}_{nx}$ により、$nx$ の構造層を $X$ 上の連接加群とみなしたものを表す。
局所スキーム $nx$ 上の任意の可逆加群は自明なので、Divisors, Remark
\ref{divisors-remark-ses-regular-section} の最初の短完全列は、この場合
$$
0 \to \mathcal{O}_X
\xrightarrow{1} \mathcal{O}_X(nx) \to \mathcal{O}_{nx} \to 0
$$
となる。$\dim_k H^0(X, \mathcal{O}_{nx}) \geq n$ であることに注意せよ。
実際、Lemma \ref{varieties-lemma-chi-tensor-finite} により
$H^0(X, \mathcal{O}_{nx}) = \mathcal{O}_{X, x}/(\pi^n)$ である。ここで $\pi$ は
$\mathcal{O}_{X, x}$ の一様化元であり、冪 $\pi^i$ は $k$ 上線形独立な元として
$\mathcal{O}_{X, x}/(\pi^n)$ に写る。ただし
$i = 0, 1, \ldots, n - 1$ である。Cohomology of Schemes, Lemma
\ref{coherent-lemma-proper-over-affine-cohomology-finite} により
$\dim_k H^1(X, \mathcal{O}_X) < \infty$ である。
$n > \dim_k H^1(X, \mathcal{O}_X)$ ならば、コホモロジー長完全列から、
$s \in \Gamma(X, \mathcal{O}_X(nx))$ が存在し、それは $\mathcal{O}_X$ の切断ではないと結論できる。
この性質をもつ $n$ を最小に取れば、$s$ は茎
$\left(\mathcal{O}_X(nx)\right)_x$ の生成元に写る。そうでなければ、
$\mathcal{O}_X((n - 1)x) \subset \mathcal{O}_X(nx)$ の切断を定めるからである。
この $n$ に対して、$s_0 = 1$ と $s_1 = s$ は可逆加群
$\mathcal{L} = \mathcal{O}_X(nx)$ を生成すると結論する。

\medskip\noindent
Constructions, Section \ref{constructions-section-projective-space} の対応する射
$f = \varphi_{\mathcal{L}, (s_0, s_1)} : X \to \mathbf{P}^1_k$ を考える。
$D_{+}(T_0)$ の逆像は $U = X \setminus \{x\}$ であることに注意せよ。
実際、切断 $s_0$ は $\mathcal{L}$ の切断として $x$ においてのみ消える。
特に $f$ は定数射ではない、すなわち $\Im(f)$ は二点以上をもつ。
従って $f$ は生成点 $\eta$（$X$ の生成点）を $\mathbf{P}^1_k$ の生成点へ写さなければならない。
従って $y \in \mathbf{P}^1_k$ が閉点ならば、$f^{-1}(\{y\})$ は
$X$ の、$\eta$ を含まない閉集合であり、それゆえ有限である。
最後に $f$ は固有である\footnote{実際、H-射影多様体は Morphisms, Lemma
\ref{morphisms-lemma-projective-is-quasi-projective-proper} により固有多様体である。
始域が固有多様体である多様体間の射は、Morphisms, Lemma
\ref{morphisms-lemma-image-proper-scheme-closed} により固有射である。}。
Cohomology of Schemes, Lemma
\ref{coherent-lemma-proper-finite-fibre-finite-in-neighbourhood}\footnote{形式関数定理に
依存するこの補題を用いないこともできる。実際、$X$ は射影的なので、
固有射 $f : X \to Y$ が、$k$ 上準射影的なスキームの間の有限なファイバーをもつ射ならば、
有限であることを示せば十分である。このため、$X$ のアフィン開集合であって、
$f$ の点 $y$ 上のファイバーを含むものを選ぶ。これは、$k$ 上準射影的なスキームの
任意の有限点集合がアフィン開集合に含まれることによる。
$Y$ を $y$ の小さなアフィン近傍へ縮小すると、アフィン間の固有射の場合へ帰着する。
このような射は Morphisms, Lemma
\ref{morphisms-lemma-integral-universally-closed} により有限である。} により、
$f$ は有限であると結論する。従って $U = f^{-1}(D_{+}(T_0))$ はアフィンである。
\end{proof}

\noindent
次の補題を Lemma \ref{varieties-lemma-proper-minus-point} と合わせると、次が分かる。
分離スキーム $X$ が次元 $1$ で $k$ 上有限型であるとき、$X \setminus Z$ は、
閉部分集合 $Z$ が $X$ の各既約成分と交わるならばアフィンである。

\begin{lemma}
\label{varieties-lemma-dim-1-nonproper-affine}
$X$ を $k$ 上有限型な分離スキームとする。$\dim(X) \leq 1$ であり、
$X$ の次元 $1$ の既約成分がいずれも固有でなければ、$X$ はアフィンである。
\end{lemma}

\begin{proof}
$X = \bigcup X_i$ を $X$ の既約成分への分解とする。$X_i$ を整スキームとみなす
（被約誘導スキーム構造を用いる。Schemes, Definition
\ref{schemes-definition-reduced-induced-scheme} を参照せよ）。特に $X_i$ は一点集合
（従ってアフィン）であるか、または曲線であり、従って Lemma
\ref{varieties-lemma-curve-affine-projective} によりアフィンである。
このとき $\coprod X_i \to X$ は有限全射であり、$\coprod X_i$ はアフィンである。
従って Cohomology of Schemes, Lemma
\ref{coherent-lemma-image-affine-finite-morphism-affine-Noetherian} により、
$X$ はアフィンである。
\end{proof}

\section{曲線上の次数}
\label{varieties-section-divisors-curves}

\noindent
体上の次元 $1$ の固有スキーム上で、可逆層、より一般には局所自由層の
次数を定義することから始める。Section \ref{varieties-section-euler} では、連接層
$\mathcal{F}$ の固有スキーム $X$ 上での Euler 標数を、基礎体 $k$ に関して、公式
$$
\chi(X, \mathcal{F}) = \sum (-1)^i \dim_k H^i(X, \mathcal{F}).
$$
によって定義した。

\begin{definition}
\label{varieties-definition-degree-invertible-sheaf}
$k$ を体とし、$X$ を次元 $\leq 1$ の固有スキームであって $k$ 上のものとし、
$\mathcal{L}$ を可逆 $\mathcal{O}_X$ 加群とする。$\mathcal{L}$ の {\it 次数}を
$$
\deg(\mathcal{L}) = \chi(X, \mathcal{L}) - \chi(X, \mathcal{O}_X)
$$
によって定義する。より一般に、$\mathcal{E}$ が階数 $n$ の局所自由層ならば、
$\mathcal{E}$ の {\it 次数}を
$$
\deg(\mathcal{E}) = \chi(X, \mathcal{E}) - n\chi(X, \mathcal{O}_X)
$$
によって定義する。
\end{definition}

\noindent
これは三つ組 $\mathcal{E}/X/k$ に依存することに注意せよ。
$X$ が非連結で、$\mathcal{E}$ が有限局所自由（ただし定階数とは限らない）ならば、
$\mathcal{E}$ の $X$ の連結成分への制限の次数を足し合わせることで、定義を修正できる。
$\mathcal{E}$ が単に連接層である場合、定義を拡張する方法は複数ある\footnote{
$X$ が固有曲線で、$\mathcal{F}$ が $X$ 上の連接層ならば、次数を
$\chi(X, \mathcal{F}) - r\chi(X, \mathcal{O}_X)$ と定義することが多い。
ここで $r = \dim_{\kappa(\xi)} \mathcal{F}_\xi$ は、$\mathcal{F}$ の生成点
$\xi$（$X$ の生成点）における階数である。}。一連の補題により、この定義が次数に期待される
すべての性質をもつことを示す。

\begin{lemma}
\label{varieties-lemma-degree-base-change}
$k'/k$ を体拡大とする。$X$ を次元 $\leq 1$ の固有スキームであって $k$ 上のものとする。
$\mathcal{E}$ を局所自由 $\mathcal{O}_X$ 加群で定階数 $n$ のものとする。このとき、
$\mathcal{E}/X/k$ の次数は $\mathcal{E}_{k'}/X_{k'}/k'$ の次数に等しい。
\end{lemma}

\begin{proof}
より正確には $X_{k'} = X \times_{\Spec(k)} \Spec(k')$ と置く。
$\mathcal{E}_{k'} = p^*\mathcal{E}$ と置く。ただし $p : X_{k'} \to X$ は射影である。
Cohomology of Schemes, Lemma \ref{coherent-lemma-flat-base-change-cohomology} により
$H^i(X_{k'}, \mathcal{E}_{k'}) = H^i(X, \mathcal{E}) \otimes_k k'$
および
$H^i(X_{k'}, \mathcal{O}_{X_{k'}}) = H^i(X, \mathcal{O}_X) \otimes_k k'$
が成り立つ。従って Euler 標数は変わらず、それゆえ次数も変わらない。
\end{proof}

\begin{lemma}
\label{varieties-lemma-degree-additive}
$k$ を体とする。$X$ を次元 $\leq 1$ の固有スキームであって $k$ 上のものとする。
$0 \to \mathcal{E}_1 \to \mathcal{E}_2 \to \mathcal{E}_3 \to 0$ を、いずれも
有限定階数の局所自由 $\mathcal{O}_X$ 加群からなる短完全列とする。このとき
$$
\deg(\mathcal{E}_2) = \deg(\mathcal{E}_1) + \deg(\mathcal{E}_3)
$$
である。
\end{lemma}

\begin{proof}
Euler 標数の加法性（Lemma \ref{varieties-lemma-euler-characteristic-additive}）と
階数の加法性から直ちに従う。
\end{proof}

\begin{lemma}
\label{varieties-lemma-degree-birational-pullback}
$k$ を体とする。$f : X' \to X$ を、次元 $\leq 1$ の固有スキームであって $k$ 上のものの間の
双有理射とする。このとき
$$
\deg(f^*\mathcal{E}) = \deg(\mathcal{E})
$$
が任意の定階数有限局所自由層について成り立つ。より一般には、$f$ が次元 $1$ の
既約成分の間の全単射を誘導し、対応する生成点における局所環の同型を誘導すれば十分である。
\end{lemma}

\begin{proof}
射 $f$ は固有であり（Morphisms, Lemma
\ref{morphisms-lemma-image-proper-scheme-closed}）、そのファイバーの次元は $\leq 0$ である。
従って $f$ は有限である（Cohomology of Schemes, Lemma
\ref{coherent-lemma-proper-finite-fibre-finite-in-neighbourhood}）。従って
$$
Rf_*f^*\mathcal{E} = f_*f^*\mathcal{E} =
\mathcal{E} \otimes_{\mathcal{O}_X} f_*\mathcal{O}_{X'}
$$
である。$f$ は次元 $1$ のすべての既約成分の生成点において局所環の同型を
誘導するので、核と余核
$$
0 \to \mathcal{K} \to \mathcal{O}_X \to f_*\mathcal{O}_{X'}
\to \mathcal{Q} \to 0
$$
の台の次元は $\leq 0$ である。$\mathcal{E}$ は局所自由なので、これを
$\mathcal{E}$ とテンソルしても完全列のままであることに注意せよ。従って
\begin{align*}
\chi(X, \mathcal{E}) - \chi(X', f^*\mathcal{E})
& =
\chi(X, \mathcal{E}) - \chi(X, f_*f^*\mathcal{E}) \\
& =
\chi(X, \mathcal{E}) - \chi(X, \mathcal{E} \otimes f_*\mathcal{O}_{X'}) \\
& =
\chi(X, \mathcal{K} \otimes \mathcal{E}) -
\chi(X, \mathcal{Q} \otimes \mathcal{E}) \\
& =
n\chi(X, \mathcal{K}) -
n\chi(X, \mathcal{Q}) \\
& =
n\chi(X, \mathcal{O}_X) - n\chi(X, f_*\mathcal{O}_{X'}) \\
& =
n\chi(X, \mathcal{O}_X) - n\chi(X', \mathcal{O}_{X'})
\end{align*}
を得て、これは求める主張を証明する。第一の等式は $f$ が有限であることから従う。
Cohomology of Schemes, Lemma \ref{coherent-lemma-relative-affine-cohomology} を参照せよ。
第二の等式は射影公式による。Cohomology, Lemma
\ref{cohomology-lemma-projection-formula} を参照せよ。第三の等式は Euler 標数の加法性による。
Lemma \ref{varieties-lemma-euler-characteristic-additive} を参照せよ。第四の等式は Lemma
\ref{varieties-lemma-chi-tensor-finite} による。
\end{proof}

\begin{lemma}
\label{varieties-lemma-degree-on-proper-curve}
$k$ を体とする。$X$ を $k$ 上の固有曲線で、生成点を $\xi$ とする。
$\mathcal{E}$ を局所自由 $\mathcal{O}_X$ 加群で階数 $n$ のものとし、$\mathcal{F}$ を
連接 $\mathcal{O}_X$ 加群とする。このとき
$$
\chi(X, \mathcal{E} \otimes \mathcal{F}) =
r \deg(\mathcal{E}) + n \chi(X, \mathcal{F})
$$
である。ここで $r = \dim_{\kappa(\xi)} \mathcal{F}_\xi$ は $\mathcal{F}$ の階数である。
\end{lemma}

\begin{proof}
$\mathcal{P}$ を、連接層 $\mathcal{F}$ が $X$ 上で補題の公式を満たすことを表す
性質とする。Cohomology of Schemes, Lemma \ref{coherent-lemma-property} の仮定 (1) と (2) が
$\mathcal{P}$ に対して成り立つと主張する。実際、Euler 標数と階数 $r$ は連接層の
短完全列に関して加法的なので、(1) は成り立つ。(2) も成り立つ：$Z = X$ ならば
$\mathcal{G} = \mathcal{O}_X$ と取ることができ、次数の定義により
$\mathcal{P}(\mathcal{O}_X)$ は真である。$i : Z \to X$ が閉点の包含ならば、
$\mathcal{G} = i_*\mathcal{O}_Z$ と取ることができ、$\mathcal{P}$ は Lemma
\ref{varieties-lemma-chi-tensor-finite} およびこの場合 $r = 0$ であることにより成り立つ。
\end{proof}

\noindent
$k$ を体とする。$X$ を $k$ 上有限型で次元 $\leq 1$ のスキームとする。
$C_i \subset X$, $i = 1, \ldots, t$ を次元 $1$ の既約成分とする。
$C_i$ には被約誘導スキーム構造を入れてスキームとみなす。$\xi_i \in C_i$ を
生成点とする。{\it $C_i$ の $X$ における重複度}を長さ
$$
m_i = \text{length}_{\mathcal{O}_{X, \xi_i}} \mathcal{O}_{X, \xi_i}
$$
として定義する。$\mathcal{O}_{X, \xi_i}$ は零次元 Noether 局所環であり、それゆえ
自分自身の上で有限長をもつので、これは意味をもつ（Algebra, Proposition
\ref{algebra-proposition-dimension-zero-ring}）。追加の情報については Chow Homology,
Section \ref{chow-section-cycle-of-closed-subscheme} を参照せよ。局所自由層の次数は、
その既約成分への制限にのみ依存することが分かる。

\begin{lemma}
\label{varieties-lemma-degree-in-terms-of-components}
$k$ を体とする。$X$ を次元 $\leq 1$ の固有スキームであって $k$ 上のものとする。
$\mathcal{E}$ を局所自由 $\mathcal{O}_X$ 加群で階数 $n$ のものとする。このとき
$$
\deg(\mathcal{E}) = \sum m_i \deg(\mathcal{E}|_{C_i})
$$
である。ここで $C_i \subset X$, $i = 1, \ldots, t$ は、被約誘導スキーム構造を入れた
次元 $1$ の既約成分であり、$m_i$ は $C_i$ の $X$ における重複度である。
\end{lemma}

\begin{proof}
$C_i \to \Spec(k)$ は次元 $1$ で固有なので、この主張は意味をもつことに注意せよ。
以下では、$X$ の任意の閉部分スキームが $k$ 上固有であることを用いる。これは Morphisms,
Lemmas \ref{morphisms-lemma-closed-immersion-proper} および
\ref{morphisms-lemma-composition-proper} による。開部分スキーム
$U_i = X \setminus (\bigcup_{j \not = i} C_j)$ を考え、$X_i \subset X$ を
$U_i$ のスキーム論的閉包とする。$X_i \cap U_i = U_i$（スキーム論的）であり、
$X_i \cap U_j = \emptyset$（集合論的）が $i \not = j$ に対して成り立つことに注意せよ。
これは Morphisms, Lemma \ref{morphisms-lemma-quasi-compact-immersion} の
スキーム論的閉包の記述から従う。従って Lemma
\ref{varieties-lemma-degree-birational-pullback} を射 $X' = \coprod X_i \to X$ に適用できる。
$C_i \subset X_i$（スキーム論的）であること、また $C_i$ の $X_i$ における重複度が
$C_i$ の $X$ における重複度に等しいことは明らかなので、次の段落で論じる場合へ帰着する。

\medskip\noindent
$X$ が生成点 $\xi$ をもつ既約スキームであると仮定する。$C = X_{red}$ の重複度を
$m$ とする。$\deg(\mathcal{E}) = m \deg(\mathcal{E}|_C)$ を示さなければならない。
$\mathcal{I} \subset \mathcal{O}_X$ を閉部分スキーム $C$ を定めるイデアルとする。
$e \geq 0$ を $\mathcal{I}^{e + 1} = 0$ となる最小のものとする（Cohomology of Schemes,
Lemma \ref{coherent-lemma-power-ideal-kills-sheaf}）。$e$ に関する帰納法で論じる。
$e = 0$ ならば $X = C$ であり、結果は直ちに従う。そうでなければ
$\mathcal{F} = \mathcal{I}^e$ と置き、これを連接 $\mathcal{O}_C$ 加群とみなす
（Cohomology of Schemes, Lemma \ref{coherent-lemma-i-star-equivalence}）。
$X' \subset X$ を連接イデアル $\mathcal{I}^e$ で切り出される閉部分スキームとし、
$m'$ を $C$ の $X'$ における重複度とする。短完全列の $\xi$ における茎を取ると
$$
0 \to \mathcal{F} \to \mathcal{O}_X \to \mathcal{O}_{X'} \to 0
$$
Algebra, Lemmas \ref{algebra-lemma-length-additive},
\ref{algebra-lemma-dimension-is-length}, および
\ref{algebra-lemma-length-independent} を用いて
$$
m = \text{length}_{\mathcal{O}_{X, \xi}} \mathcal{O}_{X, \xi}
= \dim_{\kappa(\xi)} \mathcal{F}_\xi +
\text{length}_{\mathcal{O}_{X', \xi}} \mathcal{O}_{X', \xi}
= r + m'
$$
を得る。ここで $r$ は $\mathcal{F}$ の $C$ 上の連接層としての階数である。
$\mathcal{E}$ とテンソルすると短完全列
$$
0 \to \mathcal{E}|_C \otimes \mathcal{F} \to \mathcal{E} \to
\mathcal{E} \otimes \mathcal{O}_{X'} \to 0
$$
を得る。帰納法により $\deg(\mathcal{E}|_{X'}) = m' \deg(\mathcal{E}|_C)$ である。
Lemma \ref{varieties-lemma-degree-on-proper-curve} により
$\chi(\mathcal{E}|_C \otimes \mathcal{F}) =
r \deg(\mathcal{E}|_C) + n \chi(\mathcal{F})$ である。
以上を合わせると結果を得る。
\end{proof}

\begin{lemma}
\label{varieties-lemma-degree-tensor-product}
$k$ を体とし、$X$ を次元 $\leq 1$ の固有スキームであって $k$ 上のものとし、
$\mathcal{E}$, $\mathcal{V}$ を有限定階数の局所自由 $\mathcal{O}_X$ 加群とする。このとき
$$
\deg(\mathcal{E} \otimes \mathcal{V}) =
\text{rank}(\mathcal{E}) \deg(\mathcal{V}) +
\text{rank}(\mathcal{V}) \deg(\mathcal{E})
$$
である。
\end{lemma}

\begin{proof}
Lemma \ref{varieties-lemma-degree-in-terms-of-components} と初等的な算術により、
固有曲線の場合へ帰着する。この場合は Lemma
\ref{varieties-lemma-degree-on-proper-curve} から従う。
\end{proof}

\begin{lemma}
\label{varieties-lemma-degree-and-det}
$k$ を体とし、$X$ を次元 $\leq 1$ の固有スキームであって $k$ 上のものとし、
$\mathcal{E}$ を局所自由 $\mathcal{O}_X$ 加群で階数 $n$ のものとする。このとき
$$
\deg(\mathcal{E}) = \deg(\wedge^n(\mathcal{E})) = \deg(\det(\mathcal{E}))
$$
である。
\end{lemma}

\begin{proof}
Lemma \ref{varieties-lemma-degree-in-terms-of-components} と初等的な算術により、
固有曲線の場合へ帰着する。このとき、$f : X' \to X$ という修正であって、
$f^*\mathcal{E}$ が可逆加群を逐次商にもつフィルトレーションを備えるものが存在する。
Divisors, Lemma \ref{divisors-lemma-filter-after-modification} を参照せよ。
Lemma \ref{varieties-lemma-degree-birational-pullback} により $X'$ 上で考えてよい。
従ってフィルトレーション
$$
0 = \mathcal{E}_0 \subset \mathcal{E}_1 \subset \mathcal{E}_2 \subset
\ldots \subset \mathcal{E}_n = \mathcal{E}
$$
が局所自由 $\mathcal{O}_X$ 加群からなり、
$\mathcal{L}_i = \mathcal{E}_i/\mathcal{E}_{i - 1}$ が可逆であると仮定してよい。
Modules, Lemma \ref{modules-lemma-det-ses} と帰納法により
$\det(\mathcal{E}) = \mathcal{L}_1 \otimes \ldots \otimes \mathcal{L}_n$ を得る。
従って等式は Lemma \ref{varieties-lemma-degree-tensor-product} と加法性
（Lemma \ref{varieties-lemma-degree-additive}）から従う。
\end{proof}


\begin{lemma}
\label{varieties-lemma-degree-effective-Cartier-divisor}
$k$ を体とし、$X$ を次元 $\leq 1$ の固有スキームであって $k$ 上のものとする。
$D$ を $X$ 上の有効 Cartier 因子とする。このとき $D$ は $\Spec(k)$ 上有限で、次数は
$\deg(D) = \dim_k \Gamma(D, \mathcal{O}_D)$ である。局所自由層 $\mathcal{E}$ で
階数 $n$ のものに対して
$$
\deg(\mathcal{E}(D)) = n\deg(D) + \deg(\mathcal{E})
$$
が成り立つ。ここで $\mathcal{E}(D) = \mathcal{E} \otimes_{\mathcal{O}_X} \mathcal{O}_X(D)$ である。
\end{lemma}

\begin{proof}
$D$ は $X$ の中で稠密でないので（Divisors, Lemma
\ref{divisors-lemma-complement-effective-Cartier-divisor}）、$\dim(D) \leq 0$ である。
従って Lemma \ref{varieties-lemma-algebraic-scheme-dim-0} により $D$ は $k$ 上有限である。
$k$ は体なので、射 $D \to \Spec(k)$ は有限局所自由であり、従って次数をもつ
（Morphisms, Definition \ref{morphisms-definition-finite-locally-free}）。この次数が、補題で
述べた通り $\dim_k \Gamma(D, \mathcal{O}_D)$ に等しいことは明らかである。Divisors,
Definition
\ref{divisors-definition-invertible-sheaf-effective-Cartier-divisor} により、短完全列
$$
0 \to \mathcal{O}_X \to \mathcal{O}_X(D) \to i_*i^*\mathcal{O}_X(D) \to 0
$$
が存在する。ここで $i : D \to X$ は閉埋め込みである。$\mathcal{E}$ とテンソルすると、短完全列
$$
0 \to \mathcal{E} \to \mathcal{E}(D) \to i_*i^*\mathcal{E}(D) \to 0
$$
を得る。補題の等式は Euler 標数の加法性（Lemma
\ref{varieties-lemma-euler-characteristic-additive}）と Lemma
\ref{varieties-lemma-chi-tensor-finite} から従う。
\end{proof}

\begin{lemma}
\label{varieties-lemma-divisible}
$k$ を体とする。$X$ を $k$ 上の固有スキームであって、被約かつ連結なものとする。
$\kappa = H^0(X, \mathcal{O}_X)$ と置く。このとき $\kappa/k$ は有限次体拡大であり、
$w = [\kappa : k]$ は次の各数を割り切る：
\begin{enumerate}
\item $\deg(\mathcal{E})$（任意の局所自由 $\mathcal{O}_X$ 加群 $\mathcal{E}$ に対して）、
\item $[\kappa(x) : k]$（任意の閉点 $x \in X$ に対して）、および
\item $\deg(D)$（任意の零次元閉部分スキーム $D \subset X$ に対して）。
\end{enumerate}
\end{lemma}

\begin{proof}
$\kappa$ に関する主張については Lemma
\ref{varieties-lemma-proper-geometrically-reduced-global-sections} を参照せよ。任意の準連接
$\mathcal{O}_X$ 加群について、$k$ ベクトル空間 $H^i(X, \mathcal{F})$ は
$\kappa$ ベクトル空間である。各可除性はこの主張と定義から直ちに従う。
\end{proof}

\begin{lemma}
\label{varieties-lemma-degree-pullback-map-proper-curves}
$k$ を体とする。$f : X \to Y$ を $k$ 上の固有曲線の間の非定値射とする。
$\mathcal{E}$ を局所自由 $\mathcal{O}_Y$ 加群とする。このとき
$$
\deg(f^*\mathcal{E}) = \deg(X/Y) \deg(\mathcal{E})
$$
である。
\end{lemma}

\begin{proof}
$X$ の $Y$ 上の次数は Morphisms, Definition \ref{morphisms-definition-degree} で定義されている。
従って $f_*\mathcal{O}_X$ は連接 $\mathcal{O}_Y$ 加群であり、その階数は
$\deg(X/Y)$、すなわち
$\deg(X/Y) = \dim_{\kappa(\xi)} (f_*\mathcal{O}_X)_\xi$ である。ここで $\xi$ は
$Y$ の生成点である。従って
\begin{align*}
\chi(X, f^*\mathcal{E})
& =
\chi(Y, f_*f^*\mathcal{E}) \\
& =
\chi(Y, \mathcal{E} \otimes f_*\mathcal{O}_X) \\
& =
\deg(X/Y) \deg(\mathcal{E}) + n \chi(Y, f_*\mathcal{O}_X) \\
& =
\deg(X/Y) \deg(\mathcal{E}) + n \chi(X, \mathcal{O}_X)
\end{align*}
を得て、これは求める主張である。第一の等式は $f$ が有限であることから従う。Cohomology of
Schemes, Lemma \ref{coherent-lemma-relative-affine-cohomology} を参照せよ。第二の等式は
射影公式による。Cohomology, Lemma \ref{cohomology-lemma-projection-formula} を参照せよ。
第三の等式は Lemma \ref{varieties-lemma-degree-on-proper-curve} による。
\end{proof}

\noindent
次は、上で得た次数に関する結果の自明だが重要な帰結である。

\begin{lemma}
\label{varieties-lemma-check-invertible-sheaf-trivial}
$k$ を体とする。$X$ を $k$ 上の固有曲線とする。$\mathcal{L}$ を可逆
$\mathcal{O}_X$ 加群とする。
\begin{enumerate}
\item $\mathcal{L}$ が零でない切断をもつならば、$\deg(\mathcal{L}) \geq 0$ である。
\item $\mathcal{L}$ が、ある点で消える零でない切断 $s$ をもつならば、
$\deg(\mathcal{L}) > 0$ である。
\item $\mathcal{L}$ と $\mathcal{L}^{-1}$ が零でない切断をもつならば、
$\mathcal{L} \cong \mathcal{O}_X$ である。
\item $\deg(\mathcal{L}) \leq 0$ で、かつ $\mathcal{L}$ が零でない切断をもつならば、
$\mathcal{L} \cong \mathcal{O}_X$ である。
\item $\mathcal{N} \to \mathcal{L}$ が可逆 $\mathcal{O}_X$ 加群の間の零でない射ならば、
$\deg(\mathcal{L}) \geq \deg(\mathcal{N})$ であり、等号が成り立つならばその射は同型である。
\end{enumerate}
\end{lemma}

\begin{proof}
$s$ を $\mathcal{L}$ の零でない切断とする。$X$ は曲線なので、$s$ は正則切断である。
従って、有効 Cartier 因子 $D \subset X$ と同型 $\mathcal{L} \to \mathcal{O}_X(D)$ が存在し、
この同型は $s$ を標準切断 $1$（$\mathcal{O}_X(D)$ の標準切断）に写す。Divisors, Lemma
\ref{divisors-lemma-characterize-OD} を参照せよ。すると
$\deg(\mathcal{L}) = \deg(D)$ である。これは Lemma
\ref{varieties-lemma-degree-effective-Cartier-divisor} による。$\deg(D) \geq 0$ であり、$= 0$ となるのは
$D = \emptyset$ のとき、かつそのときに限るので、(1) と (2) が従う。(3) の場合には
$\deg(\mathcal{L}) = 0$ および $D = \emptyset$ を得る。(4) についても同様である。
(5) を見るには、可逆層
$$
\mathcal{L} \otimes_{\mathcal{O}_X} \mathcal{N}^{\otimes -1} =
\SheafHom_{\mathcal{O}_X}(\mathcal{N}, \mathcal{L})
$$
に (1) と (4) を適用する。この層の次数は $\deg(\mathcal{L}) - \deg(\mathcal{N})$ である。
これは Lemma \ref{varieties-lemma-degree-tensor-product} による。
\end{proof}

\begin{lemma}
\label{varieties-lemma-no-sections-dual-nef}
$k$ を体とする。$X$ を $k$ 上の固有スキームであって、被約、連結、かつ次元 $1$ の
等次元なものとする。$\mathcal{L}$ を可逆 $\mathcal{O}_X$ 加群とする。
$\deg(\mathcal{L}|_C) \leq 0$ が、すべての既約成分 $C$、すなわち $X$ の既約成分について
成り立つならば、$H^0(X, \mathcal{L}) = 0$ または $\mathcal{L} \cong \mathcal{O}_X$ である。
\end{lemma}

\begin{proof}
$s \in H^0(X, \mathcal{L})$ を零でないものとする。$X$ は被約なので、既約成分 $C$、すなわち
$X$ の既約成分であって $s|_C \not = 0$ となるものが存在する。しかし $s|_C$ が零でなければ、
Lemma \ref{varieties-lemma-check-invertible-sheaf-trivial} により $s$ は $C$ 上のどの点でも消えない。
すると $s$ は、$X$ の既約成分で $C$ と交わるもののすべての上でもどの点でも消えない。
$X$ は連結なので、$s$ はどこでも消えないと結論でき、補題が従う。
\end{proof}


\begin{lemma}
\label{varieties-lemma-ample-curve}
$k$ を体とする。$X$ を $k$ 上の固有曲線とする。$\mathcal{L}$ を可逆
$\mathcal{O}_X$ 加群とする。このとき $\mathcal{L}$ が豊富であるための必要十分条件は
$\deg(\mathcal{L}) > 0$ である。
\end{lemma}

\begin{proof}
$\mathcal{L}$ が豊富ならば、$n > 0$ と切断 $s \in H^0(X, \mathcal{L}^{\otimes n})$ で
$X_s$ がアフィンとなるものが存在する。$X$ はアフィンでない（そうでなければ Morphisms,
Lemma \ref{morphisms-lemma-finite-proper} により $X$ は有限となる）ので、$s$ はある点で消える。
従って Lemma \ref{varieties-lemma-check-invertible-sheaf-trivial} により
$\deg(\mathcal{L}^{\otimes n}) > 0$ である。Lemma \ref{varieties-lemma-degree-tensor-product} により
$\deg(\mathcal{L}) = 1/n\deg(\mathcal{L}^{\otimes n}) > 0$ と結論する。

\medskip\noindent
$\deg(\mathcal{L}) > 0$ と仮定する。このとき
$$
\dim_k H^0(X, \mathcal{L}^{\otimes n}) \geq \chi(X, \mathcal{L}^n)
= n\deg(\mathcal{L}) + \chi(X, \mathcal{O}_X)
$$
は $n$ に関して線型に増大する。従って閉点 $x_1, \ldots, x_t$（$X$ の閉点）の
任意の有限集合に対し、$n$ を
$\dim_k H^0(X, \mathcal{L}^{\otimes n}) > \sum \dim_k \kappa(x_i)$ となるように取れる。
（Hilbert の零点定理により拡大体 $\kappa(x_i)/k$ が有限であることを思い出そう。例えば
Morphisms, Lemma \ref{morphisms-lemma-closed-point-fibre-locally-finite-type} を参照せよ。）
従って零でない $s \in H^0(X, \mathcal{L}^{\otimes n})$ であって、
$x_1, \ldots, x_t$ で消えるものを取れる。特に、$x_1, \ldots, x_t$ を、
$X \setminus \{x_1, \ldots, x_t\}$ がアフィンとなるように選べば、
$X_s$ もアフィンである（例えば Properties, Lemma
\ref{properties-lemma-affine-cap-s-open} による。ただし、有限集合を
$\mathcal{L}|_{X \setminus \{x_1, \ldots, x_t\}}$ が自明となるように選べば直ちに分かる）。
結論として、$n > 0$ と零でない切断 $s \in H^0(X, \mathcal{L}^{\otimes n})$ で
$X_s$ がアフィンとなるものを取れる。

\medskip\noindent
任意の準連接イデアル層 $\mathcal{I}$ に対し、ある $m > 0$ が存在して
$H^1(X, \mathcal{I} \otimes \mathcal{L}^{\otimes m})$ が零となることを示す。
これにより Cohomology of Schemes, Lemma
\ref{coherent-lemma-vanshing-gives-ample} から証明が完了する。これを見るため、写像
$$
\mathcal{I} \xrightarrow{s}
\mathcal{I} \otimes \mathcal{L}^{\otimes n} \xrightarrow{s}
\mathcal{I} \otimes \mathcal{L}^{\otimes 2n} \xrightarrow{s} \ldots
$$
を考える。$\mathcal{I}$ は捩れなしなので、これらの写像は単射であり、$X_s$ 上では
同型である。従って余核の $H^1$ は消える（例えば Cohomology of Schemes, Lemma
\ref{coherent-lemma-coherent-support-dimension-0} による）。ゆえにベクトル空間の写像
$$
H^1(X, \mathcal{I}) \to
H^1(X, \mathcal{I} \otimes \mathcal{L}^{\otimes n}) \to
H^1(X, \mathcal{I} \otimes \mathcal{L}^{\otimes 2n}) \to \ldots
$$
は全射である。一方、$H^1(X, \mathcal{I})$ の次元は有限であり、各元は Cohomology of
Schemes, Lemma \ref{coherent-lemma-section-affine-open-kills-classes} により、いずれ零へ写る。
従って、ある $e > 0$ に対して
$H^1(X, \mathcal{I} \otimes \mathcal{L}^{\otimes en})$ は零である。これで証明が完了する。
\end{proof}

\begin{lemma}
\label{varieties-lemma-ampleness-in-terms-of-degrees-components}
$k$ を体とする。$X$ を次元 $\leq 1$ の固有スキームであって $k$ 上のものとする。
$\mathcal{L}$ を可逆 $\mathcal{O}_X$ 加群とする。$C_i \subset X$, $i = 1, \ldots, t$ を
次元 $1$ の既約成分とする。次は同値である：
\begin{enumerate}
\item $\mathcal{L}$ は豊富である。
\item $\deg(\mathcal{L}|_{C_i}) > 0$ が $i = 1, \ldots, t$ に対して成り立つ。
\end{enumerate}
\end{lemma}

\begin{proof}
$x_1, \ldots, x_r \in X$ を孤立閉点とする。$x_i = \Spec(\kappa(x_i))$ をスキームとみなす。
スキームの射
$$
f : C_1 \amalg \ldots \amalg C_t \amalg x_1 \amalg \ldots \amalg x_r
\longrightarrow X
$$
を考える。これは $k$ 上固有なスキームの有限全射である（詳細は省略する）。従って
$\mathcal{L}$ が豊富であるための必要十分条件は $f^*\mathcal{L}$ が豊富であることである
（Cohomology of Schemes, Lemma
\ref{coherent-lemma-surjective-finite-morphism-ample}）。従って Lemma
\ref{varieties-lemma-ample-curve} により結論を得る。
\end{proof}

\begin{lemma}
\label{varieties-lemma-general-degree-g-line-bundle}
$k$ を代数閉体とする。$X$ を $k$ 上の固有曲線とする。このとき、次が存在する：
\begin{enumerate}
\item 可逆 $\mathcal{O}_X$ 加群 $\mathcal{L}$ で、
$\dim_k H^0(X, \mathcal{L}) = 1$ かつ $H^1(X, \mathcal{L}) = 0$ を満たすもの、
\item 可逆 $\mathcal{O}_X$ 加群 $\mathcal{N}$ で、
$\dim_k H^0(X, \mathcal{N}) = 0$ かつ $H^1(X, \mathcal{N}) = 0$ を満たすもの。
\end{enumerate}
\end{lemma}

\begin{proof}
閉埋め込み $i : X \to \mathbf{P}^n_k$ を選ぶ（Lemma
\ref{varieties-lemma-dim-1-proper-projective}）。$\mathcal{L} = i^*\mathcal{O}_{\mathbf{P}^n}(d)$ と
$d \gg 0$ に対して置くことで、可逆層 $\mathcal{L}$ で
$H^0(X, \mathcal{L}) \not = 0$ かつ $H^1(X, \mathcal{L}) = 0$ を満たすものが
存在することが分かる（消滅については Cohomology of Schemes, Lemma
\ref{coherent-lemma-vanshing-gives-ample} を、非消滅についてはそこで挙げられた文献を参照せよ）。
$t = \dim_k H^0(X, \mathcal{L})$ に関する下降帰納法によって (1) の証明を完了する。
基底の場合 $t = 1$ は自明である。$t > 1$ と仮定する。

\medskip\noindent
$U \subset X$ を Lemma \ref{varieties-lemma-dense-smooth-open-variety-over-perfect-field} で調べた
非特異点からなる空でない開部分集合とする。$s \in H^0(X, \mathcal{L})$ を零でないものとする。
閉点 $x \in U$ で $s$ が $x$ で消えないものが存在する。$\mathcal{I}$ を
Lemma \ref{varieties-lemma-regular-point-on-curve} における $i : x \to X$ のイデアル層とする。
短完全列
$$
0 \to \mathcal{I} \otimes_{\mathcal{O}_X} \mathcal{L} \to
\mathcal{L} \to i_*i^*\mathcal{L} \to 0
$$
を考える。$H^0(X, i_*i^*\mathcal{L}) = H^0(x, i^*\mathcal{L})$ の次元は $1$ である。
実際、$x$ は $k$ 有理点である（$k$ は代数閉）。
$s$ は $x$ で消えないので
$$
H^0(X, \mathcal{L}) \longrightarrow H^0(X, i_*i^*\mathcal{L})
$$
は全射である。従って
$\dim_k H^0(X, \mathcal{I} \otimes_{\mathcal{O}_X} \mathcal{L}) = t - 1$ である。
最後に、コホモロジーの長完全列から
$H^1(X, \mathcal{I} \otimes_{\mathcal{O}_X} \mathcal{L}) = 0$ も分かり、
これで帰納段階の証明が完了する。

\medskip\noindent
(2) のような可逆層を得るには、(1) のような可逆層 $\mathcal{L}$ を取り、
前段落の議論をもう一度行う。
\end{proof}

\begin{lemma}
\label{varieties-lemma-vanishing-degree-2g-and-1-line-bundle}
$k$ を代数閉体とする。$X$ を $k$ 上の固有曲線とする。
$g = \dim_k H^1(X, \mathcal{O}_X)$ と置く。任意の可逆 $\mathcal{O}_X$ 加群
$\mathcal{L}$ で $\deg(\mathcal{L}) \geq 2g - 1$ を満たすものに対して
$H^1(X, \mathcal{L}) = 0$ である。
\end{lemma}

\begin{proof}
Lemma \ref{varieties-lemma-general-degree-g-line-bundle} の (2) で得た可逆加群 $\mathcal{N}$ を取る。
$\mathcal{N}$ の次数は
$\chi(X, \mathcal{N}) - \chi(X, \mathcal{O}_X) = 0 - (1 - g) = g - 1$ である。
従って $\mathcal{L} \otimes \mathcal{N}^{\otimes - 1}$ の次数は
$\deg(\mathcal{L}) - (g - 1) \geq g$ である。従って
$\chi(X, \mathcal{L} \otimes \mathcal{N}^{\otimes -1}) \geq g + 1 - g = 1$ である。
ゆえに、零でない大域切断 $s$ で、その零点スキームが有効 Cartier 因子 $D$ となり、
その次数が $\deg(\mathcal{L}) - (g - 1)$ であるものが存在する。これは短完全列
$$
0 \to \mathcal{N} \xrightarrow{s} \mathcal{L} \to i_*(\mathcal{L}|_D) \to 0
$$
を与える。ここで $i : D \to X$ は包含射である。従って
$H^0(X, \mathcal{L})$ は $H^0(D, \mathcal{L}|_D)$ へ同型に写り、後者の次元は
$\deg(\mathcal{L}) - (g - 1)$ である。結果は次数の定義から従う。
\end{proof}


\section{数値的交点}
\label{varieties-section-num}

\noindent
この節では、固有スキーム上の連接層の Euler 標数を用いて、可逆加群の数値的交点数を得る。
主な道具は次の補題である。

\begin{lemma}
\label{varieties-lemma-numerical-polynomial-from-euler}
$k$ を体とする。$X$ を $k$ 上の固有スキームとする。$\mathcal{F}$ を連接
$\mathcal{O}_X$ 加群とする。$\mathcal{L}_1, \ldots, \mathcal{L}_r$ を可逆
$\mathcal{O}_X$ 加群とする。写像
$$
(n_1, \ldots, n_r) \longmapsto
\chi(X, \mathcal{F} \otimes
\mathcal{L}_1^{\otimes n_1} \otimes \ldots \otimes
\mathcal{L}_r^{\otimes n_r})
$$
は $n_1, \ldots, n_r$ の数値多項式であり、その全次数は $\mathcal{F}$ の台の次元以下である。
\end{lemma}

\begin{proof}
$\dim(\text{Supp}(\mathcal{F}))$ に関する帰納法で証明する。この数が零ならば、Lemma
\ref{varieties-lemma-chi-tensor-finite} により関数は定数で、その値は
$\dim_k \Gamma(X, \mathcal{F})$ である。$\dim(\text{Supp}(\mathcal{F})) > 0$ と仮定する。

\medskip\noindent
$\mathcal{F}$ が埋め込まれた随伴点をもつならば、Divisors, Lemma
\ref{divisors-lemma-remove-embedded-points} で構成された短完全列
$0 \to \mathcal{K} \to \mathcal{F} \to \mathcal{F}' \to 0$ を考えられる。
$\mathcal{K}$ の台の次元は真に小さいので、帰納法の仮定により、結果は $\mathcal{K}$ に
対して真であり、全次数も真に小さい。Euler 標数の加法性（Lemma
\ref{varieties-lemma-euler-characteristic-additive}）により、$\mathcal{F}'$ に対して結果を示せば十分である。
従って $\mathcal{F}$ は埋め込まれた随伴点をもたないと仮定してよい。

\medskip\noindent
$i : Z \to X$ が閉埋め込みで、$\mathcal{F} = i_*\mathcal{G}$ ならば、
$X$, $\mathcal{F}$, $\mathcal{L}_1, \ldots, \mathcal{L}_r$ に対する結果は、
$Z$, $\mathcal{G}$, $i^*\mathcal{L}_1, \ldots, i^*\mathcal{L}_r$ に対する結果と同値である
（コホモロジーが一致するため。Cohomology of Schemes, Lemma
\ref{coherent-lemma-relative-affine-cohomology} を参照せよ）。Divisors, Lemma
\ref{divisors-lemma-no-embedded-points-endos} を適用して、$X$ は埋め込まれた成分をもたず、
$X = \text{Supp}(\mathcal{F})$ であると仮定してよい。

\medskip\noindent
正則な有理型切断 $s$ を $\mathcal{L}_1$ から選ぶ。Divisors, Lemma
\ref{divisors-lemma-regular-meromorphic-section-exists} を参照せよ。
$\mathcal{I} \subset \mathcal{O}_X$ を $s$ の分母のイデアルとし、Divisors, Lemma
\ref{divisors-lemma-make-maps-regular-section} の写像
$$
\mathcal{I}\mathcal{F} \to \mathcal{F},\quad
\mathcal{I}\mathcal{F} \to \mathcal{F} \otimes \mathcal{L}_1
$$
を考える。これらは単射であり、余核 $\mathcal{Q}$, $\mathcal{Q}'$ は
$X = \text{Supp}(\mathcal{F})$ の稠密でない閉部分スキーム上に台をもつ。可逆加群
$\mathcal{L}_1^{\otimes n_1} \otimes \ldots \otimes \mathcal{L}_r^{\otimes n_r}$ と
テンソルする操作は完全であるから、再び加法性を用いると
\begin{align*}
&\chi(X, \mathcal{F} \otimes \mathcal{L}_1^{\otimes n_1} \otimes \ldots \otimes
\mathcal{L}_r^{\otimes n_r}) -
\chi(X, \mathcal{F} \otimes \mathcal{L}_1^{\otimes n_1 + 1}
\otimes \ldots \otimes \mathcal{L}_r^{\otimes n_r}) \\
& =
\chi(\mathcal{Q} \otimes \mathcal{L}_1^{\otimes n_1} \otimes \ldots \otimes
\mathcal{L}_r^{\otimes n_r}) -
\chi(\mathcal{Q}' \otimes \mathcal{L}_1^{\otimes n_1} \otimes \ldots \otimes
\mathcal{L}_r^{\otimes n_r})
\end{align*}
を得る。従って補題の関数 $P(n_1, \ldots, n_r)$ は、
$$
P(n_1 + 1, n_2, \ldots, n_r) - P(n_1, \ldots, n_r)
$$
が全次数 $<$ $\mathcal{F}$ の台の次元である数値多項式となるという性質をもつ。
もちろん対称性により、任意の $i \in \{1, \ldots, r\}$ に対して
$$
P(n_1, \ldots, n_{i - 1}, n_i + 1, n_{i + 1}, \ldots, n_r)
- P(n_1, \ldots, n_r)
$$
についても同じことが成り立つ。簡単な算術的議論により、$P$ は全次数が
$\dim(\text{Supp}(\mathcal{F}))$ 以下の数値多項式である。
\end{proof}

\noindent
次の補題は、最高次係数が連接加群の台の生成点における長さだけに依存することを大まかに示す。

\begin{lemma}
\label{varieties-lemma-numerical-polynomial-leading-term}
$k$ を体とする。$X$ を $k$ 上の固有スキームとする。$\mathcal{F}$ を連接
$\mathcal{O}_X$ 加群とする。$\mathcal{L}_1, \ldots, \mathcal{L}_r$ を可逆
$\mathcal{O}_X$ 加群とする。$d = \dim(\text{Supp}(\mathcal{F}))$ と置く。
$Z_i \subset X$ を、$\text{Supp}(\mathcal{F})$ の次元 $d$ の既約成分とする。
$\xi_i \in Z_i$ を生成点とし、
$m_i = \text{length}_{\mathcal{O}_{X, \xi_i}}(\mathcal{F}_{\xi_i})$ と置く。このとき
$$
\chi(X, \mathcal{F} \otimes \mathcal{L}_1^{\otimes n_1} \otimes \ldots \otimes
\mathcal{L}_r^{\otimes n_r}) -
\sum\nolimits_i
m_i\ \chi(Z_i, \mathcal{L}_1^{\otimes n_1} \otimes \ldots \otimes
\mathcal{L}_r^{\otimes n_r}|_{Z_i})
$$
は $n_1, \ldots, n_r$ の数値多項式であり、その全次数は $< d$ である。
\end{lemma}

\begin{proof}
組 $(\xi , Z)$ で、$Z \subset X$ が次元 $d$ の整閉部分スキームであり、$\xi$ がその生成点である
ものを考える。このとき有限 $\mathcal{O}_{X, \xi}$ 加群 $\mathcal{F}_\xi$ の台は
$\{\xi\}$ に含まれるので、長さ
$m_Z = \text{length}_{\mathcal{O}_{X, \xi}}(\mathcal{F}_\xi)$ は有限である（Algebra, Lemma
\ref{algebra-lemma-support-point}）。また、$Z = Z_i$ となる $i$ がなければこの長さは零である。
従って補題の式は
$$
E(\mathcal{F}) =
\chi(X, \mathcal{F} \otimes \mathcal{L}_1^{\otimes n_1} \otimes \ldots \otimes
\mathcal{L}_r^{\otimes n_r}) -
\sum\nolimits
m_Z\ \chi(Z, \mathcal{L}_1^{\otimes n_1} \otimes \ldots \otimes
\mathcal{L}_r^{\otimes n_r}|_Z)
$$
と書ける。ここで和は整閉部分スキーム $Z \subset X$ で次元が $d$ のもの全体にわたる。対応
$\mathcal{F} \mapsto E(\mathcal{F})$ は、短完全列
$0 \to \mathcal{F} \to \mathcal{F}' \to \mathcal{F}'' \to 0$ をなす連接
$\mathcal{O}_X$ 加群で、その台の次元が
$\leq d$ であるものに関して加法的である。
これは Euler 標数の加法性（Lemma \ref{varieties-lemma-euler-characteristic-additive}）と長さの加法性
（Algebra, Lemma \ref{algebra-lemma-length-additive}）から従う。Cohomology of Schemes,
Lemma \ref{coherent-lemma-coherent-filter} を適用して、連接部分層によるフィルトレーション
$$
0 = \mathcal{F}_0 \subset \mathcal{F}_1 \subset
\ldots \subset \mathcal{F}_m = \mathcal{F}
$$
を取る。これは、各 $j = 1, \ldots, m$ に対して整閉部分スキーム
$V_j \subset X$ と零でないイデアル層 $\mathcal{I}_j \subset \mathcal{O}_{V_j}$ が存在し、
$$
\mathcal{F}_j/\mathcal{F}_{j - 1}
\cong (V_j \to X)_* \mathcal{I}_j
$$
を満たすものである。従って $V_j \subset \text{Supp}(\mathcal{F})$ であり、ゆえに
$\dim(V_j) \leq d$ である。上で述べた加法性により、各部分商
$\mathcal{F}_j/\mathcal{F}_{j - 1}$ に対して結果を示せば十分である。従って
$\mathcal{F} = (V \to X)_*\mathcal{I}$ で、$V \subset X$ が次元 $\leq d$ の整閉部分スキーム、
$\mathcal{I} \subset \mathcal{O}_V$ が零でない連接イデアル層である場合に結果を示せば十分である。
$\dim(V) < d$ の場合、より一般には $\mathcal{F}$ の台の次元が $< d$ である場合、
$E(\mathcal{F})$ の第一項は Lemma \ref{varieties-lemma-numerical-polynomial-from-euler} により
全次数 $< d$ であり、第二項は零である。$\dim(V) = d$ ならば、短完全列
$$
0 \to (V \to X)_*\mathcal{I} \to (V \to X)_*\mathcal{O}_V
\to (V \to X)_*(\mathcal{O}_V/\mathcal{I}) \to 0
$$
を使える。中項の層に対して結果は成り立つ。実際、和に現れる $Z$ は
$Z = V$ だけで、$m_Z = 1$ であり、
$$
H^i(X, ((V \to X)_*\mathcal{O}_V) \otimes 
 \mathcal{L}_1^{\otimes n_1} \otimes \ldots \otimes
\mathcal{L}_r^{\otimes n_r}) =
H^i(V,  \mathcal{L}_1^{\otimes n_1} \otimes \ldots \otimes
\mathcal{L}_r^{\otimes n_r}|_V)
$$
が射影公式（Cohomology, Section \ref{cohomology-section-projection-formula}）および Cohomology of
Schemes, Lemma \ref{coherent-lemma-relative-affine-cohomology} により成り立つからである。
従ってこの場合は実際に $E(\mathcal{F}) = 0$ である。右項の層に対しては、その台の次元が
$< d$ なので結果が成り立つ。ゆえに左項の層に対しても結果が成り立ち、補題が証明された。
\end{proof}

\begin{definition}
\label{varieties-definition-intersection-number}
$k$ を体とする。$X$ を $k$ 上の固有スキームとする。
$i : Z \to X$ を次元 $d$ の閉部分スキームとする。
$\mathcal{L}_1, \ldots, \mathcal{L}_d$ を可逆 $\mathcal{O}_X$ 加群とする。
{\it 交点数} $(\mathcal{L}_1 \cdots \mathcal{L}_d \cdot Z)$ を、数値多項式
$$
\chi(X, i_*\mathcal{O}_Z \otimes \mathcal{L}_1^{\otimes n_1} \otimes
\ldots \otimes \mathcal{L}_d^{\otimes n_d}) =
\chi(Z, \mathcal{L}_1^{\otimes n_1} \otimes
\ldots \otimes \mathcal{L}_d^{\otimes n_d}|_Z)
$$
における $n_1 \ldots n_d$ の係数として定義する。特別な場合
$\mathcal{L}_1 = \ldots = \mathcal{L}_d = \mathcal{L}$ には
$(\mathcal{L}^d \cdot Z)$ と書く。
\end{definition}

\noindent
定義中の表示等式は射影公式
（Cohomology, Section \ref{cohomology-section-projection-formula}）および
Cohomology of Schemes, Lemma
\ref{coherent-lemma-relative-affine-cohomology} から従う。
これらの交点数についていくつかの補題を証明する。

\begin{lemma}
\label{varieties-lemma-intersection-number-integer}
Definition \ref{varieties-definition-intersection-number} の状況において、交点数
$(\mathcal{L}_1 \cdots \mathcal{L}_d \cdot Z)$ は整数である。
\end{lemma}

\begin{proof}
$e$ 次の任意の $n_1, \ldots, n_d$ の数値多項式は、$\mathbf{Z}$ 線形結合として一意に書ける。
ここで用いる関数は ${n_1 \choose k_1}{n_2 \choose k_2} \ldots {n_d \choose k_d}$ で、
$k_1 + \ldots + k_d \leq e$ を満たすものとする。
$e = d$ としてこれを適用する。詳細は演習として残す。
\end{proof}

\begin{lemma}
\label{varieties-lemma-intersection-number-additive}
Definition \ref{varieties-definition-intersection-number} の状況において、交点数
$(\mathcal{L}_1 \cdots \mathcal{L}_d \cdot Z)$ は加法的である。すなわち、
$\mathcal{L}_i = \mathcal{L}_i' \otimes \mathcal{L}_i''$ ならば
$$
(\mathcal{L}_1 \cdots \mathcal{L}_i \cdots \mathcal{L}_d \cdot Z) =
(\mathcal{L}_1 \cdots \mathcal{L}_i' \cdots \mathcal{L}_d \cdot Z) +
(\mathcal{L}_1 \cdots \mathcal{L}_i'' \cdots \mathcal{L}_d \cdot Z)
$$
が成り立つ。
\end{lemma}

\begin{proof}
これは Lemma \ref{varieties-lemma-numerical-polynomial-from-euler} により、関数
$$
(n_1, \ldots, n_{i - 1}, n_i', n_i'', n_{i + 1}, \ldots, n_d)
\mapsto
\chi(Z, \mathcal{L}_1^{\otimes n_1} \otimes
\ldots \otimes (\mathcal{L}_i')^{\otimes n_i'} \otimes
(\mathcal{L}_i'')^{\otimes n_i''} \otimes \ldots \otimes
\mathcal{L}_d^{\otimes n_d}|_Z)
$$
が全次数高々 $d$ の $d + 1$ 変数の数値多項式だからである。
\end{proof}

\begin{lemma}
\label{varieties-lemma-intersection-number-in-terms-of-components}
Definition \ref{varieties-definition-intersection-number} の状況において、
$Z_i \subset Z$ を次元 $d$ の既約成分とする。ここで
$m_i = \text{length}_{\mathcal{O}_{X, \xi_i}}(\mathcal{O}_{Z, \xi_i})$
とし、$\xi_i \in Z_i$ は生成点とする。このとき
$$
(\mathcal{L}_1 \cdots \mathcal{L}_d \cdot Z) =
\sum m_i(\mathcal{L}_1 \cdots \mathcal{L}_d \cdot Z_i)
$$
が成り立つ。
\end{lemma}

\begin{proof}
Lemma \ref{varieties-lemma-numerical-polynomial-leading-term} と定義から直ちに従う。
\end{proof}

\begin{lemma}
\label{varieties-lemma-intersection-number-and-pullback}
$k$ を体とする。$f : Y \to X$ を $k$ 上の固有スキームの射とする。
$Z \subset Y$ を次元 $d$ の整閉部分スキームとし、
$\mathcal{L}_1, \ldots, \mathcal{L}_d$ を可逆 $\mathcal{O}_X$ 加群とする。このとき
$$
(f^*\mathcal{L}_1 \cdots f^*\mathcal{L}_d \cdot Z) =
\deg(f|_Z : Z \to f(Z)) (\mathcal{L}_1 \cdots \mathcal{L}_d \cdot f(Z))
$$
が成り立つ。ここで $\deg(Z \to f(Z))$ は Morphisms, Definition
\ref{morphisms-definition-degree} によるものとするが、$0$ と置くのは
$\dim(f(Z)) < d$ の場合である。
\end{lemma}

\begin{proof}
左辺は、関数
$$
\chi(Y, \mathcal{O}_Z \otimes f^*\mathcal{L}_1^{\otimes n_1} \otimes
\ldots \otimes f^*\mathcal{L}_d^{\otimes n_d}) =
\sum (-1)^i
\chi(X, R^if_*\mathcal{O}_Z \otimes
\mathcal{L}_1^{\otimes n_1} \otimes \ldots \otimes
\mathcal{L}_d^{\otimes n_d})
$$
における $n_1 \ldots n_d$ の係数を用いて計算される。この等式は Lemma
\ref{varieties-lemma-euler-characteristic-morphism} と射影公式
（Cohomology, Lemma \ref{cohomology-lemma-projection-formula}）から従う。
$f(Z)$ の次元が $< d$ ならば、右辺は Lemma
\ref{varieties-lemma-numerical-polynomial-from-euler} により全次数 $< d$ の多項式であり、結果は真である。
$\dim(f(Z)) = d$ と仮定する。$\xi \in f(Z)$ を生成点とする。次元論
（Lemmas \ref{varieties-lemma-dimension-locally-algebraic} および
\ref{varieties-lemma-dimension-fibres-locally-algebraic} を参照）により、
$Z$ の生成点は、$Z$ の点で $\xi$ に写るものとして唯一である。
すると $f : Z \to f(Z)$ は $f(Z)$ の空でない開集合上で有限である。Morphisms, Lemma
\ref{morphisms-lemma-generically-finite} を参照せよ。
従って $\deg(f : Z \to f(Z))$ は定義され、実際これは
$f_*\mathcal{O}_Z$ の $\xi$ における茎の $\mathcal{O}_{X, \xi}$ 上の長さに等しい。さらに、
$R^if_*\mathcal{O}_X$ の $\xi$ における茎は $i > 0$ に対して零である。実際、直前に
$f|_Z$ が $\xi$ の近傍で有限であることを見たからである
（従って Cohomology of Schemes, Lemma
\ref{coherent-lemma-finite-pushforward-coherent} が消滅を与える）。
従って項 $\chi(X, R^if_*\mathcal{O}_Z \otimes
\mathcal{L}_1^{\otimes n_1} \otimes \ldots \otimes
\mathcal{L}_d^{\otimes n_d})$ は $i > 0$ に対して全次数 $< d$ であり、
$$
\chi(X, f_*\mathcal{O}_Z \otimes
\mathcal{L}_1^{\otimes n_1} \otimes \ldots \otimes
\mathcal{L}_d^{\otimes n_d})
=
\deg(f : Z \to f(Z)) \chi(f(Z),
\mathcal{L}_1^{\otimes n_1} \otimes \ldots \otimes
\mathcal{L}_d^{\otimes n_d}|_{f(Z)})
$$
は Lemma \ref{varieties-lemma-numerical-polynomial-leading-term} により、全次数 $< d$ の多項式を法として成り立つ。
求める結果が従う。
\end{proof}

\begin{lemma}
\label{varieties-lemma-numerical-intersection-effective-Cartier-divisor}
$k$ を体とする。$X$ を $k$ 上固有とする。$Z \subset X$ を次元 $d$ の閉部分スキームとする。
$\mathcal{L}_1, \ldots, \mathcal{L}_d$ を可逆 $\mathcal{O}_X$ 加群とする。
有効 Cartier 因子 $D \subset Z$ が存在し、
$\mathcal{L}_1|_Z \cong \mathcal{O}_Z(D)$ を満たすと仮定する。このとき
$$
(\mathcal{L}_1 \cdots \mathcal{L}_d \cdot Z) =
(\mathcal{L}_2 \cdots \mathcal{L}_d \cdot D)
$$
が成り立つ。
\end{lemma}

\begin{proof}
$X$ を $Z$ で、$\mathcal{L}_i$ を $\mathcal{L}_i|_Z$ で置き換えてよい。
従って $X = Z$ かつ $\mathcal{L}_1 = \mathcal{O}_X(D)$ と仮定してよい。
$\mathcal{L}_1^{-1}$ は $D$ のイデアル層であり、短完全列
$$
0 \to \mathcal{L}_1^{\otimes -1} \to \mathcal{O}_X \to \mathcal{O}_D \to 0
$$
を考えられる。
$P(n_1, \ldots, n_d) =
\chi(X, \mathcal{L}_1^{\otimes n_1} \otimes \ldots \otimes
\mathcal{L}_d^{\otimes n_d})$
および
$Q(n_1, \ldots, n_d) =
\chi(D, \mathcal{L}_1^{\otimes n_1} \otimes \ldots \otimes
\mathcal{L}_d^{\otimes n_d}|_D)$
と置く。加法性から
$$
P(n_1, \ldots, n_d) - P(n_1 - 1, n_2, \ldots, n_d) =
Q(n_1, \ldots, n_d)
$$
を得る。$P$ の全次数は高々 $d$ なので、$n_1 \ldots n_d$ の $P$ における係数は
$n_2 \ldots n_d$ の $Q$ における係数に等しい。
\end{proof}

\begin{lemma}
\label{varieties-lemma-ample-positive}
$k$ を体とする。$X$ を $k$ 上固有とする。$Z \subset X$ を次元 $d$ の閉部分スキームとする。
$\mathcal{L}_1, \ldots, \mathcal{L}_d$ が豊富ならば、
$(\mathcal{L}_1 \cdots \mathcal{L}_d \cdot Z)$ は正である。
\end{lemma}

\begin{proof}
$d$ に関する帰納法で証明する。$d = 0$ の場合は Lemma
\ref{varieties-lemma-chi-tensor-finite} から従う。$d > 0$ と仮定する。
Lemma \ref{varieties-lemma-intersection-number-in-terms-of-components} により、
$Z$ は整閉部分スキームであると仮定してよい。実際、$X$ を $Z$ で、$\mathcal{L}_i$ を
$\mathcal{L}_i|_Z$ で置き換え、$Z = X$ が次元 $d$ の固有多様体である場合に帰着してよい。
Lemma \ref{varieties-lemma-intersection-number-additive} により、$\mathcal{L}_1$ を正のテンソル冪で
置き換えてよい。従って、零でない切断 $s \in \Gamma(X, \mathcal{L}_1)$ で
$X_s$ がアフィンとなるものが存在すると仮定してよい
（ここでは豊富な可逆層の定義を用いる。Properties, Definition
\ref{properties-definition-ample} を参照せよ）。
$X$ はアフィンでないことに注意する。なぜなら、固有かつアフィンならば有限となる
（Morphisms, Lemma \ref{morphisms-lemma-finite-proper}）が、これは $d > 0$ に反するからである。
従って $s$ は空でない消滅スキーム $Z(s) \subset X$ をもつ。
$X$ は多様体なので、$s$ は $\mathcal{L}_1$ の正則切断である。ゆえに $Z(s)$ は有効 Cartier
因子であり、従って $Z(s)$ は余次元 $1$ をもつ部分スキームとして $X$ に含まれ、ゆえに
$Z(s)$ の次元は $d - 1$ である（ここでは Divisors, Sections
\ref{divisors-section-effective-Cartier-divisors},
\ref{divisors-section-effective-Cartier-invertible}, および
\ref{divisors-section-Noetherian-effective-Cartier} の内容と、Lemma
\ref{varieties-lemma-dimension-locally-algebraic} におけるような次元論を用いる）。
Lemma \ref{varieties-lemma-numerical-intersection-effective-Cartier-divisor} により
$$
(\mathcal{L}_1 \cdots \mathcal{L}_d \cdot X) =
(\mathcal{L}_2 \cdots \mathcal{L}_d \cdot Z(s))
$$
が成り立つ。帰納法により右辺は正であり、証明が完了する。
\end{proof}

\begin{definition}
\label{varieties-definition-degree}
$k$ を体とする。$X$ を $k$ 上の固有スキームとする。
$\mathcal{L}$ を豊富な可逆 $\mathcal{O}_X$ 加群とする。
任意の閉部分スキーム $Z$ の {\it $\mathcal{L}$ に関する次数} は
$\deg_\mathcal{L}(Z)$ と書き、交点数 $(\mathcal{L}^d \cdot Z)$ と定める。
ここで $d = \dim(Z)$ である。
\end{definition}

\noindent
Lemma \ref{varieties-lemma-ample-positive} により、部分スキームの次数は常に正の整数である。
次の式から分かるように、$\deg_\mathcal{L}(Z) = d$ であることと
$$
\chi(Z, \mathcal{L}^{\otimes n}|_Z) = \frac{d}{\dim(Z)!} n^{\dim(Z)} + l.o.t
$$
であることは同値である。
$$
(n_1 + \ldots + n_{\dim(Z)})^{\dim(Z)} =
\dim(Z)!\ n_1 \ldots n_{\dim(Z)} + \text{other terms}
$$

\begin{lemma}
\label{varieties-lemma-degree-finite-morphism-in-terms-degrees}
$k$ を体とする。$f : Y \to X$ を $k$ 上の固有多様体の有限支配射とする。
$\mathcal{L}$ を豊富な可逆 $\mathcal{O}_X$ 加群とする。このとき
$$
\deg_{f^*\mathcal{L}}(Y) = \deg(f) \deg_\mathcal{L}(X)
$$
が成り立つ。ここで $\deg(f)$ は Morphisms, Definition
\ref{morphisms-definition-degree} による。
\end{lemma}

\begin{proof}
主張が意味をもつのは、Morphisms, Lemma
\ref{morphisms-lemma-pullback-ample-tensor-relatively-ample} により
$f^*\mathcal{L}$ が豊富だからである。これを踏まえれば、結果は Lemma
\ref{varieties-lemma-intersection-number-and-pullback} の特別な場合である。
\end{proof}

\noindent
最後に、曲線との交点数を Section \ref{varieties-section-divisors-curves} で導入した
曲線上の可逆加群の次数の概念と関係づける。

\begin{lemma}
\label{varieties-lemma-intersection-numbers-and-degrees-on-curves}
$k$ を体とする。$X$ を $k$ 上の固有スキームとする。
$Z \subset X$ を次元 $\leq 1$ の閉部分スキームとする。
$\mathcal{L}$ を可逆 $\mathcal{O}_X$ 加群とする。このとき
$$
(\mathcal{L} \cdot Z) = \deg(\mathcal{L}|_Z)
$$
が成り立つ。ここで $\deg(\mathcal{L}|_Z)$ は Definition
\ref{varieties-definition-degree-invertible-sheaf} による。
$\mathcal{L}$ が豊富ならば、$\deg_\mathcal{L}(Z) = \deg(\mathcal{L}|_Z)$ である。
\end{lemma}

\begin{proof}
これは、関数 $n \mapsto \chi(Z, \mathcal{L}|_Z^{\otimes n})$ の次数が $1$ であり、
従って最高次係数が連続する二つの値の差であることから従う。
\end{proof}

\begin{proposition}[漸近的 Riemann--Roch]
\label{varieties-proposition-asymptotic-riemann-roch}
$k$ を体とする。$X$ を $k$ 上の次元 $d$ の固有スキームとする。
$\mathcal{L}$ を豊富な可逆 $\mathcal{O}_X$ 加群とする。このとき
$$
\dim_k \Gamma(X, \mathcal{L}^{\otimes n}) \sim c n^d + l.o.t.
$$
が成り立つ。ここで $c = \deg_\mathcal{L}(X)/d!$ は正の定数である。
\end{proposition}

\begin{proof}
これは定義、Lemma \ref{varieties-lemma-ample-positive}、および Cohomology of Schemes, Lemma
\ref{coherent-lemma-vanshing-gives-ample} における高次コホモロジーの消滅から従う。
\end{proof}

\section{埋め込み次元}
\label{varieties-section-embedding-dimension}

\noindent
埋め込み次元を定義する方法はいくつかあるが、代数閉体上の代数的スキームの閉点については、
すべての定義が次のものと同値である。

\begin{definition}
\label{varieties-definition-embed-dim}
$k$ を代数閉体とする。$X$ を局所代数的 $k$ スキームとし、$x \in X$ を閉点とする。
{\it $X$ の $x$ における埋め込み次元} とは
$\dim_k \mathfrak m_x/\mathfrak m_x^2$ のことである。
\end{definition}

\noindent
埋め込み次元に関する事実を挙げる。
$k, X, x$ は Definition \ref{varieties-definition-embed-dim} と同じものとする。
\begin{enumerate}
\item $X$ の $x$ における埋め込み次元は、接空間
$T_{X/k, x}$（Definition \ref{varieties-definition-tangent-space}）の
$k$ ベクトル空間としての次元である。
\item $X$ の $x$ における埋め込み次元は、次の全射が存在するような最小の整数
$d \geq 0$ である。
$$
k[[x_1, \ldots, x_d]] \longrightarrow \mathcal{O}_{X, x}^\wedge
$$
ここでこれは $k$ 代数の全射である。
\item $X$ の $x$ における埋め込み次元は、次を満たす最小の整数 $d \geq 0$ である。
開近傍 $U \subset X$ が $x$ を含み、閉埋め込み $U \to Y$ が存在して、
$Y$ は次元 $d$ の $k$ 上滑らかな多様体である。
\item $X$ の $x$ における埋め込み次元は、次を満たす最小の整数 $d \geq 0$ である。
開近傍 $U \subset X$ が $x$ を含み、不分岐射 $U \to \mathbf{A}^d_k$ が存在する。
\item $X \to Y$ が閉埋め込みで、$Y$ が $k$ 上滑らかであるとする。このとき
$X$ の $x$ における埋め込み次元は、次を満たす最小の整数 $d \geq 0$ である。
閉部分スキーム $Z \subset Y$ が存在して
$X \subset Z$、$Z \to \Spec(k)$ が $x$ で滑らか、かつ
$\dim_x(Z) = d$ となる。
\end{enumerate}
これらが必要になれば、正確な結果を定式化し証明を与える。

\medskip\noindent
代数閉でない基礎体または閉でない点の場合を考える。
$k$ を体とし、$X$ を局所代数的 $k$ スキームとする。
$x \in X$ が点ならば、$X$ の $x$ における埋め込み次元にはいくつかの選択肢がある。
すなわち、次を用いることができる。
\begin{enumerate}
\item $\dim_{\kappa(x)}(\mathfrak m_x/\mathfrak m_x^2)$,
\item $\dim_{\kappa(x)}(T_{X/k, x}) =
\dim_{\kappa(x)}(\Omega_{X/k, x} \otimes_{\mathcal{O}_{X, x}} \kappa(x))$
（Lemma \ref{varieties-lemma-tangent-space-cotangent-space}）,
\item 次を満たす最小の整数 $d \geq 0$。開近傍 $U \subset X$ が $x$ を含み、閉埋め込み
$U \to Y$ が存在して、$Y$ は次元 $d$ の $k$ 上滑らかな多様体である。
\end{enumerate}
標数零では (1) $=$ (2) である。ただし $x$ は閉点とする。より一般に、$\kappa(x)$ が
$k$ 上分離代数的ならばこれが成り立つ。Lemma
\ref{varieties-lemma-tangent-space-rational-point} を参照せよ。
幾何学的定義 (3) は「埋め込み次元」という語が想起させる幾何学的直観に最も近いと思われる。
(3) と (2) が同じ数を定義することを示せるので
（これは Lemma \ref{varieties-lemma-immersion-into-affine} から従う）、これを用いる。
用語上、基礎体に相対的な埋め込み次元を取っていることを明示する。

\begin{definition}
\label{varieties-definition-embedding-dimension}
$k$ を体とする。$X$ を局所代数的 $k$ スキームとする。$x \in X$ を点とする。
{\it $X/k$ の $x$ における埋め込み次元} とは
$\dim_{\kappa(x)}(T_{X/k, x})$ のことである。
\end{definition}

\noindent
$(A, \mathfrak m, \kappa)$ が Noether 局所環ならば、{\it $A$ の埋め込み次元} は
$\mathfrak m/\mathfrak m^2$ の $\kappa$ 上の次元として定義されることがある。
上で見たように、$A$ が体 $k$ 上の代数として与えられる場合には、
$\dim_\kappa(\Omega_{A/k} \otimes_A \kappa)$ を用いる方が望ましいことがある。
この量を {\it $A/k$ の埋め込み次元} と呼ぶことにする。この用語のもとで、
$$
\text{embed dim of }X/k\text{ at }x =
\text{embed dim of }\mathcal{O}_{X, x}/k =
\text{embed dim of }\mathcal{O}_{X, x}^\wedge/k
$$
が成り立つ。ただし $k, X, x$ は Definition
\ref{varieties-definition-embedding-dimension} と同じものとする。

\section{Bertini の諸定理}
\label{varieties-section-bertini}

\noindent
この節では次の形の結果を証明する。滑らかな射影多様体 $X$ が体 $k$ 上で与えられたとき、
滑らかである豊富因子 $H \subset X$ が存在する。

\begin{lemma}
\label{varieties-lemma-generate-over-complement}
$k$ を体とする。$X$ を $k$ 上の固有スキームとする。$\mathcal{L}$ を豊富な可逆
$\mathcal{O}_X$ 加群とする。$Z \subset X$ を閉部分スキームとする。このとき整数 $n_0$ が存在し、
すべての $n \geq n_0$ に対して、核 $V_n$、すなわち写像
$\Gamma(X, \mathcal{L}^{\otimes n}) \to \Gamma(Z, \mathcal{L}^{\otimes n}|_Z)$
の核は $\mathcal{L}^{\otimes n}|_{X \setminus Z}$ を生成し、標準射
$$
X \setminus Z \longrightarrow \mathbf{P}(V_n)
$$
は $k$ 上のスキームの埋め込みである。
\end{lemma}

\begin{proof}
$\mathcal{I} \subset \mathcal{O}_X$ を $Z$ の準連接イデアル層とする。包含
$\mathcal{I} \otimes_{\mathcal{O}_X} \mathcal{L}^{\otimes n} \subset
\mathcal{L}^{\otimes n}$ によって
$V_n = \Gamma(X, \mathcal{I} \otimes_{\mathcal{O}_X} \mathcal{L}^{\otimes n})$
となることに注意する。$n_1$ を、$n \geq n_1$ ならば層
$\mathcal{I} \otimes \mathcal{L}^{\otimes n}$ が大域生成されるように選ぶ。
Properties, Proposition \ref{properties-proposition-characterize-ample} を参照せよ。
従って $V_n$ は $\mathcal{L}^{\otimes n}|_{X \setminus Z}$ を
$n \geq n_1$ に対し生成する。

\medskip\noindent
$n \geq n_1$ に対し、制限写像を
$\psi_n : V_n \to
\Gamma(X \setminus Z, \mathcal{L}^{\otimes n}|_{X \setminus Z})$
と書く。Constructions, Example \ref{constructions-example-projective-space} により標準射
$$
\varphi = \varphi_{\mathcal{L}^{\otimes n}|_{X \setminus Z}, \psi_n} :
X \setminus Z
\longrightarrow
\mathbf{P}(V_n)
$$
を得る。$n_2$ を、すべての $n \geq n_2$ に対して可逆層
$\mathcal{L}^{\otimes n}$ が $X$ 上非常に豊富となるように選ぶ。
$n_0 = n_1 + n_2$ でよいと主張する。

\medskip\noindent
主張を証明する。$n \geq n_0$ とし、$n = n_1 + n'$ と書く。
$x \in X \setminus Z$ に対し、$s_1 \in V_1$ を $x$ で消えないように選べる。
$V' = \Gamma(X, \mathcal{L}^{\otimes n'})$ と置く。$n$ と $n'$ の選び方により、
対応する射 $\varphi' : X \to \mathbf{P}(V')$ は閉埋め込みである。従って
$s' \in \Gamma(X, \mathcal{L}^{\otimes n'})$ を $x$ で消えないように選べば、
$X_{s'} = (\varphi')^{-1}(D_+(s'))$（Constructions, Lemma
\ref{constructions-lemma-invertible-map-into-proj} を参照）はアフィンであり、
$X_{s'} \to D_+(s')$ は閉埋め込みである。
すると $s = s_1 \otimes s' \in V_n$ は $x$ で消えない。
$D_+(s) \subset \mathbf{P}(V_n)$ がこの射影空間の対応するアフィン開集合を表すならば、
$\varphi^{-1}(D_+(s)) = X_s \subset X \setminus Z$ である（上の参照を見よ）。
開集合 $X_s = X_{s'} \cap X_{s_1}$ はアフィンである。Properties, Lemma
\ref{properties-lemma-affine-cap-s-open} を参照せよ。
射 $X_s \to D_+(s)$ を定める環準同型
$$
\text{Sym}(V)_{(s)} \longrightarrow \mathcal{O}_X(X_s)
$$
を考える。$X_{s'} \to D_+(s')$ は閉埋め込みなので、元
$$
\frac{s_1 \otimes t'}{s_1 \otimes s'}
$$
で $t' \in V'$ であるものの像は、
$\mathcal{O}_X(X_{s'}) \to \mathcal{O}_X(X_s)$ の像を生成する。
$X_s \to X_{s'}$ は開埋め込みなので、これは $X_s \to D_+(s)$ が
アフィンスキームの埋め込みであることを意味する（下を参照）。
従って Morphisms, Lemma \ref{morphisms-lemma-check-immersion} により
$\varphi_n$ は埋め込みである。

\medskip\noindent
$a : A' \to A$ と $c : B \to A$ を環準同型とし、$\Spec(a)$ は埋め込みで
$\Im(a) \subset \Im(c)$ であると仮定する。
$B' = A' \times_A B$ と置き、射影を $b : B' \to B$ および
$c' : B' \to A'$ とする。仮定により $c'$ は全射なので、
$\Spec(c')$ は閉埋め込みである。従って
$\Spec(c') \circ \Spec(a)$ は埋め込みである
（Schemes, Lemma \ref{schemes-lemma-composition-immersion}）。
すると $\Spec(c)$ は埋め込みでなければならない。実際これは埋め込み
$\Spec(c') \circ \Spec(a) = \Spec(b) \circ \Spec(c)$ を分解するからである。
Morphisms, Lemma \ref{morphisms-lemma-immersion-permanence} を参照せよ。
\end{proof}

\begin{situation}
\label{varieties-situation-family-divisors}
$k$ を体、$X$ を $k$ 上のスキーム、$\mathcal{L}$ を可逆 $\mathcal{O}_X$ 加群、
$V$ を有限次元 $k$ ベクトル空間とし、
$\psi : V \to \Gamma(X, \mathcal{L})$ を $k$ 線形写像とする。
$\dim(V) = r$ とし、$v_1, \ldots, v_r$ が $V$ の基底であるとする。
このとき「普遍因子」
$$
H_{univ} = Z(s_{univ}) \subset  \mathbf{A}^r \times_k X
$$
を、切断
$$
s_{univ} = \sum\nolimits_{i = 1, \ldots, r} x_i \psi(v_i) \in
\Gamma(\mathbf{A}^r \times_k X, \text{pr}_2^*\mathcal{L})
$$
の零スキーム（Divisors, Definition \ref{divisors-definition-zero-scheme-s}）として得る。
体拡大 $k'/k$ に対し、$k'$ 点 $v \in \mathbf{A}^r_k(k')$ は
ベクトル $(a_1, \ldots, a_r)$ に対応し、その各成分は $k'$ の元である。
従って一方では $v$ を元
$v = \sum_{i = 1, \ldots, r} a_i v_i \in V \otimes_k k'$ とみなせ、他方では
$v$ に切断
$$
\psi(v) = \sum\nolimits_{i = 1, \ldots, r} a_i \psi(v_i) \in
\Gamma(X_{k'}, \mathcal{L}|_{X_{k'}})
$$
を対応させられる。この記法のもとで、$H_{univ}$ の
$v \in V \otimes k'$ 上のファイバーは $\psi(v)$ の零スキームであることが明らかである。
式では
$$
H_v = H_{univ, v} = Z(\psi(v))
$$
となる。この共通の値を、表示のとおり $H_v$ と書く。最後に、この状況で $P$ を
$v \in V \otimes_k k'$ の性質とする。ただし $k'/k$ は任意の体拡大である
\footnote{例えば、$H_v$ が $k'$ 上滑らかである、または $k'$ 上幾何学的に既約であるという
条件を考えられる。}。$P$ が $v \in V \otimes_k k'$ に対して
{\it 一般に成り立つ} とは、空でない Zariski 開集合
$U \subset \mathbf{A}^r_k$ が存在し、$v$ が $k'$ 点として $U$ に対応するならば、
任意の $k'/k$ に対して $P(v)$ が成り立つことをいう。
\end{situation}

\begin{lemma}
\label{varieties-lemma-bertini}
Situation \ref{varieties-situation-family-divisors} において、次を仮定する。
\begin{enumerate}
\item $X$ は $k$ 上滑らかである。
\item $\psi : V \to \Gamma(X, \mathcal{L})$ の像は $\mathcal{L}$ を生成する。
\item 対応する射
$\varphi_{\mathcal{L}, \psi} : X \to \mathbf{P}(V)$ は埋め込みである。
\end{enumerate}
このとき、一般の $v \in V \otimes_k k'$ に対してスキーム
$H_v$ は $k'$ 上滑らかである。
\end{lemma}

\begin{proof}
（$X$ は射影空間の局所閉部分スキームとして分離的かつ有限型であることに注意する。）
補題の主張の前に導入した記法を用いる。射影
$$
\xymatrix{
\mathbf{A}^r_k \times_k X \ar[d] &
H_{univ} \ar[l] \ar[ld]^p \ar[r] \ar[rd]_q &
\mathbf{A}^r_k \times_k X \ar[d] \\
X & &
\mathbf{A}^r_k
}
$$
を考える。$\Sigma \subset H_{univ}$ を射
$q : H_{univ} \to \mathbf{A}^r_k$ の特異軌跡、すなわち $q$ が滑らかでない点の集合とする。
滑らかな軌跡は定義により開なので、$\Sigma$ は閉である。
滑らかな射のファイバーは滑らかなので、$q(\Sigma)$ が
$\mathbf{A}^r_k$ の真の閉部分集合に含まれることを示せば十分である。
$\Sigma$ は被約誘導スキーム構造を入れると $k$ 上有限型のスキームなので、
$\dim(\Sigma) < r$ を示せば十分である。これは Lemma
\ref{varieties-lemma-dimension-fibres-locally-algebraic} から従う。
$k$ をより大きな体で置き換えても次元は変わらないので
（Morphisms, Lemma \ref{morphisms-lemma-dimension-fibre-after-base-change}）、
$k$ は代数閉であると仮定してよいし、そう仮定する。
次元論（Lemma \ref{varieties-lemma-dimension-fibres-locally-algebraic}）により、
閉点 $x \in X \setminus Z$ に対して
$p^{-1}(\{x\}) \cap \Sigma$ の次元が $< r - \dim(X')$ であることを示せば十分である。
ここで $X'$ は $X$ の唯一の既約成分で、$x$ を含むものである。
$X$ は $k$ 上滑らかであり、$x$ は閉点なので
$\dim(X') = \dim \mathfrak m_x/\mathfrak m_x^2$
である（Morphisms, Lemma \ref{morphisms-lemma-smooth-omega-finite-locally-free}
および Algebra, Lemma \ref{algebra-lemma-rank-omega}）。従って、すべての閉点
$x \in X$ に対して
$$
\dim p^{-1}(x) \cap \Sigma < r - \dim \mathfrak m_x/\mathfrak m_x^2
$$
ならばよい。

\medskip\noindent
$V$ が $\mathcal{L}$ を大域生成したので、$X'$ を $X$ の各既約成分とすると、
$\mathbf{A}^r$ の空でない Zariski 開集合が存在して、その上の $q$ のファイバーは
$X'$ を含まない。（例えば、$x' \in X'$ を閉点とすれば、
ベクトル $v \in V$ で $\psi(v)$ が $x'$ で消えないものに対応する開集合を取れる。
これは $\mathbf{A}^r_k$ における超平面の補集合である。）
$U \subset \mathbf{A}^r$ をこれらの開集合の空でない共通部分とする。
すると $q^{-1}(U) \to U$ のファイバーは $U \times_k X \to U$ のファイバー上の
有効 Cartier 因子である（整スキーム上の可逆加群の消えない切断は正則切断だからである）。
従って射 $q^{-1}(U) \to U$ は Divisors, Lemma
\ref{divisors-lemma-fibre-Cartier} により平坦である。
従って $x \in X$ を閉点、$v \in V = \mathbf{A}^r_k(k)$ とし、
$(x, v) \in H_{univ}$、すなわち $x \in H_v$ ならば、$q$ が $(x, v)$ で滑らかであることと
$H_v$ が $x$ において滑らかであることは同値である。Morphisms, Lemma
\ref{morphisms-lemma-smooth-at-point} を参照せよ。

\medskip\noindent
像 $\psi(v)_x$ を茎 $\mathcal{L}_x$ において考える。これは
$v \in V$ に対応する切断の像である。次が成り立つ。
$$
x \in H_v \Leftrightarrow \psi(v)_x \in \mathfrak m_x\mathcal{L}_x
$$
これが真ならば、
$$
H_v\text{ singular at }x \Leftrightarrow
\psi(v)_x \in \mathfrak m_x^2\mathcal{L}_x
$$
が成り立つ。すなわち、$\psi(v)_x$ が $\mathfrak m_x^2\mathcal{L}_x$ に含まれない
$\Leftrightarrow$
$H_v \subset X$ の $x$ における局所方程式が $\mathfrak m_x^2$ に含まれない
$\Leftrightarrow$
$\mathcal{O}_{H_v, x}$ は正則である
（Algebra, Lemma \ref{algebra-lemma-regular-ring-CM}）
$\Leftrightarrow$
$H_v$ は $x$ において $k$ 上滑らかである
（Algebra, Lemma \ref{algebra-lemma-separable-smooth}）。
従って $p^{-1}(x) \cap \Sigma$ の閉点は、$v \in V$ で
$\psi(v)_x \in \mathfrak m_x^2\mathcal{L}_x$ を満たすものに対応する。
しかし $\varphi_{\mathcal{L}, \psi}$ は埋め込みなので、写像
$$
V \longrightarrow \mathcal{L}_x/\mathfrak m_x^2\mathcal{L}_x
$$
は全射である（細部は省略する）。上の議論により、軌跡
$p^{-1}(x) \cap \Sigma$ を $V$ の部分空間とみたとき、その閉点はこの写像の核であり、
従って望むとおり次元 $r - \dim \mathfrak m_x/\mathfrak m_x^2 - 1$ をもつ。
\end{proof}

\section{Enriques--Severi--Zariski}
\label{varieties-section-vanishing-negative}

\noindent
この節では、「非常に負」の可逆加群で捻ると低次コホモロジーが消える、という形のいくつかの
結果を証明する。また、次元 $\geq 2$ の正規固有スキームの超曲面切断が連結であることも導く。

\begin{lemma}
\label{varieties-lemma-vanishin-h0-negative}
$k$ を体とする。$X$ を $k$ 上の固有スキームとする。$\mathcal{L}$ を豊富な可逆
$\mathcal{O}_X$ 加群とする。$\mathcal{F}$ を連接 $\mathcal{O}_X$ 加群とする。
$\text{Ass}(\mathcal{F})$ が閉点を含まないならば、
$\Gamma(X, \mathcal{F} \otimes_{\mathcal{O}_X} \mathcal{L}^{\otimes n}) = 0$
が $n \ll 0$ に対して成り立つ。
\end{lemma}

\begin{proof}
連接 $\mathcal{O}_X$ 加群 $\mathcal{F}$ に対し、$\mathcal{P}(\mathcal{F})$ を次の性質とする。
$n_0 \in \mathbf{Z}$ が存在し、$n \leq n_0$ ならば、すべての切断 $s$、すなわち
$\mathcal{F} \otimes_{\mathcal{O}_X} \mathcal{L}^{\otimes n}$ の切断の台が閉点だけからなる。
$\text{Ass}(\mathcal{F}) =
\text{Ass}(\mathcal{F} \otimes_{\mathcal{O}_X} \mathcal{L}^{\otimes n})$
なので、$\mathcal{P}$ が $X$ 上のすべての連接加群について成り立つことを示せば十分である。
このために Cohomology of Schemes, Lemma
\ref{coherent-lemma-property-higher-rank-cohomological} の条件 (1), (2), (3) が満たされることを示す。

\medskip\noindent
条件 (1) を見るため、連接 $\mathcal{O}_X$ 加群の短完全列
$$
0 \to \mathcal{F}_1 \to \mathcal{F} \to \mathcal{F}_2 \to 0
$$
が与えられ、$\mathcal{P}$ が $\mathcal{F}_i$, $i = 1, 2$ について成り立つと仮定する。
得られる境界値を $n_1, n_2$ とする。
$\mathcal{F}'_2 \subset \mathcal{F}_2$ を、台が閉点の有限集合である最大の連接部分加群とする。
$\mathcal{I} \subset \mathcal{O}_X$ を $\mathcal{F}'_2$ の零化イデアルとする。
$\mathcal{L}$ は豊富なので、$e > 0$ で
$\mathcal{I} \otimes_{\mathcal{O}_X} \mathcal{L}^{\otimes e}$
が大域生成されるものを取れる。
$n_0 = \min(n_2, n_1 - e)$ と置く。$n \leq n_0$ とし、$t$ を
$\mathcal{F} \otimes \mathcal{L}^{\otimes n}$ の大域切断とする。
$t$ の $\mathcal{F}_2 \otimes \mathcal{L}^{\otimes n}$ における像は
$\mathcal{F}'_2 \otimes \mathcal{L}^{\otimes n}$ に入る。これは $n \leq n_2$ だからである。
従って任意の
$s \in \Gamma(X, \mathcal{I} \otimes_{\mathcal{O}_X} \mathcal{L}^{\otimes e})$
に対して積 $t \otimes s$ は
$\mathcal{F}_1 \otimes \mathcal{L}^{\otimes n + e}$ に入る。
ゆえに $t \otimes s$ の台は、$\text{Ass}(\mathcal{F}_1)$ に含まれる閉点の有限集合に
含まれる。これは $n + e \leq n_1$ だからである。
$e$ の選び方により、$s$ を任意の点で可逆となるように選べる。ただしその点は
$\mathcal{F}'_2$ の台にないものとする。従って、
$t$ の台は、$\text{Ass}(\mathcal{F}_1)$ に含まれる閉点の有限集合と
$\text{Ass}(\mathcal{F}_2)$ に含まれる閉点の有限集合との合併に含まれる。
これで条件 (1) の証明が完了する。

\medskip\noindent
条件 (2) は直ちに成り立つ。

\medskip\noindent
条件 (3) については $\mathcal{G} = \mathcal{O}_Z$ を選ぶ。
この場合、$Z$ が $X$ の閉点ならば示すことはない。$\dim(Z) > 0$ ならば、
$\Gamma(Z, \mathcal{L}^{\otimes n}|_Z) = 0$ が $n < 0$ に対して成り立つことを示す。
実際、$s$ を $\mathcal{L}|_Z$ の負の冪の零でない切断とする。
$t$ を $\mathcal{L}|_Z$ の正の冪の零でない切断として選ぶ
（$\mathcal{L}$ が豊富なので可能である。Properties, Proposition
\ref{properties-proposition-characterize-ample} を参照せよ）。
すると $s^{\deg(t)} \otimes t^{\deg(s)}$ は $\mathcal{O}_Z$ の零でない大域切断であり
（$Z$ は整だからである）、従って単元である
（Lemma \ref{varieties-lemma-proper-geometrically-reduced-global-sections}）。
これは $t$ が $\mathcal{L}$ の正の冪を自明化する切断であることを意味する。
従って関数 $n \mapsto \dim_k \Gamma(X, \mathcal{L}^{\otimes n})$ は正整数の無限集合上で
有界となるが、これは漸近的 Riemann--Roch
（Proposition \ref{varieties-proposition-asymptotic-riemann-roch}）および
$\dim(Z) > 0$ に反する。
\end{proof}

\begin{lemma}[Enriques--Severi--Zariski]
\label{varieties-lemma-vanishin-h1-negative}
$k$ を体とする。$X$ を $k$ 上の固有スキームとする。$\mathcal{L}$ を豊富な可逆
$\mathcal{O}_X$ 加群とする。$\mathcal{F}$ を連接 $\mathcal{O}_X$ 加群とする。
すべての閉点 $x \in X$ に対し $\text{depth}(\mathcal{F}_x) \geq 2$ であると仮定する。
このとき
$H^1(X, \mathcal{F} \otimes_{\mathcal{O}_X} \mathcal{L}^{\otimes m}) = 0$
が $m \ll 0$ に対して成り立つ。
\end{lemma}

\begin{proof}
閉埋め込み $i : X \to \mathbf{P}^n_k$ で、
$i^*\mathcal{O}(1) \cong \mathcal{L}^{\otimes e}$ がある $e > 0$ に対して成り立つものを選ぶ
（Morphisms, Lemma
\ref{morphisms-lemma-finite-type-over-affine-ample-very-ample} を参照）。
すると $\mathbf{P}^n_k$ 上の
$$
\mathcal{G} =
i_*(\mathcal{F} \oplus \mathcal{F} \otimes \mathcal{L} \oplus \ldots
\oplus \mathcal{F} \otimes \mathcal{L}^{\otimes e - 1})
\quad\text{and}\quad \mathcal{O}(1)
$$
について補題を証明すれば十分である。実際、
$$
H^1(\mathbf{P}^n_k, \mathcal{G}(m)) =
\bigoplus\nolimits_{j = 0, \ldots, e - 1}
H^1(X, \mathcal{F} \otimes \mathcal{L}^{\otimes j + me})
$$
が Cohomology of Schemes, Lemma
\ref{coherent-lemma-relative-affine-cohomology} により成り立つ。
また、$y \in \mathbf{P}^n_k$ が閉点ならば、
$\text{depth}(\mathcal{G}_y) = \infty$ が $y \not \in i(X)$ の場合に成り立ち、
$\text{depth}(\mathcal{G}_y) = \text{depth}(\mathcal{F}_x)$ が $y = i(x)$ の場合に成り立つ。
後者の場合には
$\mathcal{G}_y \cong \mathcal{F}_x^{\oplus e}$ が
$\mathcal{O}_{\mathbf{P}^n_k, x}$ 上の加群として成り立ち、例えば Algebra, Lemma
\ref{algebra-lemma-depth-goes-down-finite} を使って等式を得られるからである。

\medskip\noindent
$X = \mathbf{P}^n_k$、$\mathcal{L} = \mathcal{O}(1)$ と仮定し、
$k$ は無限であるとする。完全列
$$
0 \to \mathcal{F}(-1) \xrightarrow{s} \mathcal{F} \to \mathcal{G} \to 0
$$
を定める $s \in H^0(\mathbf{P}^1_k, \mathcal{O}(1))$ を選ぶ。Lemma
\ref{varieties-lemma-exact-sequence-induction} におけるものと同じである。写像
$\mathcal{F}(-1) \to \mathcal{F}$ はアフィン局所的に $\mathcal{F}$ 上の非零因子による
乗法で与えられるので、閉点 $x \in \mathbf{P}^n_k$ に対し
$\text{depth}(\mathcal{G}_x) \geq 1$ である。Algebra, Lemma
\ref{algebra-lemma-depth-drops-by-one} を参照せよ。
従って Lemma \ref{varieties-lemma-vanishin-h0-negative} により、
$H^0(\mathcal{G}(m)) = 0$ が $m \ll 0$ に対して成り立つ。
捻った後のコホモロジー長完全列
（Remark \ref{varieties-remark-exact-sequence-induction-cohomology} を参照）を見ると、数列
$$
\dim H^1(\mathbf{P}^n_k, \mathcal{F}(m))
$$
は $m \leq m_0$ で安定する。そのような整数を $m_0$ とする。
$N$ を $m \leq m_0$ に対するこれらの空間の共通次元とする。$N = 0$ を示す必要がある。

\medskip\noindent
$d > 0$ および $m \leq m_0$ に対し、双線形写像
$$
H^0(\mathbf{P}^n_k, \mathcal{O}(d)) \times
H^1(\mathbf{P}^n_k, \mathcal{F}(m - d))
\longrightarrow
H^1(\mathbf{P}^n_k, \mathcal{F}(m))
$$
を考える。線形代数により、余次元 $\leq N^2$ の部分空間
$V_m \subset H^0(\mathbf{P}^n_k, \mathcal{O}(d))$ が存在し、
$s' \in V_m$ による乗法は
$H^1(\mathbf{P}^n_k, \mathcal{F}(m - d))$ を零化する。
$m' < m \leq m_0$ に対し、図式
$$
\xymatrix{
H^0(\mathbf{P}^n_k, \mathcal{O}(d)) \times
H^1(\mathbf{P}^n_k, \mathcal{F}(m' - d)) \ar[r] \ar[d]^{1 \times s^{m' - m}} &
H^1(\mathbf{P}^n_k, \mathcal{F}(m')) \ar[d]^{s^{m' - m}}\\
H^0(\mathbf{P}^n_k, \mathcal{O}(d)) \times
H^1(\mathbf{P}^n_k, \mathcal{F}(m - d)) \ar[r] &
H^1(\mathbf{P}^n_k, \mathcal{F}(m))
}
$$
は可換で、縦の写像は同型である。従って $V_m = V$ は
$m \leq m_0$ に依存しない。$x \in \text{Ass}(\mathcal{F})$ に対し
$Z = \overline{\{x\}}$ と置く。$d$ が十分大きければ線形写像
$$
H^0(\mathbf{P}^n_k, \mathcal{O}(d)) \to H^0(Z, \mathcal{O}(d)|_Z)
$$
の階数は $> N^2$ である。これは $\dim(Z) \geq 1$ だからである
（例えば漸近的 Riemann--Roch と豊富性 $\mathcal{O}(1)$ から従う。詳細は省略する）。
従って $s' \in V$ を、$s'$ が $\mathcal{F}$ のどの随伴点でも消えないように選べる
（随伴点の集合が有限であることを用いる）。すると
$$
0 \to \mathcal{F}(-d) \xrightarrow{s'} \mathcal{F} \to \mathcal{G}' \to 0
$$
を得て、前と同様に、$s'$ による
$H^1(\mathbf{P}^n_k, \mathcal{F}(m - d))$ 上の乗法は $m \ll 0$ に対して単射であると結論する。
これは $s'$ の選び方に反する。ただし $N = 0$ ならば反しないので、望む結果を得る。

\medskip\noindent
$k$ が有限である場合をまだ扱う必要がある。この場合、$K/k$ を任意の無限代数体拡大とする。
$\mathcal{F}_K$ および $\mathcal{L}_K$ を、$\mathcal{F}$ および $\mathcal{L}$ の
$X_K = \Spec(K) \times_{\Spec(k)} X$ への引き戻しと書く。このとき
$$
H^1(X_K, \mathcal{F}_K \otimes \mathcal{L}_K^{\otimes m}) =
H^1(X, \mathcal{F} \otimes \mathcal{L}^{\otimes m}) \otimes_k K
$$
が Cohomology of Schemes, Lemma
\ref{coherent-lemma-flat-base-change-cohomology} により成り立つ。
一方、$x_K$ を $X_K$ の閉点とすると、これは閉点 $x$ に写り、その点は $X$ に属する。
$K/k$ は代数拡大だからである。
環準同型 $\mathcal{O}_{X, x} \to \mathcal{O}_{X_K, x_K}$ は平坦である
（Lemma \ref{varieties-lemma-change-fields-flat}）。従って
$$
\text{depth}(\mathcal{F}_{x_K}) =
\text{depth}(\mathcal{F}_x \otimes_{\mathcal{O}_{X, x}}
\mathcal{O}_{X_K, x_K}) \geq
\text{depth}(\mathcal{F}_x)
$$
が Algebra, Lemma \ref{algebra-lemma-apply-grothendieck-module} により成り立つ
（実際ここでは等号が成り立つが必要ない）。
従って $k$ 上の結果は無限体 $K$ 上の結果から従い、証明が完了する。
\end{proof}

\begin{lemma}
\label{varieties-lemma-connectedness-ample-divisor}
$k$ を体とする。$X$ を $k$ 上の固有スキームとする。$\mathcal{L}$ を豊富な可逆
$\mathcal{O}_X$ 加群とする。$s \in \Gamma(X, \mathcal{L})$ とする。次を仮定する。
\begin{enumerate}
\item $s$ は正則切断である
（Divisors, Definition \ref{divisors-definition-regular-section}）。
\item すべての閉点 $x \in X$ に対して
$\text{depth}(\mathcal{O}_{X, x}) \geq 2$ である。
\item $X$ は連結である。
\end{enumerate}
このとき $Z(s)$、すなわち $s$ の零スキームは連結である。
\end{lemma}

\begin{proof}
$s$ は正則切断なので、$s^n \in \Gamma(X, \mathcal{L}^{\otimes n})$ も
すべての $n > 1$ に対し正則切断である。
さらに、包含射 $Z(s) \to Z(s^n)$ は台となる位相空間上の全単射である。
従って $Z(s)$ が非連結ならば $Z(s^n)$ も非連結である。ここで標準短完全列
$$
0 \to \mathcal{L}^{\otimes -n} \xrightarrow{s^n}
\mathcal{O}_X \to \mathcal{O}_{Z(s^n)} \to 0
$$
を考える。$k$ 代数 $R_n = \Gamma(X, \mathcal{O}_{Z(s^n)})$ を考える。
$Z(s)$ が非連結、すなわち $Z(s^n)$ が非連結ならば、
$R_n$ は $Z(s^n) = \emptyset$ の場合には零であり、または
$R_n$ は非自明な冪等元を含む。後者は $Z(s^n) = U \amalg V$ で
$U, V \subset Z(s^n)$ が開かつ空でない場合である（読者は Lemma
\ref{varieties-lemma-proper-geometrically-reduced-global-sections} を参照するとよい）。
従って写像 $\Gamma(X, \mathcal{O}_X) \to R_n$ は同型ではありえない。
よって $H^0(X, \mathcal{L}^{\otimes -n})$ または
$H^1(X, \mathcal{L}^{\otimes -n})$ のいずれかが、無限個の正整数 $n$ に対して零でない。
これは Lemma \ref{varieties-lemma-vanishin-h0-negative} または
\ref{varieties-lemma-vanishin-h1-negative} に反し、証明が完了する。
\end{proof}
